\section{Calculus of blowups along smooth centres}
\label{sec:blowup-calculus}

This section collects the facts about a single blowup along a smooth centre, and about finite
sequences of them, that the canonical resolution of marked ideals (Wlodarczyk's algorithm, in the analytic setting) uses: the exceptional divisor, the controlled transform of
an ideal, and the restriction of a blowup to an open subset. The sources are written for smooth
algebraic varieties (Wlodarczyk) or for manifolds and analytic spaces (Bierstone--Milman); here
everything is stated for complex manifolds and for the blowup \(\pi\colon Y\to X\) along a
nowhere dense closed submanifold \(C\) in the sense of Definition~\ref{ana:blowups}, i.e.\ by its
universal property (its existence is Lemma~\ref{ana:blowup-existence}). The route differs from the
sources as follows.
\begin{enumerate}
\item Ideals are ideal sheaves of \(\mathcal O_X\) recorded by their germs
(Definition~\ref{ana:ideal-sheaves}), and the controlled transform of a marked ideal
\((\mathfrak a,\mu)\) is \emph{defined} as the colon ideal \((\pi^*\mathfrak a:\mathcal I_D^{\mu})\),
which needs no hypothesis; it is the quotient \(\mathcal I_D^{-\mu}\pi^*\mathfrak a\) of the
sources as soon as \(\mathfrak a\subset\mathcal I_C^{\mu}\).
\item The hypothesis of the sources on the centre, \(C\subset\operatorname{supp}(\mathfrak a,\mu)=
\{x:\operatorname{ord}_x\mathfrak a\ge\mu\}\), enters only through its consequence
\(\mathfrak a\subset\mathcal I_C^{\mu}\) (Wlodarczyk, \S2.2, Lemma, part (1)); the statements below
take the inclusion as their hypothesis.
\item Open embeddings take the place of the smooth (\'etale) morphisms of the algebraic sources.
\end{enumerate}

\begin{lemma}[The exceptional divisor]
\label{bcalc:exceptional-divisor}
\uses{ana:blowups,ana:ideal-sheaves,ana:blowup-exceptional-image}
\source{wlodarczyk05-hironaka}
\lean{Hironaka.BlowupExceptionalDivisor.exceptionalDivisor,Hironaka.BlowupExceptionalDivisor.exceptionalDivisor_ideal,Hironaka.BlowupExceptionalDivisor.support_exceptionalDivisor}\upstreamlean{TranscendentalNakamaye.Pending.BlowupExceptionalDivisor.exceptionalLocus_subset_support_exceptionalDivisor,TranscendentalNakamaye.Pending.BlowupExceptionalDivisor.interior_support_exceptionalDivisor}\lean{Hironaka.BlowupExceptionalDivisor.interior_exceptionalLocus_eq_empty,Hironaka.BlowupExceptionalDivisor.pullback_idealPow_vanishingIdeal,Hironaka.BlowupExceptionalDivisor.isCoherent_vanishingIdeal_centre,Hironaka.BlowupExceptionalDivisor.totalTransform,Hironaka.BlowupExceptionalDivisor.support_totalTransform,Hironaka.BlowupExceptionalDivisor.zeroLocus_vanishingIdeal,Hironaka.BlowupExceptionalDivisor.zeroLocus_pullback,Hironaka.BlowupExceptionalDivisor.pullback_unitIdeal,Hironaka.BlowupExceptionalDivisor.pullback_idealPow,Hironaka.BlowupExceptionalDivisor.isInvertible_unitIdeal,Hironaka.BlowupExceptionalDivisor.isInvertible_idealPow,Hironaka.BlowupExceptionalDivisor.isCoherent_ideal}
\mixedleanok
Let \(\pi\colon Y\to X\) be the blowup of \(X\) along \(C\). The ideal \(\pi^*\mathcal I_C=\mathcal I_C\mathcal O_Y\)
is invertible; it is the ideal \(\mathcal O_Y(-D)\) of an effective divisor \(D\), the \emph{exceptional divisor}, and
\begin{enumerate}
\item \(\operatorname{supp}D=\pi^{-1}(C)\), which is nowhere dense in \(Y\);
\item \(\operatorname{Exc}\pi\subset\pi^{-1}(C)\), and \(\operatorname{Exc}\pi\) is nowhere dense;
\item \(\pi^*(\mathcal I_C^{\mu})=\mathcal O_Y(-\mu D)\) for every \(\mu\ge0\), and \(\mathcal O_Y(-\mu D)\) is invertible;
\item for an effective divisor \(H\) on \(X\) the total transform \(\pi^*H\), with ideal \(\pi^*\mathcal O_X(-H)\), is an
effective divisor with \(\operatorname{supp}\pi^*H=\pi^{-1}(\operatorname{supp}H)\).
\end{enumerate}
\end{lemma}
\begin{proof}
\mixedleanok
Invertibility of \(\pi^*\mathcal I_C\) is part of Definition~\ref{ana:blowups}. For a coherent ideal \(\mathfrak a\) and a
holomorphic map \(f\), the zero locus of \(f^*\mathfrak a\) is \(f^{-1}\) of the zero locus of \(\mathfrak a\): a germ of
\(\mathfrak a\) not vanishing at \(f(y)\) pulls back to a germ of \(f^*\mathfrak a\) not vanishing at \(y\), and conversely every
germ of \(f^*\mathfrak a\) at \(y\) is a combination of pullbacks of germs of \(\mathfrak a\) at \(f(y)\). The zero locus of
\(\mathcal I_A\) is \(A\) for a closed set \(A\): off \(A\) the constant \(1\) belongs to \(\mathcal I_A\). Since \(\mathcal I_C\) is
coherent, this gives \(\operatorname{supp}D=\pi^{-1}(C)\), and likewise (4), the ideal of an effective divisor being coherent and
total transforms of invertible ideals along maps whose exceptional locus is nowhere dense being invertible. The preimage
\(\pi^{-1}(C)\) is nowhere dense because a local generator of the invertible ideal \(\pi^*\mathcal I_C\) vanishes on it and
vanishes identically on no nonempty open set (Lemma~\ref{ana:blowup-exceptional-image}). Off \(\pi^{-1}(C)\) the map \(\pi\) is a
biholomorphism onto \(X\setminus C\) (the universal property applied to the inclusion of \(X\setminus C\), where \(\mathcal I_C\)
is the unit ideal), so \(\operatorname{Exc}\pi\subset\pi^{-1}(C)\) and (2) follows from (1). Finally
\(\pi^*(\mathcal I_C^{\mu})=(\pi^*\mathcal I_C)^{\mu}\) because pullback commutes with products of coherent ideals and sends the unit
ideal to the unit ideal, and powers of invertible ideals are invertible.
\end{proof}

\begin{lemma}[Cancelling an invertible ideal]
\label{bcalc:cancellation}
\uses{ana:ideal-sheaves,ana:invertible-colon}
\lean{Hironaka.BlowupControlledTransform.colon_mul_eq,Hironaka.BlowupControlledTransform.colon_unitIdeal,Hironaka.BlowupControlledTransform.pullback_id,Hironaka.BlowupControlledTransform.idealPow_unitIdeal,Hironaka.BlowupControlledTransform.unitIdeal_mul,Hironaka.BlowupControlledTransform.mul_mul_mul_comm,Hironaka.BlowupControlledTransform.idealPow_mul}
\leanok
Let \(\mathfrak l\) be an invertible ideal and \(\mathfrak c\) a coherent ideal. Then \((\mathfrak l\mathfrak c:\mathfrak l)=\mathfrak c\); in
particular \((\mathfrak c:\mathcal O)=\mathfrak c\). Moreover, for coherent ideals, \(\mathrm{id}^*\mathfrak c=\mathfrak c\),
\(\mathcal O^{\mu}=\mathcal O\), \(\mathcal O\mathfrak c=\mathfrak c\), \((\mathfrak a\mathfrak b)(\mathfrak c\mathfrak d)=(\mathfrak a\mathfrak c)(\mathfrak b\mathfrak d)\)
and \((\mathfrak a\mathfrak b)^{\mu}=\mathfrak a^{\mu}\mathfrak b^{\mu}\).
\end{lemma}
\begin{proof}
\leanok
If \(f\in\mathfrak c_x\) then \(gf\in\mathfrak l\mathfrak c\) for every \(g\in\mathfrak l_x\). Conversely let \(f\) be holomorphic at
\(x\) with \(\mathfrak l_xf\subset(\mathfrak l\mathfrak c)_x\), and let \(h\) generate \(\mathfrak l\) near \(x\). Then
\(hf=\sum_jc_jp_jq_j\) near \(x\) with \(p_j\in\mathfrak l_x\), \(q_j\in\mathfrak c_x\), and \(p_j=e_jh\); so
\(h\bigl(f-\sum_jc_je_jq_j\bigr)=0\) near \(x\). As \(h\) vanishes identically on no nonempty open set, the holomorphic function
\(f-\sum_jc_je_jq_j\) vanishes on a dense subset of a neighbourhood of \(x\), hence near \(x\), and \(f\in\mathfrak c_x\) since
\(\mathfrak c\) is coherent. The case \(\mathfrak l=\mathcal O\) is the second claim (\(\mathcal O\mathfrak c=\mathfrak c\)). The
remaining identities follow from commutativity and associativity of the product of coherent ideals
(Lemma~\ref{ana:invertible-colon} and its context) by induction on \(\mu\).
\end{proof}

\begin{definition}[Controlled transform]
\label{bcalc:controlled-transform-def}
\uses{bcalc:exceptional-divisor}
\source{wlodarczyk05-hironaka,bierstone-milman08-functoriality}
\lean{Hironaka.BlowupControlledTransform.controlledTransform}
\leanok
Let \(\pi\colon Y\to X\) be the blowup along \(C\) with exceptional divisor \(D\), \(\mathfrak a\) an ideal on \(X\) and
\(\mu\ge0\). The \emph{controlled transform} of the marked ideal \((\mathfrak a,\mu)\) is the colon ideal
\[
 \sigma^{c}(\mathfrak a,\mu):=\bigl(\pi^*\mathfrak a:\mathcal O_Y(-\mu D)\bigr).
\]
\end{definition}

\begin{lemma}[The controlled transform is the quotient by \(\mathcal I_D^{\mu}\)]
\label{bcalc:controlled-transform}
\uses{bcalc:controlled-transform-def,bcalc:cancellation,ana:invertible-colon}
\source{wlodarczyk05-hironaka,bierstone-milman08-functoriality}
\lean{Hironaka.BlowupControlledTransform.pullback_le_idealPow_exceptionalDivisor,Hironaka.BlowupControlledTransform.pullback_eq_mul_controlledTransform,Hironaka.BlowupControlledTransform.isCoherent_controlledTransform,Hironaka.BlowupControlledTransform.eq_controlledTransform_of_eq_mul,Hironaka.BlowupControlledTransform.controlledTransform_zero,Hironaka.BlowupControlledTransform.isInvertible_idealPow_exceptionalDivisor}
\leanok
Let \(\mathfrak a\) be coherent with \(\mathfrak a\subset\mathcal I_C^{\mu}\). Then \(\pi^*\mathfrak a\subset\mathcal O_Y(-\mu D)\),
\[
 \pi^*\mathfrak a=\mathcal O_Y(-\mu D)\cdot\sigma^{c}(\mathfrak a,\mu),
\]
\(\sigma^{c}(\mathfrak a,\mu)\) is coherent, and it is the only coherent ideal \(\mathfrak c\) with
\(\pi^*\mathfrak a=\mathcal O_Y(-\mu D)\,\mathfrak c\). For \(\mu=0\) (and \(\mathfrak a\) coherent) \(\sigma^{c}(\mathfrak a,0)=\pi^*\mathfrak a\).
\end{lemma}
\begin{proof}
\leanok
By Lemma~\ref{bcalc:exceptional-divisor}(3), \(\pi^*\mathfrak a\subset\pi^*(\mathcal I_C^{\mu})=\mathcal O_Y(-\mu D)\) (Wlodarczyk,
\S2.2, Lemma, part (2)). The ideal \(\pi^*\mathfrak a\) is coherent and \(\mathcal O_Y(-\mu D)\) is invertible, so
Lemma~\ref{ana:invertible-colon} gives the factorization with a coherent colon ideal, which is \(\sigma^{c}(\mathfrak a,\mu)\) by
definition. If also \(\pi^*\mathfrak a=\mathcal O_Y(-\mu D)\mathfrak c\) with \(\mathfrak c\) coherent, then
\(\sigma^{c}(\mathfrak a,\mu)=(\mathcal O_Y(-\mu D)\mathfrak c:\mathcal O_Y(-\mu D))=\mathfrak c\) by Lemma~\ref{bcalc:cancellation}. For
\(\mu=0\), \(\mathcal O_Y(-0\cdot D)=\mathcal O_Y\) and \((\pi^*\mathfrak a:\mathcal O_Y)=\pi^*\mathfrak a\), again by
Lemma~\ref{bcalc:cancellation}.
\end{proof}

\begin{lemma}[Restriction of a blowup to an open subset]
\label{bcalc:restrict-open}
\uses{ana:blowups,rem:open-subset-analytic,ana:ideal-sheaves}
\source{wlodarczyk05-hironaka,bierstone-milman08-functoriality}
\lean{Hironaka.BlowupRestrictOpen.isBlowupAlong_restrict}\upstreamlean{TranscendentalNakamaye.Pending.BlowupRestrictOpen.exists_isBlowupAlong_restrict}\lean{Hironaka.BlowupRestrictOpen.contMDiff_opens_mk,Hironaka.BlowupRestrictOpen.vanishingIdeal_opens_mem_iff,Hironaka.BlowupRestrictOpen.vanishingIdeal_opens_mem_comp_val,Hironaka.BlowupRestrictOpen.pullback_vanishingIdeal_opens,Hironaka.BlowupRestrictOpen.isClosedSubmanifold_preimage_val,Hironaka.BlowupRestrictOpen.interior_preimage_val_eq_empty}
\mixedleanok
Let \(\pi\colon Y\to X\) be the blowup along \(C\) and \(V\subset X\) open. Then the restriction
\(\pi^{-1}(V)\to V\) of \(\pi\) is the blowup of the complex manifold \(V\) along \(C\cap V\).
\end{lemma}
\begin{proof}
\mixedleanok
\(C\cap V\) is a closed submanifold of \(V\) (regularity is local, Lemma~\ref{rem:open-subset-analytic}) with empty interior (an open
subset of \(V\) inside \(C\) is open in \(X\)), and the restriction is holomorphic. A holomorphic map \(g\colon Z\to V\) into the
open set \(V\) is the same as a holomorphic map \(\iota\circ g\colon Z\to X\) with values in \(V\), and germs on \(V\) are germs on
\(X\) at points of \(V\) (extend by \(0\)); in particular \(g^*\mathcal I_{C\cap V}=(\iota\circ g)^*\mathcal I_C\). For \(g=\) the
restriction of \(\pi\) to \(\pi^{-1}(V)\) this is the restriction of the invertible ideal \(\pi^*\mathcal I_C\) to the open set
\(\pi^{-1}(V)\), hence invertible. If \(g^*\mathcal I_{C\cap V}\) is invertible, the universal property of \(\pi\) gives a unique
holomorphic \(h\colon Z\to Y\) with \(\pi\circ h=\iota\circ g\); it takes values in \(\pi^{-1}(V)\), so it is a holomorphic lift
of \(g\) through the restriction, and any such lift, composed with the inclusion of \(\pi^{-1}(V)\), is a lift of \(\iota\circ g\)
through \(\pi\), hence equal to \(h\).
\end{proof}

\paragraph{Towers.} A \emph{tower} over \(X\) is a finite sequence \(X=X_0\leftarrow X_1\leftarrow\dots\leftarrow X_r\) of blowups
\(\pi_i\colon X_i\to X_{i-1}\) along nowhere dense closed submanifolds \(C_{i-1}\subset X_{i-1}\), recorded with all its spaces,
centres and maps; its projection is \(\pi_1\circ\dots\circ\pi_r\colon X_r\to X\). The transform of a marked ideal along a tower,
multiple test blow-ups and resolutions are those of Bierstone--Milman, in the ambient form \((M,N,E,\mathcal I,d)\) used by the
chapter on the resolution of marked ideals.

\begin{lemma}[Restriction of a tower to an open subset]
\label{bcalc:tower-restrict}
\uses{bcalc:restrict-open,hir:blowup-tower}
\source{wlodarczyk05-hironaka,bierstone-milman08-functoriality}
\lean{Hironaka.BlowupTowerTransforms.openPreimage,Hironaka.BlowupTowerTransforms.restrict,Hironaka.BlowupTowerTransforms.proj_restrict,Hironaka.BlowupTowerTransforms.length_restrict}
\leanok
Let \(T\) be a tower over \(X\) with projection \(p\colon X_r\to X\), and \(U\subset X\) open. The open sets
\(U_i=(\pi_1\circ\dots\circ\pi_i)^{-1}(U)\subset X_i\), with the restrictions of the \(\pi_i\) and the centres \(C_{i-1}\cap U_{i-1}\),
form a tower over \(U\) of the same length, whose projection is the restriction of \(p\) to \(p^{-1}(U)\to U\).
\end{lemma}
\begin{proof}
\leanok
Induction on the length, using Lemma~\ref{bcalc:restrict-open} at each step:
\(\pi_i^{-1}(U_{i-1})=U_i\), and the restriction of \(\pi_i\) to \(U_i\to U_{i-1}\) is the blowup along \(C_{i-1}\cap U_{i-1}\).
\end{proof}

\begin{definition}[The exceptional divisor of a tower]
\label{bcalc:tower-exceptional}
\uses{bcalc:exceptional-divisor,hir:blowup-tower}
\source{wlodarczyk05-hironaka}
\lean{Hironaka.BlowupTowerTransforms.exceptionalDivisor,Hironaka.BlowupTowerTransforms.exceptionalDivisor_refl_ideal,Hironaka.BlowupTowerTransforms.exceptionalDivisor_step_ideal}
\leanok
The exceptional divisor \(E_T\) of a tower \(T\) is \(0\) for the empty tower, and \(\pi_r^*E_{T'}+D_r\) if \(T\) is \(T'\) followed by
the blowup \(\pi_r\) with exceptional divisor \(D_r\) (Lemma~\ref{bcalc:exceptional-divisor}); thus
\(E_T=\sum_i(\pi_{i+1}\circ\dots\circ\pi_r)^*D_i\).
\end{definition}

\begin{lemma}[Transforms with \(N=X\)]
\label{bcalc:transform-ambient}
\uses{bcalc:controlled-transform,bcalc:exceptional-divisor,hir:marked-ideal,hir:transform}
\source{bierstone-milman08-functoriality,wlodarczyk05-hironaka}
\upstreamlean{TranscendentalNakamaye.Pending.BlowupTowerTransforms.strictTransform_univ}\lean{Hironaka.BlowupTowerTransforms.restrictOn_univ}\upstreamlean{TranscendentalNakamaye.Pending.BlowupTowerTransforms.transform_ideal_eq_controlledTransform,TranscendentalNakamaye.Pending.BlowupTowerTransforms.transform_N_eq_univ}
\mixedleanok
Let \(\pi\colon Y\to X\) be the blowup along \(C\). The strict transform of \(X\) is \(Y\). For a marked ideal \((X,X,E,\mathfrak a,d)\)
with \(\mathfrak a\) coherent, its transform by \(\pi\) is \((Y,Y,E',\sigma^{c}(\mathfrak a,d),d)\): the ideal of the transform,
\(\mathcal I_D^{-d}\pi^*\mathfrak a\) read on the strict transform \(Y\), is the controlled transform of
Definition~\ref{bcalc:controlled-transform-def}.
\end{lemma}
\begin{proof}
\mixedleanok
The strict transform of \(X\) is the closure of \(Y\setminus\operatorname{Exc}\pi\), which is \(Y\) since \(\operatorname{Exc}\pi\)
is nowhere dense (Lemma~\ref{bcalc:exceptional-divisor}(2)). Reading an ideal on the whole space changes nothing for a coherent
ideal (a germ agreeing everywhere near the point with a combination of germs of the ideal is that combination), so the colon
along \(Y\) of \(\pi^*\mathfrak a\) by \(\mathcal O_Y(-dD)=(\pi^*\mathcal I_C)^{d}\) is the colon \(\sigma^{c}(\mathfrak a,d)\).
\end{proof}

\begin{proposition}[Resolution gives principalization]
\label{bcalc:principalization}
\uses{bcalc:transform-ambient,bcalc:controlled-transform,bcalc:tower-exceptional,bcalc:cancellation,hir:test-tower,hir:cosupport}
\source{wlodarczyk05-hironaka}
\upstreamlean{TranscendentalNakamaye.Pending.BlowupTowerTransforms.pullback_eq_mul_transformAlong,TranscendentalNakamaye.Pending.BlowupTowerTransforms.pullback_eq_exceptionalDivisor_of_isResolution,TranscendentalNakamaye.Pending.BlowupTowerTransforms.isInvertible_pullback_of_isResolution}\lean{Hironaka.BlowupTowerTransforms.le_vanishingIdeal_of_subset_zeroLocus,Hironaka.BlowupTowerTransforms.eq_unitIdeal_of_zeroLocus_eq_empty}
\mixedleanok
Let \(\mathfrak a\) be a coherent ideal on \(X\) and \(T\) a multiple test blow-up of the marked ideal \((X,X,E,\mathfrak a,1)\), with
projection \(p\) and transform \((X_r,X_r,E_r,\mathfrak a_r,1)\). Then \(\mathfrak a_r\) is coherent and
\(p^*\mathfrak a=\mathcal O(-E_T)\,\mathfrak a_r\). If \(T\) is a resolution (empty final cosupport), then \(\mathfrak a_r=\mathcal O\)
and \(p^*\mathfrak a=\mathcal O(-E_T)\) is invertible.
\end{proposition}
\begin{proof}
\mixedleanok
Induction on the length. With the mark \(1\) the cosupport of a marked ideal \((X_i,X_i,E_i,\mathfrak a_i,1)\) is the zero locus of
\(\mathfrak a_i\). An admissible centre \(C_i\) lies in it, so every germ of the coherent ideal \(\mathfrak a_i\) at a point \(x\)
vanishes on \(C_i\) near \(x\) (the germ belongs to \(\mathfrak a_i\) at every point near \(x\), by coherence), i.e.\
\(\mathfrak a_i\subset\mathcal I_{C_i}\). By Lemmas~\ref{bcalc:transform-ambient} and~\ref{bcalc:controlled-transform},
\(\mathfrak a_{i+1}=\sigma^{c}(\mathfrak a_i,1)\) is coherent and \(\pi_{i+1}^*\mathfrak a_i=\mathcal O(-D_{i+1})\,\mathfrak a_{i+1}\).
If \((\pi_1\circ\dots\circ\pi_i)^*\mathfrak a=\mathcal O(-E_i)\,\mathfrak a_i\), then pulling back by \(\pi_{i+1}\), which commutes with
products of coherent ideals, gives \((\pi_1\circ\dots\circ\pi_{i+1})^*\mathfrak a=\mathcal O(-\pi_{i+1}^*E_i-D_{i+1})\,\mathfrak a_{i+1}\).
If the final cosupport is empty, the coherent ideal \(\mathfrak a_r\) has no zero: near each point it contains a germ \(b\) with
\(b(x)\ne0\), and then every holomorphic germ \(g=(g/b)\,b\) belongs to it, so \(\mathfrak a_r=\mathcal O\) and
\(p^*\mathfrak a=\mathcal O(-E_T)\) (Lemma~\ref{bcalc:cancellation}).
\end{proof}

\paragraph{Submanifolds through the centre.} In the remaining statements \(\pi\colon Y\to X\) is the blowup along \(C\) and
\(N\supset C\) is a closed subset of \(X\), regular of dimension \(k_N\) at each of its points. Coordinates are those of
Definition~\ref{ana:blowups}: a set is \emph{flattened} by a local biholomorphism \(\Phi\) (after the preferred chart) onto
\(\mathbf C^k\times0\), the \emph{last} \(n-k\) coordinates vanishing.

\begin{lemma}[Nested flattening coordinates]
\label{bcalc:nested-flattening}
\uses{ana:blowups,ana:local-holomorphic-inverse}
\lean{Hironaka.NestedFlattening.exists_nested_flattening,Hironaka.NestedFlattening.exists_nested_flattening_model,Hironaka.NestedFlattening.nestedMap,Hironaka.NestedFlattening.nestedMap_mem_plane_iff,Hironaka.NestedFlattening.nestedMap_coordinateInclusion,Hironaka.NestedFlattening.midCoords,Hironaka.NestedFlattening.midCoords_apply,Hironaka.NestedFlattening.midCoords_coordinateInclusion,Hironaka.NestedFlattening.eq_zero_of_projection_midCoords,Hironaka.NestedFlattening.coordinateInclusion_comp,Hironaka.NestedFlattening.coordinateProjection_comp,Hironaka.NestedFlattening.sub_inclusion_projection_apply,Hironaka.BlowupStrictSubmanifold.le_of_flattening_model,Hironaka.BlowupStrictSubmanifold.le_of_regularAtDim_subset}
\leanok
Let \(C\subset N\) be regular at \(x\) of dimensions \(k_C\) and \(k_N\). Then \(k_C\le k_N\), and one local biholomorphism
\(\Phi\) at \(x\) (with \(\Phi(x)=0\)) flattens \(C\) onto \(\mathbf C^{k_C}\times0\) and \(N\) onto \(\mathbf C^{k_N}\times0\) simultaneously.
\end{lemma}
\begin{proof}
\leanok
Work in coordinates \(u\) flattening \(N\), so \(N=\mathbf C^{k_N}\times0\) near \(0\), and let \(\Theta\) flatten \(C\) there. The map
\(\varphi(v)=\Theta^{-1}(v,0)\), \(v\in\mathbf C^{k_C}\) near \(0\), parametrizes \(C\) and takes values in \(\mathbf C^{k_N}\times0\); write
\(\varphi=\iota\circ\varphi'\) with \(\varphi'\colon\mathbf C^{k_C}\to\mathbf C^{k_N}\). Since \(\Theta\circ\iota\circ\varphi'=(\cdot,0)\) near
\(0\), the differential \(D=d\varphi'(0)\) is injective; in particular \(k_C\le k_N\). Let \(J\colon\mathbf C^{k_N-k_C}\to\mathbf C^{k_N}\) be
linear and injective onto a complement of \(D(\mathbf C^{k_C})\), and write \(w=(v,s,t)\in\mathbf C^{k_C}\times\mathbf C^{k_N-k_C}\times
\mathbf C^{n-k_N}\). The map \(F(v,s,t)=(\varphi'(v)+Js,\,t)\) has the invertible differential \((v,s,t)\mapsto(Dv+Js,t)\) at \(0\), so it
has a local holomorphic inverse \(\Phi\) (Lemma~\ref{ana:local-holomorphic-inverse}). Now \(F(w)\in\mathbf C^{k_N}\times0\) iff \(t=0\),
and \(F(v,0,0)=\varphi(v)\in C\); conversely a point of \(C\) near \(0\) is \(\varphi(v)\) with \(v\) small (it is
\(\Theta^{-1}(\Theta(\cdot))\) and \(\Theta(\cdot)\in\mathbf C^{k_C}\times0\)), so \(\Phi\) maps it to \((v,0,0)\). Hence near \(0\),
\(\Phi(C)=\mathbf C^{k_C}\times0\) and \(\Phi(N)=\mathbf C^{k_N}\times0\); compose with the coordinates \(u\).
\end{proof}

\begin{lemma}[Monomial charts of an abstract blowup]
\label{bcalc:monomial-lift}
\uses{ana:blowups,bex:monomial-chart,bex:hadamard}
\source{demailly-cadg:VII.12.1}
\lean{Hironaka.BlowupMonomialCharts.monomial,Hironaka.BlowupMonomialCharts.monomialDomain,Hironaka.BlowupMonomialCharts.monomialChart,Hironaka.BlowupMonomialCharts.contMDiff_monomialChart,Hironaka.BlowupMonomialCharts.isInvertible_pullback_monomialChart,Hironaka.BlowupMonomialCharts.monomialChart_mem_flatDomain,Hironaka.BlowupLocalSurjectivity.existsUnique_monomialLift,Hironaka.BlowupLocalSurjectivity.exists_monomialLocalInverse,Hironaka.BlowupLocalSurjectivity.exists_monomialLift_eq,Hironaka.BlowupFibreDirection.fderiv_monomialInverse_ne_zero,Hironaka.BlowupStrictSubmanifold.monomial_mem_plane_iff,Hironaka.BlowupStrictSubmanifold.monomial_mem_plane_iff_of_lt,Hironaka.BlowupStrictSubmanifold.eq_zero_of_monomial_mem_plane,Hironaka.BlowupStrictSubmanifold.sub_mem_plane_iff}
\leanok
Let \(C\) be flattened near \(y_0\) by \(\Phi\) onto \(\mathbf C^{k}\times0\), and for \(j\ge k\) let \(\sigma_j(w)_i=w_iw_j\) (\(k\le i\ne j\)),
\(\sigma_j(w)_i=w_i\) otherwise, and \(\gamma_j=e^{-1}\circ\Phi^{-1}\circ\sigma_j\) on \(\sigma_j^{-1}(\Phi(\text{target}))\). Then
\(\gamma_j^*\mathcal I_C\) is generated by \(w_j\), so \(\gamma_j\) lifts uniquely to a holomorphic \(\ell_j\) with \(\pi\circ\ell_j=\gamma_j\).
Every point \(z\) of \(Y\) over the flattening domain has a neighbourhood \(O\), an index \(j\ge k\) and a holomorphic
\(W\colon O\to\mathbf C^n\) with \(\sigma_j\circ W=\Phi\circ e\circ\pi\) and \(\ell_j\circ W=\mathrm{id}\) on \(O\); the differential of \(W\)
(in a chart) is injective. Moreover \(\sigma_j(w)\in\mathbf C^{k}\times0\iff w_j=0\); for \(k\le j<m\),
\(\sigma_j(w)\in\mathbf C^{m}\times0\iff w_j=0\) or \(w\in\mathbf C^{m}\times0\); for \(j\ge m\), \(\sigma_j(w)\in\mathbf C^m\times0\) forces \(w_j=0\).
\end{lemma}
\begin{proof}
\leanok
The pullback of \(\mathcal I_C\), generated by the flat coordinates \(\tau_i\) (\(i\ge k\), Lemma~\ref{bex:hadamard}), is generated by the
\(w_iw_j\) and \(w_j\), i.e.\ by \(w_j\), which vanishes on no open set; the universal property gives \(\ell_j\). Near \(z\) the pulled back
coordinates are \(\tau_i\circ\pi=\alpha_ih\) with \(h\) a local generator of \(\pi^*\mathcal I_C\); writing \(h\) back in terms of the
\(\tau_i\) shows that some \(\alpha_i(z)\ne0\); for \(j\) with \(|\alpha_j(z)|\) maximal, \(W=(\alpha_i/\alpha_j\ (i\ge k,\ i\ne j);\ \tau_i\circ\pi
\ \text{otherwise})\) satisfies \(\sigma_j\circ W=\Phi\circ e\circ\pi\), and \(\ell_j\circ W\) and the inclusion are two lifts of \(\pi\) near
\(z\), hence equal; then \(d\ell_j\circ dW=\mathrm{id}\) shows \(dW\) injective. The statements about \(\sigma_j\) are read off its formula.
\end{proof}

\begin{theorem}[The strict transform of a submanifold through the centre]
\label{bcalc:strict-submanifold}
\uses{bcalc:nested-flattening,bcalc:monomial-lift,bcalc:exceptional-divisor,ana:regular-transport,ana:local-holomorphic-inverse}
\source{bierstone-milman08-functoriality,wlodarczyk05-hironaka}
\lean{Hironaka.BlowupStrictSubmanifold.geoStrictTransform,Hironaka.BlowupStrictSubmanifold.geoStrictTransform_subset_preimage,Hironaka.BlowupStrictSubmanifold.mem_geoStrictTransform_iff_of_notMem,Hironaka.BlowupStrictSubmanifold.regularAtDim_geoStrictTransform_of_notMem,Hironaka.BlowupStrictSubmanifold.exists_local_geoStrictTransform,Hironaka.BlowupStrictSubmanifold.regularAtDim_geoStrictTransform_of_mem,Hironaka.BlowupStrictSubmanifold.regularAtDim_geoStrictTransform}\upstreamlean{TranscendentalNakamaye.Pending.BlowupStrictSubmanifold.isClosedSubmanifold_geoStrictTransform}
\mixedleanok
The strict transform \(N'=\overline{\pi^{-1}(N\setminus C)}\) is a closed submanifold of \(Y\), regular of dimension \(k_N\) at each of
its points. In the monomial chart of Lemma~\ref{bcalc:monomial-lift} at a point \(z\) over \(C\), built on nested coordinates
(Lemma~\ref{bcalc:nested-flattening}), \(N'\) near \(z\) is \(W^{-1}(\mathbf C^{k_N}\times0)\) if \(j<k_N\), and \(N'\) does not meet the chart
if \(j\ge k_N\).
\end{theorem}
\begin{proof}
\mixedleanok
Off \(\pi^{-1}(C)\) the map \(\pi\) is a local biholomorphism (Lemma~\ref{bcalc:exceptional-divisor}(2)) and \(N'=\pi^{-1}(N)\) there
(\(\pi^{-1}(N)\) is closed, and agrees with \(\pi^{-1}(N\setminus C)\) off \(\pi^{-1}(C)\)), so \(N'\) is regular of dimension \(k_N\) there
(Lemma~\ref{ana:regular-transport}). Let \(\pi(z)\in C\); take nested coordinates for \(C\subset N\) at \(\pi(z)\) and the chart
\(j\), \(O\), \(W\) at \(z\). For \(u\in O\): \(\pi(u)\in C\iff W(u)_j=0\), and \(\pi(u)\in N\iff\sigma_j(W(u))\in\mathbf C^{k_N}\times0\). If
\(j\ge k_N\) this forces \(W(u)_j=0\), so \(\pi^{-1}(N\setminus C)\cap O=\emptyset\) and \(N'\cap O=\emptyset\). If \(j<k_N\),
\(\pi^{-1}(N\setminus C)\cap O=\{u\in O:W(u)\in\mathbf C^{k_N}\times0,\ W(u)_j\ne0\}\). Its closure in \(O\) lies in the closed set
\(\{W\in\mathbf C^{k_N}\times0\}\); conversely if \(W(u)\in\mathbf C^{k_N}\times0\) and \(W(u)_j=0\), then for \(t\ne0\) small
\(w_t=W(u)+te_j\in\mathbf C^{k_N}\times0\) (as \(j<k_N\)) with \((w_t)_j=t\ne0\), so \(\ell_j(w_t)\in\pi^{-1}(N\setminus C)\), and
\(\ell_j(w_t)\to\ell_j(W(u))=u\). So \(N'\cap O=W^{-1}(\mathbf C^{k_N}\times0)\cap O\). Since \(W\circ e_z^{-1}\) has injective differential at
\(e_z(z)\), it is a local biholomorphism there (Lemma~\ref{ana:local-holomorphic-inverse}); recentred by \(W(z)\in\mathbf C^{k_N}\times0\),
it flattens \(N'\) onto \(\mathbf C^{k_N}\times0\) at \(z\).
\end{proof}

\begin{theorem}[The ideal of the strict transform]
\label{bcalc:strict-ideal}
\uses{bcalc:strict-submanifold,bcalc:nested-flattening,bcalc:monomial-lift,bcalc:exceptional-divisor,bcalc:controlled-transform,bex:hadamard,ana:local-holomorphic-inverse}
\source{bierstone-milman08-functoriality,wlodarczyk05-hironaka}
\lean{Hironaka.BlowupStrictIdeal.pullback_vanishingIdeal_eq_mul,Hironaka.BlowupStrictIdeal.controlledTransform_vanishingIdeal,Hironaka.BlowupStrictIdeal.isClosedSubmanifold_geoStrictTransform_of_isClosedSubmanifold,Hironaka.BlowupStrictIdeal.idealPow_one_exceptionalDivisor,Hironaka.BlowupStrictIdeal.mem_pullback_iff_mem_mul_of_notMem,Hironaka.BlowupStrictIdeal.mem_pullback_iff_mem_mul_of_mem,Hironaka.BlowupStrictIdeal.pullback_vanishingIdeal_le,Hironaka.BlowupStrictIdeal.mem_pullback_vanishingIdeal_of_notMem,Hironaka.BlowupStrictIdeal.mem_pullback_vanishingIdeal_iff_of_notMem,Hironaka.BlowupStrictIdeal.mem_mul_vanishingIdeal_iff_of_notMem,Hironaka.BlowupStrictIdeal.pulledCoords,Hironaka.BlowupStrictIdeal.holomorphicAt_pulledCoords,Hironaka.BlowupStrictIdeal.mem_pullback_vanishingIdeal_iff_pulledCoords,Hironaka.BlowupStrictIdeal.isInSpanAt_pulledCoords_centre_iff,Hironaka.BlowupStrictIdeal.exceptionalDivisor_mem_iff_of_chart,Hironaka.BlowupStrictIdeal.planeCoords,Hironaka.BlowupStrictIdeal.mdifferentiableOn_planeCoords,Hironaka.BlowupStrictIdeal.exists_vanishingIdeal_mem_iff_planeCoords,Hironaka.BlowupStrictIdeal.isInSpanAt_zero,Hironaka.BlowupStrictIdeal.mem_vanishingIdeal_of_isInSpanAt',Hironaka.BlowupStrictIdeal.vanishingIdeal_mem_congr,Hironaka.BlowupStrictIdeal.vanishingIdeal_mem_iff_of_disjoint,Hironaka.BlowupStrictIdeal.holomorphicAt_coord,Hironaka.BlowupStrictIdeal.mdifferentiableOn_coord}
\leanok
Let \(N\supset C\) be a closed submanifold of \(X\) (of any, possibly varying, dimension) and \(N'=\overline{\pi^{-1}(N\setminus C)}\).
Then \(N'\) is a closed submanifold of \(Y\),
\[
 \pi^*\mathcal I_N=\mathcal O_Y(-D)\cdot\mathcal I_{N'},
 \qquad\text{equivalently}\qquad
 \sigma^{c}(\mathcal I_N,1)=\mathcal I_{N'}.
\]
\end{theorem}
\begin{proof}
\leanok
\(N'\) is a closed submanifold by Theorem~\ref{bcalc:strict-submanifold}, applied at each point with the dimension of \(N\) there.
Both sides of the identity are computed stalk by stalk at \(z\in Y\). In general \(f^*\mathcal I_A\subset\mathcal I_{f^{-1}(A)}\) for a
continuous \(f\): the generators \(a\circ f\) vanish on \(f^{-1}(A)\).

If \(\pi(z)\notin C\), then \(\pi\) has a holomorphic section \(s\) over \(X\setminus C\) with \(s(\pi(z'))=z'\) for \(\pi(z')\notin C\)
(the universal property applied to the inclusion of \(X\setminus C\), where \(\mathcal I_C\) is the unit ideal). A germ \(g\) vanishing on
\(\pi^{-1}(N)\) near \(z\) is \(a\circ\pi\) with \(a=g\circ s\), and \(a\) vanishes on \(N\) near \(\pi(z)\); so
\(\pi^*\mathcal I_N=\mathcal I_{\pi^{-1}(N)}\) at \(z\). Near \(z\), \(N'=\pi^{-1}(N)\), and \(\mathcal O_Y(-D)\) contains \(1=1\circ\pi\), so
\(\mathcal O_Y(-D)\cdot\mathcal I_{N'}=\mathcal I_{N'}\) at \(z\).

Let \(\pi(z)\in C\). Take nested coordinates \(\Phi\) at \(\pi(z)\) with \(C=\mathbf C^{k_C}\times0\) and \(N=\mathbf C^{k_N}\times0\)
(Lemma~\ref{bcalc:nested-flattening}) and the monomial chart \(j\ge k_C\), \(O\), \(W\) at \(z\) (Lemma~\ref{bcalc:monomial-lift}), so
\(\sigma_j\circ W=\Phi\circ e\circ\pi\) on \(O\). For \(A=C\) or \(A=N\), flattened onto \(\mathbf C^{m}\times0\), every germ of
\(\mathcal I_A\) at \(\pi(z')\) is \(\sum_{i\ge m}b_i\Phi_i\) (Lemma~\ref{bex:hadamard}); hence on \(O\) the ideal \(\pi^*\mathcal I_A\) is
generated by the functions \(\sigma_j(W)_i\), \(i\ge m\), each of which lies in it. For \(A=C\) these are \(W_iW_j\) (\(i\ne j\)) and \(W_j\),
so \(\mathcal O_Y(-D)=(W_j)\) on \(O\). For \(A=N\):
\begin{itemize}
\item if \(j<k_N\), all \(i\ge k_N\) differ from \(j\), so \(\pi^*\mathcal I_N=(W_jW_i:i\ge k_N)\). By Theorem~\ref{bcalc:strict-submanifold},
\(N'\cap O=W^{-1}(\mathbf C^{k_N}\times0)\cap O\). If \(z\in N'\), \(W\) has injective differential at \(z\), so it is a local
biholomorphism there (Lemma~\ref{ana:local-holomorphic-inverse}); recentred by \(W(z)\in\mathbf C^{k_N}\times0\), it flattens \(N'\), and
Hadamard's lemma in these coordinates gives \(\mathcal I_{N'}=(W_i:i\ge k_N)\) near \(z\); the product of the two presentations is
\((W_jW_i:i\ge k_N)=\pi^*\mathcal I_N\). If \(z\notin N'\), then \(\mathcal I_{N'}=\mathcal O_Y\) near \(z\), and some \(W_{i_0}\), \(i_0\ge k_N\),
does not vanish at \(z\), so \(W_j=W_{i_0}^{-1}\cdot W_jW_{i_0}\in\pi^*\mathcal I_N\) and \(\pi^*\mathcal I_N=(W_j)=\mathcal O_Y(-D)\cdot\mathcal O_Y\).
\item if \(j\ge k_N\), then \(N'\) misses \(O\), \(\mathcal I_{N'}=\mathcal O_Y\) on \(O\), and the generators of \(\pi^*\mathcal I_N\) are
\(W_j\) (\(i=j\)) and multiples \(W_iW_j\) of it: \(\pi^*\mathcal I_N=(W_j)=\mathcal O_Y(-D)\cdot\mathcal O_Y\).
\end{itemize}
Finally \(\mathcal I_N\subset\mathcal I_C\), \(\mathcal I_{N'}\) is coherent (it is the ideal of a closed submanifold) and
\(\mathcal O_Y(-D)^1=\mathcal O_Y(-D)\), so the uniqueness part of Lemma~\ref{bcalc:controlled-transform} turns the identity into
\(\sigma^{c}(\mathcal I_N,1)=\mathcal I_{N'}\).
\end{proof}

\begin{lemma}[Geometric and analytic strict transforms]
\label{bcalc:strict-geo}
\uses{bcalc:exceptional-divisor,ana:strict-transform-composition}
\lean{Hironaka.BlowupStrictIdeal.geoStrictTransform_eq_strictTransform,Hironaka.BlowupStrictIdeal.geoStrictTransform_subset_strictTransform,Hironaka.BlowupStrictIdeal.geoStrictTransform_self}
\leanok
For \(A\subset X\), \(\overline{\pi^{-1}(A\setminus C)}\subset\overline{\pi^{-1}(A)\setminus\operatorname{Exc}\pi}\), with equality when \(C\) is nowhere
dense in \(A\) (\(A\subset\overline{A\setminus C}\)). The first set is empty for \(A=C\); the second is \(C\) itself when \(C\) is a smooth
hypersurface (then \(\pi\) is a biholomorphism and \(\operatorname{Exc}\pi=\emptyset\)).
\end{lemma}
\begin{proof}
\leanok
\(\operatorname{Exc}\pi\subset\pi^{-1}(C)\) (Lemma~\ref{bcalc:exceptional-divisor}) gives the inclusion. Conversely let \(\pi(u)\in A\),
\(u\notin\operatorname{Exc}\pi\). Then \(\pi\) maps every neighbourhood \(V\) of \(u\) onto a neighbourhood of \(\pi(u)\)
(Lemma~\ref{ana:strict-transform-composition}), which meets \(A\setminus C\) because \(\pi(u)\in\overline{A\setminus C}\); so \(V\) meets
\(\pi^{-1}(A\setminus C)\). The closure of \(\pi^{-1}(A)\setminus\operatorname{Exc}\pi\) is therefore in \(\overline{\pi^{-1}(A\setminus C)}\).
\end{proof}
The second clause is why the transform of a marked ideal must use \(\overline{\pi^{-1}(A\setminus C)}\) for \(N\) and for the members of
\(E\): along a centre of codimension one that is a member \(H\) of \(E\), the other set would keep \(H\), and \(E'\) would list it twice.

\begin{lemma}[The controlled transform of a function]
\label{bcalc:function-transform}
\uses{bcalc:controlled-transform,bcalc:exceptional-divisor,bcalc:cancellation}
\source{wlodarczyk05-hironaka}
\lean{Hironaka.BlowupFunctionTransform.IsLocalEquation,Hironaka.BlowupFunctionTransform.exists_isLocalEquation,Hironaka.BlowupFunctionTransform.IsLocalEquation.holomorphicAt}\upstreamlean{TranscendentalNakamaye.Pending.BlowupFunctionTransform.IsLocalEquation.nonvanishing_pow}\lean{Hironaka.BlowupFunctionTransform.IsLocalEquation.idealPow_mem_iff,Hironaka.BlowupFunctionTransform.IsLocalEquation.pow_mem,Hironaka.BlowupFunctionTransform.exists_eq_pow_mul}\upstreamlean{TranscendentalNakamaye.Pending.BlowupFunctionTransform.eventuallyEq_of_pow_mul}\lean{Hironaka.BlowupFunctionTransform.mem_controlledTransform_of_eq_pow_mul}\upstreamlean{TranscendentalNakamaye.Pending.BlowupFunctionTransform.mem_controlledTransform_iff}\lean{Hironaka.BlowupFunctionTransform.holomorphicAt_pow,Hironaka.BlowupFunctionTransform.isInSpanAt_range_const_iff,Hironaka.BlowupFunctionTransform.holomorphicAt_of_mem_idealPow}
\mixedleanok
A \emph{local equation} of \(D\) on an open set \(U\subset Y\) is a holomorphic \(h\) on \(U\), vanishing identically on no nonempty open
subset, with \(\mathcal O_Y(-D)=(h)\) on \(U\); local equations exist around every point. On \(U\), \(\mathcal O_Y(-\mu D)=(h^{\mu})\), and for
\(z\in U\):
\begin{enumerate}
\item if \(u\in\mathcal I_C^{\mu}\) at \(\pi(z)\), then \(u\circ\pi=h^{\mu}u'\) near \(z\) for a germ \(u'\) holomorphic at \(z\), unique as a germ
(the \emph{controlled transform} \(u'=h^{-\mu}\pi^*u\) of \(u\));
\item if moreover \(u\in\mathfrak a\) at \(\pi(z)\), then \(u'\in\sigma^{c}(\mathfrak a,\mu)\) at \(z\);
\item if \(\mathfrak a\subset\mathcal I_C^{\mu}\) is coherent, \(\sigma^{c}(\mathfrak a,\mu)\) is generated at \(z\) by the controlled transforms of the
germs of \(\mathfrak a\) at \(\pi(z)\).
\end{enumerate}
\end{lemma}
\begin{proof}
\mixedleanok
Local equations are the local generators of the invertible ideal \(\mathcal O_Y(-D)\). By induction on \(\mu\), the product of the
presentations \((h^{\mu})\) and \((h)\) is \((h^{\mu+1})\). (1) \(u\circ\pi\in\pi^*(\mathcal I_C^{\mu})=\mathcal O_Y(-\mu D)=(h^{\mu})\)
(Lemma~\ref{bcalc:exceptional-divisor}); uniqueness: if \(h^{\mu}v=h^{\mu}w\) near \(z\), then \(v=w\) on the dense set \(\{h\ne0\}\), hence near
\(z\). (2) Every \(\ell\in\mathcal O_Y(-\mu D)\) at \(z\) is \(ch^{\mu}\), and \(\ell u'=c\,(u\circ\pi)\in\pi^*\mathfrak a\); so \(u'\) lies in the colon ideal
\(\sigma^{c}(\mathfrak a,\mu)\). (3) If \(g\in\sigma^{c}(\mathfrak a,\mu)\), then \(h^{\mu}g\in\pi^*\mathfrak a\), i.e.\ \(h^{\mu}g=\sum_ia_i\,(u_i\circ\pi)\) with
\(u_i\in\mathfrak a\); writing \(u_i\circ\pi=h^{\mu}u_i'\) by (1) and cancelling \(h^{\mu}\) gives \(g=\sum_ia_iu_i'\). Conversely the \(u_i'\) lie in the
coherent ideal \(\sigma^{c}(\mathfrak a,\mu)\) by (2) (Lemma~\ref{bcalc:controlled-transform}), which absorbs their combinations.
\end{proof}

\begin{corollary}[Strict transforms of hypersurfaces through the centre]
\label{bcalc:giraud-strict}
\uses{bcalc:function-transform,bcalc:strict-ideal}
\source{wlodarczyk05-hironaka}
\upstreamlean{TranscendentalNakamaye.Pending.BlowupFunctionTransform.vanishingIdeal_geoStrictTransform_mem_iff,TranscendentalNakamaye.Pending.BlowupFunctionTransform.mem_geoStrictTransform_iff_eq_zero,TranscendentalNakamaye.Pending.BlowupFunctionTransform.mem_iff_eq_zero_of_span,TranscendentalNakamaye.Pending.BlowupFunctionTransform.vanishingIdeal_le_of_subset}
\upstreamleanok
Let \(N\supset C\) be a closed submanifold, \(z\in U\), and suppose \(\mathcal I_N=(u)\) at \(\pi(z)\) with \(u\) holomorphic, and
\(u\circ\pi=h\,u'\) near \(z\). Then \(\mathcal I_{N'}=(u')\) at \(z\), where \(N'=\overline{\pi^{-1}(N\setminus C)}\); in particular \(z\in N'\) if and only
if \(u'(z)=0\). This is the statement ``\(V(u')\) is the strict transform of \(V(u)\)'' of Włodarczyk's Lemma on tangent directions.
\end{corollary}
\begin{proof}
\upstreamleanok
By Theorem~\ref{bcalc:strict-ideal}, \(\mathcal I_{N'}=\sigma^{c}(\mathcal I_N,1)\), which by Lemma~\ref{bcalc:function-transform}(3) is generated at \(z\) by the
\(v\) with \(w\circ\pi=hv\), \(w\in\mathcal I_N\) at \(\pi(z)\) (note \(\mathcal I_N\subset\mathcal I_C\) as \(C\subset N\)). Such a \(w\) is \(cu\), so
\(hv=(c\circ\pi)hu'\) and \(v=(c\circ\pi)u'\); and \(u'\) itself is one of them. For the pointwise statement: if \(z\in N'\), the generator \(u'\)
vanishes at \(z\); if \(z\notin N'\), the constant \(1\) lies in \(\mathcal I_{N'}=(u')\) at \(z\), so \(u'(z)\ne0\).
\end{proof}

\begin{lemma}[Zero sets of functions of multiplicity one]
\label{bcalc:hypersurface}
\uses{bcalc:strict-ideal,bcalc:snc-chart,hir:derivative-ideal,ana:local-holomorphic-inverse}
\source{wlodarczyk05-hironaka}
\lean{Hironaka.LocalGiraudTransform.exists_hypersurface_coords,Hironaka.LocalGiraudTransform.mem_plane_pred_iff,Hironaka.LocalGiraudTransform.vanishingIdeal_zeroSet_mem_iff}\upstreamlean{TranscendentalNakamaye.Pending.LocalGiraudTransform.regularAtDim_zeroSet,TranscendentalNakamaye.Pending.LocalGiraudTransform.isClosed_preimage_val_zeroSet}
\mixedleanok
Let \(u\) be holomorphic on an open set \(U\subset X\), of multiplicity one at its zeros in \(U\) (at each such zero some partial derivative of
\(u\) in some chart does not vanish). Then \(V(u)\cap U\) is closed in \(U\) and regular of dimension \(n-1\) at each of its points, and at every
point of \(U\) the ideal of \(V(u)\cap U\) is generated by \(u\).
\end{lemma}
\begin{proof}
\mixedleanok
Closedness in \(U\) is the continuity of \(u\). Let \(x\in U\) with \(u(x)=0\), and let \(e\) be the preferred chart at \(x\). The partial derivatives of
\(u\) in any chart are holomorphic combinations of those in \(e\) (chain rule for the transition), so some \(\partial_j(u\circ e^{-1})(e(x))\neq0\);
in particular \(n>0\). Let \(\tau\) be the transposition of \(j\) and \(n-1\) and \(W=(e_{\tau(0)},\dots,e_{\tau(n-2)},u)\). The differential of \(W\) at
\(x\) is injective: a vector killed by it has all coordinates but the \(j\)-th equal to \(0\), and then \(du(x)\) kills it only if that coordinate is
\(0\). Hence \(W-W(x)\) is a coordinate system at \(x\) (Lemma~\ref{bcalc:snc-chart}; inverse function theorem,
Lemma~\ref{ana:local-holomorphic-inverse}) in which \(V(u)\cap U=\mathbf C^{n-1}\times0\): this is the regularity. By Hadamard's lemma in the
coordinates \(W\) (as in the proof of Theorem~\ref{bcalc:strict-ideal}) the ideal of \(V(u)\cap U\) is generated near \(x\) by the coordinate
\(W_{n-1}=u\). At a point \(x\in U\) with \(u(x)\neq0\), \(u\) is a unit near \(x\) and \(V(u)\) misses a neighbourhood of \(x\): both the ideal of
\(V(u)\cap U\) and \((u)\) consist of all holomorphic germs.
\end{proof}

\begin{theorem}[Włodarczyk's lemma on tangent directions on an open set]
\label{bcalc:giraud-local}
\uses{bcalc:hypersurface,bcalc:restrict-open,bcalc:giraud-strict,bcalc:multiplicity-one,bcalc:strict-submanifold,bcalc:function-transform,mc:tangent-direction,mc:tangent-direction-support}
\source{wlodarczyk05-hironaka}
\upstreamlean{TranscendentalNakamaye.Pending.LocalGiraudTransform.vanishingIdeal_geoStrictTransform_mem_iff_of_open,TranscendentalNakamaye.Pending.LocalGiraudTransform.mem_geoStrictTransform_iff_eq_zero_of_open,TranscendentalNakamaye.Pending.LocalGiraudTransform.eventually_mem_geoStrictTransform_iff_of_open,TranscendentalNakamaye.Pending.LocalGiraudTransform.exists_chartDeriv_ne_zero_controlledTransform_of_open}\lean{Hironaka.LocalGiraudTransform.isLocalEquation_restrict,Hironaka.LocalGiraudTransform.geoStrictTransform_preimage_val,Hironaka.LocalGiraudTransform.vanishingIdeal_comp_val_iff}\upstreamlean{TranscendentalNakamaye.Pending.LocalGiraudTransform.isInSpanAt_eq_comp_val_iff,TranscendentalNakamaye.Pending.LocalGiraudTransform.IsTangentDirection.giraud_hypotheses}
\mixedleanok
Let \(u\) be holomorphic on an open set \(U\subset X\), of multiplicity one at its zeros in \(U\), with \(C\cap U\subset V(u)\), and put
\(N=V(u)\cap U\), \(N'=\overline{\pi^{-1}(N\setminus C)}\). Let \(h\) be a local equation of the exceptional divisor on \(U'\subset Y\),
\(z\in U'\cap\pi^{-1}(U)\), and \(u\circ\pi=h\,u'\) near \(z\) with \(u'\) holomorphic at \(z\). Then:
\begin{enumerate}
\item the ideal of \(N'\) at \(z\) is \((u')\); hence \(N'\) and \(V(u')\) agree near \(z\) (Włodarczyk's Lemma on tangent directions, part (3): \(V(u')\)
is the strict transform of \(V(u)\) over \(U\));
\item if \(z\in N'\), some partial derivative of \(u'\) at \(z\) in some chart does not vanish (part (2): \(u'\) has multiplicity one).
\end{enumerate}
The hypotheses on \(u\) hold for a tangent direction \(u\) of a marked ideal \((\mathfrak a,\mu)\), \(\mu\ge1\), on \(U\)
(Definition~\ref{mc:tangent-direction}) and a centre \(C\) in the cosupport of \((\mathfrak a,\mu)\) (Lemma~\ref{mc:tangent-direction-support}).
\end{theorem}
\begin{proof}
\mixedleanok
By Lemma~\ref{bcalc:restrict-open} the restriction \(\pi^{-1}(U)\to U\) of \(\pi\) is the blowup of the manifold \(U\) along \(C\cap U\). In \(U\),
\(N\) is a closed submanifold containing \(C\cap U\), of dimension \(n-1\) at each of its points, with ideal \((u)\) at every point
(Lemma~\ref{bcalc:hypersurface}). The restriction of \(h\) to \(U'\cap\pi^{-1}(U)\) is a local equation of the exceptional divisor of the restricted
blowup: germs on an open subset are germs on the ambient manifold at its points (extension by zero), and the exceptional ideal of the
restriction is the restriction of \(\pi^*\mathcal I_C\). Strict transforms commute with the restriction: the closure in the open set
\(\pi^{-1}(U)\) of a subset of it is the trace of its closure in \(Y\). Corollary~\ref{bcalc:giraud-strict} for the restricted blowup gives
\(\mathcal I_{N'}=(u')\) at \(z\), read back on \(Y\) through the same identifications; as \(N'\) is closed, \(z\in N'\) if and only if \(u'(z)=0\), and
the hypotheses persist at the points near \(z\), which gives the agreement near \(z\). For (2): if \(z\in N'\), then \(u(\pi(z))=0\), so \(n>0\) by
Lemma~\ref{bcalc:hypersurface}; \(N'\) is regular of dimension \(n-1\) at \(z\) (Theorem~\ref{bcalc:strict-submanifold} in \(\pi^{-1}(U)\), regularity
being local), and a generator of the ideal of a submanifold of dimension \(n-1<n\) has a nonvanishing partial derivative
(Lemma~\ref{bcalc:multiplicity-one}). Finally, a tangent direction \(u\) is a germ of \(\mathcal D^{\mu-1}(\mathfrak a)\) at every point of \(U\), hence
holomorphic there, of multiplicity one by definition, and \(C\cap U\subset\operatorname{cosupp}(\mathfrak a,\mu)\cap U\subset V(u)\)
(Lemma~\ref{mc:tangent-direction-support}).
\end{proof}

\begin{lemma}[A generator modulo a submanifold has a nonvanishing tangential derivative]
\label{bcalc:tangent-generator-along}
\uses{bcalc:nested-flattening,bcalc:multiplicity-one,mc:derivative-along,hir:order-along-flat}
\source{wlodarczyk05-hironaka}
\lean{Hironaka.LocalGiraudTransformAlong.exists_fieldDeriv_ne_zero_of_vanishingIdeal_le,Hironaka.LocalGiraudTransformAlong.apply_eq_zero_of_isInSpanAt,Hironaka.LocalGiraudTransformAlong.isInSpanAt_of_mul_isInSpanAt}
\leanok
Let \(A\subset B\) be regular at \(z\) of dimensions \(a<b\), and let \(f\) be holomorphic at \(z\) with \(f(z)=0\), such that every germ of
\(\mathcal I_A\) at \(z\) is a combination of germs of \(\mathcal I_B\) and of \(f\). Then some holomorphic vector field \(\xi\) tangent to \(B\) at \(z\)
(Definition~\ref{mc:derivative-along}) has \(\xi(f)(z)\ne0\). Moreover: a combination of germs vanishing at \(z\) vanishes at \(z\); and if \(h\) is
holomorphic at \(z\) and vanishes identically near no point, \(g\) is holomorphic at \(z\), and \(hg\) is a combination of products \(ht\) with \(t\) in
a family of holomorphic germs, then \(g\) is a combination of the \(t\).
\end{lemma}
\begin{proof}
\leanok
Take coordinates \(\psi\) centred at \(z\) in which \(A=\mathbf C^a\times0\) and \(B=\mathbf C^b\times0\) (Lemma~\ref{bcalc:nested-flattening}). The field
\(\xi=\partial/\partial\psi_a\) is tangent to \(B\), as \(a<b\) (Lemma~\ref{hir:order-along-flat}(2)), and \(\psi_a\in\mathcal I_A\). Write
\(\psi_a=\sum_kc_ks_k\) near \(z\) with each \(s_k\) a germ of \(\mathcal I_B\) or \(f\). By the product rule
\(1=\xi(\psi_a)(z)=\sum_k\xi(c_k)(z)s_k(z)+c_k(z)\xi(s_k)(z)\). Every \(s_k\) vanishes at \(z\in A\subset B\); if \(s_k\in\mathcal I_B\), then \(\xi(s_k)\in\mathcal I_B\)
(tangency) vanishes at \(z\). So if \(\xi(f)(z)=0\), the right side vanishes, a contradiction. (For \(B=X\) this is the argument of
Lemma~\ref{bcalc:multiplicity-one}.) The second claim is evaluation at \(z\). For the third, \(hg=\sum c_k\,ht_k\) near \(z\) gives
\(h\,(g-\sum c_kt_k)=0\) near \(z\), and a holomorphic germ killed by \(h\) vanishes near \(z\), since \(h\) vanishes identically near no point.
\end{proof}

\begin{theorem}[Włodarczyk's lemma on tangent directions along a submanifold]
\label{bcalc:giraud-local-along}
\uses{bcalc:giraud-local,bcalc:tangent-generator-along,bcalc:restrict-open,bcalc:strict-ideal,bcalc:strict-submanifold,bcalc:function-transform,bcalc:exceptional-divisor,mc:contact-hypersurface-along,mc:tangent-direction-along,mc:tangent-direction-along-basic,mc:derivative-along}
\source{wlodarczyk05-hironaka}
\lean{Hironaka.LocalGiraudTransformAlong.vanishingIdeal_geoStrictTransform_inter_zeroSet_mem_iff,Hironaka.LocalGiraudTransformAlong.eventually_mem_geoStrictTransform_inter_zeroSet_iff,Hironaka.LocalGiraudTransformAlong.exists_fieldDeriv_ne_zero_controlledTransform_along}\upstreamlean{TranscendentalNakamaye.Pending.LocalGiraudTransformAlong.giraud_hypotheses_of_isTangentDirectionAlong}\lean{Hironaka.LocalGiraudTransformAlong.isClosedSubmanifold_preimage_val_inter_zeroSet,Hironaka.LocalGiraudTransformAlong.isInSpanAt_pullback_of_mem_vanishingIdeal_geoStrictTransform,Hironaka.LocalGiraudTransformAlong.eventually_eq_zero_of_mem_geoStrictTransform,Hironaka.LocalGiraudTransformAlong.ne_zero_of_isLocalEquation,Hironaka.LocalGiraudTransformAlong.frequently_ne_zero_of_isLocalEquation}
\mixedleanok
Let \(N\supset C\) be a closed submanifold, \(u\) holomorphic on an open set \(U\subset X\) with multiplicity one along \(N\) at its zeros on \(N\) (at each
such point some holomorphic field tangent to \(N\) has \(\xi(u)\ne0\)), and \(C\cap U\subset V(u)\). Put \(W=N\cap V(u)\cap U\), and let \(W'\), \(N'\) be the strict
transforms. Let \(h\) be a local equation of the exceptional divisor on \(U'\subset Y\), \(z\in U'\cap\pi^{-1}(U)\), and \(u\circ\pi=h\,u'\) near \(z\) with \(u'\)
holomorphic at \(z\). Then:
\begin{enumerate}
\item if \(\pi(z)\in W\), then \(\mathcal I_{W'}=\mathcal I_{N'}+(u')\) at \(z\);
\item near \(z\), a point of \(N'\) lies on \(W'\) if and only if \(u'\) vanishes there (part (3) of Włodarczyk's lemma along \(N\));
\item if \(z\in W'\), some holomorphic field tangent to \(N'\) at \(z\) has \(\xi(u')(z)\ne0\): \(u'\) has multiplicity one along \(N'\) (part (2)).
\end{enumerate}
The hypotheses on \(u\) hold for a tangent direction \(u\) along \(N\) of \((\mathfrak a,\mu)\), \(\mu\ge1\), on \(U\) (Definition~\ref{mc:tangent-direction-along})
and a centre \(C\) in \(\operatorname{cosupp}(\mathfrak a|_N,\mu)\) (Lemma~\ref{mc:tangent-direction-along-basic}). For \(N=X\) this is
Theorem~\ref{bcalc:giraud-local}.
\end{theorem}
\begin{proof}
\mixedleanok
Two facts on \(h\): it vanishes nowhere on \(U'\setminus\pi^{-1}(C)\) (there \(\mathcal O(-D)\ni1\), Lemma~\ref{bcalc:exceptional-divisor}), and it vanishes identically
near no point (Lemma~\ref{bcalc:function-transform}). Hence \(u'\) vanishes on \(W'\) near \(z\): on \(\pi^{-1}(W\setminus C)\) we have \(h\,u'=u\circ\pi=0\) and
\(h\ne0\), and \(u'\) is continuous.

(1) By Lemma~\ref{mc:contact-hypersurface-along}, \(W\) is a submanifold of \(U\), closed in \(U\), of dimension one less than \(N\) at each of its points,
and near \(\pi(z)\) its ideal is generated by \(u\) and the normal coordinates of \(N\), so \(\mathcal I_W\subset\mathcal I_N+(u)\) at \(\pi(z)\). By
Lemma~\ref{bcalc:restrict-open} the restriction \(\pi^{-1}(U)\to U\) is the blowup of \(U\) along \(C\cap U\subset W\), and Theorem~\ref{bcalc:strict-ideal} for it,
read back on \(Y\) as in Theorem~\ref{bcalc:giraud-local}, gives \(h\,g\in\pi^*\mathcal I_W\) for every germ \(g\) of \(\mathcal I_{W'}\) at \(z\). Each germ of
\(\pi^*\mathcal I_W\) is a combination of \(a\circ\pi\), \(a\in\mathcal I_N\), and of \(u\circ\pi=h\,u'\); and \(\pi^*\mathcal I_N=\mathcal O(-D)\,\mathcal I_{N'}\)
(Theorem~\ref{bcalc:strict-ideal} for \(N\)), so \(a\circ\pi\) is a combination of products \(h\,t\), \(t\in\mathcal I_{N'}\). Thus \(h\,g\) is a combination of
the \(h\,t\) (\(t\in\mathcal I_{N'}\) or \(t=u'\)), and \(g\) is a combination of the \(t\) (Lemma~\ref{bcalc:tangent-generator-along}). Conversely
\(\mathcal I_{N'}\subset\mathcal I_{W'}\) as \(W'\subset N'\), and \(u'\in\mathcal I_{W'}\) at \(z\) by the first paragraph.

(2) Let \(y\in N'\) near \(z\); the hypotheses hold at \(y\). If \(y\in W'\), then \(u'(y)=0\) by the first paragraph. Let \(u'(y)=0\). If \(\pi(y)\notin C\),
then \(\pi(y)\in N\) and \(u(\pi(y))=h(y)u'(y)=0\), so \(y\in\pi^{-1}(W\setminus C)\subset W'\). If \(\pi(y)\in C\), then \(\pi(y)\in W\); were \(y\notin W'\), the
constant \(1\) would lie in \(\mathcal I_{W'}\) at \(y\), hence by (1) be a combination of germs vanishing at \(y\in N'\), which is absurd.

(3) As \(W\) is closed in \(U\), \(\pi(z)\in W\). Let \(N\) have dimension \(k\) at \(\pi(z)\); then \(k\ge1\) and \(W\) has dimension \(k-1\) there
(Lemma~\ref{mc:contact-hypersurface-along}). By Theorem~\ref{bcalc:strict-submanifold} (for \(W\) in the restricted blowup, and for \(N\)), \(W'\) and \(N'\) are
regular at \(z\) of dimensions \(k-1<k\), and \(W'\subset N'\). By (1) and \(u'(z)=0\), Lemma~\ref{bcalc:tangent-generator-along} applies with \(A=W'\), \(B=N'\),
\(f=u'\). Finally, a tangent direction \(u\) along \(N\) is a germ of \(\mathcal D_N^{\mu-1}(\mathfrak a)\) at every point of \(U\), hence holomorphic there, of
multiplicity one along \(N\) by definition, and \(C\cap U\subset\operatorname{cosupp}(\mathfrak a|_N,\mu)\cap U\subset V(u)\)
(Lemma~\ref{mc:tangent-direction-along-basic}).
\end{proof}

\begin{lemma}[Blowing up a codimension-one part of the support]
\label{bcalc:codim-one}
\uses{bcalc:function-transform,bcalc:exceptional-divisor,hir:hadamard-bridge,hir:giraud-support,hir:cosupport,hir:order-basic,hir:geo-transform-ambient,mc:maximal-order,mc:principal-power}
\source{wlodarczyk05-hironaka,wlodarczyk09-analytic-resolution:page-0018,wlodarczyk09-analytic-resolution:page-0019}
\upstreamlean{TranscendentalNakamaye.Pending.CodimOneCentre.notMem_cosupport_controlledTransform_of_codim_one}\lean{Hironaka.CodimOneCentre.notMem_cosupport_succ_of_isMaximalOrder}\upstreamlean{TranscendentalNakamaye.Pending.CodimOneCentre.cosupport_controlledTransform_inter_preimage_eq_empty,TranscendentalNakamaye.Pending.CodimOneCentre.cosupp_geoTransform_inter_preimage_eq_empty,TranscendentalNakamaye.Pending.CodimOneCentre.exists_generator_vanishingIdeal_of_regularAtDim,TranscendentalNakamaye.Pending.CodimOneCentre.VanishesToOrderOn.congr}\lean{Hironaka.CodimOneCentre.vanishesToOrderOn_pow}\upstreamlean{TranscendentalNakamaye.Pending.CodimOneCentre.apply_ne_zero_of_not_vanishesToOrderOn}
\mixedleanok
Let \(\mathfrak a\) be coherent, \(\mu\ge1\), and let the centre \(C\) of the blowup \(\pi\) lie in \(\operatorname{cosupp}(\mathfrak a,\mu)\) (orders on the whole
manifold). If \(z\in\pi^{-1}(C)\), \(C\) has dimension \(n-1\) at \(x=\pi(z)\), and \(\operatorname{ord}_x\mathfrak a\le\mu\), then
\(z\notin\operatorname{cosupp}(\sigma^{c}(\mathfrak a,\mu),\mu)\). Consequently, if \((\mathfrak a,\mu)\) is of maximal order and \(C\) has codimension one
at each of its points, the cosupport of \(\sigma^{c}(\mathfrak a,\mu)\) does not meet the exceptional divisor \(\pi^{-1}(C)\) (the transform half of
Włodarczyk's Lemma ``le:codim1'', after Villamayor); the same holds for the corrected transform of a marked ideal with \(N=M\).
\end{lemma}
\begin{proof}
\mixedleanok
Near \(x\), a flattening coordinate \(u\) generates \(\mathcal I_C\) (Hadamard's lemma in flattening coordinates), so \(\mathcal I_C^\mu=(u^\mu)\)
(Lemma~\ref{mc:principal-power}). By the Hadamard bridge (Lemma~\ref{hir:hadamard-bridge}) \(\mathfrak a\subset\mathcal I_C^\mu\). As
\(\operatorname{ord}_x\mathfrak a\le\mu\), some \(f\in\mathfrak a_x\) does not vanish to order \(\mu+1\) at \(x\); write \(f=c\,u^\mu\). If \(c(x)=0\), then
\(c\) and \(u\) vanish at \(x\) and \(f\) vanishes to order \(1+\mu\) (Lemma~\ref{hir:order-basic}); so \(c(x)\neq0\). Let \(h\) be a local equation of the
exceptional divisor at \(z\) (Lemma~\ref{bcalc:exceptional-divisor}). Both \(u\circ\pi\) and \(h\) generate \(\pi^*\mathcal I_C\) at \(z\) (the germs of
\(\pi^*\mathcal I_C\) are combinations of \(a\circ\pi\), \(a\in(u)\)), so \(u\circ\pi=d\,h\) and \(h=d'\,(u\circ\pi)\), whence \(h\,(d'd-1)=0\) and
\(d'd=1\) near \(z\) (cancellation of \(h\)), and \(d(z)\neq0\). Then \(f\circ\pi=h^\mu\,(c\circ\pi)\,d^\mu\), so \((c\circ\pi)\,d^\mu\) lies in
\(\sigma^{c}(\mathfrak a,\mu)\) at \(z\) (Lemma~\ref{bcalc:function-transform}); it does not vanish at \(z\), hence \(z\) is not in the cosupport of
\((\sigma^{c}(\mathfrak a,\mu),\mu)\) for \(\mu\ge1\) (Definition~\ref{hir:cosupport}). If \((\mathfrak a,\mu)\) is of maximal order, \(\mathcal D^\mu\mathfrak a\)
contains \(1\), so \(\operatorname{cosupp}(\mathfrak a,\mu+1)\subset\operatorname{cosupp}(\mathcal D^\mu\mathfrak a,1)=\emptyset\) (Lemma~\ref{hir:giraud-support},
Definition~\ref{mc:maximal-order}): \(\operatorname{ord}_x\mathfrak a\le\mu\) everywhere. For a marked ideal with \(N=M\) the corrected transform has
\(N'=M'\) and ideal \(\sigma^{c}(\mathfrak a,d)\) (Lemma~\ref{hir:geo-transform-ambient}).
\end{proof}

\paragraph{Codimension-one components of the support.} Włodarczyk's Lemma ``le:codim1'' has a second half: for a marked ideal of maximal order,
the codimension-one components of the support are smooth and isolated. In the analytic category the components are the irreducible analytic
components (Definition~\ref{ana:nonkahler-locus}); ``isolated'' is made precise as being a connected component of the support and meeting no
other irreducible component; and Włodarczyk's argument ``the support lies in the smooth hypersurface \(V(u)\) of a tangent direction'' is completed
by an identity theorem on \(V(u)\), Lemma~\ref{bcalc:analytic-fill}. The analytic version of the lemma (Włodarczyk 2009, Lemma 5.3.7, printed p.~48) has the same statements for germs \(M_Z\), under the
additional hypothesis that the support has codimension one, which is not used here; its proof writes \(\mathcal I=(u^\mu)\) locally, while
Lemma~\ref{bcalc:codim-one} only uses \(\mathfrak a\subset\mathcal I_C^\mu\) and \(\operatorname{ord}\mathfrak a\le\mu\), and so applies to any centre of
dimension \(n-1\) inside the support (as in Step~1aa for \(s=1\), where the centre is a member).

\begin{lemma}[Filling a regular set by an analytic set of the same local dimension]
\label{bcalc:analytic-fill}
\uses{rem:thin-dimension-drop,ana:regular-component,rem:open-subset-analytic,ana:analytic-preimage,lag:regular-point-basics,rem:dimension}
\lean{Hironaka.CodimOneSupport.exists_isOpen_inter_subset_of_locDimGE,Hironaka.CodimOneSupport.eventually_mem_of_locDimGE,Hironaka.CodimOneSupport.eventually_mem_iff_of_regularAtDim_of_locDimGE,Hironaka.CodimOneSupport.regularAtDim_of_eventually_mem_iff,Hironaka.CodimOneSupport.mem_of_locDimGE}
\leanok
Let \(A\) be an analytic subset of a complex manifold \(X\), let \(H\subset X\) be regular of dimension \(k\) at \(x\), and suppose \(A\subset H\) near \(x\)
and \(\dim_xA\ge k\) (every neighbourhood of \(x\) contains a regular point of \(A\) of dimension \(\ge k\), Remark~\ref{rem:dimension}). Then
\(A=H\) near \(x\); in particular \(A\) is regular of dimension \(k\) at \(x\).
\end{lemma}
\begin{proof}
\leanok
As \(A\) is closed and its regular points accumulate at \(x\), \(x\in A\). Let \(g_1,\dots,g_m\) be holomorphic local equations of \(A\) at \(x\) in the
chart \(e\) at \(x\), defined on an open set \(D\ni e(x)\), and let \(N_0\) be an open neighbourhood of \(x\) on which \(A\subset H\). Let \(\Phi\) flatten \(H\)
at \(x\): \(\Phi(e(x))=0\), and for \(y\) in the chart domain with \(e(y)\) in the domain of \(\Phi\), \(y\in H\iff\Phi(e(y))\in\mathbf C^k\times0\). Choose
a ball \(B\) of radius \(r\) about \(0\) in the target of \(\Phi\) with \(\Phi^{-1}(B)\subset D\cap e(N_0)\), and put \(W=e^{-1}\Phi^{-1}(B)\), an open
neighbourhood of \(x\) inside \(N_0\). Then \(H\cap W=\psi(B\cap(\mathbf C^k\times0))\) with \(\psi=e^{-1}\circ\Phi^{-1}\), every point of \(H\cap W\) is
regular of dimension \(k\) for \(H\), and \(H\cap W\) is closed in \(W\) (Lemma~\ref{ana:regular-component}).

First, \(A\cap W\) contains a nonempty relatively open part of \(H\cap W\). Otherwise, regard \(W\) as a complex manifold. In it, \(T=H\cap W\) is a
closed submanifold, hence analytic (Lemma~\ref{ana:analytic-preimage}), and all its regular points have dimension \(k\) (the dimension at a regular
point is unique, Lemma~\ref{lag:regular-point-basics}); \(T'=A\cap W\subset T\) has empty interior relative to \(T\). By the thin-dimension drop
(Remark~\ref{rem:thin-dimension-drop}) every regular point of \(T'\) has dimension \(<k\). But \(\dim_xA\ge k\) gives a point \(y\in W\) at which \(A\)
is regular of some dimension \(j\ge k\), and \(T'\) is regular of dimension \(j\) at \(y\) in \(W\) (regularity is local,
Remark~\ref{rem:open-subset-analytic}): a contradiction.

Second, the identity theorem on the flattened ball. Let \(\iota\colon\mathbf C^k\to\mathbf C^n\) be the inclusion \(w\mapsto(w,0)\); the set
\(\iota^{-1}(B)\) is open and convex, and each \(g_i\circ\Phi^{-1}\circ\iota\) is holomorphic on it. By the first step it vanishes on the nonempty open
set of parameters \(w\) with \(\psi(\iota(w))\) in the open part found there; hence it vanishes on all of \(\iota^{-1}(B)\)
(Lemma~\ref{ana:regular-component}), i.e.\ \(H\cap W=\psi(\iota(\iota^{-1}B))\subset A\). With \(A\cap W\subset H\) this gives \(A\cap W=H\cap W\), and
regularity at \(x\) only depends on the germ (Lemma~\ref{lag:regular-point-basics}).
\end{proof}

\begin{lemma}[Irreducible analytic sets: connectedness and local dimension]
\label{bcalc:irreducible-connected}
\uses{fpgeom:complex-analytic-subvariety,ana:analytic-preimage,ana:pure-dimension,lag:regular-locus-dense,rem:dimension}
\upstreamlean{TranscendentalNakamaye.Pending.CodimOneSupport.isPreconnected_of_isComplexAnalyticallyIrreducible,TranscendentalNakamaye.Pending.CodimOneSupport.locDimGE_carrier}
\upstreamleanok
An analytically irreducible subset \(V\) of a complex manifold is connected. If \(V\) is a closed irreducible analytic subvariety, then
\(\dim_xV\ge\dim V\) at every point \(x\in V\).
\end{lemma}
\begin{proof}
\upstreamleanok
Let \(P,Q\) be closed sets with \(V\subset P\cup Q\) and \(V\cap P\cap Q=\emptyset\). Near each point of \(V\cap P\) the set \(V\cap P\) agrees with \(V\)
(on the open set \(X\setminus Q\)), and \(V\cap P\) is closed, so it is analytic (Lemma~\ref{ana:analytic-preimage}); likewise \(V\cap Q\). As
\(V=(V\cap P)\cup(V\cap Q)\), irreducibility (Definition~\ref{fpgeom:complex-analytic-subvariety}) gives \(V=V\cap P\) or \(V=V\cap Q\), i.e.\ \(V\subset P\)
or \(V\subset Q\). For the second claim, every neighbourhood of \(x\in V\) contains a regular point of \(V\) (Lemma~\ref{lag:regular-locus-dense}), and
every regular point of \(V\) has dimension \(\dim V\) (Lemma~\ref{ana:pure-dimension}).
\end{proof}

\begin{lemma}[The support of a coherent marked ideal is analytic]
\label{bcalc:cosupport-analytic}
\uses{ana:zero-locus-coherent,hir:giraud-support,hir:derivative-ideal-generators}
\lean{Hironaka.CodimOneSupport.isComplexAnalyticSubset_zeroLocus}\upstreamlean{TranscendentalNakamaye.Pending.CodimOneSupport.isComplexAnalyticSubset_cosupport}
\mixedleanok
If \(\mathfrak a\) is coherent and \(\mu\ge1\), then \(\operatorname{cosupp}(\mathfrak a,\mu)\) (orders on the whole manifold) is an analytic subset.
\end{lemma}
\begin{proof}
\mixedleanok
By Lemma~\ref{hir:giraud-support}, \(\operatorname{cosupp}(\mathfrak a,\mu)=V(\mathcal D^{\mu-1}\mathfrak a)\), and \(\mathcal D^{\mu-1}\mathfrak a\) is coherent
(Lemma~\ref{hir:derivative-ideal-generators}). The zero locus of a coherent ideal is analytic: near each point it is the common zero set of finitely
many local generators (Lemma~\ref{ana:zero-locus-coherent}).
\end{proof}

\begin{lemma}[Near a point of local dimension \(n-1\) the support is the hypersurface of a tangent direction]
\label{bcalc:codim-one-local}
\uses{bcalc:analytic-fill,bcalc:hypersurface,mc:tangent-direction,mc:tangent-direction-support,mc:tangent-direction-exists,mc:maximal-order}
\source{wlodarczyk05-hironaka}
\upstreamlean{TranscendentalNakamaye.Pending.CodimOneSupport.eventually_mem_iff_of_locDimGE,TranscendentalNakamaye.Pending.CodimOneSupport.regularAtDim_of_locDimGE}
\upstreamleanok
Let \(\mathfrak a\) be coherent, \(\mu\ge1\), and \(A\subset\operatorname{cosupp}(\mathfrak a,\mu)\) analytic with \(\dim_xA\ge n-1\). If \(u\) is a tangent
direction of \((\mathfrak a,\mu)\) on an open set \(U\ni x\) (Definition~\ref{mc:tangent-direction}), then \(A=V(u)\) near \(x\). If moreover
\((\mathfrak a,\mu)\) is of maximal order, \(A\) is regular of dimension \(n-1\) at \(x\).
\end{lemma}
\begin{proof}
\upstreamleanok
\(A\) is closed and its regular points accumulate at \(x\), so \(x\in A\); as \(A\cap U\subset\operatorname{cosupp}(\mathfrak a,\mu)\cap U\subset V(u)\)
(Lemma~\ref{mc:tangent-direction-support}), \(u(x)=0\). The tangent direction \(u\) is holomorphic on \(U\) and of multiplicity one at its zeros, so
\(H=V(u)\cap U\) is regular of dimension \(n-1\) at \(x\) (Lemma~\ref{bcalc:hypersurface}), and \(A\subset H\) on \(U\). By
Lemma~\ref{bcalc:analytic-fill}, \(A=H\) near \(x\), i.e.\ \(A=V(u)\) near \(x\). If \((\mathfrak a,\mu)\) is of maximal order
(Definition~\ref{mc:maximal-order}), a tangent direction exists on a neighbourhood of every point of the support
(Lemma~\ref{mc:tangent-direction-exists}), and then \(A\) is regular of dimension \(n-1\) at \(x\) because \(H\) is.
\end{proof}

\begin{definition}[The codimension-one part of the support]
\label{bcalc:codim-one-part-def}
\uses{hir:cosupport}
\upstreamlean{TranscendentalNakamaye.Pending.CodimOneSupport.codimOnePart,TranscendentalNakamaye.Pending.CodimOneSupport.codimOnePart_subset}
\upstreamleanok
The \emph{codimension-one part} \(R(\mathfrak a,\mu)\) of the support \(S=\operatorname{cosupp}(\mathfrak a,\mu)\) is the set of points at which \(S\) is
regular of dimension \(n-1\).
\end{definition}

\begin{theorem}[The codimension-one part of the support is a smooth isolated hypersurface]
\label{bcalc:codim-one-part}
\uses{bcalc:codim-one-part-def,bcalc:codim-one-local,bcalc:cosupport-analytic,bcalc:hypersurface,bcalc:nested-flattening,lag:regular-point-basics,lag:regular-locus-dense,hir:giraud-support,mc:tangent-direction-exists,mc:tangent-direction-support,rem:dimension}
\source{wlodarczyk05-hironaka}
\upstreamlean{TranscendentalNakamaye.Pending.CodimOneSupport.mem_codimOnePart_iff,TranscendentalNakamaye.Pending.CodimOneSupport.eventually_mem_codimOnePart_iff,TranscendentalNakamaye.Pending.CodimOneSupport.isClosed_codimOnePart,TranscendentalNakamaye.Pending.CodimOneSupport.regularAtDim_codimOnePart,TranscendentalNakamaye.Pending.CodimOneSupport.isClosedSubmanifold_codimOnePart,TranscendentalNakamaye.Pending.CodimOneSupport.exists_isOpen_cosupport_inter_eq_codimOnePart,TranscendentalNakamaye.Pending.CodimOneSupport.eventually_mem_codimOnePart_iff_eq_zero,TranscendentalNakamaye.Pending.CodimOneSupport.isClosed_cosupport_diff_codimOnePart,TranscendentalNakamaye.Pending.CodimOneSupport.lt_of_regularAtDim_cosupport_diff_codimOnePart,TranscendentalNakamaye.Pending.CodimOneSupport.dimLE_cosupport_diff_codimOnePart}
\upstreamleanok
Let \((\mathfrak a,\mu)\) be coherent, of maximal order, \(\mu\ge1\), with support \(S\) and codimension-one part \(R\). Then:
\begin{enumerate}
\item \(R\) is the set of points \(x\) with \(\dim_xS\ge n-1\);
\item \(R\) is closed, open in \(S\), and regular of dimension \(n-1\) at each of its points; in particular \(R\) is a closed submanifold;
\item near each \(x\in R\), \(R=S=V(u)\) for every tangent direction \(u\) on a neighbourhood of \(x\);
\item \(S\setminus R\) is closed, and every regular point of \(S\setminus R\) has dimension \(<n-1\): \(\dim(S\setminus R)\le n-2\).
\end{enumerate}
\end{theorem}
\begin{proof}
\upstreamleanok
(1) A regular point of dimension \(n-1\) has local dimension \(n-1\) (Remark~\ref{rem:dimension}); conversely \(S\) is analytic
(Lemma~\ref{bcalc:cosupport-analytic}), and Lemma~\ref{bcalc:codim-one-local} with \(A=S\) shows that \(S\) is regular of dimension \(n-1\) at a point of
local dimension \(\ge n-1\). (2) A flattening of \(S\) at \(x\in R\) flattens \(S\) at all nearby points of \(S\), with the same dimension
(Theorem~\ref{lag:regular-locus-dense}); so \(R=S\) near \(x\). Hence \(R\) is open in \(S\) (the union of these neighbourhoods meets \(S\) in \(R\)), and
\(R\) is regular of dimension \(n-1\) at \(x\) because \(S\) is and the two agree near \(x\). A limit of points of \(R\) has local dimension \(\ge n-1\) for
\(S\), so it lies in \(R\) by (1): \(R\) is closed. (3) \(R=S\) near \(x\), and \(S=V(u)\) near \(x\) by Lemma~\ref{bcalc:codim-one-local}.
(4) If \(O\) is open with \(S\cap O=R\), then \(S\setminus R=S\setminus O\), closed since \(S\) is (Lemma~\ref{hir:giraud-support}). Let \(S\setminus R\) be
regular of dimension \(j\) at \(y\). As \(R\) is closed, \(S\setminus R\) and \(S\) agree on the open neighbourhood \(X\setminus R\) of \(y\), so \(S\) is regular
of dimension \(j\) at \(y\). With a tangent direction \(u\) on \(U\ni y\) (Lemma~\ref{mc:tangent-direction-exists}), \(S\cap U\subset V(u)\cap U\)
(Lemma~\ref{mc:tangent-direction-support}), which is regular of dimension \(n-1\) at \(y\) (Lemma~\ref{bcalc:hypersurface}); a submanifold of a
submanifold has smaller dimension (Lemma~\ref{bcalc:nested-flattening}), so \(j\le n-1\); and \(j\ne n-1\) because \(y\notin R\).
\end{proof}

\begin{theorem}[Codimension-one components of the support are smooth and isolated]
\label{bcalc:codim-one-components}
\uses{bcalc:codim-one-part,bcalc:codim-one-local,bcalc:irreducible-connected,bcalc:cosupport-analytic,bcalc:codim-one,rem:dimension-monotone,mc:tangent-direction-exists,mc:tangent-direction-support,ana:nonkahler-locus}
\source{wlodarczyk05-hironaka,wlodarczyk09-analytic-resolution:page-0018,wlodarczyk09-analytic-resolution:page-0019}
\upstreamlean{TranscendentalNakamaye.Pending.CodimOneSupport.carrier_subset_codimOnePart,TranscendentalNakamaye.Pending.CodimOneSupport.eventually_mem_carrier_iff,TranscendentalNakamaye.Pending.CodimOneSupport.regularAtDim_carrier,TranscendentalNakamaye.Pending.CodimOneSupport.eventually_mem_carrier_iff_mem_cosupport,TranscendentalNakamaye.Pending.CodimOneSupport.exists_isOpen_cosupport_inter_eq_carrier,TranscendentalNakamaye.Pending.CodimOneSupport.connectedComponentIn_eq_carrier,TranscendentalNakamaye.Pending.CodimOneSupport.disjoint_carrier,TranscendentalNakamaye.Pending.CodimOneSupport.isIrreducibleComponent_codim_one,TranscendentalNakamaye.Pending.CodimOneSupport.cosupport_controlledTransform_inter_preimage_carrier_eq_empty}
\upstreamleanok
Let \((\mathfrak a,\mu)\) be coherent, of maximal order, \(\mu\ge1\), with support \(S\), and let \(V\subset S\) be a closed irreducible analytic subvariety
with \(\dim V=n-1\) (for instance an irreducible component of \(S\) of codimension one). Then \(V\subset R(\mathfrak a,\mu)\), and:
\begin{enumerate}
\item \(V\) is regular of dimension \(n-1\) at each of its points;
\item near each \(x\in V\), \(V=V(u)\) for every tangent direction \(u\) on a neighbourhood of \(x\), and \(V=S\);
\item \(V\) is open in \(S\), and it is the connected component of \(S\) through each of its points;
\item \(V\) meets no irreducible component \(V'\neq V\) of \(S\);
\item after the blowup \(\pi\) along \(V\), the cosupport of \(\sigma^{c}(\mathfrak a,\mu)\) does not meet the exceptional divisor \(\pi^{-1}(V)\).
\end{enumerate}
Items (1)--(4) are the first half of Włodarczyk's Lemma ``le:codim1'' (after Villamayor), item (5) its second half for one component.
\end{theorem}
\begin{proof}
\upstreamleanok
At every \(x\in V\), \(\dim_xV\ge\dim V=n-1\) (Lemma~\ref{bcalc:irreducible-connected}); hence \(\dim_xS\ge n-1\) (local dimension is monotone,
Remark~\ref{rem:dimension-monotone}; \(S\) is analytic by Lemma~\ref{bcalc:cosupport-analytic}) and \(x\in R\) by Theorem~\ref{bcalc:codim-one-part}.
Lemma~\ref{bcalc:codim-one-local} with \(A=V\) gives (1) and the first half of (2). For the second half, take a tangent direction \(u\) on \(U\ni x\)
(Lemma~\ref{mc:tangent-direction-exists}): near \(x\), \(S\cap U\subset V(u)=V\subset S\) (Lemma~\ref{mc:tangent-direction-support}). (3) The union \(O\)
of the neighbourhoods of the points of \(V\) on which \(V=S\) is open, contains \(V\), and \(S\cap O=V\). By Lemma~\ref{bcalc:irreducible-connected}, \(V\)
is connected, so it lies in the connected component \(K\) of \(S\) through \(x\in V\). Conversely \(K\subset S\subset O\cup(X\setminus V)\) and \(K\) does not
meet \(O\setminus V=(S\cap O)\setminus V\) inside \(S\); so the connected set \(K\) lies in \(O\) or in \(X\setminus V\), and \(x\in K\cap V\) excludes the second:
\(K\subset S\cap O=V\). (4) If an irreducible component \(V'\) of \(S\) (Definition~\ref{ana:nonkahler-locus}) contains \(x\in V\), then \(V'\) is connected
(Lemma~\ref{bcalc:irreducible-connected}) and contained in \(S\), hence in the connected component of \(S\) through \(x\), which is \(V\) by (3); since
\(V\subset S\), maximality of \(V'\) gives \(V'=V\). (5) By (1), \(V\) is a centre of dimension \(n-1\) at each of its points inside the support;
apply Lemma~\ref{bcalc:codim-one}.
\end{proof}

\begin{corollary}[Blowing up the codimension-one part of the support]
\label{bcalc:codim-one-step}
\uses{bcalc:codim-one,bcalc:codim-one-part,hir:geo-transform-ambient}
\source{wlodarczyk05-hironaka,wlodarczyk09-analytic-resolution:page-0027}
\upstreamlean{TranscendentalNakamaye.Pending.CodimOneSupport.cosupport_controlledTransform_inter_preimage_codimOnePart_eq_empty,TranscendentalNakamaye.Pending.CodimOneSupport.cosupp_geoTransform_inter_preimage_codimOnePart_eq_empty}
\upstreamleanok
Let \((\mathfrak a,\mu)\) be coherent, of maximal order, \(\mu\ge1\), and let \(\pi\) be the blowup along its codimension-one part \(R\). Then the
cosupport of \(\sigma^{c}(\mathfrak a,\mu)\) does not meet the exceptional divisor \(\pi^{-1}(R)\); the same holds for the corrected transform of a marked
ideal with \(N=M\). This is Step 1ba of Włodarczyk's algorithm, the elimination of the codimension-one part of the support.
\end{corollary}
\begin{proof}
\upstreamleanok
\(R\) lies in the support and has dimension \(n-1\) at each of its points (Theorem~\ref{bcalc:codim-one-part}); apply Lemma~\ref{bcalc:codim-one}.
\end{proof}

\paragraph{The transform away from the centre.} Off the exceptional divisor a blowup is biholomorphic and the transform of a marked ideal is its
pullback; in particular the cosupport does not change there.

\begin{lemma}[A blowup is an isomorphism off its centre]
\label{bcalc:isomorphism-off-centre}
\uses{ana:blowups}
\lean{Hironaka.BlowupIsomorphismOff.exists_section_compl,Hironaka.BlowupIsomorphismOff.isIsomorphismOver_compl_of_isBlowupAlong,Hironaka.BlowupIsomorphismOff.exceptionalLocus_subset_preimage}
\leanok
Let \(\pi\colon Y\to X\) be the blowup along \(C\). There is a holomorphic map \(s\colon X\setminus C\to Y\) with \(\pi\circ s=\mathrm{id}\) and
\(s(\pi(z))=z\) whenever \(\pi(z)\notin C\). Hence \(\pi\) is an isomorphism over \(X\setminus C\), and \(\operatorname{Exc}(\pi)\subset\pi^{-1}(C)\).
\end{lemma}
\begin{proof}
\leanok
On the open manifold \(X\setminus C\) the pullback of \(\mathcal I_C\) by the inclusion is the unit ideal, which is invertible; by the universal property
of the blowup (Definition~\ref{ana:blowups}) the inclusion factors holomorphically as \(\pi\circ s\). On the open set \(\pi^{-1}(X\setminus C)\) both the
inclusion and \(s\circ\pi\) are holomorphic lifts of the restriction of \(\pi\), whose ideal \(\pi^{*}\mathcal I_C\) is again the unit ideal; by the
uniqueness clause they agree, i.e.\ \(s(\pi(z))=z\). So \(\pi\) maps \(\pi^{-1}(X\setminus C)\) bijectively onto \(X\setminus C\), and the chain rule for
\(\pi\circ s=\mathrm{id}\) makes \(d\pi\) surjective, hence bijective (equal dimensions), at every point over \(X\setminus C\).
\end{proof}

\begin{theorem}[The transform away from the centre]
\label{bcalc:transform-off-centre}
\uses{bcalc:isomorphism-off-centre,bcalc:controlled-transform,hir:geo-transform,hir:geo-transform-wellformed,hir:geo-transform-ambient,hir:cosupport,hir:order-basic,hir:point-power-pullback,hir:open-restriction-ideal,hir:marked-ideal-operations}
\source{wlodarczyk05-hironaka,bierstone-milman08-functoriality}
\lean{Hironaka.TransformOffCentre.eventually_mem_geoStrictTransform_iff,Hironaka.TransformOffCentre.IsInSpanOnAt.vanishesToOrderOn,Hironaka.TransformOffCentre.mem_cosupp_geoTransform_of_mem_cosupp,Hironaka.TransformOffCentre.mem_cosupp_of_mem_cosupp_geoTransform,Hironaka.TransformOffCentre.mem_cosupp_geoTransform_iff}\upstreamlean{TranscendentalNakamaye.Pending.TransformOffCentre.cosupp_geoTransform_diff_preimage}\lean{Hironaka.TransformOffCentre.cosupp_geoTransform_subset_preimage}\upstreamlean{TranscendentalNakamaye.Pending.TransformOffCentre.cosupport_controlledTransform_diff_preimage}
\mixedleanok
Let \(\underline{\mathcal I}=(M,N,E,\mathfrak a,d)\) be a marked ideal with \(N\) closed and \(\mathfrak a\) an ideal of holomorphic germs, \(\pi\) the blowup
along \(C\), and \(\underline{\mathcal I}'=(M',N',E',\mathfrak a',d)\) the transform (Definition~\ref{hir:geo-transform}). If \(\pi(z)\notin C\), then
\(N'=\pi^{-1}(N)\) near \(z\), and \(z\in\operatorname{cosupp}\underline{\mathcal I}'\) if and only if \(\pi(z)\in\operatorname{cosupp}\underline{\mathcal I}\).
Consequently \(\operatorname{cosupp}\underline{\mathcal I}'\setminus\pi^{-1}(C)=\pi^{-1}(\operatorname{cosupp}\underline{\mathcal I}\setminus C)\), and
\(\operatorname{cosupp}\underline{\mathcal I}'\subset\pi^{-1}(\operatorname{cosupp}\underline{\mathcal I})\) when \(C\subset\operatorname{cosupp}\underline{\mathcal I}\). For
\(N=M\) and \(\mathfrak a\) coherent: \(\operatorname{cosupp}(\sigma^{c}(\mathfrak a,\mu),\mu)\setminus\pi^{-1}(C)=\pi^{-1}(\operatorname{cosupp}(\mathfrak a,\mu)\setminus C)\).
\end{theorem}
\begin{proof}
\mixedleanok
\(N'=\overline{\pi^{-1}(N\setminus C)}\subset\pi^{-1}(N)\) as \(N\) is closed, and on the open set \(\pi^{-1}(X\setminus C)\ni z\) the set \(\pi^{-1}(N)\) equals
\(\pi^{-1}(N\setminus C)\subset N'\). At \(z\) the constant \(1\) lies in \(\pi^*\mathcal I_C\), hence in its \(d\)-th power
(Lemma~\ref{hir:geo-transform-wellformed}), so a germ \(f\) of \(\mathfrak a'=(\mathfrak a\mathcal O_{N'}:\pi^*\mathcal I_C^{\,d})\) satisfies \(f=1\cdot f\in\mathfrak a\mathcal O_{N'}\): on
\(N'\) near \(z\), \(f\) is a holomorphic combination of germs of \(\pi^*\mathfrak a\), themselves combinations of \(a\circ\pi\), \(a\in\mathfrak a_{\pi(z)}\);
conversely every \(a\circ\pi\) lies in \(\mathfrak a'\) at \(z\), because \(\mathfrak a\mathcal O_{N'}\) is closed under holomorphic multiples.

If \(\pi(z)\in\operatorname{cosupp}\underline{\mathcal I}\), each \(a\in\mathfrak a_{\pi(z)}\) agrees on \(N\) near \(\pi(z)\) with some \(h\in\mathfrak m_{\pi(z)}^d\)
(Definition~\ref{hir:cosupport}, Lemma~\ref{hir:order-basic}); then \(a\circ\pi=h\circ\pi\) on \(N'\subset\pi^{-1}(N)\) near \(z\), and
\(h\circ\pi\in\mathfrak m_z^d\) (Lemma~\ref{hir:point-power-pullback}). Germs vanishing to order \(d\) along \(N'\) are closed under holomorphic combinations
(Lemma~\ref{hir:marked-ideal-operations}), so every germ of \(\mathfrak a'\) at \(z\) vanishes to order \(d\) along \(N'\): \(z\in\operatorname{cosupp}\underline{\mathcal I}'\).

Conversely let \(z\in\operatorname{cosupp}\underline{\mathcal I}'\) and \(a\in\mathfrak a_{\pi(z)}\). As \(a\circ\pi\in\mathfrak a'_z\), there is \(h\in\mathfrak m_z^d\) with
\(a\circ\pi=h\) on \(N'\) near \(z\). Let \(s\) be the section of Lemma~\ref{bcalc:isomorphism-off-centre} on the open manifold \(X\setminus C\), with
\(s(\pi(z))=z\). Then \(h\circ s\in\mathfrak m_{\pi(z)}^d\) on \(X\setminus C\) (Lemma~\ref{hir:point-power-pullback}), and its extension by zero is a germ of
\(\mathfrak m_{\pi(z)}^d\) on \(X\) (Lemma~\ref{hir:open-restriction-ideal}). For \(y\in N\) near \(\pi(z)\), \(s(y)\) is near \(z\) and lies in
\(\pi^{-1}(N)=N'\) there, so \(h(s(y))=a(\pi(s(y)))=a(y)\): \(a\) vanishes to order \(d\) along \(N\) at \(\pi(z)\), and \(\pi(z)\in N\). So
\(\pi(z)\in\operatorname{cosupp}\underline{\mathcal I}\). The set identities follow pointwise; for \(N=M\) the transform has \(N'=M'\) and ideal
\(\sigma^{c}(\mathfrak a,\mu)\) (Lemma~\ref{hir:geo-transform-ambient}).
\end{proof}

\begin{corollary}[Cosupports along multiple test blow-ups]
\label{bcalc:transform-off-centre-tower}
\uses{bcalc:transform-off-centre,hir:geo-transform,hir:open-restriction-tower}
\lean{Hironaka.TransformOffCentre.cosupp_geoTransformAlong_subset}
\leanok
Let \(\underline{\mathcal I}\) have \(N\) a closed submanifold and \(\mathfrak a\) coherent. Along every multiple test blow-up \(T\) of \(\underline{\mathcal I}\), the
cosupport of the transform \(\underline{\mathcal I}_T\) lies over \(\operatorname{cosupp}\underline{\mathcal I}\).
\end{corollary}
\begin{proof}
\leanok
Induction along \(T\). Every stage has \(N\) a closed submanifold and a coherent ideal (Lemma~\ref{hir:open-restriction-tower}), and the next centre is
admissible, so it lies in the cosupport of that stage; Theorem~\ref{bcalc:transform-off-centre} puts the cosupport of the next stage over the cosupport
of that stage, which lies over \(\operatorname{cosupp}\underline{\mathcal I}\) by induction.
\end{proof}

\begin{corollary}[After blowing up the codimension-one part the support has codimension at least two]
\label{bcalc:codim-one-step-rest}
\uses{bcalc:codim-one-step,bcalc:codim-one-part,bcalc:transform-off-centre}
\source{wlodarczyk05-hironaka,wlodarczyk09-analytic-resolution:page-0027}
\upstreamlean{TranscendentalNakamaye.Pending.TransformOffCentre.cosupport_controlledTransform_subset_preimage_diff_codimOnePart}
\upstreamleanok
Let \((\mathfrak a,\mu)\) be coherent, of maximal order, \(\mu\ge1\), with support \(S\) and codimension-one part \(R\), and let \(\pi\) be the blowup along
\(R\). Then \(\operatorname{cosupp}(\sigma^{c}(\mathfrak a,\mu),\mu)\subset\pi^{-1}(S\setminus R)\), where \(S\setminus R\) is closed of dimension \(\le n-2\)
(Theorem~\ref{bcalc:codim-one-part}).
\end{corollary}
\begin{proof}
\upstreamleanok
Points over \(R\) are excluded by Corollary~\ref{bcalc:codim-one-step}; at the other points Theorem~\ref{bcalc:transform-off-centre} gives
\(\operatorname{cosupp}(\sigma^{c}(\mathfrak a,\mu),\mu)\setminus\pi^{-1}(R)=\pi^{-1}(S\setminus R)\).
\end{proof}

\paragraph{Codimension one inside a submanifold.} In the induction of the algorithm the marked ideals are \((M,N,E,\mathfrak a,d)\) with \(N\) a
submanifold of dimension \(k+1\), and Step~1ba eliminates the part of the cosupport of codimension one in \(N\). The transform half of
Lemma~\ref{bcalc:codim-one} holds inside \(N\) with the same proof, read on \(N\) and on the strict transform \(N'\).

\begin{lemma}[Blowing up a centre of codimension one in \(N\) inside the support along \(N\)]
\label{bcalc:codim-one-along}
\uses{bcalc:codim-one,bcalc:nested-flattening,bex:hadamard,hir:hadamard-bridge-along,hir:restricted-calculus,bcalc:hyperplane-ideal,mc:principal-power,bcalc:monomial-lift,bcalc:snc-chart,bcalc:strict-ideal,bcalc:exceptional-divisor,hir:geo-transform,hir:cosupport,hir:order-basic,mc:maximal-order-along}
\source{wlodarczyk05-hironaka,wlodarczyk09-analytic-resolution:page-0018,wlodarczyk09-analytic-resolution:page-0019,wlodarczyk09-analytic-resolution:page-0027}
\lean{Hironaka.CodimOneCentreAlong.isInSpanOnAt_subset,Hironaka.CodimOneCentreAlong.VanishesToOrderOn.congr_on,Hironaka.CodimOneCentreAlong.apply_ne_zero_of_not_vanishesToOrderOn_on,Hironaka.CodimOneCentreAlong.isInSpanOnAt_flatCoord_of_mem_vanishingIdeal,Hironaka.CodimOneCentreAlong.notMem_cosupp_geoTransform_of_codim_one_along}\upstreamlean{TranscendentalNakamaye.Pending.CodimOneCentreAlong.cosupp_geoTransform_inter_preimage_eq_empty_of_forall_along,TranscendentalNakamaye.Pending.CodimOneCentreAlong.cosupp_geoTransform_inter_preimage_eq_empty_along}
\mixedleanok
Let \(\underline{\mathcal I}=(M,N,E,\mathfrak a,d)\) have \(N\) a closed submanifold, \(\mathfrak a\) coherent and \(d\ge1\), and let the centre \(C\) of the
blowup \(\pi\) lie in \(\operatorname{cosupp}\underline{\mathcal I}\). Let \(z\in\pi^{-1}(C)\), \(x=\pi(z)\), and suppose that at \(x\) the centre is regular of
dimension \(k\) and \(N\) is regular of dimension \(k+1\), and that \(\operatorname{ord}_x(\mathfrak a|_N)\le d\). Then \(z\notin\operatorname{cosupp}\underline{\mathcal I}'\).
Consequently, if \(\underline{\mathcal I}\) is of maximal order along \(N\) (Definition~\ref{mc:maximal-order-along}: \(\operatorname{ord}(\mathfrak a|_N)\le d\) everywhere)
and \(C\) has codimension one in \(N\) at each of its points, the cosupport of the transform does not meet \(\pi^{-1}(C)\). For \(N=M\) this is
Lemma~\ref{bcalc:codim-one}.
\end{lemma}
\begin{proof}
\mixedleanok
Take nested flattening coordinates \(\Phi\) at \(x\) (Lemma~\ref{bcalc:nested-flattening}): \(C=\mathbf C^k\times0\), \(N=\mathbf C^{k+1}\times0\), and put
\(u=\Phi_k\), \(H=\{u=0\}\); \(\mathcal I_H=(u)\) near \(x\) (Lemma~\ref{bcalc:hyperplane-ideal}). At a point \(y\) of \(N\) near \(x\) a germ \(g\in\mathcal I_C\) is
\(\sum_{i\ge k}b_i\Phi_i\) (Hadamard's lemma in flattening coordinates, Lemma~\ref{bex:hadamard}), and \(\Phi_i=0\) on \(N\) for \(i\ge k+1\); so on \(N\) near \(y\),
\(g=b_ku\): \(\operatorname{Res}_N\mathcal I_C\subset\operatorname{Res}_N\mathcal I_H\) near \(x\), hence \(\operatorname{Res}_N\mathcal I_C^d\subset\operatorname{Res}_N\mathcal I_H^d\)
(restricted inclusions pass to powers, Lemma~\ref{hir:restricted-calculus}), and \(\mathcal I_H^d=(u^d)\) (Lemma~\ref{mc:principal-power}). As
\(C\subset\operatorname{cosupp}\underline{\mathcal I}\), \(\operatorname{Res}_N\mathfrak a\subset\operatorname{Res}_N\mathcal I_C^d\) (Lemma~\ref{hir:hadamard-bridge-along}).
Since \(\operatorname{ord}_x(\mathfrak a|_N)\le d\), some \(f\in\mathfrak a_x\) does not vanish to order \(d+1\) along \(N\) at \(x\); by the above \(f=c\,u^d\) on \(N\)
near \(x\), and \(c(x)\ne0\): otherwise \(c\,u^d\) vanishes to order \(d+1\) at \(x\) (Lemma~\ref{hir:order-basic}) and agrees with \(f\) on \(N\).
Suppose \(z\in\operatorname{cosupp}\underline{\mathcal I}'\); then \(z\in N'\). Take a monomial chart \(j\ge k\), \(W\) around \(z\) with \(\sigma_j\circ W=\Phi\circ\pi\)
(Lemma~\ref{bcalc:monomial-lift}). As \(z\in N'\), \(j<k+1\) (Lemma~\ref{bcalc:snc-chart}), so \(j=k\) and \(u\circ\pi=W_k\), which generates \(\mathcal O(-D)\) near
\(z\) (Lemmas~\ref{bcalc:strict-ideal} and~\ref{bcalc:exceptional-divisor}); so \(\mathcal O(-dD)=(W_k^d)\) near \(z\). On \(N'\subset\pi^{-1}(N)\) near \(z\),
\(f\circ\pi=(c\circ\pi)\,W_k^d\); hence for every \(l=b\,W_k^d\in\mathcal O(-dD)_z\), \(l\,(c\circ\pi)=b\,(f\circ\pi)\) on \(N'\) near \(z\), that is,
\(c\circ\pi\) lies in \(\mathfrak a'=(\operatorname{Res}_{N'}\pi^*\mathfrak a:\mathcal O(-dD))\) at \(z\) (Definition~\ref{hir:geo-transform}). It does not vanish at
\(z\in N'\), so \(\mathfrak a'\) does not vanish to order \(d\ge1\) along \(N'\) at \(z\) (Definition~\ref{hir:cosupport}): a contradiction. The consequence is the
pointwise statement at every \(z\in\pi^{-1}(C)\).
\end{proof}

\begin{lemma}[Near a point of local dimension \(\dim N-1\) the support along \(N\) is a hypersurface of maximal contact]
\label{bcalc:codim-one-local-along}
\uses{bcalc:codim-one-local,bcalc:analytic-fill,bcalc:cosupport-analytic,hir:giraud-along,mc:derivative-along-coherent,ana:analytic-preimage,fpgeom:complex-analytic-subvariety,mc:tangent-direction-along,mc:tangent-direction-along-basic,mc:tangent-direction-along-exists,mc:contact-hypersurface-along,mc:maximal-order-along,lag:regular-point-basics,rem:dimension}
\source{wlodarczyk05-hironaka}
\lean{Hironaka.CodimOneSupportAlong.isComplexAnalyticSubset_cosupport_along,Hironaka.CodimOneSupportAlong.eventually_mem_iff_of_locDimGE_along,Hironaka.CodimOneSupportAlong.exists_isTangentDirectionAlong_of_isMaximalOrderAlong,Hironaka.CodimOneSupportAlong.regularAtDim_of_locDimGE_along,Hironaka.CodimOneSupportAlong.apply_eq_zero_of_isInSpanOnAt,Hironaka.CodimOneSupportAlong.eventually_mem_iff_of_locDimGE_along_of_isInSpanOnAt}
\leanok
Let \(N\) be a closed submanifold, \(\mathfrak a\) coherent and \(\mu\ge1\). Then \(S=\operatorname{cosupp}(\mathfrak a|_N,\mu)\) is analytic. Let \(A\subset S\) be
analytic, \(x\in M\), \(N\) regular of dimension \(k+1\) at \(x\), and \(\dim_xA\ge k\). If \(u\) is a tangent direction of \((\mathfrak a|_N,\mu)\) along \(N\) on an open
set \(U\ni x\) (Definition~\ref{mc:tangent-direction-along}), then \(A=N\cap V(u)\) near \(x\). If moreover \((\mathfrak a|_N,\mu)\) is of maximal order along \(N\)
(Definition~\ref{mc:maximal-order-along}), such a \(u\) exists near every point of \(S\), and \(A\) is regular of dimension \(k\) at \(x\). The first conclusion holds for
any \(w\) holomorphic on an open \(O\ni x\) that agrees on \(N\), near each point of \(N\cap O\), with a combination of germs of \(\mathcal D_N^{\mu-1}\mathfrak a\) and has
multiplicity one along \(N\) at its zeros on \(N\cap O\) (such \(w\) are the tangent directions of the transforms along a tower, which are controlled only on
the strict transforms of \(N\)). For \(N=M\) this is Lemma~\ref{bcalc:codim-one-local}.
\end{lemma}
\begin{proof}
\leanok
\(S=N\cap V(\mathcal D_N^{\mu-1}\mathfrak a)\) (Lemma~\ref{hir:giraud-along}); \(N\) is analytic (Lemma~\ref{ana:analytic-preimage}), the zero locus of the coherent ideal
\(\mathcal D_N^{\mu-1}\mathfrak a\) (Lemma~\ref{mc:derivative-along-coherent}) is analytic (Lemma~\ref{bcalc:cosupport-analytic}), and so is the intersection
(Definition~\ref{fpgeom:complex-analytic-subvariety}). Let \(A\), \(x\), \(u\) be as stated. \(A\) is closed and its regular points accumulate at \(x\), so \(x\in A\);
\(S\cap U\subset V(u)\) (Lemma~\ref{mc:tangent-direction-along-basic}), so \(u(x)=0\), and \(H=N\cap V(u)\cap U\) is regular at \(x\) of dimension one less than \(N\)
(Lemma~\ref{mc:contact-hypersurface-along}), that is of dimension \(k\) (the dimension of a regular point is unique, Lemma~\ref{lag:regular-point-basics}).
Near \(x\), \(A\subset S\cap U\subset H\); by the filling lemma (Lemma~\ref{bcalc:analytic-fill}), \(A=H\) near \(x\), i.e.\ \(A=N\cap V(u)\) near \(x\). For \(w\) as in the last sentence of the statement the same
argument applies: \(w\) vanishes at the points \(y\in S\cap O\), being on \(N\) near \(y\) a combination of germs of \(\mathcal D_N^{\mu-1}\mathfrak a\), which vanish at \(y\);
and \(N\cap V(w)\cap O\) is regular of dimension \(k\) at \(x\) by Lemma~\ref{mc:contact-hypersurface-along}, which uses only the holomorphy of \(w\) and its
multiplicity one along \(N\) at \(x\). Under maximal order
along \(N\), \(\operatorname{ord}_x(\mathfrak a|_N)\le\mu\) at every \(x\in S\), so a tangent direction along \(N\) exists on a neighbourhood of \(x\)
(Lemma~\ref{mc:tangent-direction-along-exists}), and \(A\) is regular of dimension \(k\) at \(x\) because \(H\) is (Remark~\ref{rem:dimension}: regularity depends only
on the germ).
\end{proof}

\begin{definition}[The codimension-one part of the support along \(N\)]
\label{bcalc:codim-one-part-along-def}
\uses{hir:cosupport,bcalc:codim-one-part-def}
\lean{Hironaka.CodimOneSupportAlong.codimOnePartAlong}
\leanok
Let \(N\subset M\) be a submanifold. The \emph{codimension-one part along \(N\)} of \(S=\operatorname{cosupp}(\mathfrak a|_N,\mu)\) is the set \(R_N(\mathfrak a,\mu)\) of points
\(x\) at which \(N\) is regular of some dimension \(k+1\) and \(S\) is regular of dimension \(k\). For \(N=M\) it is the codimension-one part
(Definition~\ref{bcalc:codim-one-part-def}).
\end{definition}

\begin{theorem}[The codimension-one part of the support along \(N\) is a smooth isolated hypersurface of \(N\)]
\label{bcalc:codim-one-part-along}
\uses{bcalc:codim-one-part-along-def,bcalc:codim-one-local-along,bcalc:codim-one-part,hir:giraud-along,bcalc:nested-flattening,lag:regular-point-basics,lag:regular-locus-dense,mc:tangent-direction-along-basic,mc:tangent-direction-along-exists,mc:contact-hypersurface-along,rem:dimension}
\source{wlodarczyk05-hironaka}
\lean{Hironaka.CodimOneSupportAlong.codimOnePartAlong_subset,Hironaka.CodimOneSupportAlong.mem_codimOnePartAlong_iff,Hironaka.CodimOneSupportAlong.eventually_mem_codimOnePartAlong_iff,Hironaka.CodimOneSupportAlong.isClosed_codimOnePartAlong,Hironaka.CodimOneSupportAlong.regularAtDim_codimOnePartAlong,Hironaka.CodimOneSupportAlong.isClosedSubmanifold_codimOnePartAlong,Hironaka.CodimOneSupportAlong.exists_isOpen_cosupport_inter_eq_codimOnePartAlong}\upstreamlean{TranscendentalNakamaye.Pending.CodimOneSupportAlong.eventually_mem_codimOnePartAlong_iff_eq_zero}\lean{Hironaka.CodimOneSupportAlong.eventually_mem_codimOnePartAlong_iff_eq_zero_of_isInSpanOnAt,Hironaka.CodimOneSupportAlong.isClosed_cosupport_diff_codimOnePartAlong,Hironaka.CodimOneSupportAlong.lt_of_regularAtDim_cosupport_diff_codimOnePartAlong}
\mixedleanok
Let \(N\) be a closed submanifold and \((\mathfrak a|_N,\mu)\) coherent, of maximal order along \(N\), \(\mu\ge1\), with support \(S\) and codimension-one part \(R=R_N(\mathfrak a,\mu)\)
along \(N\). Then:
\begin{enumerate}
\item \(R\) is the set of points \(x\) at which \(N\) has some dimension \(k+1\) and \(\dim_xS\ge k\);
\item \(R\) is closed, open in \(S\), and at each of its points regular of dimension one less than \(N\); in particular \(R\) is a closed submanifold;
\item near each \(x\in R\), \(R=S=N\cap V(u)\) for every tangent direction \(u\) along \(N\) on a neighbourhood of \(x\), and for every \(w\) controlled only on
\(N\) as in Lemma~\ref{bcalc:codim-one-local-along};
\item \(S\setminus R\) is closed, and at a point where \(N\) has dimension \(k+1\) every regular point of \(S\setminus R\) has dimension \(<k\).
\end{enumerate}
For \(N=M\) this is Theorem~\ref{bcalc:codim-one-part}.
\end{theorem}
\begin{proof}
\mixedleanok
(1) A regular point of dimension \(k\) has local dimension \(\ge k\) (Remark~\ref{rem:dimension}); conversely \(S\) is analytic and Lemma~\ref{bcalc:codim-one-local-along} with
\(A=S\) makes \(S\) regular of dimension \(k\) at a point of local dimension \(\ge k\). (2) Near \(x\in R\), with \(N\) of dimension \(k+1\) and \(S\) of dimension \(k\) at \(x\),
every point of \(N\) is regular of dimension \(k+1\) and every point of \(S\) regular of dimension \(k\) (Theorem~\ref{lag:regular-locus-dense}); as \(S\subset N\), \(R=S\)
near \(x\). Hence \(R\) is open in \(S\) and regular of dimension \(k\) at \(x\). Let \(x\) be a limit of points of \(R\); \(x\in S\) since \(S\) is closed
(Lemma~\ref{hir:giraud-along}), \(N\) is regular of some dimension \(k_x\) at \(x\), hence of dimension \(k_x\) at the points of \(N\) near \(x\); at the points of \(R\) near
\(x\) it also has dimension \(k+1\) with \(S\) of dimension \(k\) there, so \(k_x=k+1\) (Lemma~\ref{lag:regular-point-basics}) and \(\dim_xS\ge k\): \(x\in R\) by (1).
(3) \(R=S\) near \(x\), and \(S=N\cap V(u)\) near \(x\) by Lemma~\ref{bcalc:codim-one-local-along}. (4) If \(O\) is open with \(S\cap O=R\), then \(S\setminus R=S\setminus O\) is
closed. Let \(S\setminus R\) be regular of dimension \(j\) at \(y\), where \(N\) has dimension \(k+1\). As \(R\) is closed, \(S\) is regular of dimension \(j\) at \(y\). A tangent
direction \(u\) along \(N\) on \(U\ni y\) exists (Lemma~\ref{mc:tangent-direction-along-exists}); \(S\cap U\subset N\cap V(u)\cap U\) (Lemma~\ref{mc:tangent-direction-along-basic}),
which is regular of dimension \(k\) at \(y\) (Lemma~\ref{mc:contact-hypersurface-along}); a submanifold of a submanifold has smaller dimension
(Lemma~\ref{bcalc:nested-flattening}), so \(j\le k\), and \(j\ne k\) because \(y\notin R\).
\end{proof}

\begin{corollary}[Blowing up the codimension-one part of the support along \(N\)]
\label{bcalc:codim-one-step-along}
\uses{bcalc:codim-one-along,bcalc:codim-one-part-along,bcalc:transform-off-centre,bcalc:codim-one-step,bcalc:codim-one-step-rest}
\source{wlodarczyk05-hironaka,wlodarczyk09-analytic-resolution:page-0027}
\upstreamlean{TranscendentalNakamaye.Pending.CodimOneSupportAlong.cosupp_geoTransform_inter_preimage_codimOnePartAlong_eq_empty,TranscendentalNakamaye.Pending.CodimOneSupportAlong.cosupp_geoTransform_subset_preimage_diff_codimOnePartAlong}
\upstreamleanok
Let \(\underline{\mathcal I}=(M,N,E,\mathfrak a,d)\) have \(N\) a closed submanifold, \(\mathfrak a\) coherent, \(d\ge1\), and be of maximal order along \(N\); let \(\pi\) be the
blowup along the codimension-one part \(R\) of its support along \(N\). Then the cosupport of the transform does not meet the exceptional divisor
\(\pi^{-1}(R)\), and it lies over \(\operatorname{cosupp}\underline{\mathcal I}\setminus R\), a closed set whose regular points have dimension at most \(\dim N-2\). This is
Step~1ba of Włodarczyk's algorithm inside \(N\). For \(N=M\) it is Corollaries~\ref{bcalc:codim-one-step} and~\ref{bcalc:codim-one-step-rest}.
\end{corollary}
\begin{proof}
\upstreamleanok
\(R\) lies in the support and has codimension one in \(N\) at each of its points (Theorem~\ref{bcalc:codim-one-part-along}(2)); Lemma~\ref{bcalc:codim-one-along} excludes the
points over \(R\). At the other points the transform has the same cosupport (Theorem~\ref{bcalc:transform-off-centre}:
\(\operatorname{cosupp}\underline{\mathcal I}'\setminus\pi^{-1}(R)=\pi^{-1}(\operatorname{cosupp}\underline{\mathcal I}\setminus R)\)); the dimension bound is
Theorem~\ref{bcalc:codim-one-part-along}(4).
\end{proof}

\paragraph{Members of the exceptional divisor under a blowup.} Let \(E\) be a simple normal crossings divisor (Definition~\ref{hir:snc}) such that
the centre \(C\) has normal crossings with \(E\). The ideals of the members of \(E\) transform monomially; this is the chart computation behind the
monomial marked ideals of the algorithm.

\begin{lemma}[Ideals of coordinate hyperplanes and principal ideals]
\label{bcalc:hyperplane-ideal}
\uses{bcalc:hypersurface,ana:local-holomorphic-inverse}
\lean{Hironaka.MemberPullback.injective_fderiv_toFun,Hironaka.MemberPullback.exists_chartDeriv_flatCoord_ne_zero,Hironaka.MemberPullback.vanishingIdeal_mem_iff_flatCoord,Hironaka.MemberPullback.exists_chartDeriv_coord_ne_zero,Hironaka.MemberPullback.pullback_mem_iff_of_generator,Hironaka.MemberPullback.mul_mem_iff_of_generators,Hironaka.MemberPullback.isInSpanAt_eq_congr,Hironaka.MemberPullback.isInSpanAt_one_iff,Hironaka.MemberPullback.isInSpanAt_range_one_iff}
\leanok
(1) If \(\Phi\) is a holomorphic coordinate system on an open set \(U\) and \(A\cap U=\{\Phi_t=0\}\), then at every point of \(U\) the ideal of \(A\) is
generated by \(\Phi_t\); the same holds for a coordinate \(W_t\) of a holomorphic map \(W\) with injective differential. (2) If an ideal \(\mathfrak a\) is
generated by \(h\) at \(f(z)\), then \(f^*\mathfrak a\) is generated by \(h\circ f\) at \(z\); if \(\mathfrak a=(p)\) and \(\mathfrak b=(q)\) at \(x\), then
\(\mathfrak a\mathfrak b=(pq)\) at \(x\).
\end{lemma}
\begin{proof}
\leanok
(1) A biholomorphism has invertible differential (chain rule for \(\Phi^{-1}\circ\Phi=\mathrm{id}\)), so the row \(d\Phi_t\) is nonzero at every point; the
same holds for \(W\), an injective endomorphism of \(\mathbf C^n\) being surjective. Hence \(\Phi_t\) (resp.\ \(W_t\)) has multiplicity one on \(U\), and
the ideal of its zero set is generated by it at every point of \(U\) (Lemma~\ref{bcalc:hypersurface}); \(A\) agrees with that zero set on \(U\).
(2) A germ of \(f^*\mathfrak a\) is a holomorphic combination of the \(a\circ f\), \(a\in\mathfrak a_{f(z)}\), and \(a=c\,h\) near \(f(z)\) gives
\(a\circ f=(c\circ f)(h\circ f)\); conversely \(h\circ f\) is one of the generators. Likewise \(ab=(cp)(dq)=(cd)\,pq\), and \(pq\) is a product of
generators.
\end{proof}

\begin{theorem}[The pullback of the ideal of a member]
\label{bcalc:member-pullback}
\uses{bcalc:hyperplane-ideal,bcalc:snc-simultaneous,bcalc:monomial-lift,bcalc:snc-chart,bcalc:strict-ideal,bcalc:exceptional-divisor,hir:snc}
\source{wlodarczyk05-hironaka,bierstone-milman08-functoriality}
\lean{Hironaka.MemberPullback.exists_chart_generators}\upstreamlean{TranscendentalNakamaye.Pending.MemberPullback.mem_pullback_member_iff_of_subset}\lean{Hironaka.MemberPullback.mem_pullback_member_iff_of_not_subset}
\mixedleanok
Let \(\pi\) be the blowup along \(C\), \(D=\pi^{-1}(C)\), \(E\) a simple normal crossings divisor with which \(C\) has normal crossings, \(H\) a member of
\(E\) and \(H'=\overline{\pi^{-1}(H\setminus C)}\). At a point \(z\) of the blowup:
\begin{enumerate}
\item if \(C\subset H\) near \(\pi(z)\), then \(\pi^*\mathcal I_H=\mathcal O(-D)\cdot\mathcal I_{H'}\) at \(z\);
\item otherwise \(\pi^*\mathcal I_H=\mathcal I_{H'}\) at \(z\).
\end{enumerate}
Statement (2) uses that \(H\) is a hypersurface: for \(A=\{z_1=z_3=0\}\) and \(C=\{z_1=z_2=0\}\), in the chart \(z_1=w_1\) the pullback of
\(\mathcal I_A\) is \((w_1,w_3)\), while the strict transform of \(A\) misses that chart.
\end{theorem}
\begin{proof}
\mixedleanok
If \(\pi(z)\notin C\), \(\pi\) is biholomorphic near \(z\), \(H'=\pi^{-1}(H)\) near \(z\) and \(\mathcal O(-D)\) contains \(1\) at \(z\), so both sides agree
(Theorem~\ref{bcalc:strict-ideal}); in case (2) this does not occur, since \(C\) is closed and so \(C\subset H\) near points off \(C\). If \(\pi(z)\in C\setminus H\),
both ideals are the unit ideal at \(z\) (\(H\) is closed and \(H'\subset\pi^{-1}(H)\)), and \(C\subset H\) near \(\pi(z)\) fails. Let \(x=\pi(z)\in C\cap H\).
Take coordinates \(\Phi\) at \(x\) in which \(C=\mathbf C^k\times0\) and \(H=\{w_t=0\}\) (Lemma~\ref{bcalc:snc-simultaneous} with \(N=X\)). Then \(C\subset H\) near \(x\)
if and only if \(t\ge k\): if \(t\ge k\) the coordinate \(w_t\) vanishes on \(\mathbf C^k\times0\); if \(t<k\), the points \(\Phi^{-1}(\Phi(x)+\varepsilon e_t)\),
\(\varepsilon\ne0\) small, lie in \(C\) but not in \(H\). Let \(j\ge k\) and \(W\) be a monomial chart around \(z\) with \(\sigma_j\circ W=\Phi\circ\pi\)
(Lemma~\ref{bcalc:monomial-lift}); \(W\) has injective differential (Lemma~\ref{bcalc:snc-chart}). By Lemma~\ref{bcalc:hyperplane-ideal}, \(\mathcal I_H=(\Phi_t)\)
near \(x\), so \(\pi^*\mathcal I_H\) is generated by \(\Phi_t\circ\pi=\sigma_j(W)_t\), which is \(W_t\) if \(t<k\), \(W_tW_j\) if \(t\ge k\), \(t\ne j\), and \(W_j\)
if \(t=j\); \(\mathcal O(-D)\) is generated by \(W_j\) (Theorem~\ref{bcalc:strict-ideal}). On the chart, \(H'=\{W_t=0\}\) if \(t\ne j\) and \(H'\) misses the chart if
\(t=j\) (Lemma~\ref{bcalc:snc-chart}); so \(\mathcal I_{H'}=(W_t)\), resp.\ the unit ideal, at \(z\) (Lemma~\ref{bcalc:hyperplane-ideal}). In case (1), \(t\ge k\)
and \(\pi^*\mathcal I_H=(W_jW_t)=(W_j)(W_t)\), resp.\ \((W_j)=(W_j)\cdot\mathcal O\); in case (2), \(t<k\le j\) and \(\pi^*\mathcal I_H=(W_t)=\mathcal I_{H'}\).
\end{proof}

\begin{lemma}[The monomial chart over a point of the centre]
\label{bcalc:monomial-chart}
\uses{bcalc:snc-simultaneous,bcalc:member-pullback,bcalc:monomial-lift,bcalc:snc-chart,bcalc:hyperplane-ideal,bcalc:strict-ideal,bcalc:exceptional-divisor,hir:monomial-factorization,ana:blowup-exceptional-image,hir:snc}
\source{wlodarczyk05-hironaka,bierstone-milman08-functoriality}
\lean{Hironaka.MonomialBlowup.exists_unit_of_isInSpanAt,Hironaka.MonomialBlowup.eventually_subset_iff_le,Hironaka.MonomialBlowup.exists_monomial_chart}\upstreamlean{TranscendentalNakamaye.Pending.MonomialBlowup.frequently_ne_zero_of_generator}
\mixedleanok
Let \(\pi\) be the blowup along \(C\), \(D=\pi^{-1}(C)\), \(E=(H_1,\dots,H_s)\) a simple normal crossings divisor with which \(C\) has only
normal crossings, \(H_i'=\overline{\pi^{-1}(H_i\setminus C)}\), and \(z\) a point with \(x=\pi(z)\in C\). There are an open \(U\ni z\), a holomorphic
map \(W\colon U\to\mathbf C^n\) with injective differential at \(z\), an index \(j\) and indices \(s_i\) (\(x\in H_i\)), pairwise distinct, such that:
\begin{enumerate}
\item \(W_j(z)=0\) and \(\mathcal O(-D)=(W_j)\) at every point of \(U\); \(W_j\) vanishes identically near no point of \(U\);
\item if \(z\in H_i'\), then \(x\in H_i\), \(s_i\ne j\), \(W_{s_i}(z)=0\) and \(\mathcal I_{H_i'}=(W_{s_i})\) at every point of \(U\);
\item if \(u_i\) generates \(\mathcal I_{H_i}\) at \(x\) (\(x\in H_i\)) and \(a_i\in\mathbb N\), then near \(z\)
\[
 \Bigl(\prod_{x\in H_i}u_i^{a_i}\Bigr)\circ\pi=c\,W_j^{\,e}\prod_{z\in H_i'}W_{s_i}^{a_i},\qquad c(z)\ne0,\qquad
 e=\sum_{x\in H_i,\ C\subset H_i\text{ near }x}a_i .
\]
\end{enumerate}
\end{lemma}
\begin{proof}
\mixedleanok
Take coordinates \(\Phi\) at \(x\) in which \(C=\mathbf C^k\times0\) and \(H_i=\{w_{s_i}=0\}\) for the members through \(x\), the \(s_i\) distinct
(Lemma~\ref{bcalc:snc-simultaneous} with \(N=X\)); as in the proof of Theorem~\ref{bcalc:member-pullback}, \(C\subset H_i\) near \(x\) if and only if
\(s_i\ge k\). Let \(j\ge k\) and \(W\) be a monomial chart around \(z\) with \(\sigma_j\circ W=\Phi\circ\pi\) (Lemma~\ref{bcalc:monomial-lift}); its
differential is injective and \(H_i'=\{W_{s_i}=0\}\) on the chart if \(s_i\ne j\), while \(H_i'\) misses the chart if \(s_i=j\) (Lemma~\ref{bcalc:snc-chart}).
Hence \(\mathcal O(-D)=(W_j)\) and \(\mathcal I_{H_i'}=(W_{s_i})\) on the chart (Theorem~\ref{bcalc:strict-ideal}, Lemma~\ref{bcalc:hyperplane-ideal}), and
\(z\in H_i'\) forces \(x\in H_i\) (\(H_i\) is closed), \(s_i\ne j\) and \(W_{s_i}(z)=0\). The zero set of \(W_j\) in \(U\) is \(D\cap U\), which has empty
interior (Lemma~\ref{ana:blowup-exceptional-image}). (3): \(\Phi_{s_i}\circ\pi=\sigma_j(W)_{s_i}\) is \(W_{s_i}W_j\) if \(k\le s_i\ne j\), \(W_j\) if \(s_i=j\) and
\(W_{s_i}\) if \(s_i<k\). Two generators \(u,\varphi\) of one principal ideal, \(\varphi\) vanishing identically near no point, differ by a unit:
\(u=a\varphi\) and \(\varphi=bu\) give \((1-ba)\varphi=0\), so \(ba=1\) near the point (Lemma~\ref{hir:monomial-factorization}). So \(u_i=c_i\Phi_{s_i}\)
near \(x\) with \(c_i(x)\ne0\) (Lemma~\ref{bcalc:hyperplane-ideal}), and the factors \(W_{s_i}\) with \(W_{s_i}(z)\ne0\), i.e.\ \(z\notin H_i'\), \(s_i\ne j\), are units
at \(z\). The exponent of \(W_j\) is the sum of the \(a_i\) with \(s_i\ge k\), that is, with \(C\subset H_i\) near \(x\).
\end{proof}

\begin{theorem}[Monomial ideals under the blowup of a centre with normal crossings]
\label{bcalc:monomial-transform}
\uses{bcalc:monomial-chart,bcalc:controlled-transform-def,bcalc:controlled-transform,bcalc:exceptional-divisor,hir:monomial-factorization,hir:geo-transform-ambient,hir:marked-ideal}
\source{wlodarczyk05-hironaka,bierstone-milman08-functoriality}
\upstreamlean{TranscendentalNakamaye.Pending.MonomialBlowup.mem_controlledTransform_iff_of_eq,TranscendentalNakamaye.Pending.MonomialBlowup.mem_geoTransform_iff_of_eq,TranscendentalNakamaye.Pending.MonomialBlowup.exists_monomial_chart_transform}
\upstreamleanok
Let \(\pi\), \(D\), \(z\) and \(\mathfrak a\) be as follows.
\begin{enumerate}
\item If \(\mathfrak a=(m)\) at \(\pi(z)\) and \(m\circ\pi=c\,h^{e}q\) near \(z\), with \(\mathcal O(-D)=(h)\) at \(z\), \(h\) vanishing identically near no point,
\(c(z)\ne0\) and \(\mu\le e\), then \(\sigma^c(\mathfrak a,\mu)=(h^{e-\mu}q)\) at \(z\); for a marked ideal \((X,X,E,\mathfrak a,d)\) with \(\mathfrak a\) coherent and
\(d\le e\), the transformed ideal is \((h^{e-d}q)\) at \(z\).
\item (Ambient case of Lemma~\ref{hir:monomial-blowup}, generator form.) In the situation of Lemma~\ref{bcalc:monomial-chart}, if
\(\mathfrak a=(\prod_{x\in H_i}u_i^{a_i})\) at \(x=\pi(z)\) and \(\mu\le e=\sum_{x\in H_i,\ C\subset H_i\text{ near }x}a_i\), then
\[
 \sigma^c(\mathfrak a,\mu)=\Bigl(W_j^{\,e-\mu}\prod_{z\in H_i'}W_{s_i}^{a_i}\Bigr)\quad\text{at }z :
\]
the transform is monomial, with exponent \(e-\mu\) on \(D=\{W_j=0\}\) and \(a_i\) on the strict transforms \(H_i'=\{W_{s_i}=0\}\) through \(z\).
\end{enumerate}
\end{theorem}
\begin{proof}
\upstreamleanok
(1) \(\pi^*\mathfrak a=(m\circ\pi)=(h^{e}q)=(h^{\mu})\cdot(h^{e-\mu}q)\) at \(z\), and \(\mathcal O(-\mu D)=(h^{\mu})\) (Lemma~\ref{bcalc:exceptional-divisor});
since \(h^{\mu}\) is a non-zero-divisor (Lemma~\ref{hir:monomial-factorization}), the germs \(g\) with \(h^{\mu}g\in\pi^*\mathfrak a\) are the multiples of
\(h^{e-\mu}q\) (Definition~\ref{bcalc:controlled-transform-def}). For a marked ideal with \(N=X\) and a coherent ideal, the transformed ideal is the controlled
transform (Lemma~\ref{hir:geo-transform-ambient}). (2) is (1) with \(h=W_j\), \(q=\prod_{z\in H_i'}W_{s_i}^{a_i}\) and the factorization of
Lemma~\ref{bcalc:monomial-chart}(3).
\end{proof}

\begin{corollary}[The transform of a monomial marked ideal in adapted coordinates]
\label{bcalc:monomial-transform-coords}
\uses{bcalc:monomial-transform,bcalc:monomial-chart,bcalc:snc-chart,bcalc:snc-transform,bcalc:hyperplane-ideal,bcalc:exceptional-divisor,hir:geo-transform}
\source{wlodarczyk05-hironaka,bierstone-milman08-functoriality}
\upstreamlean{TranscendentalNakamaye.Pending.MonomialBlowup.exists_snc_chart_geoTransform,TranscendentalNakamaye.Pending.MonomialBlowup.exists_snc_chart_geoTransform_of_flat}\lean{Hironaka.MonomialBlowup.mem_centre_iff_of_generator,Hironaka.MonomialBlowup.mem_iff_of_vanishingIdeal_eq,Hironaka.MonomialBlowup.prod_map_append_singleton}
\mixedleanok
Let \((X,X,E,\mathfrak a,d)\) be a marked ideal with \(\mathfrak a\) coherent, \(E=(H_1,\dots,H_s)\) a simple normal crossings divisor, \(C\) a centre with only
normal crossings with \(E\), \(\pi\) the blowup along \(C\) and \(z\) a point with \(x=\pi(z)\in C\). Suppose \(\mathfrak a=(\prod_{x\in H_i}u_i^{a_i})\) at \(x\), with
\(u_i\) generating \(\mathcal I_{H_i}\) at \(x\) (for instance the coordinates \(\Phi_{\sigma_0(i)}\) of a flattening \(\Phi\) at \(x\) in which the members through
\(x\) are the hyperplanes \(\{\Phi_{\sigma_0(i)}=0\}\)), and \(d\le e=\sum_{x\in H_i,\ C\subset H_i\text{ near }x}a_i\). Let \(E'=(H_1',\dots,H_s',D)\). Then there
are coordinates \(\Theta\) centred at \(z\), indices \(\sigma(i)\) for the members of \(E'\) through \(z\), pairwise distinct, and exponents \(a'_i=a_i\)
(\(i\le s\)), \(a'_D=e-d\), such that every member \(H\) of \(E'\) through \(z\) is the hyperplane \(\{\Theta_{\sigma(H)}=0\}\) near \(z\) and the transformed ideal is
\[
 \mathfrak a'=\Bigl(\prod_{z\in H\in E'}\Theta_{\sigma(H)}^{\,a'_H}\Bigr)\quad\text{at }z .
\]
\end{corollary}
\begin{proof}
\mixedleanok
Take the monomial chart \(W\) of Lemma~\ref{bcalc:monomial-chart} and \(\Theta=W-W(z)\), a coordinate system centred at \(z\) (Lemma~\ref{bcalc:snc-chart});
\(\sigma(H_i')=s_i\), \(\sigma(D)=j\). The coordinates \(W_j\), \(W_{s_i}\) (\(z\in H_i'\)) vanish at \(z\), so they agree with \(\Theta_j\), \(\Theta_{s_i}\). Where
\(\mathcal O(-D)=(W_j)\), the zero set of \(W_j\) is the support \(\pi^{-1}(C)\) of \(D\) (Lemma~\ref{bcalc:exceptional-divisor}); where \(\mathcal I_{H'}=(W_{s_i})\) and
\(H'\) is closed, the zero set of \(W_{s_i}\) is \(H'\) (off \(H'\) the ideal contains \(1\)). The \(s_i\) are distinct and differ from \(j\)
(Lemma~\ref{bcalc:monomial-chart}(2)). The transformed ideal is \((W_j^{e-d}\prod_{z\in H_i'}W_{s_i}^{a_i})\) at \(z\) (Theorem~\ref{bcalc:monomial-transform}), and the
product over the members of \(E'\) through \(z\) splits into the strict transforms and the last member \(D\), which passes through \(z\) (the indexing of \(E'\),
Theorem~\ref{bcalc:snc-transform}; Definition~\ref{hir:geo-transform}). Flat coordinates generate the ideals of their hyperplanes
(Lemma~\ref{bcalc:hyperplane-ideal}), which gives the variant with \(u_i=\Phi_{\sigma_0(i)}\).
\end{proof}

\begin{theorem}[Monomial ideals under the blowup, along a submanifold]
\label{bcalc:monomial-transform-along}
\uses{bcalc:monomial-chart,bcalc:monomial-transform,hir:marked-operations-along-step,hir:geo-transform-wellformed,hir:geo-transform-coherent,hir:monomial-factorization,bcalc:strict-submanifold,bcalc:exceptional-divisor,hir:geo-transform,hir:marked-ideal}
\source{wlodarczyk05-hironaka,bierstone-milman08-functoriality}
\lean{Hironaka.MonomialBlowupAlong.isInSpanOnAt_geoTransform_iff_of_eq}\upstreamlean{TranscendentalNakamaye.Pending.MonomialBlowupAlong.exists_monomial_chart_transform_along}\lean{Hironaka.MonomialBlowupAlong.isInSpanOnAt_comp_of_mem_pullback,Hironaka.MonomialBlowupAlong.isInSpanOnAt_pullback_comp}
\mixedleanok
Let \((M,N,E,\mathfrak a,d)\) be a marked ideal with \(N\) a closed submanifold and \(\mathfrak a\) coherent, \(C\) an admissible centre, \(\pi\) the blowup along \(C\),
\(N'=\overline{\pi^{-1}(N\setminus C)}\), \(\mathfrak a'\) the transformed ideal, and \(z\in N'\) with \(\pi(z)\in C\). Write \(\operatorname{Res}_W\mathfrak b=\operatorname{Res}_W\mathfrak c\)
at a point when the same functions agree on \(W\) near the point with combinations of germs of \(\mathfrak b\), resp.\ \(\mathfrak c\).
\begin{enumerate}
\item If \(\operatorname{Res}_N\mathfrak a=\operatorname{Res}_N(m)\) at \(\pi(z)\) and \(m\circ\pi=c\,h^{e}q\) near \(z\), with \(\mathcal O(-D)=(h)\) at \(z\), \(c(z)\ne0\) and \(d\le e\),
then \(\operatorname{Res}_{N'}\mathfrak a'=\operatorname{Res}_{N'}(h^{e-d}q)\) at \(z\).
\item (Lemma~\ref{hir:monomial-blowup} along \(N'\), generator form.) In the monomial chart of Lemma~\ref{bcalc:monomial-chart}, if
\(\operatorname{Res}_N\mathfrak a=\operatorname{Res}_N(\prod_{x\in H_i}u_i^{a_i})\) at \(x=\pi(z)\), with \(u_i\) generating \(\mathcal I_{H_i}\) at \(x\), and
\(d\le e=\sum_{x\in H_i,\ C\subset H_i\text{ near }x}a_i\), then \(\operatorname{Res}_{N'}\mathfrak a'=\operatorname{Res}_{N'}(W_j^{\,e-d}\prod_{z\in H_i'}W_{s_i}^{a_i})\) at \(z\).
\end{enumerate}
No transversality of \(E\) to \(N\) enters: the statement concerns restrictions of ambient generators. Reading \(W_j|_{N'}\) and \(W_{s_i}|_{N'}\)
as coordinates of \(N'\) (the monomial form of \(E'|_{N'}\)) needs a chart adapted to \(N'\) as well, which is not part of this statement; that
form is Corollary~\ref{bcalc:monomial-transform-along-coords}, under transversality of \(E\) to \(N\).
\end{theorem}
\begin{proof}
\mixedleanok
Restricted generators and restricted memberships pull back from \(N\) to \(N'\), since \(\pi(N')\subset N\) (Lemma~\ref{bcalc:strict-submanifold}). By
Lemma~\ref{hir:marked-operations-along-step}, \(\operatorname{Res}_{N'}\pi^*\mathfrak a=\operatorname{Res}_{N'}\mathcal O(-dD)\mathfrak a'\); near \(z\), \(\mathcal O(-dD)=(h_0)\) with
\(h_0\) vanishing identically on no relatively open piece of \(N'\) (Lemma~\ref{hir:geo-transform-wellformed}), and \(h^d=\kappa h_0\) with \(\kappa(z)\ne0\), two
generators of one principal ideal differing by a unit (Lemma~\ref{bcalc:monomial-chart}; \(h^d\in\mathcal O(-dD)\) and the germs of \(\mathcal O(-dD)\) are multiples of
\(h^d\), Lemma~\ref{hir:monomial-factorization}). So \(m\circ\pi=h_0\cdot c\kappa h^{e-d}q\). (1) A germ \(t\) of \(\mathfrak a'\) gives \(h_0t\in\mathcal O(-dD)\mathfrak a'\), which on
\(N'\) is a multiple \(A\,(m\circ\pi)=h_0\cdot Ac\kappa h^{e-d}q\); cancelling \(h_0\) along \(N'\), \(t=Ac\kappa h^{e-d}q\) on \(N'\). Conversely \(m\circ\pi\) is on \(N'\) a
combination of germs of \(\mathcal O(-dD)\mathfrak a'\) (\(\mathfrak a'\) is coherent, Lemma~\ref{hir:geo-transform-coherent}), each \(h_0\) times a germ of \(\mathfrak a'\);
cancelling \(h_0\), \(c\kappa h^{e-d}q\) is on \(N'\) a combination of germs of \(\mathfrak a'\), and \(c\kappa\) is a unit at \(z\). (2) is (1) with \(h=W_j\),
\(q=\prod_{z\in H_i'}W_{s_i}^{a_i}\) and Lemma~\ref{bcalc:monomial-chart}(3).
\end{proof}

\begin{lemma}[The monomial chart adapted to a submanifold through the centre]
\label{bcalc:monomial-chart-adapted}
\uses{bcalc:monomial-chart,bcalc:snc-chart,bcalc:snc-transform,bcalc:monomial-lift,bcalc:strict-submanifold}
\source{wlodarczyk05-hironaka,bierstone-milman08-functoriality}
\lean{Hironaka.MonomialBlowupAlongChart.exists_monomial_chart_of_isAdaptedAt}
\leanok
Let \(\pi\) be the blowup along \(C\), \(C\subset N\), \(E=(H_1,\dots,H_s)\) a divisor with closed members, and \(z\) a point with \(x=\pi(z)\in C\).
Suppose that \(C\subset N\) and \(E\) admit adapted coordinates at \(x\) (Theorem~\ref{bcalc:snc-transform}): \(C=\mathbf C^{k_C}\times0\),
\(N=\mathbf C^{k_N}\times0\), and the members through \(x\) are the hyperplanes \(\{w_{s_i}=0\}\), \(s_i<k_N\), the \(s_i\) distinct. Then there are
\(U\), \(W\), \(j\) and the \(s_i\) with all the properties (1)--(3) of Lemma~\ref{bcalc:monomial-chart}, and moreover, with
\(N'=\overline{\pi^{-1}(N\setminus C)}\),
\[
 N'\cap U=\{y\in U:\ W_t(y)=0\text{ for all }t\ge k_N\}\ \text{ if }j<k_N,\qquad N'\cap U=\emptyset\ \text{ if }j\ge k_N .
\]
In particular \(j<k_N\) if \(z\in N'\).
\end{lemma}
\begin{proof}
\leanok
Run the proof of Lemma~\ref{bcalc:monomial-chart} with the given adapted coordinates \(\Phi\) at \(x\) in place of coordinates adapted to \(C\) and \(E\)
alone: it uses only that \(C=\mathbf C^{k_C}\times0\) and that the members through \(x\) are the hyperplanes \(\{w_{s_i}=0\}\) with distinct \(s_i\), and
that the members are closed. The monomial chart \(W\) has index \(j\ge k_C\) (Lemma~\ref{bcalc:monomial-lift}). Since \(N=L_T\) with
\(T=\{t:t\ge k_N\}\) in the coordinates \(\Phi\), Lemma~\ref{bcalc:snc-chart} gives, for \(y\in U\), \(y\in N'\iff j\notin T\) and \(W(y)\in L_T\), which
is the description of \(N'\cap U\); \(N'\) is a strict transform of a closed set (Lemma~\ref{bcalc:strict-submanifold}).
\end{proof}

\begin{corollary}[Monomial ideals under the blowup, along a submanifold, in coordinates adapted to the strict transform]
\label{bcalc:monomial-transform-along-coords}
\uses{bcalc:monomial-chart-adapted,bcalc:monomial-transform-along,bcalc:monomial-transform-coords,bcalc:snc-simultaneous,bcalc:snc-chart,bcalc:snc-transform,bcalc:hyperplane-ideal,bcalc:exceptional-divisor,hir:geo-transform,hir:marked-ideal,hir:snc}
\source{wlodarczyk05-hironaka,bierstone-milman08-functoriality}
\lean{Hironaka.MonomialBlowupAlongChart.exists_snc_chart_geoTransform_along,Hironaka.MonomialBlowupAlongChart.exists_snc_chart_geoTransform_along_of_flat,Hironaka.MonomialBlowupAlongChart.isInSpanOnAt_eq_congr}
\leanok
Let \((M,N,E,\mathfrak a,d)\) be a marked ideal with \(N\) a closed submanifold, \(\mathfrak a\) coherent, \(E=(H_1,\dots,H_s)\) a simple normal crossings
divisor transverse to \(N\), \(C\) an admissible centre, \(\pi\) the blowup along \(C\), \(N'=\overline{\pi^{-1}(N\setminus C)}\), \(\mathfrak a'\) the
transformed ideal, \(E'=(H_1',\dots,H_s',D)\), and \(z\in N'\) with \(x=\pi(z)\in C\). Suppose
\(\operatorname{Res}_N\mathfrak a=\operatorname{Res}_N(\prod_{x\in H_i}u_i^{a_i})\) at \(x\) (notation of Theorem~\ref{bcalc:monomial-transform-along}), with \(u_i\)
generating \(\mathcal I_{H_i}\) at \(x\) (for instance the coordinates \(\Phi_{\sigma_0(i)}\) of a flattening \(\Phi\) at \(x\) in which the members through \(x\)
are the hyperplanes \(\{\Phi_{\sigma_0(i)}=0\}\); \(\Phi\) need not be adapted to \(N\)), and \(d\le e=\sum_{x\in H_i,\ C\subset H_i\text{ near }x}a_i\). Then there
are \(k\), coordinates \(\Theta\) centred at \(z\), indices \(\sigma(H)<k\) for the members \(H\) of \(E'\) through \(z\), pairwise distinct, and exponents
\(a'_{H_i'}=a_i\), \(a'_D=e-d\), such that near \(z\)
\[
 N'=\{\Theta_t=0,\ t\ge k\},\qquad H=\{\Theta_{\sigma(H)}=0\}\ (z\in H\in E'),\qquad
 \operatorname{Res}_{N'}\mathfrak a'=\operatorname{Res}_{N'}\Bigl(\prod_{z\in H\in E'}\Theta_{\sigma(H)}^{\,a'_H}\Bigr)\ \text{at }z .
\]
So \(\Theta_0,\dots,\Theta_{k-1}\) restrict to coordinates of \(N'\) at \(z\), the traces \(H\cap N'\) of the members through \(z\) are coordinate hyperplanes
of \(N'\), and \(\mathfrak a'|_{N'}\) is the monomial \(\prod\Theta_{\sigma(H)}^{a'_H}\) in these coordinates (the form of Lemma~\ref{hir:monomial-blowup} along \(N'\)).
\end{corollary}
\begin{proof}
\leanok
As \(C\subset\operatorname{cosupp}\subset N\), \(C\ne\emptyset\) and \(E\) is transverse to \(N\), \(C\subset N\) and \(E\) admit adapted coordinates at \(x\)
(Lemma~\ref{bcalc:snc-simultaneous}). Take the chart \(W\) of Lemma~\ref{bcalc:monomial-chart-adapted}; \(z\in N'\) gives \(j<k_N\) and \(W_t(z)=0\) for
\(t\ge k_N\). Put \(k=k_N\) and \(\Theta=W-W(z)\), a coordinate system centred at \(z\) (Lemma~\ref{bcalc:snc-chart}). Near \(z\), \(y\in N'\iff W_t(y)=0\)
for \(t\ge k\iff\Theta_t(y)=0\) for \(t\ge k\). As in Corollary~\ref{bcalc:monomial-transform-coords}, \(\sigma(H_i')=s_i\) and \(\sigma(D)=j\) are distinct, the
members of \(E'\) through \(z\) are the hyperplanes \(\{\Theta_{\sigma(H)}=0\}\), and the coordinates \(W_j\), \(W_{s_i}\) (\(z\in H_i'\)) agree with
\(\Theta_j\), \(\Theta_{s_i}\); moreover \(s_i<k\) (\(x\in H_i\)) and \(j<k\). By Theorem~\ref{bcalc:monomial-transform-along}(2),
\(\operatorname{Res}_{N'}\mathfrak a'=\operatorname{Res}_{N'}(W_j^{e-d}\prod_{z\in H_i'}W_{s_i}^{a_i})\) at \(z\); this generator agrees near \(z\) with
\(\prod_{z\in H\in E'}\Theta_{\sigma(H)}^{a'_H}\) (the product over the members of \(E'\) through \(z\) splits into the strict transforms and \(D\), which passes
through \(z\); Definition~\ref{hir:geo-transform}, Theorem~\ref{bcalc:snc-transform}), and whether a function is on \(N'\) near \(z\) a multiple of a
generator depends only on the germ of the generator. Flat coordinates generate the ideals of their hyperplanes (Lemma~\ref{bcalc:hyperplane-ideal}),
which gives the variant with \(u_i=\Phi_{\sigma_0(i)}\).
\end{proof}

\begin{lemma}[Monomial ideals along \(N'\) off the centre]
\label{bcalc:monomial-transform-along-off-centre}
\uses{bcalc:monomial-transform-along,hir:monomial-step,bcalc:isomorphism-off-centre,bcalc:transform-off-centre,bcalc:strict-submanifold,bcalc:snc-transform,hir:geo-transform,hir:marked-ideal}
\source{wlodarczyk05-hironaka,bierstone-milman08-functoriality}
\lean{Hironaka.MonomialBlowupAlongChart.isInSpanOnAt_geoTransform_iff_pullback_of_notMem,Hironaka.MonomialBlowupAlongChart.exists_snc_chart_geoTransform_along_of_notMem}
\leanok
Let \((M,N,E,\mathfrak a,d)\) be a marked ideal with \(N\) closed and \(\mathfrak a\) coherent, \(E=(H_1,\dots,H_s)\) a simple normal crossings divisor,
\(\pi\) the blowup along \(C\), \(N'\), \(\mathfrak a'\), \(E'=(H_1',\dots,H_s',D)\) as in Corollary~\ref{bcalc:monomial-transform-along-coords}, and \(z\) a point with
\(x=\pi(z)\notin C\).
\begin{enumerate}
\item \(\operatorname{Res}_{N'}\mathfrak a'=\operatorname{Res}_{N'}\pi^*\mathfrak a\) at \(z\).
\item Let \(n>0\), and let \(\Phi\) be coordinates centred at \(x\) in which \(N=\{\Phi_t=0,\ t\ge k\}\) and the members through \(x\) are the hyperplanes
\(\{\Phi_{\sigma_0(i)}=0\}\), \(\sigma_0(i)<k\) pairwise distinct, and suppose
\(\operatorname{Res}_N\mathfrak a=\operatorname{Res}_N(\prod_{x\in H_i}\Phi_{\sigma_0(i)}^{a_i})\) at \(x\). Then \(\Theta=\Phi\circ\pi\) are coordinates centred at \(z\) in
which, near \(z\), \(N'=\{\Theta_t=0,\ t\ge k\}\), the members of \(E'\) through \(z\) are the hyperplanes \(\{\Theta_{\sigma_0(i)}=0\}\) of the strict
transforms \(H_i'\), \(x\in H_i\) (\(D\) does not pass through \(z\)), and
\(\operatorname{Res}_{N'}\mathfrak a'=\operatorname{Res}_{N'}(\prod_{z\in H_i'}\Theta_{\sigma_0(i)}^{a_i})\) at \(z\).
\end{enumerate}
With Corollary~\ref{bcalc:monomial-transform-along-coords} this gives the form of Lemma~\ref{hir:monomial-blowup} along \(N'\) at every point of \(N'\).
\end{lemma}
\begin{proof}
\leanok
(1) Near \(z\) the ideal \(\pi^*\mathcal I_C\) contains \(1\), and so does its \(d\)-th power; the transformed ideal is the colon along \(N'\) of \(\pi^*\mathfrak a\)
by this power (Definition~\ref{hir:geo-transform}), so its germs are, on \(N'\) near \(z\), combinations of germs of \(\pi^*\mathfrak a\) (multiply by \(1\)), and
every germ \(t\) of \(\pi^*\mathfrak a\) lies in it (\(ht=h\cdot t\)). (2) Off \(\pi^{-1}(C)\) the blowup is a local biholomorphism, so \(\Phi\circ\pi\), read in
charts, is a coordinate system centred at \(z\) (Lemmas~\ref{bcalc:isomorphism-off-centre} and~\ref{hir:monomial-step}); its domain lies off \(\pi^{-1}(C)\).
There \(N'=\pi^{-1}(N)\) and \(H_i'=\pi^{-1}(H_i)\), as \(N\) and the members are closed (Lemmas~\ref{bcalc:strict-submanifold}
and~\ref{bcalc:transform-off-centre}; the indexing of \(E'\), Theorem~\ref{bcalc:snc-transform}), which gives the coordinate descriptions, and the members of
\(E'\) through \(z\) are the \(H_i'\) with \(x\in H_i\). By (1) and the pullback of restricted generators along \(\pi\) (Theorem~\ref{bcalc:monomial-transform-along},
\(\pi(N')\subset N\)), \(\operatorname{Res}_{N'}\mathfrak a'=\operatorname{Res}_{N'}(m\circ\pi)\) with \(m=\prod_{x\in H_i}\Phi_{\sigma_0(i)}^{a_i}\), and
\(m\circ\pi=\prod_{z\in H_i'}\Theta_{\sigma_0(i)}^{a_i}\) near \(z\).
\end{proof}

\paragraph{The boundary count.} Let \(E=(H_1,\dots,H_s)\) be a simple normal crossings divisor (Definition~\ref{hir:snc}). For a set \(S\) of
indices let \(H_S=\bigcap_{i\in S}H_i\), and let \(s(x)=\#\{i:x\in H_i\}\) be the number of members through \(x\). Along a blowup \(\pi\) the
count \(s'(z)\) is that of the strict transforms \(H_i'=\overline{\pi^{-1}(H_i\setminus C)}\) of the members; the new exceptional divisor is not
counted (Włodarczyk identifies the strict transforms of \(E\) with \(E\)). Step 1a of the algorithm blows up maximal intersections \(H_S\),
which lowers \(s\). Włodarczyk intersects irreducible components of the members; here the members themselves are intersected, and the
maximality of \(s\) becomes the hypothesis that no member outside \(S\) meets \(H_S\).

\begin{lemma}[Intersections of members]
\label{bcalc:member-intersection}
\uses{hir:snc,hir:transform,bcalc:snc-simultaneous,bcalc:snc-transform-nc,lag:regular-locus-dense,ana:blowups}
\source{wlodarczyk05-hironaka,bierstone-milman08-functoriality}
\lean{Hironaka.BoundaryCount.memberInter,Hironaka.BoundaryCount.isClosed_memberInter,Hironaka.BoundaryCount.mem_coordSubspace_biUnion,Hironaka.BoundaryCount.complexCoordinatePlane_min,Hironaka.BoundaryCount.hasNormalCrossings_memberInter,Hironaka.BoundaryCount.isClosedSubmanifold_memberInter}\upstreamlean{TranscendentalNakamaye.Pending.BoundaryCount.isAdmissibleCentre_memberInter}
\mixedleanok
Let \(E\) be a simple normal crossings divisor on \(X\) and \(S\) a set of indices of members. Then \(H_S\) is closed and has only normal
crossings with \(E\). If \(n>0\), \(H_S\) is a closed submanifold. If \(E\) is the divisor of a marked ideal \(\underline{\mathcal I}\) and
\(H_S\subset\operatorname{cosupp}\underline{\mathcal I}\), then \(H_S\) is an admissible centre for \(\underline{\mathcal I}\).
\end{lemma}
\begin{proof}
\mixedleanok
\(H_S\) is an intersection of closed sets. Let \(x\in H_S\) and take coordinates at \(x\) in which each member \(H_i\) through \(x\) is
\(L_{T_i}\) (Definition~\ref{hir:snc}). Every member with index in \(S\) passes through \(x\), so near \(x\)
\(H_S=\bigcap_{i\in S}L_{T_i}=L_{\bigcup_{i\in S}T_i}\) in the same coordinates; together with the \(L_{T_i}\) this is the normal crossings of
\(H_S\) with \(E\). Let \(n>0\). The divisor \(E\) is transverse to \(N=X\) (Corollary~\ref{bcalc:snc-transform-nc}), so
Lemma~\ref{bcalc:snc-simultaneous} gives coordinates at \(x\) in which \(H_S=\mathbf C^{k}\times0\) near \(x\); replacing \(k\) by
\(\min(k,n)\) does not change this plane. Hence \(H_S\) is a submanifold of dimension \(\min(k,n)\) at every point near \(x\)
(Theorem~\ref{lag:regular-locus-dense}). The last claim is Definition~\ref{hir:transform}.
\end{proof}

\begin{lemma}[Intersections of members with a transverse submanifold]
\label{bcalc:member-intersection-along}
\uses{bcalc:member-intersection,hir:snc,bcalc:snc-simultaneous,bcalc:snc-transform-nc,lag:regular-locus-dense,hir:transform,hir:marked-ideal}
\source{wlodarczyk05-hironaka,bierstone-milman08-functoriality}
\lean{Hironaka.MemberInterAlong.hasNormalCrossings_inter_memberInter,Hironaka.MemberInterAlong.isClosedSubmanifold_inter_memberInter}\upstreamlean{TranscendentalNakamaye.Pending.MemberInterAlong.isAdmissibleCentre_inter_memberInter,TranscendentalNakamaye.Pending.MemberInterAlong.isAdmissibleCentre_inter_memberInter_of_isWellFormed}
\mixedleanok
Let \(E\) be a simple normal crossings divisor on \(X\) transverse to a closed set \(N\), and \(S\) a set of indices of members. Then
\(N\cap H_S\) is closed and has only normal crossings with \(E\). If \(n>0\), \(N\cap H_S\) is a closed submanifold. If \(N\) and \(E\) are those of a
marked ideal \(\underline{\mathcal I}\) (for instance a well formed one) and \(N\cap H_S\subset\operatorname{cosupp}\underline{\mathcal I}\), then \(N\cap H_S\) is an
admissible centre for \(\underline{\mathcal I}\). (For \(N=X\) this is Lemma~\ref{bcalc:member-intersection}; along \(N\) it gives the admissibility of
intersections of the traces \(H\cap N\) of members, as in the monomial case along \(N\).)
\end{lemma}
\begin{proof}
\mixedleanok
\(N\cap H_S\) is an intersection of closed sets. Let \(x\in N\cap H_S\) and take coordinates at \(x\) in which \(N=L_{S_0}\) and each member \(H_i\)
through \(x\) is \(L_{T_i}\), \(T_i\) a single index not in \(S_0\) (transversality, Definition~\ref{hir:snc}). Every member with index in \(S\) passes
through \(x\), so near \(x\) \(N\cap H_S=L_{S_0}\cap\bigcap_{i\in S}L_{T_i}=L_{S_0\cup\bigcup_{i\in S}T_i}\) in the same coordinates; together with the
\(L_{T_i}\) this is the normal crossings of \(N\cap H_S\) with \(E\). Let \(n>0\). As in Lemma~\ref{bcalc:member-intersection}, \(E\) is transverse to
\(X\) (Corollary~\ref{bcalc:snc-transform-nc}), so Lemma~\ref{bcalc:snc-simultaneous} gives coordinates at \(x\) in which \(N\cap H_S=\mathbf C^{k}\times0\)
near \(x\), and \(N\cap H_S\) is a submanifold of dimension \(\min(k,n)\) near \(x\) (Theorem~\ref{lag:regular-locus-dense}). The last claim is
Definition~\ref{hir:transform}; a well formed marked ideal has \(N\) closed and \(E\) simple normal crossings and transverse to \(N\)
(Definition~\ref{hir:marked-ideal}).
\end{proof}

\begin{lemma}[The boundary count]
\label{bcalc:boundary-count}
\uses{hir:snc,bcalc:strict-submanifold,hir:test-tower,hir:test-tower-basic,hir:geo-transform,hir:geo-transform-basic}
\source{wlodarczyk05-hironaka,bierstone-milman08-functoriality,wlodarczyk09-analytic-resolution:page-0025}
\lean{Hironaka.BoundaryCount.boundaryCount,Hironaka.BoundaryCount.eventually_boundaryCount_le,Hironaka.BoundaryCount.boundaryCount_map,Hironaka.BoundaryCount.boundaryCount_map_geoStrictTransform_le}\upstreamlean{TranscendentalNakamaye.Pending.BoundaryCount.isClosed_of_mem_take_E_transformAlong,TranscendentalNakamaye.Pending.BoundaryCount.boundaryCount_take_E_transformAlong_le}\lean{Hironaka.BoundaryCount.isClosed_of_mem_take_E_geoTransformAlong}\upstreamlean{TranscendentalNakamaye.Pending.BoundaryCount.boundaryCount_take_E_geoTransformAlong_le}
\mixedleanok
Let \(E\) be a finite sequence of closed subsets of \(X\). (1) \(s\) is upper semicontinuous: \(s(y)\le s(x)\) for all \(y\) near \(x\).
(2) For a continuous map \(\pi\colon Y\to X\) and any \(C\), \(s'(z)\le s(\pi(z))\) for every \(z\in Y\), where \(s'(z)\) counts the
\(i\) with \(z\in\overline{\pi^{-1}(H_i\setminus C)}\).
(3) Along a tower of blowups \(T\) with projection \(p\) (either transform of Definitions~\ref{hir:test-tower} and~\ref{hir:geo-transform}), a
point \(z\) of the last space lies on at most \(s(p(z))\) of the first \(\#E\) members of the transformed divisor, which are the iterated strict
transforms of the original members.
\end{lemma}
\begin{proof}
\mixedleanok
(1) The complements of the finitely many members not through \(x\) are open, so near \(x\) every member through \(y\) passes through \(x\).
(2) \(\overline{\pi^{-1}(H_i\setminus C)}\subset\pi^{-1}(H_i)\), the latter being closed (Theorem~\ref{bcalc:strict-submanifold}); so every \(i\)
counted at \(z\) is counted at \(\pi(z)\). The count along the sequence \((H_1',\dots,H_s')\) is re-indexed by the length of the sequence.
(3) Induction on the tower. At each stage the divisor is \(E_{i+1}=(E_i\text{ strict-transformed})\) followed by the new exceptional divisor
and has \(\#E+i+1\) members (Lemma~\ref{hir:test-tower-basic}, Lemma~\ref{hir:geo-transform-basic}), so its first \(\#E\) members are the strict
transforms of the first \(\#E\) members of \(E_i\). These are closed (the original members by hypothesis, later ones as closures), and (2) at
each stage composes with the induction hypothesis.
\end{proof}

\begin{theorem}[Blowing up a maximal intersection separates the members]
\label{bcalc:boundary-count-drop}
\uses{bcalc:member-intersection,bcalc:boundary-count,bcalc:snc-simultaneous,bcalc:monomial-lift,bcalc:snc-chart,bcalc:strict-submanifold}
\source{wlodarczyk05-hironaka,bierstone-milman08-functoriality,wlodarczyk09-analytic-resolution:page-0025}
\upstreamlean{TranscendentalNakamaye.Pending.BoundaryCount.boundaryCount_map_geoStrictTransform_lt,TranscendentalNakamaye.Pending.BoundaryCount.card_filter_mem_geoStrictTransform_lt,TranscendentalNakamaye.Pending.BoundaryCount.boundaryCount_map_geoStrictTransform_lt_boundaryCount}
\upstreamleanok
Let \(E\) be a simple normal crossings divisor on \(X\), \(S\) a set of indices such that no member \(H_i\) with \(i\notin S\) meets
\(C=H_S\), and \(\pi\) the blowup along \(C\). Then \(s'(z)<\#S\) for every \(z\in\pi^{-1}(C)\). (Włodarczyk, Step 1aa: ``by blowing up
\(C=H^s_\alpha\) we separate divisors contributing to \(H^s_\alpha\), thus creating new points all with \(s(x)<s\)''.) No hypothesis
\(S\neq\emptyset\) is needed: \(H_\emptyset=X\), and a blowup centre has empty interior.

Local form. Let \(C\) be any blowup centre with only normal crossings with \(E\), and \(z\in\pi^{-1}(C)\). If \(C\supset H_S\) near
\(\pi(z)\), then \(z\) does not lie on all the \(H_i'\), \(i\in S\). In particular, if near \(x=\pi(z)\) the centre contains the intersection of
the members through \(x\), then \(s'(z)<s(x)\); no global maximality is needed. This is the form used in Step 1aa: near each of its points
the centre \(H^s_\alpha\) is the intersection of the \(s\) members through that point.
\end{theorem}
\begin{proof}
\upstreamleanok
Let \(x=\pi(z)\in C\). If \(z\in H_i'\) then \(x\in H_i\) (Lemma~\ref{bcalc:boundary-count}), so \(i\in S\); hence \(s'(z)\le\#S\), and equality
means \(z\in H_i'\) for every \(i\in S\). Suppose this. Then \(n>0\) (Lemma~\ref{bcalc:snc-simultaneous}), and \(C\) has only normal crossings with
\(E\) (Lemma~\ref{bcalc:member-intersection}). Take coordinates \(\Phi\) at \(x\) in which \(C=\mathbf C^k\times0\) and \(H_i=\{w_{s_i}=0\}\) for the
members through \(x\) (Lemma~\ref{bcalc:snc-simultaneous} with \(N=X\)), and a monomial chart \(j\ge k\), \(W\) around \(z\) with
\(\sigma_j\circ W=\Phi\circ\pi\) (Lemma~\ref{bcalc:monomial-lift}). By Lemma~\ref{bcalc:snc-chart}, \(z\in H_i'\) gives \(s_i\neq j\) and
\(W_{s_i}(z)=0\) for \(i\in S\); and \(W_j(z)=0\) since \(\pi(z)\in C\). For small \(\varepsilon\neq0\) the point \(w_\varepsilon=\sigma_j(W(z)+\varepsilon e_j)\)
lies in the target of \(\Phi\). For \(i\in S\) its coordinate \(s_i\) is \(W_{s_i}(z)\) or \(W_{s_i}(z)\,\varepsilon\), hence \(0\), so
\(\Phi^{-1}(w_\varepsilon)\in H_i\) for all \(i\in S\), that is \(\Phi^{-1}(w_\varepsilon)\in C\). But the \(j\)-th coordinate of \(w_\varepsilon\) is
\(\varepsilon\neq0\) with \(j\ge k\), so \(w_\varepsilon\notin\mathbf C^k\times0\): a contradiction.

Local form: the same argument, with \(C\) adapted directly (it has normal crossings with \(E\)) and the members of \(S\) passing through \(x\)
because \(z\) lies on their strict transforms. As \(\varepsilon\to0\), \(\Phi^{-1}(w_\varepsilon)\to\Phi^{-1}(\sigma_j(W(z)))=x\), so for small
\(\varepsilon\) the point lies in the neighbourhood of \(x\) where \(H_S\subset C\). For the count, let \(S\) be the set of members through \(x\): every
\(H_i'\) through \(z\) has \(i\in S\), and not all of them pass through \(z\).
\end{proof}

\begin{theorem}[Blowing up an intersection of members along a transverse submanifold separates them on the strict transform]
\label{bcalc:boundary-count-drop-along}
\uses{bcalc:boundary-count-drop,bcalc:boundary-count,bcalc:member-intersection-along,bcalc:snc-simultaneous,bcalc:monomial-lift,bcalc:snc-chart,bcalc:strict-submanifold}
\source{wlodarczyk05-hironaka,bierstone-milman08-functoriality,wlodarczyk09-analytic-resolution:page-0025}
\lean{Hironaka.BoundaryCountAlong.card_filter_mem_geoStrictTransform_lt_along,Hironaka.BoundaryCountAlong.boundaryCount_map_geoStrictTransform_lt_boundaryCount_along}\upstreamlean{TranscendentalNakamaye.Pending.BoundaryCountAlong.boundaryCount_map_geoStrictTransform_lt_along}
\mixedleanok
Let \(E\) be a simple normal crossings divisor on \(X\) transverse to a set \(N\), \(C\subset N\) the centre of the blowup \(\pi\), with only normal
crossings with \(E\), \(N'=\overline{\pi^{-1}(N\setminus C)}\) the strict transform, and \(S\) a set of indices of members. Let \(z\in N'\) with
\(x=\pi(z)\in C\).
\begin{enumerate}
\item If \(C\supset N\cap H_S\) near \(x\), then \(z\) does not lie on all the \(H_i'\), \(i\in S\).
\item If near \(x\) the centre contains \(N\cap\bigcap_{x\in H_i}H_i\), then \(s'(z)<s(x)\).
\item If \(C=N\cap H_S\) and no member \(H_i\) with \(i\notin S\) meets \(C\), then \(s'(z)<\#S\).
\end{enumerate}
This is Theorem~\ref{bcalc:boundary-count-drop} inside \(N\): Step~1aa of Włodarczyk's algorithm for a marked ideal along \(N\), and the
drop of the monomial invariant along \(N\) (Lemma~\ref{hir:monomial-case}(2)), where the centre is \(N\cap H_I\) and the cosupport lies in \(N'\).
For \(N=X\) it is the local form of Theorem~\ref{bcalc:boundary-count-drop}; for \(N\ne X\) the hypothesis \(C\supset H_S\) of that theorem fails, and
the point \(z\in N'\) is what replaces it.
\end{theorem}
\begin{proof}
\mixedleanok
(1) Suppose \(z\in H_i'\) for every \(i\in S\); then \(x\in H_i\) for \(i\in S\) (Lemma~\ref{bcalc:boundary-count}), and \(n>0\) since \(C\) is a
nonempty centre. Take coordinates \(\Phi\) at \(x\) adapted to \(C\subset N\) and \(E\) (Lemma~\ref{bcalc:snc-simultaneous}): \(C=\mathbf C^{k_C}\times0\),
\(N=\mathbf C^{k_N}\times0\), \(H_i=\{w_{s_i}=0\}\) for the members through \(x\); and a monomial chart \(j\ge k_C\), \(W\) around \(z\) with
\(\sigma_j\circ W=\Phi\circ\pi\) (Lemma~\ref{bcalc:monomial-lift}). By Lemma~\ref{bcalc:snc-chart}: \(z\in N'\) gives \(j<k_N\) and \(W_t(z)=0\) for
\(t\ge k_N\); \(z\in H_i'\) gives \(s_i\ne j\) and \(W_{s_i}(z)=0\); and \(W_j(z)=0\) since \(x\in C\). For small \(\varepsilon\ne0\) put
\(w_\varepsilon=\sigma_j(W(z)+\varepsilon e_j)\). Its coordinate \(t\ge k_N\) is \(W_t(z)(W_j(z)+\varepsilon)=0\) (as \(t\ne j\)), so
\(\Phi^{-1}(w_\varepsilon)\in N\); its coordinate \(s_i\), \(i\in S\), is \(W_{s_i}(z)\) or \(W_{s_i}(z)\varepsilon\), hence \(0\), so
\(\Phi^{-1}(w_\varepsilon)\in H_S\). As \(\varepsilon\to0\), \(\Phi^{-1}(w_\varepsilon)\to\Phi^{-1}(\sigma_j(W(z)))=x\), so for small \(\varepsilon\) the point lies
in the neighbourhood of \(x\) where \(N\cap H_S\subset C\). But its \(j\)-th coordinate is \(\varepsilon\ne0\) with \(j\ge k_C\), so it is not in
\(\mathbf C^{k_C}\times0\): a contradiction.
(2) Let \(S\) be the set of members through \(x\). Every \(H_i'\) through \(z\) has \(i\in S\) (Lemma~\ref{bcalc:boundary-count}), and by (1) not all of them
pass through \(z\).
(3) \(C\) is closed with normal crossings with \(E\) (Lemma~\ref{bcalc:member-intersection-along}); the members through \(x\in C\) are those of \(S\), so
near \(x\) the centre is \(N\cap\bigcap_{x\in H_i}H_i\), and (2) gives \(s'(z)<s(x)\le\#S\).
\end{proof}

\begin{lemma}[The members outside an intersection, restricted to a transverse submanifold]
\label{bcalc:member-intersection-drop-along}
\uses{bcalc:member-intersection-along,bcalc:snc-simultaneous,bcalc:sub-divisor,bcalc:nc-comparison,hir:snc,hir:marked-ideal}
\source{wlodarczyk05-hironaka,bierstone-milman08-functoriality}
\lean{Hironaka.BoundaryCountAlong.isTransverse_inter_memberInter_of_forall,Hironaka.BoundaryCountAlong.isTransverse_inter_memberInter_drop,Hironaka.BoundaryCountAlong.isSNCDivisor_drop}\upstreamlean{TranscendentalNakamaye.Pending.BoundaryCountAlong.inter_memberInter_subset_of_notMem_drop}\lean{Hironaka.BoundaryCountAlong.isWellFormed_inter_memberInter_drop}
\mixedleanok
Let \(E\) be a simple normal crossings divisor on \(X\) transverse to \(N\), \(S\) a set of indices of members, and \(E_S\) the subsequence of the
members with index not in \(S\). Then \(E_S\) is a simple normal crossings divisor, transverse to \(N\cap H_S\) (more generally, any finite sequence of
members with indices not in \(S\) is transverse to \(N\cap H_S\)), and every member of \(E\) missing from \(E_S\) contains \(N\cap H_S\). Hence, if
\((M,N,E,\mathfrak a,d)\) is well formed, \(n>0\), and \(\mathfrak b\) is coherent, then \((M,N\cap H_S,E_S,\mathfrak b,\mu)\) is well formed. This is the
restriction of Step~1a of Włodarczyk's algorithm inside \(N\): the marked ideal on \(S=N\cap H^s_\alpha\) with the divisor without the \(s\) members
containing it (Lemmas~\ref{hir:drop-containing-members} and~\ref{hir:coefficient-restriction}).
\end{lemma}
\begin{proof}
\mixedleanok
\(E_S\) is a subsequence of \(E\) (Lemma~\ref{bcalc:sub-divisor}(2)). A member of \(E\) missing from \(E_S\) is \(H_i\) with \(i\in S\), which contains \(H_S\).
Transversality: let \(x\in N\cap H_S\) and take coordinates at \(x\) in which \(N=L_{S_0}\) and every member \(H_i\) through \(x\) is \(L_{T_i}\), \(T_i\) a single
index not in \(S_0\) (transversality of \(E\) to \(N\), Definition~\ref{hir:snc}). Every member with index in \(S\) passes through \(x\), so near \(x\)
\(N\cap H_S=L_{S_0\cup\bigcup_{i\in S}T_i}\) (Lemma~\ref{bcalc:member-intersection-along}). Let \(H_j\), \(j\notin S\), pass through \(x\). Its index \(T_j\) is
not in \(S_0\); and \(T_j\ne T_i\) for \(i\in S\): otherwise \(H_j=L_{T_j}=L_{T_i}=H_i\) near \(x\), and two members of a simple normal crossings divisor
through \(x\) that agree near \(x\) are the same member (Lemma~\ref{bcalc:snc-simultaneous}, first paragraph of the proof), so \(j=i\in S\). Hence
\(T_j\) is disjoint from \(S_0\cup\bigcup_{i\in S}T_i\), which is the transversality of \(E_S\) to \(N\cap H_S\) at \(x\). For the marked ideal,
\(N\cap H_S\) is a closed submanifold (Lemma~\ref{bcalc:member-intersection-along}, \(N\) closed and \(n>0\)), and the conditions of
Definition~\ref{hir:marked-ideal} are those just proved.
\end{proof}

\begin{lemma}[Strict transforms of intersections of members]
\label{bcalc:member-intersection-strict}
\uses{bcalc:trace-strict-transform,bcalc:snc-simultaneous,bcalc:snc-chart,bcalc:monomial-lift,bcalc:member-intersection-along,hir:snc}
\source{bierstone-milman08-functoriality}
\lean{Hironaka.MemberInterTowerAlong.geoStrictTransform_inter_memberInter_eq}
\leanok
Let \(\pi\) be the blowup along \(C\), \(N\supset C\) closed, \(E\) a simple normal crossings divisor transverse to \(N\) with which \(C\) has only normal
crossings, and \(S\) a set of indices of members. Then \(\overline{\pi^{-1}((N\cap H_S)\setminus C)}=N'\cap\bigcap_{i\in S}H_i'\).
\end{lemma}
\begin{proof}
\leanok
As for one member (Lemma~\ref{bcalc:trace-strict-transform}). \(\subset\): strict transforms are monotone. \(\supset\): let \(z\in N'\cap\bigcap_{i\in S}H_i'\), so
\(\pi(z)\in N\cap H_S\). If \(\pi(z)\notin C\), \(z\in\pi^{-1}((N\cap H_S)\setminus C)\). If \(\pi(z)\in C\), in coordinates adapted to \(C\subset N\) and \(E\)
(Lemma~\ref{bcalc:snc-simultaneous}) \(N=L_{T_N}\) and \(H_i=L_{\{s_i\}}\), so \(N\cap H_S=L_{T_N\cup\{s_i:i\in S\}}\). In a monomial chart \(j\) through \(z\)
(Lemma~\ref{bcalc:monomial-lift}), Lemma~\ref{bcalc:snc-chart} turns \(z\in N'\) into \(j\notin T_N\), \(W(z)\in L_{T_N}\), and \(z\in H_i'\) into \(j\ne s_i\),
\(W(z)\in L_{\{s_i\}}\); so \(j\notin T_N\cup\{s_i:i\in S\}\) and \(W(z)\in L_{T_N\cup\{s_i:i\in S\}}\), that is \(z\in(N\cap H_S)'\) by the same lemma.
\end{proof}

\begin{lemma}[Intersections of members along multiple test blow-ups]
\label{bcalc:member-intersection-tower}
\uses{bcalc:member-intersection-strict,mc:strict-along,hir:wellformed-tower,hir:test-tower,hir:geo-transform,hir:blowup-tower,bcalc:boundary-count}
\source{wlodarczyk09-analytic-resolution:page-0025,bierstone-milman08-functoriality}
\lean{Hironaka.MemberInterTowerAlong.strictAlong_mono,Hironaka.MemberInterTowerAlong.isClosed_strictAlong,Hironaka.MemberInterTowerAlong.strictAlong_subset_preimage,Hironaka.MemberInterTowerAlong.disjoint_strictAlong,Hironaka.MemberInterTowerAlong.take_E_geoTransformAlong,Hironaka.MemberInterTowerAlong.length_E_le_length_E_geoTransformAlong,Hironaka.MemberInterTowerAlong.get_E_geoTransformAlong,Hironaka.MemberInterTowerAlong.setOf_forall_strictAlong_eq_memberInter,Hironaka.MemberInterTowerAlong.strictAlong_inter_memberInter}\upstreamlean{TranscendentalNakamaye.Pending.MemberInterTowerAlong.mem_strictAlong_inter_memberInter_iff,TranscendentalNakamaye.Pending.MemberInterTowerAlong.mem_strictAlong_inter_memberInter_self}
\mixedleanok
Let \(T\) be a tower over \(X\), and write \(A_T\) for the iterated strict transform of \(A\subset X\) (Definition~\ref{mc:strict-along}). Then \(A\mapsto A_T\) is
monotone; for \(A\) closed, \(A_T\) is closed and lies over \(A\), so the strict transforms of disjoint closed sets are disjoint. For a marked ideal
\((M,N,E,\mathfrak a,d)\), the first \(\#E\) members of the transformed divisor \(E_T\) are the strict transforms \((H_i)_T\) of the members of \(E\), in
order (the exceptional divisors follow them). If the marked ideal is well formed and \(T\) is a multiple test blow-up, then for every set \(S\) of indices
\[
(N\cap H_S)_T=N_T\cap\bigcap_{i\in S}(H_i)_T .
\]
In particular a point \(z\in N_T\) lies on the strict transform of \(N\cap H_{S(z)}\), \(S(z)\) the set of original members whose strict transforms pass
through \(z\). This justifies, inside \(N\), the identification of the strict transforms of \(E\) with \(E\) in Step~1a of Włodarczyk's algorithm: a point with
\(s(z)=s\) (Lemma~\ref{bcalc:boundary-count}) lies on the strict transform of the intersection it lies on at the start of Step~1.
\end{lemma}
\begin{proof}
\mixedleanok
Monotonicity, closedness and \(A_T\subset T^{-1}(A)\) by induction on the length, each step being a closure of a preimage (Definition~\ref{mc:strict-along}).
The members: \(E\) after one more blowup \(\pi\) along \(C\) is \(E_T'=((H_1)_T',\dots,(H_{s+r})_T',\pi^{-1}(C))\) (Definition~\ref{hir:geo-transform}), whose first
\(\#E\) entries are the strict transforms of the first \(\#E\) entries of \(E_T\). The intersection by induction on the length: for the empty tower it is the
definition of \(H_S\). If \(T'\) is \(T\) followed by the blowup along \(C\), the transform along \(T\) is well formed (Lemma~\ref{hir:wellformed-tower}), so \(N_T\)
is closed and \(E_T\) is simple normal crossings and transverse to \(N_T\), and \(C\subset N_T\) has only normal crossings with \(E_T\) (admissibility,
Definition~\ref{hir:test-tower}). By induction \((N\cap H_S)_T=N_T\cap\bigcap_{i\in S}(H_i)_T\), an intersection of \(N_T\) with members of \(E_T\), and
Lemma~\ref{bcalc:member-intersection-strict} applied to \(N_T\), \(E_T\), \(C\) gives \((N\cap H_S)_{T'}=N_T'\cap\bigcap_{i\in S}(H_i)_T'\), with
\(N_{T'}=N_T'\) and \((H_i)_{T'}=(H_i)_T'\).
\end{proof}

\begin{lemma}[Localising the centres of a tower over an open set]
\label{bcalc:centres-localise}
\uses{bcalc:member-intersection-tower,mc:strict-along,hir:open-restriction-tower,hir:test-tower,mc:coefficient-restriction-drop-along}
\source{wlodarczyk09-analytic-resolution:page-0025,bierstone-milman08-functoriality}
\lean{Hironaka.MemberInterTowerAlong.strictAlong_union}\upstreamlean{TranscendentalNakamaye.Pending.MemberInterTowerAlong.centresIn_restrict_of_centresIn_union}\lean{Hironaka.MemberInterTowerAlong.centresIn_restrict_of_forall}
\mixedleanok
Let \(T\) be a tower over \(X\) and \(U\subset X\) open, and write \(T|_U\) for its restriction (Lemma~\ref{hir:open-restriction-tower}). Then \((A\cup B)_T=A_T\cup B_T\).
(1) If the centres of \(T\) lie in the strict transforms of \(A\cup B\), with \(B\) closed and disjoint from \(U\), then the centres of \(T|_U\) lie in the strict
transforms of \(A\cap U\). (2) If every admissible centre \(C\) of every multiple test blow-up \(T\) of a marked ideal satisfies \(C\cap T^{-1}(U)\subset A_T\), then
the centres of the restriction to \(U\) of every multiple test blow-up lie in the strict transforms of \(A\cap U\).
In Step~1a inside \(N\) the centres lie in the strict transforms of the union of the intersections \(S_\alpha=N\cap H^s_\alpha\), whose traces on the cosupport
are disjoint; (1) with \(A=S_\alpha\), \(B=\bigcup_{\alpha'\ne\alpha}(S_{\alpha'}\cap\operatorname{cosupp})\) and \(U=X\setminus B\supset S_\alpha\cap\operatorname{cosupp}\) puts the
restricted tower in the setting of the restricted coefficient ideal with \(E_{S_\alpha}\) (Corollary~\ref{mc:coefficient-restriction-drop-along}) on \(U\). In Step~1b, (2) applies to the
strict transforms of a hypersurface of maximal contact \(N\cap V(u)\cap U\).
\end{lemma}
\begin{proof}
\mixedleanok
The union: by induction on the length, as \(\overline{\pi^{-1}((A\cup B)\setminus C)}=\overline{\pi^{-1}(A\setminus C)}\cup\overline{\pi^{-1}(B\setminus C)}\).
The centres of \(T|_U\) are the traces of those of \(T\) on the open preimages of \(U\), and the strict transform of \(A\cap U\) along \(T|_U\) is the trace of
\(A_T\) (Definition~\ref{mc:strict-along}). (1) A point \(p\) of a centre over \(U\) lies in \(A_T\cup B_T\); it does not lie in \(B_T\), which lies over \(B\)
(Lemma~\ref{bcalc:member-intersection-tower}), while \(p\) lies over \(U\). (2) Each centre of \(T\) is an admissible centre of a multiple test blow-up
(Definition~\ref{hir:test-tower}).
\end{proof}

\begin{lemma}[The exceptional divisors are transverse to the strict transform of a submanifold containing the centres]
\label{bcalc:exceptional-transverse-tower}
\uses{bcalc:member-intersection-tower,bcalc:snc-transform,bcalc:strict-submanifold,bcalc:sub-divisor,hir:wellformed-tower,hir:test-tower,hir:snc-basic,mc:strict-along}
\source{wlodarczyk09-analytic-resolution:page-0025,bierstone-milman08-functoriality}
\lean{Hironaka.MemberInterTowerAlong.drop_E_geoTransformAlong_step,Hironaka.MemberInterTowerAlong.isTransverse_drop_E_geoTransformAlong}\upstreamlean{TranscendentalNakamaye.Pending.MemberInterTowerAlong.isTransverse_E_geoTransformAlong,TranscendentalNakamaye.Pending.MemberInterTowerAlong.isWellFormed_strictAlong_drop}\upstreamunbound{TranscendentalNakamaye.Pending.MemberInterTowerAlong.isTransverse_drop_E_geoTransformAlong',TranscendentalNakamaye.Pending.MemberInterTowerAlong.isWellFormed_strictAlong_drop'}
\historicalleanok
Let \((M,N,E,\mathfrak a,d)\) be well formed, \(W\) a closed submanifold and \(T\) a multiple test blow-up whose centres lie in the strict transforms of \(W\).
Then \(W_T\) is a closed submanifold, and the members of \(E_T\) after the first \(\#E\) (the exceptional divisors of \(T\)) form a simple normal crossings
divisor transverse to \(W_T\); hence \((M_T,W_T,E_T\setminus E,\mathfrak b,\mu)\) is well formed for every coherent \(\mathfrak b\). If \(E\) is transverse to \(W\),
all of \(E_T\) is transverse to \(W_T\). This is Włodarczyk's observation that the new divisors, born from centres in the hypersurface of maximal contact,
are transversal to it: with \(W=N\cap V(u)\cap U\) and Lemma~\ref{bcalc:centres-localise}(2) it makes Step~1b's marked ideal on the strict transform of the
hypersurface of maximal contact, with the old members dropped in Step~1a, well formed.
\end{lemma}
\begin{proof}
\historicalleanok
Induction on the length. For the empty tower there are no exceptional members (\(\emptyset\) is transverse to every closed submanifold,
Lemma~\ref{hir:snc-basic}). Let \(T'\) be \(T\) followed by the blowup \(\pi\) along \(C\subset W_T\). The exceptional members of \(T'\) are the strict transforms of
those of \(T\) followed by \(\pi^{-1}(C)\) (Definition~\ref{hir:geo-transform}). \(C\) has only normal crossings with \(E_T\) (admissibility), hence with its
sub-sequence of exceptional members (Lemma~\ref{bcalc:sub-divisor}), which is simple normal crossings and, by induction, transverse to the closed
submanifold \(W_T\). Theorem~\ref{bcalc:snc-transform} (with \(W_T\supset C\) in place of \(N\)) shows that the transformed sequence is simple normal crossings and
transverse to \(W_T'=W_{T'}\), a closed submanifold by Theorem~\ref{bcalc:strict-submanifold}. For all of \(E_T\) the same induction runs with \(E_T\), simple normal
crossings since the transform along \(T\) is well formed (Lemma~\ref{hir:wellformed-tower}).
\end{proof}

\begin{lemma}[Normal crossings and transversality are local]
\label{bcalc:normal-crossings-local}
\uses{hir:snc,hir:snc-invariance}
\source{bierstone-milman08-functoriality}
\lean{Hironaka.CodimOneAdmissibleAlong.isCoordinateFamilyAt_congr,Hironaka.CodimOneAdmissibleAlong.isTransverse_of_forall_eventually,Hironaka.CodimOneAdmissibleAlong.hasNormalCrossings_of_forall_eventually,Hironaka.CodimOneAdmissibleAlong.isTransverse_inter_of_preimage_val}
\leanok
Adapted coordinates for a family at \(x\) (Definition~\ref{hir:snc}) are adapted coordinates for every family agreeing with it near \(x\). Hence if near each
of its points a set \(R\) agrees with a set \(W\) transverse to \(E\) (resp.\ having only normal crossings with \(E\)), then \(R\) is transverse to \(E\)
(resp.\ has only normal crossings with \(E\)). If \(O\) is open and the trace of \(W\) on \(O\) is transverse, in \(O\), to the traces of the members of \(E\), then
\(W\cap O\) is transverse to \(E\).
\end{lemma}
\begin{proof}
\leanok
For the first sentence apply Lemma~\ref{hir:snc-invariance} to the identity, whose differential is bijective, on an open set on which the families agree.
The second follows at each point \(x\in R\) from the coordinates of \(W\) at \(x\). For the last, adapted coordinates at \(z\in W\cap O\) on \(O\) are adapted
coordinates on the ambient manifold for \(W\) and the members (Lemma~\ref{hir:snc-invariance}, open inclusions), and \(W\cap O\) agrees with \(W\) near \(z\).
\end{proof}

\begin{lemma}[The exceptional divisors over an open set]
\label{bcalc:exceptional-transverse-open}
\uses{bcalc:exceptional-transverse-tower,bcalc:centres-localise,bcalc:normal-crossings-local,hir:open-restriction-tower,hir:open-restriction-marked,hir:giraud-tower-along}
\source{wlodarczyk09-analytic-resolution:page-0027,bierstone-milman08-functoriality}
\lean{Hironaka.CodimOneAdmissibleAlong.isTransverse_drop_E_inter_preimage}\upstreamlean{TranscendentalNakamaye.Pending.CodimOneAdmissibleAlong.isTransverse_drop_E_of_forall_eventually}
\mixedleanok
Let \((M,N,E,\mathfrak a,d)\) be well formed, \(U\subset M\) open and \(A\) a set whose trace on \(U\) is a closed submanifold of \(U\), and assume that every
admissible centre \(C\) of every multiple test blow-up \(T\) satisfies \(C\cap T^{-1}(U)\subset A_T\) (for a tangent direction \(u\) along \(N\) on \(U\) and
\(A=N\cap V(u)\cap U\) this is Theorem~\ref{hir:giraud-tower-along}). Then along every multiple test blow-up \(T\) the exceptional members of \(E_T\) are transverse
to \(A_T\cap T^{-1}(U)\), and so to every set \(R\) which near each of its points lies over \(U\) and agrees with \(A_T\).
\end{lemma}
\begin{proof}
\mixedleanok
Restrict to \(U\) (Lemma~\ref{hir:open-restriction-tower}): the restricted tower is a multiple test blow-up of the well-formed restriction
(Lemma~\ref{hir:open-restriction-marked}), its centres lie in the strict transforms of \(A\cap U\) (Lemma~\ref{bcalc:centres-localise}(2)), so its exceptional
members are transverse to the strict transform of \(A\cap U\) (Lemma~\ref{bcalc:exceptional-transverse-tower}). The transform along the restricted tower is the
restriction of the transform, and the strict transform of \(A\cap U\) the trace of \(A_T\); Lemma~\ref{bcalc:normal-crossings-local} returns to \(M_T\), and gives the
statement for \(R\).
\end{proof}

\begin{corollary}[Step~1ba's centre inside \(N\) is admissible]
\label{bcalc:codim-one-admissible-along}
\uses{bcalc:exceptional-transverse-open,bcalc:codim-one-part-along,bcalc:codim-one-local-along,hir:marked-ideal,hir:giraud-tower-along,mc:tangent-direction-along-exists,bcalc:transform-off-centre-tower,bcalc:giraud-local-along,hir:open-restriction-tower,hir:geo-transform-basic,bcalc:normal-crossings-local}
\source{wlodarczyk09-analytic-resolution:page-0027}
\upstreamlean{TranscendentalNakamaye.Pending.CodimOneAdmissibleAlong.isAdmissibleCentre_codimOnePartAlong}\lean{Hironaka.CodimOneAdmissibleAlong.isTransverse_codimOnePartAlong_drop_E}\upstreamlean{TranscendentalNakamaye.Pending.CodimOneAdmissibleAlong.isAdmissibleCentre_codimOnePartAlong_stage}
\mixedleanok
Let \((M,N,E,\mathfrak a,d)\) have \(N\) a closed submanifold and \(\mathfrak a\) coherent, be of maximal order along \(N\), \(d\ge1\). If the codimension-one part along \(N\) of its cosupport
(Definition~\ref{bcalc:codim-one-part-along-def}) is transverse to \(E\), it is an admissible centre. If \((M,N,E,\mathfrak a,d)\) is well formed, of
maximal order along \(N\), \(d\ge1\), then at every stage \(T\) of a multiple test blow-up the codimension-one part along \(N_T\) of the cosupport of the
transform is transverse to the exceptional members of \(E_T\) (the old ones dropped in Step~1a); hence, if the transform is of maximal order along \(N_T\), it is
an admissible centre of \((M_T,N_T,E_T\setminus E,\mathfrak a_T,d)\). Indeed near each of its points the codimension-one part is the strict transform of
the hypersurface of maximal contact \(N\cap V(u)\cap U\) (Theorem~\ref{bcalc:codim-one-part-along} for the tangent direction of the transform,
Theorem~\ref{hir:giraud-tower-along}), to which the exceptional members are transverse (Lemma~\ref{bcalc:exceptional-transverse-open}). This is Włodarczyk's
statement in Step~1ba (2009, p.~57) that the divisors are transversal to the codimension-one components.
\end{corollary}
\begin{proof}
\mixedleanok
The codimension-one part is a closed submanifold contained in the cosupport (Theorem~\ref{bcalc:codim-one-part-along}), and transversality implies normal
crossings (Definition~\ref{hir:snc}). At a stage \(T\): let \(x\) lie in the codimension-one part \(R_T\). Then \(T(x)\) lies in the cosupport of the
original marked ideal (Lemma~\ref{bcalc:transform-off-centre-tower}), where by maximal order there is a tangent direction \(u\) along \(N\) on an open \(U\ni T(x)\)
(Lemma~\ref{mc:tangent-direction-along-exists}). Theorem~\ref{hir:giraud-tower-along} gives, on an open \(O\ni x\) over \(U\), a function \(w\) controlled on
\(N_T\) with \(N_T\cap V(w)\cap O=A_T\cap O\), \(A=N\cap V(u)\cap U\), and with every admissible centre over \(U\) inside \(A_T\). By
Theorem~\ref{bcalc:codim-one-part-along} (its form for \(w\) controlled on \(N_T\)), \(R_T=N_T\cap V(w)\) near \(x\), so \(R_T\) agrees near \(x\) with
\(A_T\cap T^{-1}(U)\), to which the exceptional members are transverse (Lemma~\ref{bcalc:exceptional-transverse-open}; \(A\cap U\) is a closed submanifold of \(U\),
Lemma~\ref{bcalc:giraud-local-along}). Lemma~\ref{bcalc:normal-crossings-local} concludes.
\end{proof}

\paragraph{Canonicity of the codimension-one part and of the boundary count.} The invariant of the maximal-order step
(Proposition~\ref{hir:maximal-order-step}) gives the points of the codimension-one part of the support the value \((0,\infty)\) and counts the original members
through a point; its properties (3)--(7) (Definition~\ref{hir:centre-function}) need these two ingredients to commute with restriction to open sets and with
biholomorphisms, to be unchanged off the centre of a blowup and by blowups along \(\emptyset\), and to depend only on the cosupports of equivalent marked ideals.

\begin{lemma}[Regular points under biholomorphisms and local sections]
\label{bcalc:regular-point-transport}
\uses{ana:regular-transport,rem:open-subset-analytic,hir:biholomorphism-invariance,bcalc:isomorphism-off-centre}
\source{wlodarczyk05-hironaka}
\lean{Hironaka.CodimOnePartCanonical.regularAtDim_preimage_iff,Hironaka.CodimOnePartCanonical.regularAtDim_preimage_iff_of_section}
\leanok
Let \(\Phi\colon M'\to M\) be a biholomorphism. Then \(x\) is a regular point of \(\Phi^{-1}(A)\) of dimension \(k\) if and only if \(\Phi(x)\) is a regular point
of \(A\) of dimension \(k\). More generally, let \(\pi\colon Y\to X\) be holomorphic, \(W\subset X\) open, \(s\colon W\to Y\) holomorphic with \(\pi\circ s=\mathrm{id}_W\),
and \(z\in Y\) with \(\pi(z)\in W\) and \(s(\pi(z))=z\) (for instance a blowup along \(C\), \(W=X\setminus C\) and \(\pi(z)\notin C\),
Lemma~\ref{bcalc:isomorphism-off-centre}). Then \(z\) is a regular point of \(\pi^{-1}(A)\) of dimension \(k\) if and only if \(\pi(z)\) is a regular point of \(A\)
of dimension \(k\).
\end{lemma}
\begin{proof}
\leanok
A biholomorphism has bijective differential everywhere (Lemma~\ref{hir:biholomorphism-invariance}), so Lemma~\ref{ana:regular-transport} applied to \(\Phi\) and
\(A\), and to \(\Phi^{-1}\) and \(\Phi^{-1}(A)\), gives the two directions. In the second statement, the chain rule for \(\pi\circ s=\mathrm{id}_W\) at
\(x=\pi(z)\) gives \(d\pi_z\circ ds_x=\mathrm{id}\) (\(s(x)=z\)); in equal finite dimensions both differentials are then bijective. Lemma~\ref{ana:regular-transport}
for \(\pi\) at \(z\) gives one direction; for \(s\) at \(x\), with \(s^{-1}(\pi^{-1}(A))=A\cap W\), it gives that \(x\) is a regular point of \(A\cap W\) in \(W\), that is
of \(A\) (Remark~\ref{rem:open-subset-analytic}).
\end{proof}

\begin{proposition}[The codimension-one part along \(N\) is canonical]
\label{bcalc:codim-one-part-canonical}
\uses{bcalc:codim-one-part-along-def,bcalc:regular-point-transport,hir:biholomorphism-transform,hir:open-restriction-marked,hir:open-restriction-tower,hir:tower-transport,hir:tower-transport-resolution,hir:restricted-equivalence,hir:restricted-equivalence-basic,hir:marked-equivalence,hir:marked-equivalence-basic,mc:strict-along,hir:cosupport,rem:dimension,rem:open-subset-analytic}
\source{wlodarczyk05-hironaka,bierstone-milman08-functoriality}
\lean{Hironaka.CodimOnePartCanonical.mem_codimOnePartAlong_iff_of_eventually}\upstreamlean{TranscendentalNakamaye.Pending.CodimOnePartCanonical.mem_codimOnePartAlong_congr,TranscendentalNakamaye.Pending.CodimOnePartCanonical.codimOnePartAlong_congr}\lean{Hironaka.CodimOnePartCanonical.cosupport_pullback,Hironaka.CodimOnePartCanonical.codimOnePartAlong_pullback,Hironaka.CodimOnePartCanonical.codimOnePartAlong_pullback_val}\upstreamlean{TranscendentalNakamaye.Pending.CodimOnePartCanonical.codimOnePartAlong_geoTransformAlong_transport,TranscendentalNakamaye.Pending.CodimOnePartCanonical.codimOnePartAlong_geoTransformAlong_restrict,TranscendentalNakamaye.Pending.CodimOnePartCanonical.codimOnePartAlong_geoTransformAlong_of_restrictEquiv,TranscendentalNakamaye.Pending.CodimOnePartCanonical.codimOnePartAlong_geoTransformAlong_of_wEquiv}
\mixedleanok
Write \(R_N(\mathfrak a,\mu)\) for the codimension-one part along \(N\) of \(\operatorname{cosupp}(\mathfrak a|_N,\mu)\) (Definition~\ref{bcalc:codim-one-part-along-def}), and
\(R(\underline{\mathcal I})=R_N(\mathfrak a,d)\) for \(\underline{\mathcal I}=(M,N,E,\mathfrak a,d)\). Let \(\mathfrak a\) consist of holomorphic germs.
\begin{enumerate}
\item (Locality) Whether \(x\in R_N(\mathfrak a,\mu)\) depends only on the germs at \(x\) of \(N\) and of \(\operatorname{cosupp}(\mathfrak a|_N,\mu)\); equal cosupports along the
same \(N\) give equal codimension-one parts.
\item (Biholomorphisms) For a biholomorphism \(\Phi\colon M'\to M\), \(\operatorname{cosupp}((\Phi^*\mathfrak a)|_{\Phi^{-1}N},\mu)=\Phi^{-1}\operatorname{cosupp}(\mathfrak a|_N,\mu)\)
and \(R_{\Phi^{-1}N}(\Phi^*\mathfrak a,\mu)=\Phi^{-1}R_N(\mathfrak a,\mu)\).
\item (Open sets) For \(U\subset M\) open, \(R_{N\cap U}(\mathfrak a|_U,\mu)=R_N(\mathfrak a,\mu)\cap U\).
\item (Transported towers) If \(\mathfrak a\) is coherent, \(T\) a tower over \(M\) and \(\Phi^*T\) its transport with last-stage biholomorphism \(\Phi_T\)
(Definition~\ref{hir:tower-transport}), then \(R(\Phi^*\underline{\mathcal I}_{\Phi^*T})=\Phi_T^{-1}R(\underline{\mathcal I}_T)\).
\item (Restricted towers) If \(N\) is a closed submanifold, \(\mathfrak a\) coherent, \(T\) a multiple test blow-up of \(\underline{\mathcal I}\) and \(T|_U\) its restriction
(Lemma~\ref{hir:open-restriction-tower}), then \(R((\underline{\mathcal I}|_U)_{T|_U})\) is the trace of \(R(\underline{\mathcal I}_T)\) on the open preimage of \(U\).
\item (Representatives) If \(\underline{\mathcal I}\) and \(\underline{\mathcal J}\) are restricted-equivalent (Definition~\ref{hir:restricted-equivalence}; \(N\) closed, coherent
ideals), then \(R(\underline{\mathcal I}_T)=R(\underline{\mathcal J}_T)\) for every tower \(T\); if they are equivalent (Definition~\ref{hir:marked-equivalence}) with the same \(N\),
then \(R(\underline{\mathcal I}_T)=R(\underline{\mathcal J}_T)\) for every multiple test blow-up \(T\) of \(\underline{\mathcal I}\).
\end{enumerate}
\end{proposition}
\begin{proof}
\mixedleanok
(1) Regularity of a given dimension at \(x\) depends only on the germ of the set at \(x\) (Remark~\ref{rem:dimension}). (2) The cosupport identity is that of
Lemma~\ref{hir:biholomorphism-transform} for the marked ideal \((M,N,\emptyset,\mathfrak a,\mu)\); then apply Lemma~\ref{bcalc:regular-point-transport} to \(N\) and to the
cosupport. (3) The cosupport restricts (Lemma~\ref{hir:open-restriction-marked}), and regular points of \(A\cap U\) in \(U\) are those of \(A\) in \(U\)
(Remark~\ref{rem:open-subset-analytic}). (4) For the empty tower this is (2); from the first blowup on the two transforms are equal and \(\Phi_T\) is the identity
(Lemma~\ref{hir:tower-transport-resolution}). (5) Along a multiple test blow-up \(N_T\) stays a closed submanifold and \(\mathfrak a_T\) coherent, and the transform of
\(\underline{\mathcal I}|_U\) along \(T|_U\) is the restriction of \(\underline{\mathcal I}_T\) to the open preimage of \(U\) (Lemma~\ref{hir:open-restriction-tower}); apply (3).
(6) Restricted-equivalent marked ideals have the same \(N\) and cosupport, and equal transforms from the first blowup on
(Lemmas~\ref{hir:restricted-equivalence-basic} and~\ref{hir:marked-equivalence-basic}); equivalent marked ideals have equal cosupports along their multiple test
blow-ups (Definition~\ref{hir:marked-equivalence}), and \(N_T\) is the iterated strict transform of \(N\), the same for both (Lemma~\ref{mc:strict-along}); apply (1).
\end{proof}

\begin{lemma}[The codimension-one part off the centre]
\label{bcalc:codim-one-part-off-centre}
\uses{bcalc:codim-one-part-canonical,bcalc:regular-point-transport,bcalc:isomorphism-off-centre,bcalc:transform-off-centre,hir:open-restriction-tower,hir:geo-transform}
\source{wlodarczyk05-hironaka}
\lean{Hironaka.CodimOnePartCanonical.mem_codimOnePartAlong_geoTransform_iff}\upstreamlean{TranscendentalNakamaye.Pending.CodimOnePartCanonical.codimOnePartAlong_geoTransform_empty}\lean{Hironaka.CodimOnePartCanonical.mem_codimOnePartAlong_geoTransformAlong_step_iff}
\mixedleanok
Let \(\underline{\mathcal I}=(M,N,E,\mathfrak a,d)\) have \(N\) closed and \(\mathfrak a\) consisting of holomorphic germs, \(\pi\) the blowup along \(C\) and
\(\underline{\mathcal I}'\) the transform. If \(\pi(z)\notin C\), then \(z\in R(\underline{\mathcal I}')\) if and only if \(\pi(z)\in R(\underline{\mathcal I})\). In particular a
blowup along \(\emptyset\) gives \(R(\underline{\mathcal I}')=\pi^{-1}R(\underline{\mathcal I})\), and along a multiple test blow-up \(T\) of \(\underline{\mathcal I}\) (\(N\) a
closed submanifold, \(\mathfrak a\) coherent) followed by the blowup \(\pi\) along \(C\), \(z\in R(\underline{\mathcal I}_{T\pi})\) if and only if
\(\pi(z)\in R(\underline{\mathcal I}_T)\) whenever \(\pi(z)\notin C\).
\end{lemma}
\begin{proof}
\mixedleanok
Near \(z\) (on the open set \(\pi^{-1}(M\setminus C)\)) the strict transform \(N'\) is \(\pi^{-1}(N)\) and the cosupport of \(\underline{\mathcal I}'\) is the preimage of
that of \(\underline{\mathcal I}\) (Theorem~\ref{bcalc:transform-off-centre}). By Proposition~\ref{bcalc:codim-one-part-canonical}(1), \(z\in R(\underline{\mathcal I}')\) if and only
if for some \(k\) the sets \(\pi^{-1}(N)\) and \(\pi^{-1}(\operatorname{cosupp}\underline{\mathcal I})\) are regular at \(z\) of dimensions \(k+1\) and \(k\). The section of
\(\pi\) over \(M\setminus C\) through \(z\) (Lemma~\ref{bcalc:isomorphism-off-centre}) and Lemma~\ref{bcalc:regular-point-transport} turn this into the same condition for
\(N\) and \(\operatorname{cosupp}\underline{\mathcal I}\) at \(\pi(z)\). The tower form applies this to \(\underline{\mathcal I}_T\), whose \(N_T\) is closed and ideal
coherent (Lemma~\ref{hir:open-restriction-tower}).
\end{proof}

\begin{lemma}[The boundary count is canonical]
\label{bcalc:boundary-count-canonical}
\uses{bcalc:boundary-count,bcalc:transform-off-centre,hir:tower-transport,hir:tower-transport-resolution,hir:open-restriction-tower,hir:geo-transform}
\source{wlodarczyk05-hironaka,bierstone-milman08-functoriality}
\lean{Hironaka.CodimOnePartCanonical.boundaryCount_map_preimage,Hironaka.CodimOnePartCanonical.boundaryCount_map_geoStrictTransform_of_notMem,Hironaka.CodimOnePartCanonical.boundaryCount_take_E_geoTransformAlong_step_of_notMem}\upstreamlean{TranscendentalNakamaye.Pending.CodimOnePartCanonical.boundaryCount_take_E_geoTransformAlong_transport,TranscendentalNakamaye.Pending.CodimOnePartCanonical.boundaryCount_take_E_geoTransformAlong_restrict}
\mixedleanok
For a map \(f\) and a sequence \(E\) of sets, the number of members of \((f^{-1}H)_{H\in E}\) through \(z\) is the number of members of \(E\) through \(f(z)\). Let
\(\underline{\mathcal I}=(M,N,E,\mathfrak a,d)\) have closed members, and write \(s_T(z)\) for the number of the first \(\#E\) members of \(E_T\) (the strict transforms of the
members of \(E\)) through \(z\) (Definition~\ref{bcalc:boundary-count}). Then:
\begin{enumerate}
\item for a blowup \(\pi\) along a closed \(C\) and \(\pi(z)\notin C\), the number of strict transforms of members through \(z\) is the number of members through \(\pi(z)\);
in particular \(s_{T\pi}(z)=s_T(\pi(z))\) for a tower \(T\) followed by the blowup \(\pi\) along \(C\);
\item if \(\mathfrak a\) is coherent, along the transport \(\Phi^*T\) of a tower by a biholomorphism, \(s_{\Phi^*T}=s_T\circ\Phi_T\);
\item if \(N\) is a closed submanifold and \(\mathfrak a\) coherent, along the restriction \(T|_U\) of a multiple test blow-up, \(s_{T|_U}\) is the restriction of \(s_T\) to the
open preimage of \(U\).
\end{enumerate}
\end{lemma}
\begin{proof}
\mixedleanok
The first statement re-indexes the members. (1) Off the centre the strict transform of a closed \(H\) is \(\pi^{-1}(H)\) near \(z\) (Theorem~\ref{bcalc:transform-off-centre}); the first
\(\#E\) members of the transform of \(\underline{\mathcal I}_T\) are the strict transforms of the first \(\#E\) members of \(E_T\), which are closed
(Lemma~\ref{bcalc:boundary-count}). (2) For the empty tower the members are \(\Phi^{-1}H\); from the first blowup on the transforms are equal and \(\Phi_T\) is the
identity (Lemma~\ref{hir:tower-transport-resolution}). (3) The transform along \(T|_U\) is the restriction of the transform along \(T\)
(Lemma~\ref{hir:open-restriction-tower}), whose members are the traces of the members.
\end{proof}

\paragraph{The codimension-one part along multiple test blow-ups.} Step~1ba of the maximal-order step (Proposition~\ref{hir:maximal-order-step}) blows up the
codimension-one part of the support along \(N\); Step~1bb evaluates the invariant of the restriction of the coefficient ideal to a hypersurface of maximal
contact, and the inductive invariant (Definition~\ref{hir:centre-function}) is only defined for marked ideals whose ideal vanishes identically on no nonempty
open subset of their submanifold. Włodarczyk states that the blowup of the codimension-one components ``reduces the situation to the case when
\(\operatorname{supp}\) is of codimension \(\ge2\)'' and then works at the points off \(\operatorname{codim}(1)(\operatorname{supp})\); the three statements below make
this precise inside \(N\): the codimension-one part never reappears along the segment, and off it the restriction to a hypersurface of maximal contact is
nonzero. As in Proposition~\ref{bcalc:codim-one-part-canonical}, \(R(\underline{\mathcal I})\) is the codimension-one part along \(N\) of the cosupport of
\(\underline{\mathcal I}=(M,N,E,\mathfrak a,d)\).

\begin{theorem}[The codimension-one part along a multiple test blow-up lies over the original one]
\label{bcalc:codim-one-part-tower}
\uses{bcalc:codim-one-part-along-def,bcalc:codim-one-part-along,bcalc:codim-one-local-along,bcalc:codim-one-part-off-centre,bcalc:transform-off-centre-tower,hir:giraud-tower-along,mc:tangent-direction-along-exists,mc:strict-along,hir:wellformed-tower,hir:test-tower}
\source{wlodarczyk05-hironaka,wlodarczyk09-analytic-resolution:page-0027}
\lean{Hironaka.CodimOnePartTowerAlong.codimOnePartAlong_geoTransformAlong_subset}\upstreamlean{TranscendentalNakamaye.Pending.CodimOnePartTowerAlong.codimOnePartAlong_geoTransformAlong_eq_empty}
\mixedleanok
Let \(\underline{\mathcal I}=(M,N,E,\mathfrak a,d)\) be well formed, of maximal order along \(N\), with \(d\ge1\), and let \(T\) be a multiple test blow-up of
\(\underline{\mathcal I}\) with projection \(p\). Then \(R(\underline{\mathcal I}_T)\subset p^{-1}(R(\underline{\mathcal I}))\). In particular, if \(R(\underline{\mathcal I})=\emptyset\),
then \(R(\underline{\mathcal I}_T)=\emptyset\) for every multiple test blow-up \(T\).
\end{theorem}
\begin{proof}
\mixedleanok
Induction on \(T\); for the empty tower there is nothing to prove. Let \(T'\) be \(T\) followed by the blowup \(\pi\) along \(C\), with projection \(p'=p\circ\pi\),
and \(x\in R(\underline{\mathcal I}_{T'})\). Then \(x\) lies in the cosupport of \(\underline{\mathcal I}_{T'}\), so \(p'(x)\) lies in that of \(\underline{\mathcal I}\)
(Corollary~\ref{bcalc:transform-off-centre-tower}), and by maximal order there is a tangent direction \(u\) along \(N\) on an open \(U\ni p'(x)\)
(Proposition~\ref{mc:tangent-direction-along-exists}). Put \(A=N\cap V(u)\cap U\). By Theorem~\ref{hir:giraud-tower-along} there are an open \(O\ni x\) and a function \(w\)
on \(O\), controlled on \(N_{T'}\) as in the last sentence of Lemma~\ref{bcalc:codim-one-local-along}, whose zero set in \(N_{T'}\cap O\) is the strict transform
\(A_{T'}\) (Definition~\ref{mc:strict-along}). By Theorem~\ref{bcalc:codim-one-part-along}(3), \(R(\underline{\mathcal I}_{T'})=N_{T'}\cap V(w)\) near \(x\); hence near \(x\)
the points of \(A_{T'}\cap O\) lie in \(R(\underline{\mathcal I}_{T'})\), and \(x\in A_{T'}=\overline{\pi^{-1}(A_T\setminus C)}\). So every neighbourhood of \(x\) contains a point
\(y\in\pi^{-1}(A_T\setminus C)\) so close to \(x\) that \(y\in A_{T'}\cap O\) lies in \(R(\underline{\mathcal I}_{T'})\). As \(\pi(y)\notin C\), \(\pi(y)\in R(\underline{\mathcal I}_T)\)
(Lemma~\ref{bcalc:codim-one-part-off-centre}), and \(p'(y)=p(\pi(y))\in R(\underline{\mathcal I})\) by induction. Since \(R(\underline{\mathcal I})\) is closed
(Theorem~\ref{bcalc:codim-one-part-along}(2)) and \(p'\) continuous, \(x\in p'^{-1}(R(\underline{\mathcal I}))\). The transforms along multiple test blow-ups stay well
formed (Lemma~\ref{hir:wellformed-tower}), which the cited results need at each stage.
\end{proof}

\begin{corollary}[After Step 1ba the codimension-one part is empty]
\label{bcalc:codim-one-part-empty}
\uses{bcalc:codim-one-part-tower,bcalc:codim-one-step-along,bcalc:codim-one-part-off-centre,hir:geo-transform-basic,hir:wellformed-tower,hir:along-transform}
\source{wlodarczyk05-hironaka,wlodarczyk09-analytic-resolution:page-0027}
\upstreamlean{TranscendentalNakamaye.Pending.CodimOnePartTowerAlong.codimOnePartAlong_geoTransform_codimOnePartAlong_eq_empty,TranscendentalNakamaye.Pending.CodimOnePartTowerAlong.codimOnePartAlong_geoTransformAlong_append_eq_empty}
\upstreamleanok
Let \(\underline{\mathcal I}=(M,N,E,\mathfrak a,d)\) have \(N\) a closed submanifold, \(\mathfrak a\) coherent, \(d\ge1\), and be of maximal order along \(N\); let \(\pi\) be
the blowup along \(R(\underline{\mathcal I})\) and \(\underline{\mathcal I}'\) the transform. Then \(R(\underline{\mathcal I}')=\emptyset\). If \(\underline{\mathcal I}\) is well
formed, then \(R(\underline{\mathcal I}_T)=\emptyset\) for every multiple test blow-up \(T\) of \(\underline{\mathcal I}\) whose first blowup is \(\pi\).
\end{corollary}
\begin{proof}
\upstreamleanok
The cosupport of \(\underline{\mathcal I}'\) does not meet \(\pi^{-1}(R(\underline{\mathcal I}))\) (Corollary~\ref{bcalc:codim-one-step-along}), and at a point \(z\) with
\(\pi(z)\notin R(\underline{\mathcal I})\), \(z\in R(\underline{\mathcal I}')\) would give \(\pi(z)\in R(\underline{\mathcal I})\) (Lemma~\ref{bcalc:codim-one-part-off-centre}). A multiple
test blow-up whose first blowup is \(\pi\) is \(\pi\) followed by a multiple test blow-up of \(\underline{\mathcal I}'\) (Lemma~\ref{hir:geo-transform-basic}), and
\(\underline{\mathcal I}'\) is well formed (Lemma~\ref{hir:wellformed-tower}), of maximal order along \(N'\) (Lemma~\ref{hir:along-transform}), with the same mark; apply
Theorem~\ref{bcalc:codim-one-part-tower} to \(\underline{\mathcal I}'\).
\end{proof}

\begin{proposition}[Off the codimension-one part the restriction to a hypersurface of maximal contact is nonzero]
\label{bcalc:contact-restriction-nonzero}
\uses{bcalc:codim-one-part-canonical,bcalc:codim-one-part-tower,mc:contact-equivalence-open-along,mc:contact-hypersurface-along,mc:tangent-direction-along,mc:homogenized-along-basic,hir:open-restriction-ideal,hir:marked-equivalence,hir:exponent-along,hir:cosupport,hir:wellformed-tower,ana:ideal-sheaves}
\source{wlodarczyk05-hironaka,wlodarczyk09-analytic-resolution:page-0027}
\lean{Hironaka.CodimOnePartTowerAlong.eventually_forall_eq_zero_of_not_isNonzeroAlongAt,Hironaka.CodimOnePartTowerAlong.isNonzeroAlongAt_of_notMem_codimOnePartAlong,Hironaka.CodimOnePartTowerAlong.IsTangentDirectionAlong.exists_regularAtDim_succ}\upstreamlean{TranscendentalNakamaye.Pending.CodimOnePartTowerAlong.isNonzeroAlongAt_contact_open,TranscendentalNakamaye.Pending.CodimOnePartTowerAlong.isNonzeroAlongAt_contact_open_geoTransformAlong}
\mixedleanok
\begin{enumerate}
\item Let \(N\) be regular of dimension \(k+1\) and \(A\) regular of dimension \(k\) at \(y\), \(\mathfrak b\) coherent, and suppose that near \(y\)
\(\operatorname{cosupp}(\mathfrak a|_N,\mu)\subset A\) and \(\operatorname{cosupp}(\mathfrak b|_A,\nu)\subset\operatorname{cosupp}(\mathfrak a|_N,\mu)\). If \(y\notin R_N(\mathfrak a,\mu)\),
then some germ of \(\mathfrak b\) at \(y\) does not vanish identically on \(A\) near \(y\) (Lemma~\ref{hir:exponent-along}).
\item Let \(\underline{\mathcal I}=(M,N,E,\mathfrak a,d)\) have \(N\) a closed submanifold, \(\mathfrak a\) coherent, \(d\ge1\), let \(u\) be a tangent direction along \(N\) on an
open \(U\) (Definition~\ref{mc:tangent-direction-along}), and let \(\underline{\mathcal S}_u=(U,N\cap V(u),E|_U,\mathcal C_N(\mathfrak a)|_U,\mu_{\mathcal C})\) be the restricted
marked ideal with \(\underline{\mathcal I}|_U\simeq\underline{\mathcal S}_u\) (Lemma~\ref{mc:contact-equivalence-open-along}). At every point \(y\) of \(N\cap V(u)\cap U\), \(N\) is
regular of some dimension \(k+1\) and \(N\cap V(u)\cap U\) of dimension \(k\); if \(y\notin R(\underline{\mathcal I})\), some germ of the ideal of \(\underline{\mathcal S}_u\) at \(y\)
does not vanish identically on \(N\cap V(u)\) near \(y\).
\item If \(\underline{\mathcal I}\) is moreover well formed and of maximal order along \(N\) with \(R(\underline{\mathcal I})=\emptyset\) (for instance the transform after
Step~1ba, Corollary~\ref{bcalc:codim-one-part-empty}), then for every multiple test blow-up \(T\) and every tangent direction \(w\) of \(\underline{\mathcal I}_T\) along \(N_T\) on
an open \(U\), the conclusion of (2) holds for \(\underline{\mathcal S}_w\) at every point of \(N_T\cap V(w)\cap U\). In particular the ideal of \(\underline{\mathcal S}_w\)
vanishes identically on no nonempty open subset of \(N_T\cap V(w)\cap U\), which is the hypothesis under which the inductive invariant of
Definition~\ref{hir:centre-function} is defined, as Step~1bb of Proposition~\ref{hir:maximal-order-step} uses it.
\end{enumerate}
\end{proposition}
\begin{proof}
\mixedleanok
(1) Suppose every germ of \(\mathfrak b\) at \(y\) vanishes identically on \(A\) near \(y\). Take generators \(s_1,\dots,s_r\) of \(\mathfrak b\) on a neighbourhood of \(y\)
(Definition~\ref{ana:ideal-sheaves}); each \(s_i\) is a germ of \(\mathfrak b\) at \(y\), so all of them vanish on \(A\) near \(y\), and then at every point \(z\) near \(y\) every germ
of \(\mathfrak b_z\), a combination of the \(s_i\), vanishes on \(A\) near \(z\). Such a germ agrees on \(A\) with \(0\), which lies in every power of the maximal ideal; so
every point of \(A\) near \(y\) lies in \(\operatorname{cosupp}(\mathfrak b|_A,\nu)\) (Definition~\ref{hir:cosupport}), hence in \(\operatorname{cosupp}(\mathfrak a|_N,\mu)\subset A\).
Thus \(\operatorname{cosupp}(\mathfrak a|_N,\mu)=A\) near \(y\), regular of dimension \(k\) at \(y\), where \(N\) has dimension \(k+1\): \(y\in R_N(\mathfrak a,\mu)\)
(Proposition~\ref{bcalc:codim-one-part-canonical}(1)).
(2) By Lemma~\ref{mc:contact-hypersurface-along}, \(N\) is regular at \(y\) of some dimension \(k'\) and \(N\cap V(u)\cap U\) of dimension \(k'-1\); and \(k'\ge1\), since a field
tangent to \(N\) at \(y\) with \(\xi(u)(y)\ne0\) gives a nonzero partial derivative of \(u\) in a tangential coordinate of index \(<k'\) (proof of that lemma). The equivalence
\(\underline{\mathcal I}|_U\simeq\underline{\mathcal S}_u\) gives equal cosupports (Definition~\ref{hir:marked-equivalence}, the empty tower), and the cosupport of
\(\underline{\mathcal S}_u\) lies in its submanifold. The ideal of \(\underline{\mathcal S}_u\) is coherent (Lemma~\ref{mc:homogenized-along-basic} and
Lemma~\ref{hir:open-restriction-ideal}), and \(R(\underline{\mathcal I}|_U)=R(\underline{\mathcal I})\cap U\) (Proposition~\ref{bcalc:codim-one-part-canonical}(3)). Apply (1) on \(U\).
(3) Along \(T\), \(N_T\) is a closed submanifold, \(\mathfrak a_T\) coherent and the mark unchanged (Lemma~\ref{hir:wellformed-tower}), and \(R(\underline{\mathcal I}_T)=\emptyset\)
(Theorem~\ref{bcalc:codim-one-part-tower}); apply (2) to \(\underline{\mathcal I}_T\). If the ideal vanished identically on a nonempty open subset of
\(N_T\cap V(w)\cap U\), all its germs at a point of that subset would vanish there.
\end{proof}

\paragraph{The member locus: Step 1aa inside \(N\).} In Step~1a of the maximal-order step (Proposition~\ref{hir:maximal-order-step}) the points of an intersection
\(N\cap H^s_\alpha\) of \(s\) members contained in the cosupport get the value \((s,\infty)\), and Włodarczyk notes that ``all \(H^s_\alpha\subset\operatorname{supp}\)
will be blown up first and we reduce the situation to the case where no \(H^s_\alpha\) is contained in \(\operatorname{supp}\)''. Włodarczyk intersects irreducible
components of the members; the invariant has to be local (Definition~\ref{hir:centre-function}(5)), and an intersection \(N\cap H_S\) of members transverse
to \(N\) is a submanifold, whose connected component through \(x\) lies in an analytic set as soon as a neighbourhood of \(x\) in it does. So the value set
is defined pointwise, by the germ at \(x\) of the intersection of \(N\) with the members through \(x\).

\begin{definition}[The member locus]
\label{bcalc:member-locus-def}
\uses{hir:cosupport,bcalc:boundary-count,hir:snc}
\source{wlodarczyk05-hironaka,wlodarczyk09-analytic-resolution:page-0025}
\lean{Hironaka.MemberInterLocusAlong.memberInterLocus}
\leanok
Let \(N,S\subset M\) and let \(E\) be a finite sequence of subsets of \(M\). The \emph{member locus} \(L_N(E,S)\) is the set of points \(x\in S\) near which every
point of \(N\) lying on all the members of \(E\) through \(x\) lies in \(S\): near \(x\), \(N\cap H_{S(x)}\subset S\), where \(S(x)\) is the set of members through
\(x\). For a marked ideal \(\underline{\mathcal I}=(M,N,E,\mathfrak a,d)\) and a tower \(T\) we write \(L(\underline{\mathcal I}_T)\) for \(L_{N_T}(E_T^{\le\#E},\operatorname{cosupp}\underline{\mathcal I}_T)\),
where \(E_T^{\le\#E}\) is the sequence of the first \(\#E\) members of \(E_T\), the strict transforms of the members of \(E\) (Lemma~\ref{bcalc:boundary-count}); its
points with \(s(x)=s\) are the points of value \((s,\infty)\) of Step~1aa.
\end{definition}

\begin{proposition}[The member locus is canonical]
\label{bcalc:member-locus-canonical}
\uses{bcalc:member-locus-def,bcalc:isomorphism-off-centre,bcalc:transform-off-centre,bcalc:codim-one-part-canonical,hir:biholomorphism-transform,hir:tower-transport,hir:tower-transport-resolution,hir:open-restriction-marked,hir:open-restriction-tower,hir:restricted-equivalence,hir:restricted-equivalence-basic,hir:marked-equivalence,hir:marked-equivalence-basic,mc:strict-along,bcalc:boundary-count,hir:geo-transform,hir:geo-transform-basic}
\source{wlodarczyk05-hironaka,bierstone-milman08-functoriality}
\lean{Hironaka.MemberInterLocusAlong.forall_get_iff,Hironaka.MemberInterLocusAlong.mem_memberInterLocus_iff}\upstreamlean{TranscendentalNakamaye.Pending.MemberInterLocusAlong.memberInterLocus_subset}\lean{Hironaka.MemberInterLocusAlong.setOf_forall_mem_eq_memberInter,Hironaka.MemberInterLocusAlong.mem_memberInterLocus_map_congr}\upstreamlean{TranscendentalNakamaye.Pending.MemberInterLocusAlong.mem_memberInterLocus_map_congr_id}\lean{Hironaka.MemberInterLocusAlong.mem_memberInterLocus_preimage_iff,Hironaka.MemberInterLocusAlong.memberInterLocus_preimage}\upstreamlean{TranscendentalNakamaye.Pending.MemberInterLocusAlong.mem_memberInterLocus_append_iff}\lean{Hironaka.MemberInterLocusAlong.map_nhds_eq_of_notMem,Hironaka.MemberInterLocusAlong.mem_memberInterLocus_geoTransform_iff}\upstreamlean{TranscendentalNakamaye.Pending.MemberInterLocusAlong.memberInterLocus_geoTransform_empty}\lean{Hironaka.MemberInterLocusAlong.mem_memberInterLocus_geoTransformAlong_step_iff}\upstreamlean{TranscendentalNakamaye.Pending.MemberInterLocusAlong.memberInterLocus_geoTransformAlong_transport,TranscendentalNakamaye.Pending.MemberInterLocusAlong.memberInterLocus_geoTransformAlong_restrict,TranscendentalNakamaye.Pending.MemberInterLocusAlong.memberInterLocus_geoTransformAlong_of_restrictEquiv,TranscendentalNakamaye.Pending.MemberInterLocusAlong.memberInterLocus_geoTransformAlong_of_wEquiv}
\mixedleanok
\begin{enumerate}
\item (Locality) Whether \(x\in L_N(E,S)\) depends only on the germs at \(x\) of \(N\), of \(S\) and of the members of \(E\); members not through \(x\) play no role.
\item (Preimages) If \(f\colon M'\to M\) maps the neighbourhoods of \(z\) onto a neighbourhood basis of \(f(z)\) (\(f_*\mathcal N_z=\mathcal N_{f(z)}\): a biholomorphism, the
inclusion of an open set, a blowup at a point off its centre), then \(z\in L_{f^{-1}N}((f^{-1}H)_{H\in E},f^{-1}S)\) if and only if \(f(z)\in L_N(E,S)\).
\item (Off the centre) Let \(\underline{\mathcal I}=(M,N,E,\mathfrak a,d)\) have \(N\) closed and \(\mathfrak a\) consist of holomorphic germs, \(E_0\) a sequence of closed sets, \(\pi\) the
blowup along \(C\) and \(\pi(z)\notin C\). Then \(z\) lies in the member locus of \(\operatorname{cosupp}\underline{\mathcal I}'\) along \(N'\) and the strict transforms of \(E_0\) if and
only if \(\pi(z)\in L_N(E_0,\operatorname{cosupp}\underline{\mathcal I})\). In particular a blowup along \(\emptyset\) pulls \(L\) back, and along a multiple test blow-up \(T\)
followed by \(\pi\), \(z\in L(\underline{\mathcal I}_{T\pi})\) if and only if \(\pi(z)\in L(\underline{\mathcal I}_T)\) whenever \(\pi(z)\notin C\) (\(N\) a closed submanifold,
\(\mathfrak a\) coherent, closed members).
\item (Transported and restricted towers) If \(\mathfrak a\) is coherent, \(L(\Phi^*\underline{\mathcal I}_{\Phi^*T})=\Phi_T^{-1}L(\underline{\mathcal I}_T)\) for every tower \(T\) and
biholomorphism \(\Phi\) (Definition~\ref{hir:tower-transport}); if moreover \(N\) is a closed submanifold and \(T\) a multiple test blow-up, \(L((\underline{\mathcal I}|_U)_{T|_U})\)
is the trace of \(L(\underline{\mathcal I}_T)\) on the open preimage of \(U\).
\item (Representatives) Restricted-equivalent marked ideals (\(N\) closed, coherent ideals) have the same \(L\) along every tower; equivalent marked ideals with
the same \(N\) have the same \(L\) along every multiple test blow-up.
\end{enumerate}
\end{proposition}
\begin{proof}
\mixedleanok
(1) Quantifying over the indices of the members through \(x\) is quantifying over these members; if near \(x\) the sets agree, the defining condition at \(x\)
reads the same for both, and a member not through \(x\) never enters it. (2) Membership of \(z\) in the preimages is membership of \(f(z)\), and a property
holds near \(z\) for the points \(f(y)\) if and only if it holds near \(f(z)\), since \(f_*\mathcal N_z=\mathcal N_{f(z)}\). A biholomorphism and an open inclusion have
this property; so does a blowup at \(z\) with \(\pi(z)\notin C\): it is continuous, and its holomorphic section \(s\) over \(M\setminus C\) with \(s(\pi(z))=z\)
(Lemma~\ref{bcalc:isomorphism-off-centre}) maps a neighbourhood of \(\pi(z)\) into any given neighbourhood \(V\) of \(z\), so \(\pi(V)\) is a neighbourhood of \(\pi(z)\).
(3) Near \(z\) the strict transforms of \(N\) and of the members are their preimages, and the cosupport of \(\underline{\mathcal I}'\) is the preimage of that of
\(\underline{\mathcal I}\) (Theorem~\ref{bcalc:transform-off-centre}); apply (1) and (2). The first \(\#E\) members of the transform along \(T\pi\) are the strict transforms
of those along \(T\) (Lemma~\ref{bcalc:boundary-count}, Lemma~\ref{hir:geo-transform-basic}). (4) For the empty tower this is (2) with the cosupport identity of
Lemma~\ref{hir:biholomorphism-transform} (as in Proposition~\ref{bcalc:codim-one-part-canonical}(2)); from the first blowup on the transforms are equal and \(\Phi_T\) is
the identity (Lemma~\ref{hir:tower-transport-resolution}). The transform along \(T|_U\) is the restriction of the transform along \(T\) to the open preimage of \(U\)
(Lemmas~\ref{hir:open-restriction-tower} and~\ref{hir:open-restriction-marked}), with the traces of the members; apply (2). (5) Restricted-equivalent marked ideals have
the same \(N\), \(E\) and cosupport and equal transforms from the first blowup on (Lemmas~\ref{hir:restricted-equivalence-basic} and~\ref{hir:marked-equivalence-basic});
equivalent ones have the same \(E\), equal cosupports along their multiple test blow-ups (Definition~\ref{hir:marked-equivalence}), and the same \(N_T\) and \(E_T\)
(Definition~\ref{mc:strict-along}).
\end{proof}

\begin{theorem}[Blowing up the level sets of the member locus]
\label{bcalc:member-locus-level}
\uses{bcalc:member-locus-def,bcalc:member-locus-canonical,bcalc:boundary-count,bcalc:boundary-count-drop-along,bcalc:boundary-count-canonical,bcalc:member-intersection-tower,bcalc:sub-divisor,mc:strict-along,hir:wellformed-tower,hir:geo-transform-basic,hir:test-tower,hir:marked-ideal}
\source{wlodarczyk05-hironaka,wlodarczyk09-analytic-resolution:page-0025}
\lean{Hironaka.MemberInterLocusAlong.boundaryCount_eq_of_forall_iff,Hironaka.MemberInterLocusAlong.forall_iff_of_forall_imp_of_boundaryCount_le,Hironaka.MemberInterLocusAlong.eventually_mem_memberInterLocus_level}\upstreamlean{TranscendentalNakamaye.Pending.MemberInterLocusAlong.memberInterLocus_level_geoTransform_eq_empty}\lean{Hironaka.MemberInterLocusAlong.strictAlong_append,Hironaka.MemberInterLocusAlong.strictAlong_inter_forall_mem,Hironaka.MemberInterLocusAlong.memberInterLocus_level_geoTransformAlong_subset}\upstreamlean{TranscendentalNakamaye.Pending.MemberInterLocusAlong.memberInterLocus_level_geoTransformAlong_eq_empty,TranscendentalNakamaye.Pending.MemberInterLocusAlong.memberInterLocus_level_geoTransformAlong_append_eq_empty}
\mixedleanok
Let \(\underline{\mathcal I}=(M,N,E,\mathfrak a,d)\) be well formed, \(E_0\) a sequence of members of \(E\) (in the algorithm \(E_0=E\), or at a later stage the strict transforms of
the original members), \(s_0\in\mathbf N\), and \(C=\{x\in L_N(E_0,\operatorname{cosupp}\underline{\mathcal I}):s(x)=s_0\}\), \(s\) counting the members of \(E_0\).
\begin{enumerate}
\item Near each \(x\in C\), every point of \(N\) on all the members of \(E_0\) through \(x\) lies in \(C\): \(C\supset N\cap H_{S(x)}\) near \(x\).
\item If \(E_0\) is a subsequence of \(E\), \(C\) is an admissible centre and \(\pi\) the blowup along \(C\), then no point of the transform with \(s'=s_0\) lies in the member
locus of \(\operatorname{cosupp}\underline{\mathcal I}'\) along \(N'\) and the strict transforms of \(E_0\).
\item Along every multiple test blow-up \(T\) with projection \(p\), the points of the member locus of \(\operatorname{cosupp}\underline{\mathcal I}_T\) along \(N_T\) and the strict
transforms of \(E_0\) with \(s_T=s_0\) lie in \(p^{-1}(\overline C)\). Hence if \(C=\emptyset\) this set is empty along every multiple test blow-up, and after the blowup of (2)
it is empty along every multiple test blow-up that starts with that blowup.
\end{enumerate}
\end{theorem}
\begin{proof}
\mixedleanok
(1) The members of \(E_0\) not through \(x\) are closed, so near \(x\) a point on all the members through \(x\) lies on exactly these members; its count is \(s_0\), and
near it \(N\cap H_{S(x)}\subset\operatorname{cosupp}\underline{\mathcal I}\) still holds, since this holds on a neighbourhood of \(x\).
(2) Let \(z\) be such a point. If \(\pi(z)\notin C\), \(\pi(z)\) lies in \(L_N(E_0,\operatorname{cosupp}\underline{\mathcal I})\) with \(s(\pi(z))=s'(z)=s_0\)
(Proposition~\ref{bcalc:member-locus-canonical}(3), Lemma~\ref{bcalc:boundary-count-canonical}(1)), that is \(\pi(z)\in C\), a contradiction. If \(x=\pi(z)\in C\), then
\(z\in N'\) and near \(x\) the centre contains \(N\cap\bigcap_{x\in H}H\) by (1); a subsequence of a simple normal crossings divisor transverse to \(N\) is one, with which
\(C\) has normal crossings (Lemma~\ref{bcalc:sub-divisor}), so \(s'(z)<s(x)=s_0\) (Theorem~\ref{bcalc:boundary-count-drop-along}(2)), a contradiction.
(3) Induction on \(T\); for the empty tower the set is \(C\). Let \(T'\) be \(T\) followed by the blowup \(\pi\) along \(C_T\), and \(z\) a point of the set at stage \(T'\). The
intersection of \(N_{T'}\) with the strict transforms of the members of \(E_0\) through \(z\) is the strict transform of the intersection \(A\) of \(N\) with these members
(Lemma~\ref{bcalc:member-intersection-tower}), that is the closure of \(\pi^{-1}(A_T\setminus C_T)\) (Definition~\ref{mc:strict-along}), and it contains \(z\). Take a point
\(y\in\pi^{-1}(A_T\setminus C_T)\) in a given neighbourhood of \(z\), so close to \(z\) that \(s_{T'}(y)\le s_{T'}(z)\) (Lemma~\ref{bcalc:boundary-count}) and that the defining condition
of the member locus at \(z\) holds near \(y\). Then \(y\) lies on all the members through \(z\) and on no more, so \(y\) is in the member locus with \(s_{T'}(y)=s_0\); as
\(\pi(y)\notin C_T\), \(\pi(y)\) is a point of the set at stage \(T\) (Proposition~\ref{bcalc:member-locus-canonical}(3), Lemma~\ref{bcalc:boundary-count-canonical}(1)), so
\(p'(y)\in\overline C\) by induction. As \(p'^{-1}(\overline C)\) is closed, \(z\in p'^{-1}(\overline C)\). For the last claim, a multiple test blow-up starting with \(\pi\) is \(\pi\)
followed by a multiple test blow-up of the well-formed transform \(\underline{\mathcal I}'\) (Lemmas~\ref{hir:geo-transform-basic} and~\ref{hir:wellformed-tower}), the strict
transforms of the members of \(E_0\) are members of \(E'\), strict transforms along concatenated towers compose, and the set of (2) is empty.
\end{proof}

\begin{proposition}[The Step 1aa centre at the maximal level is admissible]
\label{bcalc:member-locus-centre}
\uses{bcalc:member-locus-def,bcalc:member-locus-level,bcalc:member-intersection-along,bcalc:sub-divisor,bcalc:boundary-count,bcalc:normal-crossings-local,bcalc:analytic-fill,bcalc:codim-one-local-along,hir:giraud-along,ana:analytic-preimage,fpgeom:complex-analytic-subvariety,lag:regular-locus-dense,rem:dimension,hir:transform}
\source{wlodarczyk05-hironaka,wlodarczyk09-analytic-resolution:page-0025}
\lean{Hironaka.MemberInterLocusAlong.setOf_forall_mem_eq_memberInter_of_mem,Hironaka.MemberInterLocusAlong.eventually_mem_level_iff,Hironaka.MemberInterLocusAlong.isAdmissibleCentre_memberInterLocus_level}
\leanok
Let \(\underline{\mathcal I}=(M,N,E,\mathfrak a,d)\) be well formed with \(d\ge1\) and \(n>0\), \(E_0\) a subsequence of \(E\), and \(s_0\) an upper bound of the number \(s\) of
members of \(E_0\) through the points of the cosupport. Then \(C=\{x\in L_N(E_0,\operatorname{cosupp}\underline{\mathcal I}):s(x)=s_0\}\) is an admissible centre: near
each of its points \(x\) it is \(N\cap H_{S(x)}\), the intersection of \(N\) with the members of \(E_0\) through \(x\). (For an open \(W\) on which \(s\le s_0\) on the cosupport,
applied to \(\underline{\mathcal I}|_W\), this is property (1) of Definition~\ref{hir:centre-function} for the value \((s_0,\infty)\), and it discharges the admissibility
hypothesis of Theorem~\ref{bcalc:member-locus-level}(2).)
\end{proposition}
\begin{proof}
\leanok
Near \(x\in C\), every point of \(N\cap H_{S(x)}\) is in \(C\) (Theorem~\ref{bcalc:member-locus-level}(1)); conversely a point \(y\in C\) near \(x\) lies on no member missing \(x\)
(the members are closed) and has \(s(y)=s_0=s(x)\), so it lies on exactly the members through \(x\) and \(y\in N\cap H_{S(x)}\). The set \(N\cap H_{S(x)}\) is a closed
submanifold with normal crossings with \(E\) (Lemma~\ref{bcalc:member-intersection-along}, with \(E_0\) simple normal crossings and transverse to \(N\),
Lemma~\ref{bcalc:sub-divisor}); both properties are local (Lemma~\ref{bcalc:normal-crossings-local}, Remark~\ref{rem:dimension}), so they pass to \(C\) once \(C\) is closed.
\(C\subset\operatorname{cosupp}\underline{\mathcal I}\) by definition. Closedness: let \(x\) be a limit of points of \(C\). Then \(x\) lies in the cosupport, which is closed
(Lemma~\ref{hir:giraud-along}); \(s\) is upper semicontinuous (Lemma~\ref{bcalc:boundary-count}), so \(s(x)\ge s_0\), and \(s(x)\le s_0\) by hypothesis. As above, the points
\(y\in C\) near \(x\) lie on exactly the members through \(x\), and near each of them \(N\cap H_{S(x)}=N\cap H_{S(y)}\subset\operatorname{cosupp}\underline{\mathcal I}\). Let
\(H=N\cap H_{S(x)}\), regular of some dimension \(k\) at \(x\), hence at its points near \(x\) (Theorem~\ref{lag:regular-locus-dense}). The set
\(A=\operatorname{cosupp}\underline{\mathcal I}\cap H\) is analytic (Lemma~\ref{bcalc:codim-one-local-along}, Lemma~\ref{ana:analytic-preimage},
Definition~\ref{fpgeom:complex-analytic-subvariety}), contained in \(H\), and agrees with \(H\) near the points \(y\), which accumulate at \(x\); so \(\dim_xA\ge k\), and the filling
lemma (Lemma~\ref{bcalc:analytic-fill}) gives \(A=H\) near \(x\): \(N\cap H_{S(x)}\subset\operatorname{cosupp}\underline{\mathcal I}\) near \(x\), that is \(x\in C\). The admissible
centre is Definition~\ref{hir:transform}.
\end{proof}

\begin{proposition}[Off the member locus the restriction to the intersection of the members is nonzero]
\label{bcalc:member-restriction-nonzero}
\uses{bcalc:member-locus-def,bcalc:member-locus-canonical,bcalc:contact-restriction-nonzero,mc:coefficient-restriction-along,mc:homogenized-along-basic,hir:exponent-along,hir:cosupport}
\source{wlodarczyk05-hironaka,wlodarczyk09-analytic-resolution:page-0025}
\lean{Hironaka.MemberInterLocusAlong.isNonzeroAlongAt_of_notMem_memberInterLocus,Hironaka.MemberInterLocusAlong.isNonzeroAlongAt_coeffMarkedAlong_of_notMem_memberInterLocus}
\leanok
Let \(y\in N\), \(E\) a finite sequence of sets, \(A=N\cap H_{S(y)}\) the intersection of \(N\) with the members through \(y\), \(\mathfrak b\) coherent, and suppose that near \(y\)
\(\operatorname{cosupp}(\mathfrak b|_A,\nu)\subset S\). If \(y\notin L_N(E,S)\), some germ of \(\mathfrak b\) at \(y\) does not vanish identically on \(A\) near \(y\). In particular, for
\(\underline{\mathcal I}=(M,N,E,\mathfrak a,d)\) with \(N\) a closed submanifold, \(\mathfrak a\) coherent and \(d\ge1\), and \(y\in N\) off \(L_N(E_0,\operatorname{cosupp}\underline{\mathcal I})\),
the coefficient ideal along \(N\) restricted to \(A\) (Lemma~\ref{mc:coefficient-restriction-along}, the restricted marked ideal of Step~1ab with any sequence of members) is
nonzero along \(A\) near \(y\): together with Theorem~\ref{bcalc:member-locus-level}(3), at every stage after Step~1aa the restricted marked ideals of Step~1ab at the points
with \(s=s_0\) satisfy the hypothesis of Definition~\ref{hir:centre-function}.
\end{proposition}
\begin{proof}
\leanok
If every germ of \(\mathfrak b\) at \(y\) vanished on \(A\) near \(y\), then, \(\mathfrak b\) being coherent, at every point \(z\) near \(y\) every germ of \(\mathfrak b\) would vanish on
\(A\) near \(z\) (proof of Proposition~\ref{bcalc:contact-restriction-nonzero}(1)), so every point of \(A\) near \(y\) would lie in \(\operatorname{cosupp}(\mathfrak b|_A,\nu)\subset S\)
(Definition~\ref{hir:cosupport}); with \(y\in A\) this is \(y\in L_N(E,S)\). For the coefficient ideal, its cosupport on \(A\subset N\) is \(\operatorname{cosupp}\underline{\mathcal I}\cap A\)
(Lemma~\ref{mc:coefficient-restriction-along}) and it is coherent (Lemma~\ref{mc:homogenized-along-basic}).
\end{proof}

\begin{proposition}[The loci of Step 1 along isomorphic towers]
\label{bcalc:locus-tower-isomorphism}
\uses{hir:tower-isomorphism,mc:tower-isomorphism-transform,bcalc:codim-one-part-canonical,bcalc:boundary-count-canonical,bcalc:member-locus-def,bcalc:member-locus-canonical,hir:biholomorphism-invariance,hir:open-restriction-tower}
\source{wlodarczyk05-hironaka,bierstone-milman08-functoriality}
\upstreamlean{TranscendentalNakamaye.Pending.LocusTowerIsomorphism.codimOnePartAlong_geoTransformAlong_isTowerIso,TranscendentalNakamaye.Pending.LocusTowerIsomorphism.boundaryCount_take_E_geoTransformAlong_isTowerIso,TranscendentalNakamaye.Pending.LocusTowerIsomorphism.memberInterLocus_geoTransformAlong_isTowerIso,TranscendentalNakamaye.Pending.LocusTowerIsomorphism.memberInterLocus_level_geoTransformAlong_isTowerIso,TranscendentalNakamaye.Pending.LocusTowerIsomorphism.isAdmissibleCentre_geoTransformAlong_isTowerIso_iff}
\upstreamleanok
Let \(\underline{\mathcal I}=(M,N,E,\mathfrak a,d)\) have \(N\) a closed submanifold and \(\mathfrak a\) coherent, \(T\) a multiple test blow-up of \(\underline{\mathcal I}\), \(\Phi\colon M'\to M\) a
biholomorphism and \(\chi\colon T'\to T\) an isomorphism of towers over \(\Phi\) with last-stage biholomorphism \(\chi_r\) (Definition~\ref{hir:tower-isomorphism}). For
\(k\in\mathbf N\) write \(s_T\) for the number of the first \(k\) members of \(E_T\) through a point and \(L_k(\underline{\mathcal I}_T)\) for the member locus of
\(\operatorname{cosupp}\underline{\mathcal I}_T\) along \(N_T\) and these members (Definition~\ref{bcalc:member-locus-def}; in the algorithm \(k=\#E\)). Then, with
\(\underline{\mathcal I}'=\Phi^*\underline{\mathcal I}\),
\begin{enumerate}
\item \(R(\underline{\mathcal I}'_{T'})=\chi_r^{-1}R(\underline{\mathcal I}_T)\) for the codimension-one parts along \(N'_{T'}\) and \(N_T\) (Proposition~\ref{bcalc:codim-one-part-canonical});
\item \(s_{T'}=s_T\circ\chi_r\);
\item \(L_k(\underline{\mathcal I}'_{T'})=\chi_r^{-1}L_k(\underline{\mathcal I}_T)\), and for every \(s_0\) the level-\(s_0\) parts (the centres of Step~1aa) correspond:
\(L_k(\underline{\mathcal I}'_{T'})\cap\{s_{T'}=s_0\}=\chi_r^{-1}(L_k(\underline{\mathcal I}_T)\cap\{s_T=s_0\})\);
\item for a subset \(C\) of the last stage of \(T\), \(\chi_r^{-1}C\) is an admissible centre of \(\underline{\mathcal I}'_{T'}\) if and only if \(C\) is one of \(\underline{\mathcal I}_T\).
\end{enumerate}
These are the blowup-calculus halves of property (6) of Definition~\ref{hir:centre-function} for the values \((0,\infty)\) and \((s,\infty)\) of Steps~1ba and~1aa and for
the count \(s\) of Proposition~\ref{hir:maximal-order-step}; the transport \(\Phi^*T\) is the case \(\chi_i=\mathrm{id}\) for \(i\ge1\)
(Propositions~\ref{bcalc:codim-one-part-canonical}(4) and~\ref{bcalc:member-locus-canonical}(4), Lemma~\ref{bcalc:boundary-count-canonical}(2), which hold even without the
hypothesis that \(T\) be a multiple test blow-up). Applied to the transform at a later stage \(T_0\) of the algorithm, the count \(k\) stays the number of the original members.
\end{proposition}
\begin{proof}
\upstreamleanok
By Lemma~\ref{mc:tower-isomorphism-transform}, \(\underline{\mathcal I}'_{T'}=\chi_r^*(\underline{\mathcal I}_T)\): its submanifold is \(\chi_r^{-1}N_T\), its members are the preimages
\(\chi_r^{-1}H\) of the members of \(E_T\) in the same order, and its ideal is \(\chi_r^*\mathfrak a_T\), coherent since \(\mathfrak a_T\) is (Lemma~\ref{hir:open-restriction-tower}).
(1) is Proposition~\ref{bcalc:codim-one-part-canonical}(2) for \(\chi_r\). (2) is the re-indexing of Lemma~\ref{bcalc:boundary-count-canonical}: the first \(k\) members of
\(\chi_r^{-1}E_T\) are the preimages of the first \(k\) members of \(E_T\). (3) The cosupport of \(\chi_r^*(\underline{\mathcal I}_T)\) is \(\chi_r^{-1}\operatorname{cosupp}\underline{\mathcal I}_T\)
(Lemma~\ref{hir:biholomorphism-invariance}), and a biholomorphism maps neighbourhood filters onto neighbourhood filters, so Proposition~\ref{bcalc:member-locus-canonical}(2) applies;
the level sets follow with (2). (4) is the invariance of admissible centres under biholomorphisms (Lemma~\ref{hir:biholomorphism-invariance}) for \(\chi_r\) and
\(\underline{\mathcal I}_T\).
\end{proof}

\paragraph{The domain of a resolution invariant.} The inductive invariant of Definition~\ref{hir:centre-function} is defined on the well-formed marked
ideals \(\underline{\mathcal I}=(M,N,E,\mathfrak a,d)\) with \(\dim N\le k\), \(d\ge1\) and \(\mathfrak a|_N\) vanishing identically on no nonempty open subset of \(N\). The
general step forms the companion ideal from the transform at a later stage of a canonical segment (Włodarczyk: ``Then we repeat the procedure for
\(({\mathcal I}_{r_1},\mu)\)''), Step~1 forms its restricted marked ideals there, and the canonicity properties (5) and (6) are applied to restrictions and
pullbacks. The two statements below show that the domain is stable under these operations; the point is that a transform stays nonzero along the
strict transform over the centre, for every admissible centre.

\begin{lemma}[A transform stays nonzero along the strict transform]
\label{bcalc:nonzero-along-transform}
\uses{hir:exponent-along,hir:geo-transform,hir:geo-transform-coherent,hir:test-tower,hir:open-restriction-tower,bcalc:strict-submanifold,bcalc:isomorphism-off-centre,bcalc:monomial-transform-along-off-centre,bcalc:contact-restriction-nonzero,hir:nonmonomial-step-along,hir:companion-reduction}
\source{wlodarczyk05-hironaka}
\lean{Hironaka.InvDomainTowerAlong.isNonzeroAlongAt_geoTransform_of_notMem,Hironaka.InvDomainTowerAlong.isNonzeroAlongAt_geoTransform_of_isAdmissibleCentre,Hironaka.InvDomainTowerAlong.isNonzeroAlongAt_geoTransformAlong}
\leanok
Let \(\underline{\mathcal I}=(M,N,E,\mathfrak a,d)\) have \(N\) a closed submanifold and \(\mathfrak a\) coherent and nonzero along \(N\) at every point of \(N\) (some germ of
\(\mathfrak a\) at the point does not vanish identically on \(N\) near it, Definition~\ref{hir:exponent-along}). Let \(\pi\) be the blowup along an admissible centre \(C\)
and \(\underline{\mathcal I}'=(M',N',E',\mathfrak a',d)\) the transform. Then \(\mathfrak a'\) is nonzero along \(N'\) at every point of \(N'\); at a point \(z\) with
\(\pi(z)\notin C\) this needs only that \(\mathfrak a\) be nonzero along \(N\) at \(\pi(z)\). Consequently, along every multiple test blow-up \(T\) of \(\underline{\mathcal I}\),
\(\mathfrak a_T\) is nonzero along \(N_T\) at every point of \(N_T\).
\end{lemma}
\begin{proof}
\leanok
Let \(\pi(z)\notin C\) and let \(a\) be a germ of \(\mathfrak a\) at \(\pi(z)\) not vanishing identically on \(N\) near \(\pi(z)\). Then \(a\circ\pi\) is a germ of
\(\pi^*\mathfrak a\) at \(z\) (Lemma~\ref{hir:companion-reduction}), and off the centre \(\operatorname{Res}_{N'}\mathfrak a'=\operatorname{Res}_{N'}\pi^*\mathfrak a\) at \(z\)
(Lemma~\ref{bcalc:monomial-transform-along-off-centre}): on \(N'\) near \(z\), \(a\circ\pi\) is a combination of germs of \(\mathfrak a'\) at \(z\). The holomorphic section
\(s\) of \(\pi\) over \(M\setminus C\) with \(s(\pi(z))=z\) (Lemma~\ref{bcalc:isomorphism-off-centre}) maps the points of \(N\setminus C\) near \(\pi(z)\) where \(a\ne0\), which
accumulate at \(\pi(z)\), to points of \(\pi^{-1}(N\setminus C)\subset N'\) near \(z\) where \(a\circ\pi\ne0\). If every germ of \(\mathfrak a'\) at \(z\) vanished on \(N'\) near
\(z\), so would the combination \(a\circ\pi\); hence \(\mathfrak a'\) is nonzero along \(N'\) at \(z\) (Lemma~\ref{hir:nonmonomial-step-along}, whose proof uses the same
step). Let now \(\pi(z)\in C\), and suppose every germ of \(\mathfrak a'\) at \(z\) vanishes identically on \(N'\) near \(z\). The ideal \(\mathfrak a'\) is coherent
(Lemma~\ref{hir:geo-transform-coherent}); taking finitely many sections generating it near \(z\), every germ of \(\mathfrak a'\) at every point \(z'\) near \(z\) vanishes on
\(N'\) near \(z'\) (proof of Proposition~\ref{bcalc:contact-restriction-nonzero}(1)). But \(N'=\overline{\pi^{-1}(N\setminus C)}\) (Theorem~\ref{bcalc:strict-submanifold}), so
there is such a \(z'\) with \(\pi(z')\in N\setminus C\), at which \(\mathfrak a'\) is nonzero along \(N'\) by the first case: a contradiction. For a multiple test blow-up,
induction on its length: along it \(N_T\) stays a closed submanifold and \(\mathfrak a_T\) coherent (Lemma~\ref{hir:open-restriction-tower}), and each centre is admissible
(Definition~\ref{hir:test-tower}).
\end{proof}

\begin{proposition}[The domain of a resolution invariant is stable]
\label{bcalc:inv-domain-stable}
\uses{hir:centre-function,bcalc:nonzero-along-transform,bcalc:strict-submanifold,hir:wellformed-tower,hir:geo-transform-basic,hir:test-tower,hir:open-restriction-tower,hir:biholomorphism-invariance,hir:open-restriction-marked,hir:open-restriction-ideal,mc:transform-source-naturality,bcalc:regular-point-transport,rem:open-subset-analytic,hir:exponent-along}
\source{wlodarczyk05-hironaka,bierstone-milman08-functoriality}
\lean{Hironaka.InvDomainTowerAlong.dimLE_geoStrictTransform,Hironaka.InvDomainTowerAlong.dimLE_geoTransformAlong,Hironaka.InvDomainTowerAlong.isInvDomain_geoTransformAlong,Hironaka.InvDomainTowerAlong.isInvDomain_pullback,Hironaka.InvDomainTowerAlong.isInvDomain_pullback_val}
\leanok
Say that \(\underline{\mathcal I}=(M,N,E,\mathfrak a,d)\) lies in the \emph{domain in dimension \(\le k\)} if it lies in the domain of a resolution invariant in dimension \(\le k\)
(Definition~\ref{hir:centre-function}): well formed, \(N\) regular of dimension \(\le k\) at each of its points, \(d\ge1\), and \(\mathfrak a\) nonzero along \(N\) at every point
of \(N\).
\begin{enumerate}
\item If \(N\) is closed, \(C\subset N\), \(\pi\) is the blowup along \(C\) and \(N\) is regular of dimension \(\le k\) at each of its points, so is the strict transform \(N'\): at
\(z\in N'\) of the dimension of \(N\) at \(\pi(z)\).
\item If \(\underline{\mathcal I}\) lies in the domain in dimension \(\le k\), so does \(\underline{\mathcal I}_T\) for every multiple test blow-up \(T\) of \(\underline{\mathcal I}\).
\item If \(\underline{\mathcal I}\) lies in the domain in dimension \(\le k\), so do \(\Phi^*\underline{\mathcal I}\) for a biholomorphism \(\Phi\colon M'\to M\) and
\(\underline{\mathcal I}|_U\) for an open \(U\subset M\).
\end{enumerate}
\end{proposition}
\begin{proof}
\leanok
(1) Let \(N\) be regular of dimension \(j\le k\) at \(\pi(z)\in N\) (\(N'\subset\pi^{-1}(N)\) as \(N\) is closed). If \(\pi(z)\notin C\), \(\pi\) is a local biholomorphism at \(z\) and
\(N'\) is \(\pi^{-1}(N)\) near \(z\); if \(\pi(z)\in C\), the centre is regular at \(\pi(z)\) of a dimension \(k_C\le j\) (a submanifold of a submanifold), and \(N'\) is regular of
dimension \(j\) at \(z\) (Theorem~\ref{bcalc:strict-submanifold}, in its pointwise form). (2) Induction on \(T\): well-formedness passes along multiple test blow-ups
(Lemma~\ref{hir:wellformed-tower}), the mark does not change (Lemma~\ref{hir:geo-transform-basic}), \(N_T\) stays a closed submanifold (Lemma~\ref{hir:open-restriction-tower}) and
the next centre lies in the cosupport, hence in \(N_T\) (Definition~\ref{hir:test-tower}), so (1) applies; nonzero along \(N_T\) is Lemma~\ref{bcalc:nonzero-along-transform}. (3)
Well-formedness is invariant under biholomorphisms and restricts to open sets (Lemmas~\ref{hir:biholomorphism-invariance} and~\ref{hir:open-restriction-marked}); regular points and
their dimensions correspond (Lemma~\ref{bcalc:regular-point-transport}, Remark~\ref{rem:open-subset-analytic}); and a germ \(a\) of \(\mathfrak a\) at \(\Phi(p)\) not vanishing
identically on \(N\) near \(\Phi(p)\) pulls back to the germ \(a\circ\Phi\) of \(\Phi^*\mathfrak a\) (Lemma~\ref{mc:transform-source-naturality}), which does not vanish identically on
\(\Phi^{-1}(N)\) near \(p\) since \(\Phi\) is a homeomorphism; likewise for the inclusion of \(U\), an open embedding (Lemma~\ref{hir:open-restriction-ideal}).
\end{proof}
These are the inputs that let the general step (Proposition~\ref{hir:general-step}) and Step~1 (Proposition~\ref{hir:maximal-order-step}) apply an invariant to marked ideals
built at a later stage of a canonical segment.

\begin{proposition}[The ideal of a smooth centre inside a submanifold]
\label{bcalc:smooth-centre-locus}
\uses{hir:cosupport,hir:cosupport-basic,hir:transform,hir:snc-basic,bcalc:codim-one-part-along-def,bcalc:codim-one-part-along,lag:regular-point-basics,bcalc:analytic-fill,ana:cartan-coherence,hir:exponent-along,bcalc:contact-restriction-nonzero,hir:centre-function,bcalc:inv-domain-stable}
\source{wlodarczyk05-hironaka}
\lean{Hironaka.SmoothCentreLocus.holomorphicAt_of_mem_vanishingIdeal,Hironaka.SmoothCentreLocus.zeroLocus_vanishingIdeal,Hironaka.SmoothCentreLocus.cosupport_vanishingIdeal_one,Hironaka.SmoothCentreLocus.cosupp_vanishingIdeal}\upstreamlean{TranscendentalNakamaye.Pending.SmoothCentreLocus.isAdmissibleCentre_vanishingIdeal,TranscendentalNakamaye.Pending.SmoothCentreLocus.codimOnePartAlong_vanishingIdeal_subset,TranscendentalNakamaye.Pending.SmoothCentreLocus.mem_codimOnePartAlong_vanishingIdeal_iff,TranscendentalNakamaye.Pending.SmoothCentreLocus.codimOnePartAlong_vanishingIdeal_eq_self,TranscendentalNakamaye.Pending.SmoothCentreLocus.codimOnePartAlong_vanishingIdeal_eq_empty,TranscendentalNakamaye.Pending.SmoothCentreLocus.frequently_mem_diff,TranscendentalNakamaye.Pending.SmoothCentreLocus.isNonzeroAlongAt_vanishingIdeal,TranscendentalNakamaye.Pending.SmoothCentreLocus.isInvDomain_vanishingIdeal}
\mixedleanok
Let \(Z\subset N\) be closed submanifolds of \(M\) and \(\underline{\mathcal I}_Z=(M,N,\emptyset,\mathcal I_Z,1)\), with \(\mathcal I_Z\) the ideal of germs of holomorphic
functions vanishing on \(Z\).
\begin{enumerate}
\item \(\operatorname{cosupp}(\mathcal I_Z|_N,1)=N\cap Z=Z\), and \(Z\) is an admissible centre of \(\underline{\mathcal I}_Z\).
\item A point \(x\in Z\) where \(Z\) and \(N\) are regular of dimensions \(j\) and \(j'\) lies in the codimension-one part \(Z_1\) of \(\operatorname{cosupp}(\mathcal I_Z|_N,1)\) along
\(N\) (Definition~\ref{bcalc:codim-one-part-along-def}) if and only if \(j'=j+1\). Hence \(Z_1=Z\) if \(Z\) has codimension one in \(N\) at each of its points, and
\(Z_1=\emptyset\) if it has codimension at least two at each of its points.
\item If \(\dim_xZ<\dim_xN\) at every \(x\in Z\), then \(\mathcal I_Z\) is nonzero along \(N\) at every point of \(N\); if moreover \(\dim N\le k\), then
\(\underline{\mathcal I}_Z\) lies in the domain of a resolution invariant in dimension \(\le k\) (Definition~\ref{hir:centre-function}, Proposition~\ref{bcalc:inv-domain-stable}).
\end{enumerate}
\end{proposition}
\begin{proof}
\mixedleanok
(1) The germs of \(\mathcal I_Z\) are holomorphic, so the cosupport with mark \(1\) is the zero locus on \(N\) (Lemma~\ref{hir:cosupport-basic}). Every germ vanishes at the
points of \(Z\); at \(x\notin Z\) the constant \(1\) is a germ of \(\mathcal I_Z\), \(Z\) being closed. So the zero locus is \(Z\). \(Z\) is a closed submanifold contained in
the cosupport, and has normal crossings with the empty divisor (Lemma~\ref{hir:snc-basic}): an admissible centre (Definition~\ref{hir:transform}).
(2) By (1), \(x\in Z_1\) if and only if for some \(k\) the set \(N\) is regular of dimension \(k+1\) and \(Z\) of dimension \(k\) at \(x\); a set is regular of at most one
dimension at a point (Lemma~\ref{lag:regular-point-basics}), so this means \(j'=j+1\). Points of \(Z_1\) lie in the cosupport \(Z\) (Theorem~\ref{bcalc:codim-one-part-along}).
(3) At \(y\in N\setminus Z\) the constant \(1\) is a germ of \(\mathcal I_Z\) not vanishing on \(N\). Let \(x\in Z\). If \(N\subset Z\) near \(x\), then \(N\) and \(Z\) agree near \(x\)
and are regular of the same dimension there (Lemma~\ref{bcalc:analytic-fill}, regularity depends only on the germ), which is excluded; so points of \(N\setminus Z\)
accumulate at \(x\). The ideal \(\mathcal I_Z\) is coherent (Theorem~\ref{ana:cartan-coherence}; for a submanifold, in flattening coordinates it is generated by the
normal coordinates). If every germ of \(\mathcal I_Z\) at \(x\) vanished on \(N\) near \(x\), then every germ at every point near \(x\) would vanish on \(N\) near that point
(proof of Proposition~\ref{bcalc:contact-restriction-nonzero}(1)), in particular \(1\) at a nearby point of \(N\setminus Z\): a contradiction (Definition~\ref{hir:exponent-along}).
\(\underline{\mathcal I}_Z\) is well formed (closed submanifold \(N\), coherent ideal, the empty divisor is simple normal crossings and transverse to \(N\),
Lemma~\ref{hir:snc-basic}), its mark is \(1\), and \(\dim N\le k\).
\end{proof}
In Lemma~\ref{hir:smooth-centre} this is applied first with \(N=U\) and then, along the induction on the codimension, with \(N=V(u)\cap U\); Włodarczyk's remark
on the weak embedded desingularization (``Its support \(\operatorname{supp}(\mathcal I_Y,1)\) is equal to \(Y\)'', and at a smooth point Steps 2a and 1b pass to a
hypersurface until the algorithm ends in Step 1ba) is the same computation.

\paragraph{The domain of Step~1ab's restricted marked ideals.} In Step~1a of the maximal-order step (Proposition~\ref{hir:maximal-order-step}), at a point
\(y\) of the cosupport on \(s_0\ge1\) of the original members, the value is \((s_0,\operatorname{inv}(x))\) with the invariant of the restriction of the coefficient ideal
to the intersection of \(N\) with the members through \(y\). Włodarczyk forms these restrictions after Step~1aa: ``After the blow-ups in Step 1aa we arrive at
\(X_p\) for which no \(H^s_\alpha\) is contained in \(\operatorname{supp}(\overline{\mathcal J}_p,\overline\mu)\) [\dots] Construct the canonical resolutions of
\(\overline{\mathcal J}_{p|H^s_\alpha}\).'' The two statements below show that these restrictions lie in the domain of the inductive invariant
(Definition~\ref{hir:centre-function}) in dimension \(\le k-1\).

\begin{lemma}[An intersection of members drops the dimension]
\label{bcalc:member-intersection-dim}
\uses{bcalc:snc-simultaneous,bcalc:member-intersection-along,hir:snc,lag:regular-point-basics}
\source{wlodarczyk05-hironaka,bierstone-milman08-functoriality}
\lean{Hironaka.InvDomainMemberRestriction.exists_regularAtDim_inter_memberInter_lt,Hironaka.InvDomainMemberRestriction.dimLE_inter_memberInter}
\leanok
Let \(E\) be a simple normal crossings divisor on \(M\) transverse to the closed set \(N\) (\(n>0\)), \(S\) a nonempty set of indices of members and
\(x\in N\cap H_S\). If \(N\) is regular of dimension \(j\) at \(x\), then \(N\cap H_S\) is regular at \(x\) of some dimension \(j'<j\). Hence if \(N\) is regular of
dimension \(\le k\) at each of its points, \(N\cap H_S\) is regular of dimension \(\le k-1\) at each of its points.
\end{lemma}
\begin{proof}
\leanok
Take coordinates at \(x\) adapted to \(N\cap H_S\subset N\) and \(E\) (Lemma~\ref{bcalc:snc-simultaneous}; \(N\cap H_S\) has normal crossings with \(E\),
Lemma~\ref{bcalc:member-intersection-along}): near \(x\), \(N\cap H_S=\mathbf C^{k_C}\times0\), \(N=\mathbf C^{k_N}\times0\) and each member \(H_i\) through \(x\) is
\(\{w_{s_i}=0\}\) with \(s_i<k_N\). So \(N\cap H_S\) and \(N\) are regular at \(x\) of dimensions \(\min(k_C,n)\) and \(\min(k_N,n)\), and \(j=\min(k_N,n)\) since a set is regular of at
most one dimension at a point (Lemma~\ref{lag:regular-point-basics}). If \(\min(k_C,n)\ge\min(k_N,n)\), let \(c\) be the coordinate vector of \(x\), \(i\in S\) and
\(v=c+\varepsilon e_{s_i}\) with \(\varepsilon>0\) so small that \(v\) lies in the coordinate domain. As \(c\in\mathbf C^{k_N}\times0\), \(c_{s_i}=0\) and \(s_i<k_N\), the vector
\(v\) lies in \(\mathbf C^{k_N}\times0\subset\mathbf C^{k_C}\times0\), so the point with coordinates \(v\) lies in \(N\cap H_S\subset H_i\); but \(v_{s_i}=\varepsilon\ne0\), so it does not
lie on \(H_i\): a contradiction. The second statement follows at each point.
\end{proof}

\begin{proposition}[Step~1ab's restricted marked ideals lie in the domain]
\label{bcalc:member-restriction-domain}
\uses{hir:centre-function,bcalc:member-intersection-dim,bcalc:member-intersection-drop-along,bcalc:member-restriction-nonzero,bcalc:member-locus-def,bcalc:member-locus-level,bcalc:boundary-count,bcalc:member-intersection-tower,bcalc:inv-domain-stable,mc:coefficient-restriction-along,mc:coefficient-restriction-drop-along,mc:homogenized-along-basic,hir:open-restriction-marked,hir:open-restriction-ideal,hir:wellformed-tower,hir:geo-transform-basic,hir:exponent-along}
\source{wlodarczyk05-hironaka,wlodarczyk09-analytic-resolution:page-0025}
\lean{Hironaka.InvDomainMemberRestriction.isInvDomain_pullback_val_of_forall,Hironaka.InvDomainMemberRestriction.memberRestrict}\upstreamlean{TranscendentalNakamaye.Pending.InvDomainMemberRestriction.memberRestrict_eq}\lean{Hironaka.InvDomainMemberRestriction.memberRestrict_N,Hironaka.InvDomainMemberRestriction.memberRestrict_E,Hironaka.InvDomainMemberRestriction.memberRestrict_ideal,Hironaka.InvDomainMemberRestriction.memberRestrict_d,Hironaka.InvDomainMemberRestriction.one_le_memberRestrict_d,Hironaka.InvDomainMemberRestriction.isWellFormed_memberRestrict,Hironaka.InvDomainMemberRestriction.dimLE_memberRestrict,Hironaka.InvDomainMemberRestriction.cosupp_memberRestrict,Hironaka.InvDomainMemberRestriction.isNonzeroAlongAt_memberRestrict,Hironaka.InvDomainMemberRestriction.membersThrough,Hironaka.InvDomainMemberRestriction.mem_membersThrough,Hironaka.InvDomainMemberRestriction.memberInter_membersThrough,Hironaka.InvDomainMemberRestriction.membersThrough_nonempty,Hironaka.InvDomainMemberRestriction.exists_opens_forall_mem_imp}\upstreamlean{TranscendentalNakamaye.Pending.InvDomainMemberRestriction.isInvDomain_memberRestrict_membersThrough}\lean{Hironaka.InvDomainMemberRestriction.mem_E_geoTransformAlong_of_mem_map}\upstreamlean{TranscendentalNakamaye.Pending.InvDomainMemberRestriction.isInvDomain_memberRestrict_geoTransformAlong,TranscendentalNakamaye.Pending.InvDomainMemberRestriction.exists_opens_isInvDomain_memberRestrict_geoTransformAlong}
\mixedleanok
Let \(\underline{\mathcal I}=(M,N,E,\mathfrak a,d)\) be well formed with \(d\ge1\), \(n>0\) and \(N\) regular of dimension \(\le k\) at each of its points. For a set \(S\) of indices of
members let \(\underline{\mathcal S}_S=(M,N\cap H_S,E_S,\mathcal C_N(\mathfrak a),\mu_{\mathcal C})\) be the restriction of the coefficient ideal along \(N\) to \(N\cap H_S\), the members of
\(S\) deleted from \(E\) (Corollary~\ref{mc:coefficient-restriction-drop-along}).
\begin{enumerate}
\item If \(\underline{\mathcal I}\) is nonzero along \(N\) at the points of \(N\) in an open set \(U\), then \(\underline{\mathcal I}|_U\) lies in the domain in dimension \(\le k\).
\item \(\underline{\mathcal S}_S\) is well formed, \(\mu_{\mathcal C}\ge1\), \(\operatorname{cosupp}\underline{\mathcal S}_S=\operatorname{cosupp}\underline{\mathcal I}\cap H_S\), and if \(S\ne\emptyset\)
then \(N\cap H_S\) is regular of dimension \(\le k-1\) at each of its points. If \(E_0\) is a finite sequence of sets and \(z\in N\) is a point near which \(H_S\) is the
intersection of the members of \(E_0\) through \(z\), and \(z\) is not in the member locus \(L_N(E_0,\operatorname{cosupp}\underline{\mathcal I})\), then the ideal of
\(\underline{\mathcal S}_S\) is nonzero along \(N\cap H_S\) at \(z\).
\item Let \(E_0\) be a sequence of members of \(E\), \(y\) a point on \(s_0\ge1\) members of \(E_0\), \(S(y)\) the set of members of \(E_0\) through \(y\), and assume that no point
of \(L_N(E_0,\operatorname{cosupp}\underline{\mathcal I})\) lies on exactly \(s_0\) members of \(E_0\). Then for every open \(U\) whose points lie on no member of \(E_0\) missing \(y\)
(for instance the complement of the union of these members, an open neighbourhood of \(y\)), \(\underline{\mathcal S}_{S(y)}|_U\) lies in the domain in dimension \(\le k-1\).
\item The same holds at every stage of a multiple test blow-up \(T\) of \(\underline{\mathcal I}\) along which no point of the member locus of the cosupport along \(N_T\) and the
strict transforms of \(E_0\) lies on exactly \(s_0\) of these strict transforms, for \(\underline{\mathcal S}\) formed from \(\underline{\mathcal I}_T\), at a point on \(s_0\ge1\) of
them; in particular at every stage after the blowup of Step~1aa at the level \(s_0\) (Theorem~\ref{bcalc:member-locus-level}(3)).
\end{enumerate}
\end{proposition}
\begin{proof}
\mixedleanok
(1) As in Proposition~\ref{bcalc:inv-domain-stable}(3): well-formedness restricts to open sets (Lemma~\ref{hir:open-restriction-marked}), so do the dimensions, and a germ \(g\)
of \(\mathfrak a\) at a point of \(U\) not vanishing identically on \(N\) near it restricts to a germ of the restricted ideal with the same property
(Lemma~\ref{hir:open-restriction-ideal}). (2) Well formed: Lemma~\ref{bcalc:member-intersection-drop-along}, \(\mathcal C_N(\mathfrak a)\) being coherent
(Lemma~\ref{mc:homogenized-along-basic}); \(\mu_{\mathcal C}=\prod_{j<d}(d-j)\ge1\); the cosupport is Lemma~\ref{mc:coefficient-restriction-along}; the dimension is
Lemma~\ref{bcalc:member-intersection-dim}. Being nonzero along a set at \(z\) depends only on the germ of the set at \(z\) (Definition~\ref{hir:exponent-along}), so the last
claim is Proposition~\ref{bcalc:member-restriction-nonzero}. (3) Let \(z\in U\cap N\cap H_{S(y)}\). It lies on every member of \(E_0\) through \(y\) and, being in \(U\), on no
other; so \(H_{S(y)}\) is the intersection of the members of \(E_0\) through \(z\), and \(z\) lies on exactly \(s_0\) members of \(E_0\) (Lemma~\ref{bcalc:boundary-count}),
hence not in the member locus. Now (2) and (1) apply, with \(S(y)\ne\emptyset\). (4) Along \(T\), \(\underline{\mathcal I}_T\) is well formed (Lemma~\ref{hir:wellformed-tower}) with
the same mark (Lemma~\ref{hir:geo-transform-basic}) and \(\dim N_T\le k\) (Proposition~\ref{bcalc:inv-domain-stable}(2)), and the strict transforms of the members of \(E_0\) are
members of \(E_T\) (Lemma~\ref{bcalc:member-intersection-tower}), closed; apply (3) to \(\underline{\mathcal I}_T\).
\end{proof}
These are the domain hypotheses of the inductive invariant for Step~1ab, the analogue for Step~1a of Proposition~\ref{bcalc:contact-restriction-nonzero} for Step~1bb.

\paragraph{The domain of the companion ideal.} In the general step (Proposition~\ref{hir:general-step}), at a point where the order \(o\) of the nonmonomial
part is positive, the value is computed from the invariant of the maximal-order step for the companion ideal (Definition~\ref{hir:companion-ideal}), formed
from the transform at the stage where \(o\) was first reached and restricted to an open set (Włodarczyk, Step~2a: ``By Step 1 we can resolve the marked
ideal of maximal order \((\mathcal J,\mu_{\mathcal J}):=O(\mathcal I,\mu)\) [\dots] Then we repeat the procedure for \((\mathcal I_{r_1},\mu)\)''). The statement below
shows that the companion ideal lies in the domain of that invariant.

\begin{proposition}[The companion ideal lies in the domain]
\label{bcalc:companion-domain}
\uses{hir:centre-function,hir:monomial-part-along,hir:companion-reduction-along,hir:marked-sum-product,hir:order-power,hir:exponent-along,ana:ideal-sheaves,bcalc:inv-domain-stable}
\source{wlodarczyk05-hironaka,bierstone-milman08-functoriality}
\lean{Hironaka.InvDomainCompanion.isNonzeroAlongAt_idealPow,Hironaka.InvDomainCompanion.isNonzeroAlongAt_sumIdeal_left,Hironaka.InvDomainCompanion.isNonzeroAlongAt_nonmonomialPartAlong,Hironaka.InvDomainCompanion.isInvDomain_companionAlong,Hironaka.InvDomainCompanion.isInvDomain_companionAlong_pullback_val,Hironaka.InvDomainCompanion.isInvDomain_pullback_val_companionAlong,Hironaka.InvDomainCompanion.isInvDomain_companionAlong_geoTransformAlong}
\leanok
Let \(\underline{\mathcal I}=(M,N,E,\mathfrak a,d)\) lie in the domain of a resolution invariant in dimension \(\le k\) (Definition~\ref{hir:centre-function}), \(n>0\), and
\(o\ge1\). Write \(\mathcal N_N(\mathfrak a)\) and \(\mathcal M_N(\mathfrak a)\) for the nonmonomial and monomial parts along \(N\) (Proposition~\ref{hir:monomial-part-along}).
\begin{enumerate}
\item If \(\mathfrak b\) is nonzero along \(N\) at \(y\), so are \(\mathfrak b^j\) and \(\mathfrak b+\mathfrak c\); and \(\mathcal N_N(\mathfrak a)\) is nonzero along \(N\) at every point
of \(N\).
\item The marked ideal \(\underline{\mathcal O}=(M,N,E,\mathcal N_N(\mathfrak a),o)+(M,N,E,\mathcal M_N(\mathfrak a),d-o)\) if \(o<d\), and
\(\underline{\mathcal O}=(M,N,E,\mathcal N_N(\mathfrak a),o)\) if \(o\ge d\) (the companion ideal of Lemma~\ref{hir:companion-reduction-along} with
\(\mathfrak b=\mathcal N_N(\mathfrak a)\), \(\mathfrak c=\mathcal M_N(\mathfrak a)\)), lies in the domain in dimension \(\le k\).
\item The same holds for the companion ideal of \(\underline{\mathcal I}_T\) along every multiple test blow-up \(T\), for the companion ideal of \(\underline{\mathcal I}|_W\)
and for \(\underline{\mathcal O}|_W\), \(W\subset M\) open.
\end{enumerate}
\end{proposition}
\begin{proof}
\leanok
(1) If \(g\in\mathfrak b_y\) does not vanish identically on \(N\) near \(y\), neither does \(g^j\in(\mathfrak b^j)_y\), and \(g\in(\mathfrak b+\mathfrak c)_y\)
(Definition~\ref{hir:exponent-along}, Lemma~\ref{hir:order-power}). Let \(y\in N\) and \(g\in\mathfrak a_y\) not vanishing identically on \(N\) near \(y\). On \(N\) near \(y\),
\(g=m\,c\) with \(m\) the local generator of \(\mathcal M_N(\mathfrak a)\) and \(c\in\mathcal N_N(\mathfrak a)_y\) (Proposition~\ref{hir:monomial-part-along}); at the points
of \(N\) near \(y\) where \(g\ne0\), also \(c\ne0\), so \(c\) does not vanish identically on \(N\) near \(y\).
(2) \(\underline{\mathcal O}\) has the submanifold \(N\) and the divisor \(E\) of \(\underline{\mathcal I}\), so it is well formed as soon as its ideal is coherent:
\(\mathcal N_N(\mathfrak a)\) and \(\mathcal M_N(\mathfrak a)\) are coherent (Proposition~\ref{hir:monomial-part-along}), and so are powers and sums of coherent ideals
(Definition~\ref{ana:ideal-sheaves}). Its mark is \(o(d-o)\ge1\) (Definition~\ref{hir:marked-sum-product}), resp. \(o\ge1\), and \(\dim N\le k\). Its ideal is
\(\mathcal N_N(\mathfrak a)^{d-o}+\mathcal M_N(\mathfrak a)^{o}\), resp. \(\mathcal N_N(\mathfrak a)\), nonzero along \(N\) at every point of \(N\) by (1).
(3) \(\underline{\mathcal I}_T\) and \(\underline{\mathcal I}|_W\) lie in the domain (Proposition~\ref{bcalc:inv-domain-stable}); apply (2), resp.\ restrict the conclusion of (2) to \(W\)
(Proposition~\ref{bcalc:inv-domain-stable}(3)).
\end{proof}
The maximal order of \(\underline{\mathcal O}\) along \(N\) and its cosupport are Lemmas~\ref{hir:companion-basic} and~\ref{hir:companion-reduction-along}; together with the
proposition above this is what the general step needs of \(\underline{\mathcal O}\) to apply the invariant of Proposition~\ref{hir:maximal-order-step} to it.

\paragraph{Restricted towers with empty centres.} The values of the invariant at a point are computed along the part of a canonical segment after the
stage at which the relevant order (the boundary count, the codimension-one part, the order of the nonmonomial part) was reached near the point; before
that stage the centres lie over a closed set missing a neighbourhood of the point, and over that neighbourhood the blowups are blowups along the empty
set, which the re-basing property of the invariant (Definition~\ref{hir:centre-function}, Lemma~\ref{hir:tower-split}) removes.

\begin{lemma}[A restricted tower whose centres miss the open set is empty]
\label{bcalc:restrict-empty-tower}
\uses{hir:blowup-tower,hir:open-restriction-tower,hir:tower-split,hir:centre-function}
\source{wlodarczyk05-hironaka}
\lean{Hironaka.RestrictEmptyTower.CentresAvoid,Hironaka.RestrictEmptyTower.CentresOver,Hironaka.RestrictEmptyTower.centresAvoid_refl,Hironaka.RestrictEmptyTower.centresOver_refl,Hironaka.RestrictEmptyTower.centresAvoid_step_iff}\upstreamlean{TranscendentalNakamaye.Pending.RestrictEmptyTower.centresOver_step_iff}\lean{Hironaka.RestrictEmptyTower.preimage_val_eq_empty_iff,Hironaka.RestrictEmptyTower.isEmptyTower_restrict_iff,Hironaka.RestrictEmptyTower.centresAvoid_of_centresOver}\upstreamlean{TranscendentalNakamaye.Pending.RestrictEmptyTower.isEmptyTower_restrict_of_centresOver,TranscendentalNakamaye.Pending.RestrictEmptyTower.centresAvoid_append_iff,TranscendentalNakamaye.Pending.RestrictEmptyTower.centresOver_append_iff,TranscendentalNakamaye.Pending.RestrictEmptyTower.centresAvoid_head,TranscendentalNakamaye.Pending.RestrictEmptyTower.centresOver_head,TranscendentalNakamaye.Pending.RestrictEmptyTower.isEmptyTower_restrict_head}
\mixedleanok
Let \(T\) be a tower over \(X\) with centres \(C_i\subset X_i\) and projections \(p_i\colon X_i\to X\), and \(U\subset X\) open. The restriction \(T|_U\)
(Lemma~\ref{hir:open-restriction-tower}) has only empty centres if and only if \(C_i\cap p_i^{-1}(U)=\emptyset\) for every \(i\). In particular, if every
\(C_i\) lies in \(p_i^{-1}(F)\) for a set \(F\) with \(U\cap F=\emptyset\), then \(T|_U\) has only empty centres. Both conditions pass to concatenations
\(T=T_0T_1\) (for \(T_1\), with \(U\) and \(F\) replaced by their preimages in the last stage of \(T_0\)), hence to the head \(T_0\) of a split of \(T\)
(Lemma~\ref{hir:tower-split}).
\end{lemma}
\begin{proof}
\mixedleanok
The centres of \(T|_U\) are the traces \(C_i\cap p_i^{-1}(U)\) of the centres of \(T\) on the open preimages of \(U\) (Lemma~\ref{hir:open-restriction-tower});
a trace on an open subset is empty iff the set misses it. If \(C_i\subset p_i^{-1}(F)\) and \(U\cap F=\emptyset\), a point of \(C_i\cap p_i^{-1}(U)\) would map into
\(U\cap F\). For \(T_0T_1\) the centres are those of \(T_0\) followed by those of \(T_1\), and the projections of the stages of \(T_1\) to \(X\) are those to the
last stage of \(T_0\) followed by \(p_{T_0}\) (Definition~\ref{hir:blowup-tower}). The head of a split followed by its tail is \(T\)
(Lemma~\ref{hir:tower-split}).
\end{proof}

\paragraph{Re-basing Step~1a's restricted marked ideal.} In Step~1a of the maximal-order step (Proposition~\ref{hir:maximal-order-step}) the restricted
marked ideal of Lemma~\ref{bcalc:member-restriction-domain} is formed at a stage depending on the point and restricted to a neighbourhood \(W\) of the image
of the point. A nearby point whose stage is later is compared by re-basing across the steps between the two stages, which are empty over \(W\)
(Lemma~\ref{bcalc:restrict-empty-tower}). The marked ideal formed at the earlier stage and transformed along these steps is then the one formed at the later
stage, up to restricted equivalence (Definition~\ref{hir:restricted-equivalence}), which the invariant does not see (Definition~\ref{hir:centre-function}(7)).
This is the lemma of the source on the commutativity of Step~1 with étale morphisms (the blow-ups whose centres ``do not hit the images'' pull
back to isomorphisms, and \(\phi^*(\overline{\mathcal J}_{i_l|H^s_{\alpha i_l}})=\overline{\mathcal J}'_{j_{l'}|(H')^s_{\alpha j_{l'}}}\)), with steps of empty centre in place of
isomorphisms.

\begin{lemma}[Re-basing Step~1a's restricted marked ideal]
\label{bcalc:member-restriction-rebase}
\uses{bcalc:member-restriction-domain,bcalc:restrict-empty-tower,bcalc:snc-simultaneous,mc:coefficient-empty-step,mc:empty-blowup-inverse,mc:homogenized-along-etale,mc:strict-along,bcalc:isomorphism-off-centre,hir:open-restriction-marked,hir:open-restriction-tower,hir:restricted-equivalence,hir:restricted-equivalence-basic,hir:geo-transform-basic,bcalc:member-intersection-tower,hir:tower-transport-resolution,hir:canonical-segments}
\source{wlodarczyk05-hironaka}
\lean{Hironaka.InvDomainRebase.eq_of_eventually_mem_iff,Hironaka.InvDomainRebase.map_get_filter_finRange,Hironaka.InvDomainRebase.memberInter_membersThrough_eq,Hironaka.InvDomainRebase.memberRestrict_membersThrough_E,Hironaka.InvDomainRebase.memberRestrict_membersThrough_N,Hironaka.InvDomainRebase.filter_map_members,Hironaka.InvDomainRebase.surjective_of_isBlowupAlong_empty,Hironaka.InvDomainRebase.restrictEquiv_memberRestrict_geoTransform_empty,Hironaka.InvDomainRebase.mem_of_preimage_val_mem_map,Hironaka.InvDomainRebase.pullback_val_memberRestrict_membersThrough,Hironaka.InvDomainRebase.restrictEquiv_memberRestrict_geoTransformAlong_of_isEmptyTower,Hironaka.InvDomainRebase.restrictEquiv_memberRestrict_rebase}
\leanok
Let \(\underline{\mathcal I}=(M,N,E,\mathfrak a,d)\) be well formed with \(d\ge1\), \(n>0\), and \(E_0\) a list of members of \(E\). For a point \(y\) let \(S(y)\) be
the set of members of \(E_0\) through \(y\) and \(\underline{\mathcal S}(\underline{\mathcal I},E_0,y)=(M,N\cap H_{S(y)},E_{S(y)},\mathcal C_N(\mathfrak a),\mu_{\mathcal C})\)
the restricted marked ideal of Lemma~\ref{bcalc:member-restriction-domain}.
\begin{enumerate}
\item \(H_{S(y)}\) is the intersection of the members of \(E\) which belong to \(E_0\) and pass through \(y\), and \(E_{S(y)}\) is \(E\) without these members,
in its order.
\item For \(U\subset M\) open and \(y\in U\), \(\underline{\mathcal S}(\underline{\mathcal I},E_0,y)|_U=\underline{\mathcal S}(\underline{\mathcal I}|_U,E_0|_U,y)\), where \(E_0|_U\) is the list of the
traces of the members of \(E_0\) on \(U\).
\item For the blowup \(\pi\colon M'\to M\) along the empty set and \(z\in M'\), \(\underline{\mathcal S}(\underline{\mathcal I}',E_0',z)\) is restricted-equivalent to
\(\underline{\mathcal S}(\underline{\mathcal I},E_0,\pi(z))'\), where primes denote transforms and \(E_0'\) is the list of the strict transforms of the members of \(E_0\).
The same holds along a tower \(T\) with only empty centres: \(\underline{\mathcal S}(\underline{\mathcal I}_T,E_{0,T},z)\) is restricted-equivalent to
\(\underline{\mathcal S}(\underline{\mathcal I},E_0,p_T(z))_T\).
\item Let \(T\) be a multiple test blow-up of \(\underline{\mathcal I}\) whose centres miss the preimages of an open set \(W\subset M\), and \(p\) a point of the last
stage over \(W\). Then \(\underline{\mathcal S}(\underline{\mathcal I}_T,E_{0,T},p)|_{p_T^{-1}(W)}\) is restricted-equivalent to
\((\underline{\mathcal S}(\underline{\mathcal I},E_0,p_T(p))|_W)_{T|_W}\).
\end{enumerate}
\end{lemma}
\begin{proof}
\leanok
(1) A member \(H_i\) has \(i\in S(y)\) exactly when \(H_i\) is a member of \(E_0\) and \(y\in H_i\); \(E_{S(y)}\) keeps the other members in their order.
(2) The coefficient ideal along \(N\) commutes with restriction to open sets (Lemma~\ref{mc:homogenized-along-etale}), and so do intersections with members and
the deletion of members, provided a member through \(y\) is deleted on \(U\) exactly when it is deleted on \(M\). Let \(H\in E\) pass through \(y\) and have the
trace on \(U\) of a member \(H_0\) of \(E_0\). Then \(H\) and \(H_0\) agree on the neighbourhood \(U\) of \(y\) and both pass through \(y\), so they are the same
member (Lemma~\ref{bcalc:snc-simultaneous}, first paragraph of the proof), and \(H\in E_0\).
(3) A blowup along the empty set is a biholomorphism (Lemma~\ref{mc:empty-blowup-inverse}); it is surjective, by its section over the complement of the centre
(Lemma~\ref{bcalc:isomorphism-off-centre}). The strict transform of a closed set is its preimage, and \(E'=(\pi^{-1}H_1,\dots,\pi^{-1}H_s,\emptyset)\)
(Lemma~\ref{hir:geo-transform-basic}). Preimages under a surjection determine the set, so \(\pi^{-1}H\in E_0'\) iff \(H\in E_0\); and \(z\in\pi^{-1}H\) iff
\(\pi(z)\in H\), while the last member misses \(z\). By (1), \(N'\cap H'_{S(z)}=\pi^{-1}(N\cap H_{S(\pi(z))})\), the strict transform of \(N\cap H_{S(\pi(z))}\), and
\(E'_{S(z)}=(E_{S(\pi(z))})'\). The coefficient ideal of \(\underline{\mathcal I}'\) restricted to this strict transform is restricted-equivalent to the transform of the
restricted coefficient ideal (Lemma~\ref{mc:coefficient-empty-step}). Along \(T\) argue by induction on its length: the stages are well formed
(Lemma~\ref{hir:tower-transport-resolution}; a tower with only empty centres is a multiple test blow-up, Lemma~\ref{hir:canonical-segments}), the strict transforms of the
members of \(E_0\) are members of \(E_T\) (Lemma~\ref{bcalc:member-intersection-tower}), and restricted-equivalent marked ideals have the same transform
(Lemma~\ref{hir:restricted-equivalence-basic}), so the restricted equivalence of the previous stage passes to the transforms.
(4) \(T|_W\) has only empty centres (Lemma~\ref{bcalc:restrict-empty-tower}). Apply (3) to \(\underline{\mathcal I}|_W\), \(E_0|_W\) and \(T|_W\) at \(p\). Here
\((\underline{\mathcal I}|_W)_{T|_W}=\underline{\mathcal I}_T|_{p_T^{-1}(W)}\) (Lemma~\ref{hir:open-restriction-tower}), the strict transforms of traces are the traces of strict transforms
(Lemma~\ref{mc:strict-along}), and (2) at both ends (with \(\underline{\mathcal I}\) and \(\underline{\mathcal I}_T\) well formed, Lemma~\ref{hir:open-restriction-marked}) takes the
restriction outside.
\end{proof}

\begin{lemma}[Re-basing the values of a resolution invariant]
\label{bcalc:inv-rebase}
\uses{hir:centre-function,hir:canonical-segments,hir:restricted-equivalence,hir:marked-equivalence-basic,hir:nonmonomial-step-along,bcalc:inv-domain-stable,hir:invariant-driver,bcalc:member-restriction-rebase,bcalc:restrict-empty-tower,bcalc:companion-domain}
\source{wlodarczyk05-hironaka}
\lean{Hironaka.InvDomainRebase.restrictEquiv_symm,Hironaka.InvDomainRebase.isInvDomain_of_restrictEquiv,Hironaka.InvDomainRebase.isInvDomain_geoTransformAlong_of_isEmptyTower,Hironaka.InvDomainRebase.isCanonical_of_restrictEquiv,Hironaka.InvDomainRebase.inv_rebase,Hironaka.InvDomainRebase.memberRestrict_rebase}\upstreamlean{TranscendentalNakamaye.Pending.InvDomainRebase.isInvDomain_companionAlong_rebase,TranscendentalNakamaye.Pending.InvDomainRebase.companionAlong_rebase}
\mixedleanok
Let \(\operatorname{inv}\) be a resolution invariant in dimension \(\le k\) (Definition~\ref{hir:centre-function}).
\begin{enumerate}
\item If \(\underline{\mathcal I}\) lies in the domain and \(\underline{\mathcal J}\), with a coherent ideal, is restricted-equivalent to it, then \(\underline{\mathcal J}\) lies in the
domain. Along a tower with only empty centres the transform of a marked ideal of the domain lies in the domain.
\item If \(\underline{\mathcal I}\) and \(\underline{\mathcal J}\) lie in the domain and are restricted-equivalent, every canonical segment of \(\underline{\mathcal I}\) is a canonical segment of
\(\underline{\mathcal J}\).
\item Let \(\underline{\mathcal I}\) lie in the domain, \(T_0\) be a tower with only empty centres, and \(\underline{\mathcal J}\), with a coherent ideal, be restricted-equivalent to
\(\underline{\mathcal I}_{T_0}\). If \(T_0T_1\) is a canonical segment of \(\underline{\mathcal I}\), then \(\underline{\mathcal J}\) lies in the domain, \(T_1\) is a canonical segment of
\(\underline{\mathcal J}\), and \(\operatorname{inv}_{\underline{\mathcal I}}(T_0T_1)=\operatorname{inv}_{\underline{\mathcal J}}(T_1)\) on \(\operatorname{cosupp}\underline{\mathcal I}_{T_0T_1}\).
\item (Step~1a) In the situation of Lemma~\ref{bcalc:member-restriction-rebase}(4), let \(\underline{\mathcal S}_0=\underline{\mathcal S}(\underline{\mathcal I},E_0,p_T(p))|_W\) lie in the
domain and \(T|_WT_1\) be a canonical segment of it, and let \(\underline{\mathcal S}_1=\underline{\mathcal S}(\underline{\mathcal I}_T,E_{0,T},p)|_{p_T^{-1}(W)}\). Then
\(\underline{\mathcal S}_1\) lies in the domain, \(T_1\) is a canonical segment of it, its transforms along \(T_1\) lie in the domain, and
\(\operatorname{inv}_{\underline{\mathcal S}_0}(T|_WT_1)=\operatorname{inv}_{\underline{\mathcal S}_1}(T_1)\) on the cosupport.
\item (The general step) Let \(\underline{\mathcal I}\) lie in the domain, \(n>0\), \(o\ge1\), \(\underline{\mathcal O}\) its companion ideal (Proposition~\ref{bcalc:companion-domain}),
\(W\) open and \(T\) a tower whose centres miss the preimages of \(W\). Then \((\underline{\mathcal O}|_W)_{T|_W}\) lies in the domain; and if \(T|_WT_1\) is a canonical
segment of \(\underline{\mathcal O}|_W\), then \(T_1\) is one of \((\underline{\mathcal O}|_W)_{T|_W}\), the transforms along \(T|_WT_1\) lie in the domain, and
\(\operatorname{inv}_{\underline{\mathcal O}|_W}(T|_WT_1)=\operatorname{inv}_{(\underline{\mathcal O}|_W)_{T|_W}}(T_1)\) on the cosupport.
\end{enumerate}
\end{lemma}
\begin{proof}
\mixedleanok
(1) \(\underline{\mathcal J}\) has the submanifold, divisor and mark of \(\underline{\mathcal I}\), so it is well formed with \(\dim N\le k\) and \(d\ge1\). At \(y\in N\) let \(g\in\mathfrak a_y\)
not vanish identically on \(N\) near \(y\); on \(N\) near \(y\), \(g=\sum_ia_ib_i\) with \(b_i\) in the stalk of the ideal of \(\underline{\mathcal J}\), and if every \(b_i\) vanished
identically on \(N\) near \(y\), so would \(g\) (Lemma~\ref{hir:nonmonomial-step-along}). A tower with only empty centres is a multiple test blow-up
(Lemma~\ref{hir:canonical-segments}(4)), along which the domain is stable (Proposition~\ref{bcalc:inv-domain-stable}).
(2) Induction along the segment. At every stage the transforms of \(\underline{\mathcal I}\) and \(\underline{\mathcal J}\) have the same cosupport
(Lemma~\ref{hir:marked-equivalence-basic}), and on it their values agree by Definition~\ref{hir:centre-function}(7) along the canonical prefix; so their maximum loci
agree, and a centre which is empty or the maximum locus for \(\underline{\mathcal I}\) is so for \(\underline{\mathcal J}\).
(3) \(\underline{\mathcal I}_{T_0}\) lies in the domain by (1), and \(T_1\) is a canonical segment of it (Lemma~\ref{hir:canonical-segments}(4)). Restricted equivalence is
symmetric, so (1) puts \(\underline{\mathcal J}\) in the domain and (2) makes \(T_1\) a canonical segment of \(\underline{\mathcal J}\). For the values, re-base by
Definition~\ref{hir:centre-function}(4') and then apply (7) along \(T_1\).
(4) Lemma~\ref{bcalc:member-restriction-rebase}(4) gives the restricted equivalence of \(\underline{\mathcal S}_1\) with \((\underline{\mathcal S}_0)_{T|_W}\); \(T|_W\) has only empty centres
(Lemma~\ref{bcalc:restrict-empty-tower}), and \(\underline{\mathcal S}_1\) has a coherent ideal (Lemma~\ref{bcalc:member-restriction-domain}). Apply (3); the transforms along a
canonical segment of a marked ideal of the domain lie in the domain (Theorem~\ref{hir:invariant-driver}).
(5) \(\underline{\mathcal O}|_W\) lies in the domain (Proposition~\ref{bcalc:companion-domain}(3)) and \(T|_W\) has only empty centres (Lemma~\ref{bcalc:restrict-empty-tower});
apply (1), Lemma~\ref{hir:canonical-segments}(4), Theorem~\ref{hir:invariant-driver} and Definition~\ref{hir:centre-function}(4').
\end{proof}
The restricted equivalence of the companion ideals formed at two stages, the analogue of Lemma~\ref{bcalc:member-restriction-rebase}(4) for the general step, is not
part of these statements.

\paragraph{Centres from local data.} The invariant built in the inductive step (Propositions~\ref{hir:maximal-order-step} and~\ref{hir:general-step})
is given near each point by local data: on a neighbourhood of the point its values are \((s_0,b)\), with \(b\) the invariant of a lower-dimensional
marked ideal on the cosupport of that marked ideal, while the remaining points of the cosupport there have smaller values. The centre is the
locus where the global invariant attains its maximum (Theorem~\ref{hir:resolution-driver}; Włodarczyk: ``The blow-ups will be performed at the
centers \(C\) [\dots] for which the invariant [\dots] attains its maximum''). The two statements below compute such a centre from the local data,
and record how centres behave under restriction to open sets and under isomorphisms of towers.

\begin{lemma}[Maximum loci from local data]
\label{bcalc:max-locus-local}
\uses{hir:resolution-driver}
\source{wlodarczyk05-hironaka}
\lean{Hironaka.MaxLocusLocal.maxLocus_comp_strictMono,Hironaka.MaxLocusLocal.preimage_maxLocus_subset,Hironaka.MaxLocusLocal.preimage_maxLocus_eq_or}\upstreamlean{TranscendentalNakamaye.Pending.MaxLocusLocal.preimage_maxLocus_eq_or_of_eq,TranscendentalNakamaye.Pending.MaxLocusLocal.preimage_level_eq_of_eq,TranscendentalNakamaye.Pending.MaxLocusLocal.preimage_level_eq_empty,TranscendentalNakamaye.Pending.MaxLocusLocal.le_of_comp_le}\lean{Hironaka.MaxLocusLocal.strictMono_toLex_coe}
\mixedleanok
Let \(\Lambda\), \(\Lambda'\) be totally ordered, \(f\colon Y\to\Lambda\), \(s\subset Y\), \(\iota\colon Y'\to Y\) a map, \(s'\subset\iota^{-1}(s)\), \(g\colon Y'\to\Lambda'\) and
\(\varphi\colon\Lambda'\to\Lambda\) strictly increasing, with \(f(\iota(p))=\varphi(g(p))\) for \(p\in s'\). Write \(\max_sf\) for the set of points of \(s\) at which \(f\)
attains its maximum on \(s\).
\begin{enumerate}
\item \(\max_t(\varphi\circ g)=\max_tg\) for every \(t\).
\item If every point of \(\iota^{-1}(s)\setminus s'\) has a value of \(f\) strictly below the value at some point of \(s'\), then \(\iota^{-1}(\max_sf)\) is
empty or equal to \(\max_{s'}g\). In particular, if \(f(\iota(p))=g(p)\) on \(\iota^{-1}(s)\), then \(\iota^{-1}(\max_sf)\) is empty or \(\max_{\iota^{-1}(s)}g\).
\item If no point of \(\iota^{-1}(s)\setminus s'\) has the value \(l\), then \(\iota^{-1}\{y\in s:f(y)=l\}\) is \(\{p\in s':g(p)=\mu\}\) if \(l=\varphi(\mu)\), and empty
if \(l\) is not a value of \(\varphi\); and \(f\circ\iota\le\varphi(\mu)\) on \(s'\) gives \(g\le\mu\) on \(s'\).
\item For a totally ordered \(A\) and \(a\in A\), \(b\mapsto(a,b)\) is strictly increasing from \(\Lambda'\) to \(A\times(\Lambda'\cup\{\infty\})\), ordered
lexicographically.
\end{enumerate}
\end{lemma}
\begin{proof}
\mixedleanok
(1) \(\varphi(b)\le\varphi(b')\) iff \(b\le b'\). (2) A point \(p\) of \(\iota^{-1}(\max_sf)\) outside \(s'\) would lie strictly below a point \(\iota(q)\), \(q\in s'\subset\iota^{-1}(s)\),
against the maximality of \(f(\iota(p))\); so \(\iota^{-1}(\max_sf)\subset s'\). Let it contain \(p_0\). If \(p\in\iota^{-1}(\max_sf)\), then \(p\in s'\) and
\(\varphi(g(q))=f(\iota(q))\le f(\iota(p))=\varphi(g(p))\) for \(q\in s'\), so \(p\in\max_{s'}g\). Conversely, for \(p\in\max_{s'}g\) and \(y\in s\),
\(f(y)\le f(\iota(p_0))=\varphi(g(p_0))\le\varphi(g(p))=f(\iota(p))\). The special case is \(\varphi=\mathrm{id}\), \(s'=\iota^{-1}(s)\). (3) On \(s'\),
\(f(\iota(p))=l\) iff \(\varphi(g(p))=\varphi(\mu)\) iff \(g(p)=\mu\) (\(\varphi\) is injective); off \(s'\) the value \(l\) does not occur. (4) The second components are
compared when the first agree, and \(\Lambda'\to\Lambda'\cup\{\infty\}\) is strictly increasing.
\end{proof}

\begin{lemma}[Centres on open sets and along isomorphisms of towers]
\label{bcalc:centre-local}
\uses{bcalc:max-locus-local,hir:invariant-driver,hir:biholomorphism-invariance,hir:open-restriction-ideal,hir:open-restriction-marked,bcalc:restrict-open,hir:centre-function,hir:canonical-segments,hir:tower-isomorphism,mc:tower-isomorphism-transform}
\source{wlodarczyk05-hironaka}
\lean{Hironaka.MaxLocusLocal.nestedOpens,Hironaka.MaxLocusLocal.nestedOpensDiffeo,Hironaka.MaxLocusLocal.nestedOpensDiffeo_apply,Hironaka.MaxLocusLocal.pullback_val_pullback_val,Hironaka.MaxLocusLocal.isAdmissibleCentre_pullback_val_of_forall_exists_opens}\upstreamlean{TranscendentalNakamaye.Pending.MaxLocusLocal.preimage_maxLocus_restrict_eq_or,TranscendentalNakamaye.Pending.MaxLocusLocal.maxLocus_isTowerIso}
\mixedleanok
Let \(\underline{\mathcal I}=(M,N,E,\mathfrak a,d)\) have a coherent ideal.
\begin{enumerate}
\item For open sets \(V\subset W\subset M\), the open subset \(W\cap V\) of \(W\) is biholomorphic to \(V\) by the identity on points, and under this
biholomorphism \((\underline{\mathcal I}|_W)|_{W\cap V}\) corresponds to \(\underline{\mathcal I}|_V\).
\item Let \(W\subset M\) be open and \(C\subset M\). If every point of \(W\) has an open neighbourhood \(V\subset W\) such that \(C\cap V\) is an admissible centre of
\(\underline{\mathcal I}|_V\), then \(C\cap W\) is an admissible centre of \(\underline{\mathcal I}|_W\).
\item Let \(\operatorname{inv}\) be a resolution invariant (Definition~\ref{hir:centre-function}), \(\underline{\mathcal I}\) in its domain and \(T\) a canonical segment. For
\(U\subset M\) open, the trace on the preimage of \(U\) of the maximum locus of \(\operatorname{inv}_{\underline{\mathcal I}}(T)\) on \(\operatorname{cosupp}\underline{\mathcal I}_T\) is
empty or the maximum locus of \(\operatorname{inv}_{\underline{\mathcal I}|_U}(T|_U)\) on \(\operatorname{cosupp}(\underline{\mathcal I}|_U)_{T|_U}\). For an isomorphism
\(\chi\colon T'\to T\) of towers over a biholomorphism \(\Phi\) (Definition~\ref{hir:tower-isomorphism}), the maximum locus of
\(\operatorname{inv}_{\Phi^*\underline{\mathcal I}}(T')\) on its cosupport is \(\chi_r^{-1}\) of that of \(\operatorname{inv}_{\underline{\mathcal I}}(T)\).
\end{enumerate}
\end{lemma}
\begin{proof}
\mixedleanok
(1) The identity on points and its inverse are holomorphic (Lemma~\ref{bcalc:restrict-open}). Restricting the members, \(N\) and the ideal to \(W\) and then
to \(W\cap V\) is restricting them to \(V\) along the composite inclusion, which is the inclusion of \(V\) after the biholomorphism; pullbacks of coherent
ideals compose (Lemma~\ref{hir:open-restriction-ideal}). (2) Admissible centres are local (Theorem~\ref{hir:invariant-driver}): apply this to
\(\underline{\mathcal I}|_W\) with the open sets \(W\cap V\) of \(W\); by (1) and the invariance of admissible centres under biholomorphisms
(Lemma~\ref{hir:biholomorphism-invariance}), \(C\cap W\cap V\) is admissible for \((\underline{\mathcal I}|_W)|_{W\cap V}\). (3) The canonical segment is a multiple test
blow-up (Theorem~\ref{hir:invariant-driver}), so the cosupport along \(T|_U\) is the trace of the cosupport along \(T\) (Lemma~\ref{hir:canonical-segments}) and,
along \(T'\), its preimage under \(\chi_r\) (Lemma~\ref{mc:tower-isomorphism-transform}); the values agree there by Definition~\ref{hir:centre-function}(5),
resp.\ (6). Apply Lemma~\ref{bcalc:max-locus-local}(2) with \(\varphi=\mathrm{id}\), resp.\ the pullback of a maximum locus along the surjection \(\chi_r\)
(Lemma~\ref{hir:canonical-segments}).
\end{proof}

\begin{lemma}[A resolution invariant on a smaller open set]
\label{bcalc:restrict-nested}
\uses{bcalc:centre-local,hir:open-restriction-tower,hir:tower-isomorphism,hir:centre-function,hir:canonical-segments,bcalc:inv-domain-stable,mc:tower-isomorphism-transform,bcalc:restrict-open}
\source{wlodarczyk05-hironaka}
\lean{Hironaka.MaxLocusLocal.nestedDiffeo,Hironaka.MaxLocusLocal.openPreimage_mono,Hironaka.MaxLocusLocal.mem_openPreimage_restrict_nestedOpens,Hironaka.MaxLocusLocal.isTowerIso_restrict_nestedOpens,Hironaka.MaxLocusLocal.pullback_symm_nestedOpensDiffeo,Hironaka.MaxLocusLocal.restrict_nestedOpens}
\leanok
Let \(T\) be a tower over \(X\), \(V\subset W\subset X\) open, and \(\underline{\mathcal I}\) a marked ideal on \(X\) with a coherent ideal.
\begin{enumerate}
\item Restricting \(T\) to \(W\) and then to the open subset \(W\cap V\) of \(W\) gives a tower isomorphic to \(T|_V\) over the biholomorphism \(W\cap V\cong V\) of
Lemma~\ref{bcalc:centre-local}(1), with the identity on points at every stage (Definition~\ref{hir:tower-isomorphism}).
\item Let \(\operatorname{inv}\) be a resolution invariant (Definition~\ref{hir:centre-function}), \(\underline{\mathcal I}|_W\) in its domain and \(T|_W\) a canonical segment of
it. Then \(\underline{\mathcal I}|_V\) lies in the domain, \(T|_V\) is a canonical segment of it, and at the points of
\(\operatorname{cosupp}(\underline{\mathcal I}|_V)_{T|_V}\) the values \(\operatorname{inv}_{\underline{\mathcal I}|_V}(T|_V)\) and
\(\operatorname{inv}_{\underline{\mathcal I}|_W}(T|_W)\) agree.
\end{enumerate}
\end{lemma}
\begin{proof}
\leanok
(1) Induction along \(T\): the stages of both towers are the open preimages of \(V\), with the same points, and the identity on points between open
sets with the same points is biholomorphic (Lemma~\ref{bcalc:restrict-open}); the centres of both are the traces of the centres of \(T\)
(Lemma~\ref{hir:open-restriction-tower}), and the blowup maps commute with the identity on points. (2) \((\underline{\mathcal I}|_W)|_{W\cap V}\) lies in the domain
(Proposition~\ref{bcalc:inv-domain-stable}) and \((T|_W)|_{W\cap V}\) is a canonical segment of it (Lemma~\ref{hir:canonical-segments}(1)); by
Lemma~\ref{bcalc:centre-local}(1) this marked ideal corresponds to \(\underline{\mathcal I}|_V\) under the biholomorphism, and by (1) and
Lemma~\ref{hir:canonical-segments}(2) \(T|_V\) is a canonical segment of \(\underline{\mathcal I}|_V\), which lies in the domain (Proposition~\ref{bcalc:inv-domain-stable}).
The cosupports correspond (Lemma~\ref{mc:tower-isomorphism-transform}), and the values agree by Definition~\ref{hir:centre-function}(6) along the isomorphism of (1)
and (5) for the restriction from \(W\) to \(W\cap V\).
\end{proof}

\paragraph{The monomial and nonmonomial parts under local biholomorphisms.} In the general step (Proposition~\ref{hir:general-step}) the companion ideal
(Definition~\ref{hir:companion-ideal}) is formed from the monomial and nonmonomial parts along \(N\) of the transform at the stage where the order of the
nonmonomial part was first reached near a point, restricted to an open set; nearby points are compared through restrictions to open sets, biholomorphisms and
steps with empty centre (Definition~\ref{hir:centre-function}(4'), (5), (6)). All three are \emph{local biholomorphisms}: holomorphic maps \(\varphi\colon M'\to M\)
with a holomorphic local inverse at every point (Lemma~\ref{mc:field-transfer}): the inclusion of an open set, a biholomorphism, a blowup along \(\emptyset\)
(Lemma~\ref{mc:empty-blowup-inverse}). The lemma below shows that both parts commute with local biholomorphisms, depend only on the restriction of the ideal to
\(N\), and ignore empty members of \(E\); this is how Włodarczyk's statement that the resolution procedure commutes with étale morphisms reaches the companion ideal.

\begin{lemma}[Monomial and nonmonomial parts under local biholomorphisms]
\label{bcalc:monomial-part-etale}
\uses{hir:monomial-part,hir:monomial-part-along,hir:divisor-exponent,hir:monomial-invariance-along,hir:restricted-calculus,hir:point-power-pullback,hir:marked-ideal-operations,mc:field-transfer,mc:empty-blowup-inverse,ana:restricted-colon,ana:ideal-sheaves}
\source{wlodarczyk05-hironaka,bierstone-milman08-functoriality}
\lean{Hironaka.MonomialPartEtale.comp_mem_idealPow_vanishingIdeal_preimage,Hironaka.MonomialPartEtale.comp_mem_idealPow_vanishingIdeal,Hironaka.MonomialPartEtale.map_nhds_eq_of_local,Hironaka.MonomialPartEtale.vanishesToOrderOn_comp,Hironaka.MonomialPartEtale.cosupport_pullback_of_local,Hironaka.MonomialPartEtale.monomialPartAlong_mem_iff',Hironaka.MonomialPartEtale.monomialPartAlong_mem_of_isInSpanAt,Hironaka.MonomialPartEtale.monomialPartAlong_append_empty,Hironaka.MonomialPartEtale.divExpAlong_congr,Hironaka.MonomialPartEtale.monomialPartAlong_congr,Hironaka.MonomialPartEtale.nonmonomialPartAlong_congr,Hironaka.MonomialPartEtale.nonmonomialPartAlong_append_empty,Hironaka.MonomialPartEtale.pullback_monomialPartAlong,Hironaka.MonomialPartEtale.mem_pullback_iff_of_local,Hironaka.MonomialPartEtale.mem_restrictOn_pullback_iff_of_local,Hironaka.MonomialPartEtale.restrictOn_mem_congr,Hironaka.MonomialPartEtale.pullback_colonAlong_of_local,Hironaka.MonomialPartEtale.pullback_nonmonomialPartAlong}
\leanok
Let \(\underline{\mathcal I}=(M,N,E,\mathfrak a,d)\) be as in Lemma~\ref{hir:divisor-exponent}, with \(\mathfrak a\) an ideal of holomorphic germs, and write
\(\mathcal M_N(\mathfrak a)\), \(\mathcal N_N(\mathfrak a)\) for its monomial and nonmonomial parts along \(N\) (Definition~\ref{hir:monomial-part}).
\begin{enumerate}
\item Let \(\varphi\colon M'\to M\) be holomorphic with a holomorphic local inverse at every point. Then
\(\operatorname{cosupp}((\varphi^*\mathfrak a)|_{\varphi^{-1}N},e)=\varphi^{-1}\operatorname{cosupp}(\mathfrak a|_N,e)\) for every \(e\), and
\(\varphi^*\mathcal M_N(\mathfrak a)=\mathcal M_{\varphi^{-1}N}(\varphi^*\mathfrak a)\), the right-hand side formed with the members \(\varphi^{-1}H_i\). For coherent
\(\mathfrak b\), \(\mathfrak l\) with \((\widetilde{\mathfrak b}^W:\mathfrak l)\) coherent (Lemma~\ref{ana:restricted-colon}),
\(\varphi^*(\widetilde{\mathfrak b}^W:\mathfrak l)=(\widetilde{\varphi^*\mathfrak b}^{\varphi^{-1}W}:\varphi^*\mathfrak l)\); hence, if \(\mathfrak a\), \(\mathcal M_N(\mathfrak a)\) and
\(\mathcal N_N(\mathfrak a)\) are coherent, \(\varphi^*\mathcal N_N(\mathfrak a)=\mathcal N_{\varphi^{-1}N}(\varphi^*\mathfrak a)\).
\item If \(N\) is closed and \(\operatorname{Res}_N\mathfrak a=\operatorname{Res}_N\mathfrak a'\), then \(\mathcal M_N(\mathfrak a')=\mathcal M_N(\mathfrak a)\) and
\(\mathcal N_N(\mathfrak a')=\mathcal N_N(\mathfrak a)\).
\item Appending the empty set to \(E\) as a last member changes neither part.
\end{enumerate}
\end{lemma}
\begin{proof}
\leanok
(1) Let \(y\in M'\) and \(\psi\) a local inverse of \(\varphi\) at \(y\). Composition with \(\varphi\) maps germs vanishing on \(H\) at \(\varphi(y)\) to germs vanishing on
\(\varphi^{-1}H\) at \(y\), and composition with \(\psi\) maps germs vanishing on \(\varphi^{-1}H\) at \(y\) to germs vanishing on \(H\) at \(\varphi(y)\)
(Lemma~\ref{mc:field-transfer}); by induction on \(e\) the same holds for the powers \(\mathcal I_H^e\), whose germs are combinations of products of a germ of
\(\mathcal I_H^{e-1}\) and one of \(\mathcal I_H\). Orders along \(N\) pull back along maps holomorphic near the point (Lemma~\ref{hir:point-power-pullback}), and
combinations of germs vanishing to order \(e\) along \(N\) vanish to order \(e\) (Lemma~\ref{hir:marked-ideal-operations}): a germ of \(\varphi^*\mathfrak a\) at \(y\) is a
combination of germs \(a\circ\varphi\), \(a\in\mathfrak a_{\varphi(y)}\); conversely, if \(a\circ\varphi\) vanishes to order \(e\) along \(\varphi^{-1}N\) at \(y\), then
\(a=(a\circ\varphi)\circ\psi\) near \(\varphi(y)\) vanishes to order \(e\) along \(N\), as \(\varphi\circ\psi=\mathrm{id}\) near \(\varphi(y)\). So the cosupports pull back; since
\(\varphi\) maps the neighbourhoods of \(y\) onto those of \(\varphi(y)\), the exponents pull back too, \(e_{\varphi^{-1}H}(\varphi^*\mathfrak a)(y)=e_H(\mathfrak a)(\varphi(y))\)
(Lemma~\ref{hir:monomial-invariance-along}). Membership in \(\mathcal M_N(\mathfrak a)_y\) is a condition on each member through \(y\) and is closed under holomorphic
combinations. A generator \(b\circ\varphi\), \(b\in\mathcal M_N(\mathfrak a)_{\varphi(y)}\), satisfies it at \(y\) by the transfer of the powers along \(\varphi\); conversely,
for \(g\in\mathcal M_{\varphi^{-1}N}(\varphi^*\mathfrak a)_y\) the transfer along \(\psi\) gives \(g\circ\psi\in\mathcal M_N(\mathfrak a)_{\varphi(y)}\), and \(g=(g\circ\psi)\circ\varphi\)
near \(y\). For a coherent \(\mathfrak c\), \(f\in(\varphi^*\mathfrak c)_y\) if and only if \(f\circ\psi\in\mathfrak c_{\varphi(y)}\): a coherent ideal is closed under holomorphic
combinations and its membership only depends on the germ (Definition~\ref{ana:ideal-sheaves}), and \((b\circ\varphi)\circ\psi=b\) near \(\varphi(y)\). The same holds for the
restricted germs of \(\widetilde{\mathfrak b}^W\) on \(W\) and of \(\widetilde{\varphi^*\mathfrak b}^{\varphi^{-1}W}\) on \(\varphi^{-1}W\), since \(\psi\) maps the points of
\(W\) near \(\varphi(y)\) into \(\varphi^{-1}W\); so \(f\) lies in the quotient at \(y\) if and only if \(f\circ\psi\) does at \(\varphi(y)\). The nonmonomial part is the quotient
\((\widetilde{\mathfrak a}^N:\mathcal M_N(\mathfrak a))\) (Definition~\ref{hir:monomial-part}), and \(\varphi^*\mathcal M_N(\mathfrak a)=\mathcal M_{\varphi^{-1}N}(\varphi^*\mathfrak a)\).
(2) The exponents only depend on the cosupports along \(N\), and \(\operatorname{Res}_N\mathfrak a=\operatorname{Res}_N\mathfrak a'\) gives the same cosupports and the same
extensions \(\widetilde{\mathfrak a}^N=\widetilde{\mathfrak a'}^N\) (Lemma~\ref{hir:restricted-calculus}). (3) No point lies on the empty member, so its condition in the
definition of \(\mathcal M_N\) is vacuous, and the nonmonomial part is the quotient by the same monomial part.
\end{proof}

\paragraph{Re-basing the companion ideal.} In the general step (Proposition~\ref{hir:general-step}) the value at a point where the order \(o\) of the
nonmonomial part is positive comes from the companion ideal formed at the stage where \(o\) was first reached near the point, restricted to an open set, and
nearby points whose stages differ are compared by re-basing across the steps between the two stages, which are empty near them
(Definition~\ref{hir:centre-function}(4'), Lemma~\ref{bcalc:restrict-empty-tower}). This is the general-step twin of Lemma~\ref{bcalc:member-restriction-rebase}.
For \(\underline{\mathcal I}=(M,N,E,\mathfrak a,d)\) and \(o\ge1\) write
\(\mathcal O(\underline{\mathcal I},o)=(M,N,E,\mathcal N_N(\mathfrak a),o)+(M,N,E,\mathcal M_N(\mathfrak a),d-o)\) if \(o<d\) and \((M,N,E,\mathcal N_N(\mathfrak a),o)\) if \(o\ge d\),
the companion ideal of Lemma~\ref{hir:companion-reduction-along} at the monomial and nonmonomial parts, as in Proposition~\ref{bcalc:companion-domain}.

\begin{lemma}[The companion ideal under local biholomorphisms, empty steps and re-basing]
\label{bcalc:companion-rebase}
\uses{bcalc:monomial-part-etale,bcalc:companion-domain,hir:companion-reduction-along,hir:marked-sum-product,hir:controlled-operations,bcalc:exceptional-divisor,mc:empty-blowup-inverse,hir:restricted-equivalence,hir:restricted-equivalence-basic,hir:open-restriction-marked,hir:open-restriction-tower,bcalc:restrict-empty-tower,bcalc:inv-domain-stable,bcalc:inv-rebase,hir:centre-function,hir:invariant-driver}
\source{wlodarczyk05-hironaka,bierstone-milman08-functoriality}
\lean{Hironaka.CompanionRebase.companionAlong_ideal,Hironaka.CompanionRebase.companionAlong_d,Hironaka.CompanionRebase.pullback_companionAlong,Hironaka.CompanionRebase.companionOf,Hironaka.CompanionRebase.companionOf_congr,Hironaka.CompanionRebase.pullback_companionOf_of_local,Hironaka.CompanionRebase.pullback_val_companionOf,Hironaka.CompanionRebase.isLocalInverseAt_diffeo,Hironaka.CompanionRebase.pullback_diffeo_companionOf,Hironaka.CompanionRebase.restrictEquiv_companionOf_geoTransform_empty,Hironaka.CompanionRebase.restrictEquiv_companionOf_geoTransformAlong_of_isEmptyTower,Hironaka.CompanionRebase.restrictEquiv_companionOf_rebase}\upstreamlean{TranscendentalNakamaye.Pending.CompanionRebase.companionOf_rebase}
\mixedleanok
Let \(n>0\) and \(o\ge1\).
\begin{enumerate}
\item For coherent \(\mathfrak b\), \(\mathfrak c\) and a holomorphic \(\varphi\colon M'\to M\), the pullback of \((M,N,E,\mathfrak b,o)+(M,N,E,\mathfrak c,d-o)\) is
\((M',\varphi^{-1}N,\varphi^{-1}E,\varphi^*\mathfrak b,o)+(M',\varphi^{-1}N,\varphi^{-1}E,\varphi^*\mathfrak c,d-o)\), and likewise for the single term if \(o\ge d\).
\item Let \(\underline{\mathcal I}\) be well formed with \(\mathfrak a\) nonzero along \(N\) at every point of \(N\). For a holomorphic \(\varphi\) with a holomorphic local inverse at
every point, \(\varphi^*\mathcal O(\underline{\mathcal I},o)=\mathcal O(\varphi^*\underline{\mathcal I},o)\); in particular \(\mathcal O(\underline{\mathcal I},o)|_U=\mathcal O(\underline{\mathcal I}|_U,o)\)
for \(U\subset M\) open and \(\Phi^*\mathcal O(\underline{\mathcal I},o)=\mathcal O(\Phi^*\underline{\mathcal I},o)\) for a biholomorphism \(\Phi\). For the blowup \(\pi\) along
\(\emptyset\), \(\mathcal O(\underline{\mathcal I}',o)\simeq\mathcal O(\underline{\mathcal I},o)'\) (restricted equivalence, Definition~\ref{hir:restricted-equivalence}).
\item If \(\underline{\mathcal I}\simeq\underline{\mathcal J}\) and \(N\) is closed, then \(\mathcal O(\underline{\mathcal I},o)=\mathcal O(\underline{\mathcal J},o)\).
\item Let \(\underline{\mathcal I}\) lie in the domain of resolution invariants in dimension \(\le k\) (Definition~\ref{hir:centre-function}). Along a tower \(T\) with only empty
centres, \(\mathcal O(\underline{\mathcal I}_T,o)\simeq\mathcal O(\underline{\mathcal I},o)_T\). If \(T\) is a multiple test blow-up of \(\underline{\mathcal I}\) whose centres miss the open
\(W\subset M\), then \(\mathcal O(\underline{\mathcal I}_T,o)|_{T^{-1}W}\simeq(\mathcal O(\underline{\mathcal I},o)|_W)_{T|_W}\).
\item In the situation of (4), let \(\operatorname{inv}\) be a resolution invariant in dimension \(\le k\) and \(T|_W\,T_1\) a canonical segment of \(\mathcal O(\underline{\mathcal I},o)|_W\). Then
\(\mathcal O(\underline{\mathcal I}_T,o)|_{T^{-1}W}\) lies in the domain, \(T_1\) is a canonical segment of it, its transforms along \(T_1\) lie in the domain, and at the points of the
cosupport along \(T|_W\,T_1\)
\[\operatorname{inv}_{\mathcal O(\underline{\mathcal I},o)|_W}(T|_W\,T_1)=\operatorname{inv}_{\mathcal O(\underline{\mathcal I}_T,o)|_{T^{-1}W}}(T_1).\]
\end{enumerate}
\end{lemma}
\begin{proof}
\mixedleanok
(1) Pullback along a holomorphic map commutes with sums and powers of ideals of holomorphic germs (Lemmas~\ref{hir:controlled-operations}
and~\ref{bcalc:exceptional-divisor}), and the operations of Definition~\ref{hir:marked-sum-product} only combine the ideals and the marks. (2) The parts are
coherent (Proposition~\ref{bcalc:companion-domain}), so by (1) and Lemma~\ref{bcalc:monomial-part-etale}(1)
\(\varphi^*\mathcal O(\underline{\mathcal I},o)\) is the companion ideal of \(\varphi^*\underline{\mathcal I}\); the inclusion of an open set and a biholomorphism have holomorphic local
inverses. For \(\pi\) along \(\emptyset\): both sides have the submanifold \(\pi^{-1}N\) and the divisor \(\pi^{-1}E\) followed by the empty exceptional member, and the same mark
(Lemma~\ref{mc:empty-blowup-inverse}). On \(\pi^{-1}N\) the transform \(\mathfrak a'\) has the restriction of \(\pi^*\mathfrak a\) (Lemma~\ref{mc:empty-blowup-inverse}), so
\(\mathcal M_{\pi^{-1}N}(\mathfrak a')=\mathcal M_{\pi^{-1}N}(\pi^*\mathfrak a)=\pi^*\mathcal M_N(\mathfrak a)\) and likewise for \(\mathcal N\) (Lemma~\ref{bcalc:monomial-part-etale}(1)--(3),
\(\pi\) having holomorphic local inverses); hence the ideal of \(\mathcal O(\underline{\mathcal I}',o)\) is \(\pi^*\) of the ideal of \(\mathcal O(\underline{\mathcal I},o)\) by (1), which has
on \(\pi^{-1}N\) the restriction of the ideal of \(\mathcal O(\underline{\mathcal I},o)'\) (Lemma~\ref{mc:empty-blowup-inverse} again). (3) The parts only depend on
\(\operatorname{Res}_N\mathfrak a\) (Lemma~\ref{bcalc:monomial-part-etale}(2)), and \(N\), \(E\), \(d\) agree. (4) Induction along \(T\) with (2): restricted-equivalent marked ideals have
the same transform (Lemma~\ref{hir:restricted-equivalence-basic}; the stages and their companion ideals lie in the domain, Propositions~\ref{bcalc:inv-domain-stable}
and~\ref{bcalc:companion-domain}). For the second statement, \(T|_W\) has only empty centres (Lemma~\ref{bcalc:restrict-empty-tower}), \((\underline{\mathcal I}|_W)_{T|_W}\) is
\(\underline{\mathcal I}_T|_{T^{-1}W}\) (Lemma~\ref{hir:open-restriction-tower}), and the companion ideal commutes with restriction to open sets by (2), at both stages.
(5) Lemma~\ref{bcalc:inv-rebase} with \(T_0=T|_W\) and the restricted equivalence of (4); the transforms along a canonical segment lie in the domain
(Theorem~\ref{hir:invariant-driver}).
\end{proof}

\paragraph{The companion ideal on an open set.} The general step (Proposition~\ref{hir:general-step}) takes, at a point where the order \(o\) of the
nonmonomial part is positive, an open set \(W\) on which this order is at most \(o\) at the points of the cosupport, and the companion ideal of
\(\underline{\mathcal I}|_W\). Lemmas~\ref{hir:companion-reduction-along} and~\ref{hir:companion-maximal-order-along} are stated for an abstract factorization
\((M,N,E,\mathfrak b\mathfrak c,d)\) with the bound on all of \(M\); the statement below is their form for the marked ideal itself, on an open set, which is the
open-set statement Lemma~\ref{hir:companion-basic} records as missing. We write \(\mathcal O(\underline{\mathcal I},o)\) as in Lemma~\ref{bcalc:companion-rebase}.

\begin{lemma}[The companion ideal on an open set]
\label{bcalc:companion-open}
\uses{bcalc:companion-rebase,bcalc:monomial-part-etale,hir:monomial-part-along,hir:companion-reduction-along,hir:companion-maximal-order-along,mc:maximal-order-along,hir:giraud-along,hir:marked-equivalence-basic,hir:restricted-equivalence-basic,hir:open-restriction-marked,bcalc:inv-domain-stable,mc:derivative-along-etale}
\source{wlodarczyk05-hironaka,bierstone-milman08-functoriality}
\lean{Hironaka.CompanionOpen.restrictEquiv_mul_parts,Hironaka.CompanionOpen.cosupport_mul_parts,Hironaka.CompanionOpen.isClosed_cosupport_nonmonomialPartAlong,Hironaka.CompanionOpen.isMaximalOrderAlong_companionOf,Hironaka.CompanionOpen.companionOf_reduction,Hironaka.CompanionOpen.bound_pullback_val,Hironaka.CompanionOpen.isMaximalOrderAlong_companionOf_pullback_val}\upstreamlean{TranscendentalNakamaye.Pending.CompanionOpen.companionOf_reduction_pullback_val}
\mixedleanok
Let \(n>0\) and \(\underline{\mathcal I}=(M,N,E,\mathfrak a,d)\) be well formed with \(\mathfrak a\) nonzero along \(N\) at every point of \(N\); write \(\mathcal N=\mathcal N_N(\mathfrak a)\),
\(\mathcal M=\mathcal M_N(\mathfrak a)\).
\begin{enumerate}
\item \(\underline{\mathcal I}\simeq(M,N,E,\mathcal N\mathcal M,d)\) (restricted equivalence); in particular both have the cosupport of \(\underline{\mathcal I}\). For every \(q\) the
set of points of \(N\) where \(\operatorname{ord}(\mathcal N|_N)\ge q\) is closed.
\item Let \(o\ge1\) with \(\operatorname{ord}_x(\mathcal N|_N)\le o\) at every \(x\in\operatorname{cosupp}\underline{\mathcal I}\). Then \(\mathcal O(\underline{\mathcal I},o)\) is of maximal
order along \(N\) (Definition~\ref{mc:maximal-order-along}); every multiple test blow-up \(T\) of \(\mathcal O(\underline{\mathcal I},o)\) is one of \(\underline{\mathcal I}\) and of
\((M,N,E,\mathcal N,o)\); along it \(\operatorname{ord}_z(\mathcal N_T|_{N_T})\le o\) at every \(z\in\operatorname{cosupp}\underline{\mathcal I}_T\), and
\(\operatorname{cosupp}\mathcal O(\underline{\mathcal I},o)_T=\operatorname{cosupp}(\mathcal N_T,o)\cap\operatorname{cosupp}\underline{\mathcal I}_T\).
\item Let \(W\subset M\) be open, \(o\ge1\), and \(\operatorname{ord}_x(\mathcal N|_N)\le o\) only at the points \(x\in\operatorname{cosupp}\underline{\mathcal I}\cap W\). Then the
conclusions of (2) hold for \(\underline{\mathcal I}|_W\) and \(\mathcal O(\underline{\mathcal I}|_W,o)=\mathcal O(\underline{\mathcal I},o)|_W\) (if \(\underline{\mathcal I}\) lies in the
domain of resolution invariants, Definition~\ref{hir:centre-function}).
\end{enumerate}
\end{lemma}
\begin{proof}
\mixedleanok
(1) On \(N\), \(\mathfrak a=\mathcal M\mathcal N\) (Proposition~\ref{hir:monomial-part-along}), so the restrictions agree; restricted-equivalent marked ideals have the same cosupport
(Lemma~\ref{hir:restricted-equivalence-basic}). \(\mathcal N\) is coherent (Proposition~\ref{hir:monomial-part-along}) and the cosupport along a closed submanifold of a coherent
ideal is closed (Lemma~\ref{hir:giraud-along}). (2) By (1) the bound holds on \(\operatorname{cosupp}((\mathcal N\mathcal M)|_N,d)\), and
Lemmas~\ref{hir:companion-maximal-order-along} and~\ref{hir:companion-reduction-along} apply with \(\mathfrak b=\mathcal N\), \(\mathfrak c=\mathcal M\). They give the statements
for \((M,N,E,\mathcal N\mathcal M,d)\), which is equivalent to \(\underline{\mathcal I}\) by (1) (Lemma~\ref{hir:marked-equivalence-basic}: the same multiple test blow-ups and the
same cosupports along them). (3) \(\underline{\mathcal I}|_W\) is well formed and nonzero along \(N\cap W\) (Proposition~\ref{bcalc:inv-domain-stable}); its cosupport is
\(\operatorname{cosupp}\underline{\mathcal I}\cap W\) (Lemma~\ref{hir:open-restriction-marked}), its nonmonomial part is \(\mathcal N|_W\) and the cosupports along \(N\cap W\) of
\(\mathcal N|_W\) are the traces of those of \(\mathcal N\) (Lemma~\ref{bcalc:monomial-part-etale}(1), the inclusion of \(W\) having holomorphic local inverses,
Proposition~\ref{mc:derivative-along-etale}). So the bound holds on \(\operatorname{cosupp}(\underline{\mathcal I}|_W)\), and (2) applies to \(\underline{\mathcal I}|_W\); the last
equality is Lemma~\ref{bcalc:companion-rebase}(2).
\end{proof}

\paragraph{Order zero of the nonmonomial part.} In the general step (Proposition~\ref{hir:general-step}), at a point where the order of the nonmonomial
part is \(0\) the marked ideal is monomial near the point and the invariant is that of the monomial case (Lemma~\ref{hir:monomial-case}). The lemma below
gives the open set on which this holds.

\begin{lemma}[Where the nonmonomial part has order zero]
\label{bcalc:nonmonomial-order-zero}
\uses{bcalc:companion-open,bcalc:monomial-part-etale,hir:monomial-part,hir:monomial-part-along,ana:restricted-colon,ana:restricted-calculus,hir:order-basic,mc:derivative-along-etale}
\source{wlodarczyk05-hironaka,bierstone-milman08-functoriality}
\lean{Hironaka.NonmonomialOrderZero.restrictLe_monomialPartAlong_of_notMem}\upstreamlean{TranscendentalNakamaye.Pending.NonmonomialOrderZero.isOpen_compl_cosupport_one}\lean{Hironaka.NonmonomialOrderZero.isMonomialAlong_pullback_val}\upstreamlean{TranscendentalNakamaye.Pending.NonmonomialOrderZero.isMonomialAlong_pullback_val_compl}
\mixedleanok
Let \(\underline{\mathcal I}=(M,N,E,\mathfrak a,d)\) with \(\mathfrak a\) coherent, and write \(\mathcal M=\mathcal M_N(\mathfrak a)\), \(\mathcal N=\mathcal N_N(\mathfrak a)\).
\begin{enumerate}
\item If \(y\in N\) and \(\operatorname{ord}_y(\mathcal N|_N)=0\), then \(\operatorname{Res}_N\mathcal M\subset\operatorname{Res}_N\mathfrak a\) at \(y\).
\item If \(W\subset M\) is open and \(\operatorname{ord}_y(\mathcal N|_N)=0\) at every \(y\in N\cap W\), then \(\underline{\mathcal I}|_W\) is monomial along \(N\cap W\)
(Definition~\ref{hir:monomial-part}).
\item If \(n>0\) and \(\underline{\mathcal I}\) is well formed with \(\mathfrak a\) nonzero along \(N\) at every point of \(N\), the set \(W\) of points that are off \(N\) or where
\(\operatorname{ord}(\mathcal N|_N)=0\) is open, and \(\underline{\mathcal I}|_W\) is monomial along \(N\cap W\).
\end{enumerate}
\end{lemma}
\begin{proof}
\mixedleanok
(1) The order is \(0\), so some \(c\in\mathcal N_y\) does not vanish at \(y\) (Lemma~\ref{hir:order-basic}). As \(\mathcal N\) is the quotient
\((\widetilde{\mathfrak a}^N:\mathcal M)\) (Definition~\ref{hir:monomial-part}, Lemma~\ref{ana:restricted-colon}), for \(g\in\mathcal M_y\) the restriction of \(gc\) to \(N\) lies
in \(\operatorname{Res}_N\mathfrak a\) near \(y\); \(c\) is holomorphic and nonzero near \(y\), so it can be cancelled from this relation. (2) By (1) the restricted inclusion holds
at the points of \(N\cap W\); it pulls back along the inclusion of \(W\) (Lemma~\ref{ana:restricted-calculus}), and \(\mathcal M_{N\cap W}(\mathfrak a|_W)=\mathcal M|_W\)
(Lemma~\ref{bcalc:monomial-part-etale}(1), the inclusion having holomorphic local inverses, Proposition~\ref{mc:derivative-along-etale}). (3) The points of \(N\) where
\(\operatorname{ord}(\mathcal N|_N)\ge1\) form a closed set (Lemma~\ref{bcalc:companion-open}(1)); its complement is \(W\); apply (2).
\end{proof}

\paragraph{The order of the nonmonomial part along a tower.} In the general step (Proposition~\ref{hir:general-step}) the value at a point
\(y\) of the cosupport along a tower \(T\) begins with the order at \(y\) of the nonmonomial part of the transform \(\underline{\mathcal I}_T\), taken
with its whole divisor \(E_T\). The lemma below collects the properties of this order that do not depend on the invariant.

\begin{lemma}[The order of the nonmonomial part]
\label{bcalc:nonmonomial-order-local}
\uses{hir:monomial-part-along,hir:nonmonomial-transform-along,hir:exponent-along,hir:cosupport-basic,bcalc:companion-open,bcalc:monomial-part-etale,bcalc:companion-rebase,mc:empty-blowup-inverse,mc:derivative-along-etale,bcalc:restrict-empty-tower,bcalc:inv-rebase,bcalc:inv-domain-stable,hir:open-restriction-tower,hir:open-restriction-marked,bcalc:tower-restrict,hir:tower-isomorphism,mc:tower-isomorphism-transform,hir:restricted-calculus,hir:companion-reduction-along,hir:centre-function,hir:test-tower}
\source{wlodarczyk05-hironaka,bierstone-milman08-functoriality}
\lean{Hironaka.NonmonomialOrderLocal.nmPart,Hironaka.NonmonomialOrderLocal.ordGe,Hironaka.NonmonomialOrderLocal.nmOrd,Hironaka.NonmonomialOrderLocal.ordGe_anti,Hironaka.NonmonomialOrderLocal.ordGe_subset_N,Hironaka.NonmonomialOrderLocal.nmOrd_congr,Hironaka.NonmonomialOrderLocal.nmOrd_le_of_notMem,Hironaka.NonmonomialOrderLocal.mem_ordGe_of_succ_le,Hironaka.NonmonomialOrderLocal.succ_le_nmOrd_of_mem,Hironaka.NonmonomialOrderLocal.exists_notMem_ordGe,Hironaka.NonmonomialOrderLocal.mem_ordGe_iff,Hironaka.NonmonomialOrderLocal.isClosed_ordGe,Hironaka.NonmonomialOrderLocal.eventually_nmOrd_le,Hironaka.NonmonomialOrderLocal.mem_ordGe_pullback_iff_of_local}\upstreamlean{TranscendentalNakamaye.Pending.NonmonomialOrderLocal.nmOrd_pullback_of_local}\lean{Hironaka.NonmonomialOrderLocal.mem_ordGe_pullback_val_iff,Hironaka.NonmonomialOrderLocal.nmOrd_pullback_val,Hironaka.NonmonomialOrderLocal.mem_ordGe_pullback_diffeo_iff}\upstreamlean{TranscendentalNakamaye.Pending.NonmonomialOrderLocal.nmOrd_pullback_diffeo}\lean{Hironaka.NonmonomialOrderLocal.mem_ordGe_geoTransform_empty_iff,Hironaka.NonmonomialOrderLocal.mem_ordGe_geoTransformAlong_of_isEmptyTower,Hironaka.NonmonomialOrderLocal.mem_ordGe_geoTransformAlong_of_centresAvoid}\upstreamlean{TranscendentalNakamaye.Pending.NonmonomialOrderLocal.nmOrd_geoTransformAlong_of_centresAvoid}\lean{Hironaka.NonmonomialOrderLocal.mem_ordGe_geoTransform_of_notMem,Hironaka.NonmonomialOrderLocal.nmOrd_geoTransform_of_notMem,Hironaka.NonmonomialOrderLocal.mem_ordGe_geoTransformAlong_isTowerIso,Hironaka.NonmonomialOrderLocal.nmOrd_geoTransformAlong_isTowerIso,Hironaka.NonmonomialOrderLocal.companionOf_tower_facts,Hironaka.NonmonomialOrderLocal.notMem_ordGe_geoTransform_of_mem,Hironaka.NonmonomialOrderLocal.nmOrd_geoTransform_le_of_mem}
\mixedleanok
Let \(n\ge1\). For a well-formed marked ideal \(\underline{\mathcal I}=(M,N,E,\mathfrak a,d)\) with \(\mathfrak a\) nonzero along \(N\) at every point of \(N\), write
\(\nu_{\underline{\mathcal I}}(y)=\operatorname{ord}_y(\mathcal N_N(\mathfrak a)|_N)\) for \(y\in N\) (Proposition~\ref{hir:monomial-part-along}).
\begin{enumerate}
\item \(\nu_{\underline{\mathcal I}}(y)\) is finite for every \(y\in N\), the sets \(\{\nu_{\underline{\mathcal I}}\ge q\}\) are closed, and \(\nu_{\underline{\mathcal I}}\) is upper
semicontinuous on \(N\).
\item If \(\varphi\colon M'\to M\) is a local biholomorphism (for instance the inclusion of an open set, or a biholomorphism), then
\(\nu_{\varphi^*\underline{\mathcal I}}=\nu_{\underline{\mathcal I}}\circ\varphi\). If \(\pi\) is the blowup along \(\emptyset\) and \(\underline{\mathcal I}'\) the transform, then
\(\nu_{\underline{\mathcal I}'}=\nu_{\underline{\mathcal I}}\circ\pi\).
\item Let \(\underline{\mathcal I}\) moreover lie in the domain of Definition~\ref{hir:centre-function}, and let \(T\) be a multiple test blow-up of it. If the centres of
\(T\) miss an open set \(W\subset M\), then \(\nu_{\underline{\mathcal I}_T}=\nu_{\underline{\mathcal I}}\circ p_T\) at the points over \(W\); in particular, for a test blowup
\(\sigma\) along \(C\), \(\nu_{\underline{\mathcal I}_\sigma}(z)=\nu_{\underline{\mathcal I}}(\sigma(z))\) whenever \(\sigma(z)\notin C\). If \(\chi\colon T'\to T\) is an
isomorphism of towers over a biholomorphism \(\Phi\) (Definition~\ref{hir:tower-isomorphism}), then
\(\nu_{(\Phi^*\underline{\mathcal I})_{T'}}=\nu_{\underline{\mathcal I}_T}\circ\chi\).
\item Let \(o\ge1\) with \(\nu_{\underline{\mathcal I}}\le o\) on \(\operatorname{cosupp}\underline{\mathcal I}\), and let \(T\) be a multiple test blow-up of the companion ideal
\(\mathcal O(\underline{\mathcal I},o)\) (Lemma~\ref{hir:companion-reduction-along}). Then \(T\) is a multiple test blow-up of \(\underline{\mathcal I}\),
\(\nu_{\underline{\mathcal I}_T}\le o\) on \(\operatorname{cosupp}\underline{\mathcal I}_T\), and
\(\operatorname{cosupp}\mathcal O(\underline{\mathcal I},o)_T\subset\{\nu_{\underline{\mathcal I}_T}\ge o\}\cap\operatorname{cosupp}\underline{\mathcal I}_T\).
\item Let \(\underline{\mathcal I}\) lie in the domain and \(\sigma\) be a test blowup of it along \(C\) with \(\nu_{\underline{\mathcal I}}=q\ge1\) on \(C\). Then
\(\nu_{\underline{\mathcal I}_\sigma}\le q\) at the points of \(\operatorname{cosupp}\underline{\mathcal I}_\sigma\) over \(C\).
\end{enumerate}
\end{lemma}
\begin{proof}
\mixedleanok
(1) \(\mathfrak a\subset\mathcal N_N(\mathfrak a)\) (Theorem~\ref{hir:nonmonomial-transform-along}), and a germ of \(\mathfrak a_y\) not vanishing on \(N\) near \(y\) has finite
order along \(N\): through a local retraction onto \(N\) a germ of infinite order along \(N\) vanishes on \(N\) near \(y\) (Lemma~\ref{hir:exponent-along}). The sets
\(\{\nu\ge q\}\) are closed by Lemma~\ref{bcalc:companion-open}(1); hence near \(y\in N\) the closed set \(\{\nu\ge\nu(y)+1\}\) is avoided, and \(\nu\le\nu(y)\)
there. The order is the least \(q\) with \(y\notin\{\nu\ge q+1\}\), so it is determined by the sets \(\{\nu\ge q\}\) (and \(\{\nu\ge0\}=N\),
Lemma~\ref{hir:cosupport-basic}).

(2) \(\varphi^*\mathcal N_N(\mathfrak a)=\mathcal N_{\varphi^{-1}N}(\varphi^*\mathfrak a)\), and cosupports along \(N\) pull back along local biholomorphisms
(Lemma~\ref{bcalc:monomial-part-etale}; inclusions of open sets and biholomorphisms have holomorphic local inverses, Proposition~\ref{mc:derivative-along-etale},
Lemma~\ref{bcalc:companion-rebase}). Across a blowup along \(\emptyset\), \(N'=\pi^{-1}N\), the divisor is \(\pi^{-1}E\) followed by an empty member,
\(\operatorname{Res}_{N'}\mathfrak a'=\operatorname{Res}_{N'}\pi^*\mathfrak a\), and \(\pi\) is a local biholomorphism (Lemma~\ref{mc:empty-blowup-inverse}); the nonmonomial
part only depends on the restriction to \(N'\) and ignores an appended empty member (Lemma~\ref{bcalc:monomial-part-etale}), so it is
\(\pi^*\mathcal N_N(\mathfrak a)\), and the first assertion applies.

(3) Over \(W\) the tower restricts to one with only empty centres (Lemma~\ref{bcalc:restrict-empty-tower}), along which (2) applies stage by stage, the stages
staying in the domain (Lemma~\ref{bcalc:inv-rebase}); the transform of \(\underline{\mathcal I}|_W\) along \(T|_W\) is the restriction of \(\underline{\mathcal I}_T\) to the
preimage of \(W\) (Lemma~\ref{hir:open-restriction-tower}), restriction to open sets does not change the order by (2), and \(p_{T|_W}=p_T\) on points
(Lemma~\ref{bcalc:tower-restrict}). For one blowup take \(W=M\setminus C\). For isomorphisms of towers,
\((\Phi^*\underline{\mathcal I})_{T'}=\chi^*\underline{\mathcal I}_T\) (Lemma~\ref{mc:tower-isomorphism-transform}) with \(\underline{\mathcal I}_T\) in the domain
(Lemma~\ref{bcalc:inv-domain-stable}); apply (2) to \(\chi\).

(4) Theorem~\ref{hir:nonmonomial-transform-along} gives the first two assertions, and \(\operatorname{Res}_{N_T}\mathcal N_{N_T}(\mathfrak a_T)\subset\operatorname{Res}_{N_T}\mathcal N_T\)
for the transform \(\mathcal N_T\) of \((M,N,E,\mathcal N_N(\mathfrak a),o)\), which lives on \(N_T\). By Lemma~\ref{bcalc:companion-open}(2),
\(\operatorname{cosupp}\mathcal O(\underline{\mathcal I},o)_T=\operatorname{cosupp}(\mathcal N_T|_{N_T},o)\cap\operatorname{cosupp}\underline{\mathcal I}_T\), and by the restricted
inclusion \(\operatorname{cosupp}(\mathcal N_T|_{N_T},o)\subset\{\nu_{\underline{\mathcal I}_T}\ge o\}\) (Lemma~\ref{hir:restricted-calculus}).

(5) The set \(W=\{\nu_{\underline{\mathcal I}}\le q\}\) is open by (1) and contains \(C\); on \(W\) the order is at most \(q\). By (2) and
Lemma~\ref{bcalc:companion-open}(2) (for the empty tower), \(C\) lies in \(\{\nu\ge q\}\cap\operatorname{cosupp}\underline{\mathcal I}|_W=\operatorname{cosupp}\mathcal O(\underline{\mathcal I}|_W,q)\);
the companion ideal has the submanifold and the divisor of \(\underline{\mathcal I}|_W\), and \(C\) is admissible for \(\underline{\mathcal I}|_W\) (Lemma~\ref{hir:open-restriction-tower}),
so the restriction of \(\sigma\) to \(\sigma^{-1}W\) is a test blowup of \(\mathcal O(\underline{\mathcal I}|_W,q)\) (Definition~\ref{hir:test-tower}). By (4),
\(\nu\le q\) on the cosupport of its transform, which is the restriction of \(\underline{\mathcal I}_\sigma\) to \(\sigma^{-1}W\) and contains the cosupport points over \(C\);
conclude with (2) for the inclusion of \(\sigma^{-1}W\).
\end{proof}

\paragraph{The stage at which a point reached its order.} In the general step the companion ideal at a point \(y\) of order \(o\ge1\) is
formed at the stage at which the images of \(y\) first had order \(o\), on a neighbourhood on which the order is at most \(o\); nearby points are
compared across the steps between their stages. The lemma below is the analogue for the general step of Lemmas~\ref{hir:regime-run}
and~\ref{hir:step-one-mask-tower} for Step~1.

\begin{lemma}[The regime of the order of the nonmonomial part]
\label{bcalc:nonmonomial-order-regime}
\uses{bcalc:nonmonomial-order-local,hir:regime-run,bcalc:restrict-empty-tower,bcalc:companion-rebase,hir:open-restriction-tower,hir:open-restriction-marked,bcalc:tower-restrict,hir:tower-isomorphism,mc:tower-isomorphism-transform,bcalc:transform-off-centre,hir:geo-transform-basic,hir:test-tower,bcalc:inv-domain-stable,hir:companion-reduction-along,hir:centre-function}
\source{wlodarczyk05-hironaka,bierstone-milman08-functoriality}
\lean{Hironaka.NonmonomialOrderRegime.regimeRun_restrict_of,Hironaka.NonmonomialOrderRegime.regimeRun_isTowerIso_of,Hironaka.NonmonomialOrderRegime.OrdRegime,Hironaka.NonmonomialOrderRegime.ordRegime_iff,Hironaka.NonmonomialOrderRegime.ordRegime_step_iff_of_notMem,Hironaka.NonmonomialOrderRegime.ordRegime_restrict_iff,Hironaka.NonmonomialOrderRegime.regimeRun_ordRegime_restrict,Hironaka.NonmonomialOrderRegime.ordRegime_isTowerIso_iff,Hironaka.NonmonomialOrderRegime.regimeRun_ordRegime_isTowerIso,Hironaka.NonmonomialOrderRegime.notMem_ordGe_of_mask,Hironaka.NonmonomialOrderRegime.mem_ordGe_of_mem_centre_of_mask,Hironaka.NonmonomialOrderRegime.ordRegime_and_le_regimeRun,Hironaka.NonmonomialOrderRegime.CentresOrdGt,Hironaka.NonmonomialOrderRegime.centresOrdGt_step_iff,Hironaka.NonmonomialOrderRegime.centresAvoid_of_centresOrdGt}\upstreamlean{TranscendentalNakamaye.Pending.NonmonomialOrderRegime.companionOf_rebase_of_centresOrdGt}
\mixedleanok
Let \(n\ge1\), \(\underline{\mathcal I}\) a marked ideal in the domain of Definition~\ref{hir:centre-function}, \(\nu\) the order of the nonmonomial part of
Lemma~\ref{bcalc:nonmonomial-order-local}, and \(o\in\mathbf N\). A point \(y\) of the last stage of a tower \(T\) is \emph{in the regime \(o\)} if
\(y\in\operatorname{cosupp}\underline{\mathcal I}_T\) and \(\nu_{\underline{\mathcal I}_T}(y)=o\); write \(r_o(T,y)\) for its run (Lemma~\ref{hir:regime-run}).
\begin{enumerate}
\item The runs of Lemma~\ref{hir:regime-run}(2), (3) under restriction to open sets and along isomorphisms of towers hold as soon as the regimes correspond
along the towers with a property inherited by prefixes. Along multiple test blow-ups of \(\underline{\mathcal I}\), the regime \(o\) and its run correspond under
restriction to an open \(U\subset X\) (for \(\underline{\mathcal I}|_U\) and \(T|_U\)) and along an isomorphism of towers over a biholomorphism \(\Phi\) (for
\(\Phi^*\underline{\mathcal I}\)); off the centre of a test blowup \(\sigma\), \(z\) is in the regime \(o\) iff \(\sigma(z)\) is.
\item (The mask) Let \(o\ge1\), \(T_0T_1\) a multiple test blow-up of \(\underline{\mathcal I}\), and \(W\) an open set of the last stage of \(T_0\) with
\(\nu_{\underline{\mathcal I}_{T_0}}\le o\) on \(\operatorname{cosupp}\underline{\mathcal I}_{T_0}\cap W\), such that \(T_1|_W\) is a multiple test blow-up of the
companion ideal \(\mathcal O(\underline{\mathcal I}_{T_0}|_W,o)\). Then every point \(z\) of \(\operatorname{cosupp}\underline{\mathcal I}_{T_0T_1}\) over \(W\) with
\(\nu(z)\ge o\) is in the regime \(o\) and \(r_o(T_0T_1,z)\ge\operatorname{length}T_1\).
\item Let \(W\subset X\) be open with \(\nu_{\underline{\mathcal I}}\le o\) on \(\operatorname{cosupp}\underline{\mathcal I}\cap W\), and \(T\) a multiple test blow-up of
\(\underline{\mathcal I}\) each of whose centres lies in \(\{\nu\ge o+1\}\) at its stage. Then the centres of \(T\) miss the preimages of \(W\), and
\(\nu_{\underline{\mathcal I}_T}\le o\) at the points of \(\operatorname{cosupp}\underline{\mathcal I}_T\) over \(W\). Consequently, for \(o\ge1\), the companion ideal
\(\mathcal O(\underline{\mathcal I}_T,o)|_{p_T^{-1}W}\) re-bases to \((\mathcal O(\underline{\mathcal I},o)|_W)_{T|_W}\) with the conclusions of
Lemma~\ref{bcalc:companion-rebase}(5).
\end{enumerate}
\end{lemma}
\begin{proof}
\mixedleanok
(1) Induction along the tower as in Lemma~\ref{hir:regime-run}, using the property only for the prefixes. Under restriction, the transform of
\(\underline{\mathcal I}|_U\) along \(T|_U\) is the restriction of \(\underline{\mathcal I}_T\) (Lemma~\ref{hir:open-restriction-tower}), cosupports restrict
(Lemma~\ref{hir:open-restriction-marked}) and so does \(\nu\) (Lemma~\ref{bcalc:nonmonomial-order-local}(2)); along an isomorphism \(\chi\), cosupports
correspond (Lemma~\ref{mc:tower-isomorphism-transform}) and so does \(\nu\) (Lemma~\ref{bcalc:nonmonomial-order-local}(3)); off the centre the cosupport
(Lemma~\ref{bcalc:transform-off-centre}) and \(\nu\) (Lemma~\ref{bcalc:nonmonomial-order-local}(3)) are those of the image.

(2) Write \(\underline{\mathcal I}_0=\underline{\mathcal I}_{T_0}\), in the domain (Lemma~\ref{bcalc:inv-domain-stable}), and \(\underline{\mathcal I}_{T_0T_1}=(\underline{\mathcal I}_0)_{T_1}\)
(Lemma~\ref{hir:geo-transform-basic}). For every prefix \(S\) of \(T_1\), \(S|_W\) is a multiple test blow-up of \(\mathcal O(\underline{\mathcal I}_0|_W,o)\), so by
Lemma~\ref{bcalc:nonmonomial-order-local}(4) on \(W\) (with (2) of that lemma for the restrictions) \(\nu\le o\) at the cosupport points over \(W\), and the
centre of the next step, over \(W\), lies in the cosupport of the transform of the companion ideal, hence in \(\{\nu\ge o\}\). Induction along \(T_1\): \(z\) is
in the regime \(o\) (\(\nu(z)\ge o\) and \(\le o\)); its image \(w=\sigma(z)\) lies in the cosupport (in the centre, or off it by
Lemma~\ref{bcalc:transform-off-centre}) and has \(\nu(w)\ge o\) (in the centre by the above, off it by Lemma~\ref{bcalc:nonmonomial-order-local}(3)); so
\(w\) is in the regime, \(r_o(z)=r_o(w)+1\), and by induction \(r_o(w)\ge\operatorname{length}S\).

(3) Induction along \(T\). A centre \(C\) lies in the cosupport, where over \(W\) the order is \(\le o\) by induction, and in \(\{\nu\ge o+1\}\); so it misses
the preimage of \(W\). Then a cosupport point over \(W\) is off the centre, its image lies in the cosupport (Lemma~\ref{bcalc:transform-off-centre}) and has
the same order (Lemma~\ref{bcalc:nonmonomial-order-local}(3)), at most \(o\). The re-basing is Lemma~\ref{bcalc:companion-rebase}(5), whose hypothesis is that the
centres miss \(W\) (Lemma~\ref{bcalc:restrict-empty-tower}).
\end{proof}

\paragraph{The order of the nonmonomial part along canonical segments.} In the general step (Proposition~\ref{hir:general-step}) the value at a point
\(y\) of the cosupport begins with the order \(\nu(y)\) of the nonmonomial part, and along a canonical segment every centre is empty or the maximum locus of
the value. The lemma below collects what follows from this first entry alone, for every value function whose order dominates \(\nu\); it supplies the
hypotheses of Lemma~\ref{bcalc:nonmonomial-order-regime}(3) at the start of the run of a point.

\begin{lemma}[The order of the nonmonomial part along canonical segments]
\label{bcalc:nonmonomial-order-canonical}
\uses{bcalc:nonmonomial-order-local,bcalc:nonmonomial-order-regime,hir:nonmonomial-step-along,hir:nonmonomial-transform-along,hir:monomial-factorization-unique-along,hir:monomial-invariant-along,hir:restricted-equivalence-basic,hir:geo-transform-coherent,bcalc:transform-wellformed,bcalc:transform-off-centre,bcalc:inv-domain-stable,hir:centre-function,hir:invariant-order,hir:regime-run,hir:tower-split,hir:test-tower,hir:geo-transform-basic,bcalc:companion-rebase,hir:cosupport-basic}
\source{wlodarczyk05-hironaka,bierstone-milman08-functoriality}
\lean{Hironaka.NonmonomialOrderCanonical.notMem_ordGe_one_geoTransform_of_mem,Hironaka.NonmonomialOrderCanonical.nmOrd_geoTransform_le_of_level,Hironaka.NonmonomialOrderCanonical.OrdMonotone,Hironaka.NonmonomialOrderCanonical.ordMonotone_toLex,Hironaka.NonmonomialOrderCanonical.nmOrd_le_of_mem_maxLocus,Hironaka.NonmonomialOrderCanonical.nmOrd_eq_of_mem_maxLocus,Hironaka.NonmonomialOrderCanonical.subset_ordGe_of_canonical,Hironaka.NonmonomialOrderCanonical.nmOrd_step_le_of_canonical,Hironaka.NonmonomialOrderCanonical.nmOrd_le_nmOrd_proj,Hironaka.NonmonomialOrderCanonical.isInvUSCOn_of_ordMonotone}\upstreamlean{TranscendentalNakamaye.Pending.NonmonomialOrderCanonical.isInvUSCOn_toLex}\lean{Hironaka.NonmonomialOrderCanonical.mem_cosupp_and_succ_le_before_start,Hironaka.NonmonomialOrderCanonical.centresOrdGt_of_canonical,Hironaka.NonmonomialOrderCanonical.centresAvoid_of_canonical}\upstreamlean{TranscendentalNakamaye.Pending.NonmonomialOrderCanonical.restrictEquiv_companionOf_rebase_of_canonical}
\mixedleanok
Let \(n\ge1\), \(\underline{\mathcal I}\) a marked ideal in the domain of Definition~\ref{hir:centre-function} and \(\nu\) the order of the nonmonomial part of
Lemma~\ref{bcalc:nonmonomial-order-local}. Let \(f\) assign to every tower \(T\) over \(X\) a function \(f(T)\) on its last stage with values in a totally ordered set,
such that along every multiple test blow-up \(T\) of \(\underline{\mathcal I}\), \(f(T)(y)\le f(T)(y')\) implies \(\nu_{\underline{\mathcal I}_T}(y)\le\nu_{\underline{\mathcal I}_T}(y')\)
for \(y,y'\in\operatorname{cosupp}\underline{\mathcal I}_T\); for instance \(f(T)(y)=(\nu_{\underline{\mathcal I}_T}(y),g(T)(y))\) ordered lexicographically, for any \(g\).
\begin{enumerate}
\item Let \(\sigma\) be a test blowup of \(\underline{\mathcal I}\) along \(C\) with \(\nu_{\underline{\mathcal I}}=q\) on \(C\), for any \(q\ge0\). Then
\(\nu_{\underline{\mathcal I}_\sigma}\le q\) at the points of \(\operatorname{cosupp}\underline{\mathcal I}_\sigma\) over \(C\). If \(q=0\), \(\nu_{\underline{\mathcal I}_\sigma}=0\) at
every point of the strict transform \(N'\) over \(C\).
\item Let \(T\) be a multiple test blow-up of \(\underline{\mathcal I}\). The maximum locus of \(f(T)\) on \(\operatorname{cosupp}\underline{\mathcal I}_T\) lies in one level set of
\(\nu_{\underline{\mathcal I}_T}\), on which \(\nu\) takes its largest value on the cosupport; hence if \(C\) is empty or this maximum locus and \(w\) is a point of the
cosupport with \(\nu(w)\ge q\), then \(C\subset\{\nu\ge q\}\). If \(T_0M\) is a multiple test blow-up of \(\underline{\mathcal I}\) and a canonical segment of \(f\)
(Definition~\ref{hir:centre-function}), every point \(z\) of \(\operatorname{cosupp}\underline{\mathcal I}_{T_0M}\) lies over \(\operatorname{cosupp}\underline{\mathcal I}_{T_0}\) and
\(\nu(z)\le\nu(p_M(z))\): the order never goes up along a canonical segment.
\item Let \(T\) be a multiple test blow-up of \(\underline{\mathcal I}\). If near every point \(y\) of \(\operatorname{cosupp}\underline{\mathcal I}_T\) the points \(y'\) of the cosupport
with \(\nu(y')=\nu(y)\) have \(f(T)(y')\le f(T)(y)\), then \(f(T)\) is upper semicontinuous on the cosupport (Definition~\ref{hir:invariant-order}).
\item Let \(T\) be a multiple test blow-up of \(\underline{\mathcal I}\) and a canonical segment of \(f\), and \(y\in\operatorname{cosupp}\underline{\mathcal I}_T\) with
\(\nu(y)=o\) and run \(r=r_o(T,y)<\operatorname{length}T\) in the regime \(o\) (Lemma~\ref{bcalc:nonmonomial-order-regime}). Then the image of \(y\) at the stage
\(\operatorname{length}T-r-1\) lies in the cosupport and has order \(\ge o+1\).
\item Let \(T_0M\sigma\) (\(\sigma\) a blowup along \(C\)) be a multiple test blow-up of \(\underline{\mathcal I}\) and a canonical segment of \(f\), and \(w\) a point of
\(\operatorname{cosupp}\underline{\mathcal I}_{T_0M}\) with \(\nu(w)\ge o+1\). Then every centre of \(M\sigma\) lies in \(\{\nu\ge o+1\}\) at its stage. Consequently, for an
open set \(W\) of the last stage of \(T_0\) with \(\nu\le o\) on \(\operatorname{cosupp}\underline{\mathcal I}_{T_0}\cap W\), the centres of \(M\sigma\) miss the preimages
of \(W\), \(\nu\le o\) at the cosupport points over \(W\), and for \(o\ge1\)
\(\mathcal O(\underline{\mathcal I}_{T_0M\sigma},o)|_{(M\sigma)^{-1}W}\simeq(\mathcal O(\underline{\mathcal I}_{T_0},o)|_W)_{(M\sigma)|_W}\).
\end{enumerate}
\end{lemma}
\begin{proof}
\mixedleanok
(1) For \(q\ge1\) this is Lemma~\ref{bcalc:nonmonomial-order-local}(5). Let \(q=0\), and write \(\mathcal N=\mathcal N_N(\mathfrak a)\), \(\mathcal M=\mathcal M_N(\mathfrak a)\).
By Theorem~\ref{hir:nonmonomial-transform-along}, \(\operatorname{Res}_N\mathfrak a=\operatorname{Res}_N(\mathcal N\mathcal M)\), both parts are coherent, \(\mathcal M\) is nonzero
and monomial along \(N\), and \(\mathcal N\) is nonzero along \(N\) and, restricted to \(N\), contained in the ideal of no member at the points of \(N\); restricted-equivalent
marked ideals have the same cosupport and the same transform (Lemma~\ref{hir:restricted-equivalence-basic}). Now \(C\subset\operatorname{cosupp}(\mathcal N\mathcal M|_N,d)\),
\(C\subset N=\operatorname{cosupp}(\mathcal N|_N,0)\) (Lemma~\ref{hir:cosupport-basic}), and \(\operatorname{ord}(\mathcal N|_N)=0\) on \(C\). So Lemma~\ref{hir:nonmonomial-step-along}
with \(o=0\) gives \(\operatorname{Res}_{N'}\mathfrak a'=\operatorname{Res}_{N'}(\mathfrak b'\mathcal M')\), where \(\mathcal M'\) is principal monomial along \(N'\) at every point of
\(N'\), hence nonzero and monomial along \(N'\) (Lemma~\ref{hir:monomial-invariant-along}), and \(\mathfrak b'=(\pi^*\mathcal N:\mathcal O)\) is the ideal of the transform of
\((M,N,E,\mathcal N,0)\), coherent (Lemma~\ref{hir:geo-transform-coherent}) and, restricted to \(N'\), in the ideal of no member of \(E'\). The transform is well formed
(Lemma~\ref{bcalc:transform-wellformed}), so by Lemma~\ref{hir:monomial-factorization-unique-along} \(\mathcal M_{N'}(\mathfrak a')=\mathcal M_{N'}(\mathcal M')\) and
\(\mathfrak b'\subset\mathcal N_{N'}(\mathfrak a')\). Let \(z\in N'\) with \(\sigma(z)\in C\). Since \(\nu(\sigma(z))=0\), some germ \(c\in\mathcal N_{\sigma(z)}\) has \(c(\sigma(z))\ne0\);
then \(c\circ\sigma\) is a germ of \(\pi^*\mathcal N\subset\mathfrak b'\subset\mathcal N_{N'}(\mathfrak a')\) at \(z\) which does not vanish at \(z\), so \(\nu(z)=0\).

(2) For \(c\) in the maximum locus and \(y\) in the cosupport, \(f(T)(y)\le f(T)(c)\), so \(\nu(y)\le\nu(c)\); two points of the maximum locus have the same value,
hence the same order. If \(\nu(w)\ge q\), a point \(c\) of the maximum locus has \(\nu(c)\ge q\), and \(c\in N_T\), so \(c\in\{\nu\ge q\}\)
(Lemma~\ref{bcalc:nonmonomial-order-local}(1); the stages lie in the domain, Proposition~\ref{bcalc:inv-domain-stable}). For one canonical step \(\sigma\) along \(C\) and
\(z\) in the new cosupport, \(\sigma(z)\) lies in the old cosupport (in \(C\), or off \(C\) by Lemma~\ref{bcalc:transform-off-centre}). If \(\sigma(z)\notin C\),
\(\nu(z)=\nu(\sigma(z))\) (Lemma~\ref{bcalc:nonmonomial-order-local}(3)); if \(\sigma(z)\in C\), then \(C\) is the maximum locus, which lies in the level set of
\(\nu(\sigma(z))\), and (1) applies. Induction along \(M\).

(3) Near \(y\), \(\nu\le\nu(y)\) (Lemma~\ref{bcalc:nonmonomial-order-local}(1): the set \(\{\nu\ge\nu(y)+1\}\) is closed). At a nearby point \(y'\) of the cosupport
with \(\nu(y')<\nu(y)\), \(f(T)(y)\le f(T)(y')\) would give \(\nu(y)\le\nu(y')\); so \(f(T)(y')<f(T)(y)\). At the points with \(\nu(y')=\nu(y)\) the hypothesis
applies.

(4) Split \(T=T_0M\) after its first \(\operatorname{length}T-r-1\) blowups (Lemma~\ref{hir:tower-split}) and let \(w=p_M(y)\). By the maximality of the run
(Lemma~\ref{hir:regime-run}) \(w\) is not in the regime \(o\). By (2) for \(T_0M\), \(w\) lies in the cosupport and \(\nu(w)\ge\nu(y)=o\); hence \(\nu(w)\ge o+1\).

(5) Induction along \(M\). The centre \(C\) is empty or the maximum locus, so \(C\subset\{\nu\ge o+1\}\) by (2). If \(M=M'\tau\), the image of \(w\) under \(\tau\)
lies in \(\operatorname{cosupp}\underline{\mathcal I}_{T_0M'}\) and has order \(\ge\nu(w)\ge o+1\) by (2), and the induction hypothesis gives the centres of \(M'\tau\).
The tower \(M\sigma\) is a multiple test blow-up of \(\underline{\mathcal I}_{T_0}\) (Lemma~\ref{hir:geo-transform-basic}), with transforms
\((\underline{\mathcal I}_{T_0})_{M'}=\underline{\mathcal I}_{T_0M'}\); the rest is Lemma~\ref{bcalc:nonmonomial-order-regime}(3) and
Lemma~\ref{bcalc:companion-rebase}(4) for \(\underline{\mathcal I}_{T_0}\).
\end{proof}

\paragraph{The companion phase after the start of a run.} In the general step the value at a point \(y\) of order \(o\ge1\) is read off the
companion ideal formed at the start of the run of \(y\) in the regime \(o\), on the open set where \(\nu\le o\), along the rest of the segment. The lemma
below shows that this rest, restricted to that open set, is a multiple test blow-up of the companion ideal (the hypothesis of
Lemma~\ref{bcalc:nonmonomial-order-regime}(2)), and that nearby points of the same order started no later.

\begin{lemma}[The restricted tail of the companion phase]
\label{bcalc:nonmonomial-order-tail}
\uses{bcalc:nonmonomial-order-canonical,bcalc:nonmonomial-order-local,bcalc:nonmonomial-order-regime,bcalc:companion-open,hir:nonmonomial-transform-along,hir:companion-reduction-along,hir:open-restriction-tower,hir:open-restriction-marked,hir:test-tower,hir:geo-transform-basic,hir:regime-run,hir:tower-split,bcalc:inv-domain-stable,hir:centre-function}
\source{wlodarczyk05-hironaka,bierstone-milman08-functoriality}
\lean{Hironaka.NonmonomialOrderTail.E_geoTransformAlong_congr,Hironaka.NonmonomialOrderTail.cosupp_companionOf_geoTransformAlong,Hironaka.NonmonomialOrderTail.isGeoTestTower_companion_of_canonical,Hironaka.NonmonomialOrderTail.ordMask,Hironaka.NonmonomialOrderTail.mem_ordMask_iff,Hironaka.NonmonomialOrderTail.notMem_ordGe_of_mem_ordMask,Hironaka.NonmonomialOrderTail.regime_start,Hironaka.NonmonomialOrderTail.proj_tail_mem_ordMask,Hironaka.NonmonomialOrderTail.isGeoTestTower_companion_tail_of_canonical}\upstreamlean{TranscendentalNakamaye.Pending.NonmonomialOrderTail.eventually_ordRegime_and_le_regimeRun}
\mixedleanok
Let \(n\ge1\), \(\underline{\mathcal I}\) a marked ideal in the domain of Definition~\ref{hir:centre-function}, \(\nu\) the order of the nonmonomial part, \(o\ge1\), and
\(f\) a value function whose order dominates \(\nu\) as in Lemma~\ref{bcalc:nonmonomial-order-canonical}. Write \(W_o(\underline{\mathcal J})\) for the open set of the points
outside the closure of \(\{\nu_{\underline{\mathcal J}}\ge o+1\}\); on the domain it is \(\{\nu_{\underline{\mathcal J}}\le o\}\).
\begin{enumerate}
\item Along a multiple test blow-up \(T\) of a companion ideal \(\mathcal O(\underline{\mathcal J},o)\) (\(\underline{\mathcal J}\) in the domain, \(\nu\le o\) on its cosupport),
\(\operatorname{cosupp}\mathcal O(\underline{\mathcal J},o)_T=\{\nu_{\underline{\mathcal J}_T}\ge o\}\cap\operatorname{cosupp}\underline{\mathcal J}_T\). Transforms of marked ideals with the
same divisor have the same divisor.
\item Let \(T_0T_1\) be a multiple test blow-up of \(\underline{\mathcal I}\) and a canonical segment of \(f\), \(W\) an open set of the last stage of \(T_0\) with
\(\nu\le o\) on \(\operatorname{cosupp}\underline{\mathcal I}_{T_0}\cap W\), and suppose some point has run \(\ge\operatorname{length}T_1\) in the regime \(o\)
(Lemma~\ref{bcalc:nonmonomial-order-regime}). Then \(T_1|_W\) is a multiple test blow-up of \(\mathcal O(\underline{\mathcal I}_{T_0}|_W,o)\).
\item Let \(T\) be a multiple test blow-up of \(\underline{\mathcal I}\) and a canonical segment of \(f\), \(y\in\operatorname{cosupp}\underline{\mathcal I}_T\) with
\(\nu(y)=o\), \(r\) its run, and \(T=T_0T_1\) the split with \(\operatorname{length}T_1=r\) (Lemma~\ref{hir:tower-split}). The image of \(y\) at the last stage of
\(T_0\) is in the regime \(o\) and in \(W=W_o(\underline{\mathcal I}_{T_0})\), and \(T_1|_W\) is a multiple test blow-up of \(\mathcal O(\underline{\mathcal I}_{T_0}|_W,o)\).
\item In the situation of (3), every point \(y'\) of \(\operatorname{cosupp}\underline{\mathcal I}_T\) near \(y\) with \(\nu(y')\ge o\) is in the regime \(o\), and its run is at
least \(r\).
\end{enumerate}
\end{lemma}
\begin{proof}
\mixedleanok
(1) The inclusion \(\subset\) is Lemma~\ref{bcalc:nonmonomial-order-local}(4). Conversely, by Lemma~\ref{bcalc:companion-open}(2),
\(\operatorname{cosupp}\mathcal O(\underline{\mathcal J},o)_T=\operatorname{cosupp}(\mathcal N_T|_{N_T},o)\cap\operatorname{cosupp}\underline{\mathcal J}_T\), where \(\mathcal N_T\) is the
transform of \((M,N,E,\mathcal N_N(\mathfrak a),o)\), which lives on \(N_T\) and is contained in \(\mathcal N_{N_T}(\mathfrak a_T)\)
(Theorem~\ref{hir:nonmonomial-transform-along}); a point where every germ of \(\mathcal N_{N_T}(\mathfrak a_T)\) vanishes to order \(o\) along \(N_T\) is therefore in
\(\operatorname{cosupp}(\mathcal N_T|_{N_T},o)\). The divisor of a transform is the list of strict transforms of the members followed by the exceptional member
(Lemma~\ref{hir:geo-transform-basic}); induction along the tower.

(2) Induction along \(T_1=S\sigma\), \(\sigma\) along \(C\). The run of the point \(y\) at the last stage is at least \(\operatorname{length}S+1\ge1\), so its image
\(w\) at the last stage of \(T_0S\) is in the regime \(o\) and has run \(\ge\operatorname{length}S\) (Lemma~\ref{hir:regime-run}); by induction \(S|_W\) is a multiple test
blow-up of \(\mathcal O=\mathcal O(\underline{\mathcal I}_{T_0}|_W,o)\). The centre \(C\) is empty or the maximum locus of \(f\); as \(w\) lies in the cosupport with
\(\nu(w)=o\), \(C\subset\{\nu\ge o\}\) (Lemma~\ref{bcalc:nonmonomial-order-canonical}(2)). The trace of \(C\) over \(W\) is an admissible centre of the transform of
\(\underline{\mathcal I}_{T_0}|_W\) along \(S|_W\), which is \(\underline{\mathcal I}_{T_0S}\) restricted to the preimage of \(W\) (Lemma~\ref{hir:open-restriction-tower},
the tower \(S\) being a multiple test blow-up of \(\underline{\mathcal I}_{T_0}\), Lemma~\ref{hir:geo-transform-basic}). The stages lie in the domain
(Proposition~\ref{bcalc:inv-domain-stable}) and the bound holds on \(W\) (Lemma~\ref{bcalc:nonmonomial-order-local}(2)), so by (1) the cosupport of
\(\mathcal O_{S|_W}\) is \(\{\nu\ge o\}\cap\operatorname{cosupp}\) of that transform, which contains the trace of \(C\) (restriction to open sets does not change \(\nu\) or the
cosupport, Lemmas~\ref{bcalc:nonmonomial-order-local}(2) and~\ref{hir:open-restriction-marked}); the divisors agree by (1), since \(\mathcal O\) has the divisor of
\(\underline{\mathcal I}_{T_0}|_W\) (Lemma~\ref{hir:companion-reduction-along}). So the trace of \(C\) is an admissible centre of \(\mathcal O_{S|_W}\)
(Definition~\ref{hir:test-tower}).

(3) If \(r=0\) the image is \(y\) itself; otherwise the image at the start of the run is in the regime by Lemma~\ref{hir:regime-run}. On the domain the set
\(\{\nu\ge o+1\}\) is closed (Lemma~\ref{bcalc:nonmonomial-order-local}(1)), so a point in the regime \(o\) lies in \(W\) and the bound holds on \(W\). Apply (2)
with the point \(y\), whose run is \(r=\operatorname{length}T_1\).

(4) By (3) and Lemma~\ref{bcalc:nonmonomial-order-regime}(2), every point \(y'\) of the cosupport over \(W\) with \(\nu(y')\ge o\) is in the regime \(o\) with run
\(\ge\operatorname{length}T_1=r\); the points of the last stage over \(W\) form an open neighbourhood of \(y\) by (3).
\end{proof}

\paragraph{The order, its regime and its mask under the operations of the general step.} Properties (4'), (5), (6) and (7) of the general
step's value (Definition~\ref{hir:centre-function}) compare values across an empty prefix, a restriction to an open set, an isomorphism of towers and a
change of representative. The lemma below gives these comparisons for the first entry of the value and for the data it determines: the regime \(o\), its
run, and the mask \(W_o\) of Lemma~\ref{bcalc:nonmonomial-order-tail}.

\begin{lemma}[The order of the nonmonomial part under the operations of the general step]
\label{bcalc:nonmonomial-order-transport}
\uses{bcalc:nonmonomial-order-local,bcalc:nonmonomial-order-regime,bcalc:nonmonomial-order-tail,bcalc:monomial-part-etale,hir:restricted-equivalence,hir:restricted-equivalence-basic,hir:open-restriction-tower,hir:open-restriction-marked,hir:tower-isomorphism,mc:tower-isomorphism-transform,hir:geo-transform-basic,hir:regime-run,bcalc:inv-domain-stable,mc:smooth-centre-invariant,hir:centre-function}
\source{wlodarczyk05-hironaka,bierstone-milman08-functoriality}
\lean{Hironaka.NonmonomialOrderTransport.ordGe_congr,Hironaka.NonmonomialOrderTransport.ordGe_geoTransformAlong_congr,Hironaka.NonmonomialOrderTransport.nmOrd_geoTransformAlong_congr,Hironaka.NonmonomialOrderTransport.ordRegime_congr,Hironaka.NonmonomialOrderTransport.regimeRun_ordRegime_congr,Hironaka.NonmonomialOrderTransport.ordMask_geoTransformAlong_congr,Hironaka.NonmonomialOrderTransport.mem_ordMask_restrict_iff,Hironaka.NonmonomialOrderTransport.mem_ordMask_isTowerIso_iff,Hironaka.NonmonomialOrderTransport.ordRegime_append_iff,Hironaka.NonmonomialOrderTransport.ordRegime_isEmptyTower_iff,Hironaka.NonmonomialOrderTransport.regimeRun_ordRegime_isEmptyTower,Hironaka.NonmonomialOrderTransport.ordMask_isEmptyTower}
\leanok
Let \(n\ge1\), \(\nu\) the order of the nonmonomial part, \(o\in\mathbf N\), \(r_o\) the run in the regime \(o\) (Lemma~\ref{bcalc:nonmonomial-order-regime}) and
\(W_o\) the mask of Lemma~\ref{bcalc:nonmonomial-order-tail}.
\begin{enumerate}
\item (Representatives) If \(\underline{\mathcal I}\simeq\underline{\mathcal J}\) (Definition~\ref{hir:restricted-equivalence}) with \(N\) closed and both ideals coherent, then along
every tower \(T\) the sets \(\{\nu\ge q\}\), the orders, the regimes \(o\), their runs and the masks \(W_o\) of \(\underline{\mathcal I}_T\) and \(\underline{\mathcal J}_T\) agree.
\item (Open sets) Let \(\underline{\mathcal I}\) lie in the domain of Definition~\ref{hir:centre-function}, \(T\) a multiple test blow-up of it and \(U\subset X\) open. A point of
the preimage of \(U\) lies in \(W_o\) of the transform of \(\underline{\mathcal I}|_U\) along \(T|_U\) iff it lies in \(W_o(\underline{\mathcal I}_T)\).
\item (Isomorphisms) Let \(\underline{\mathcal I}\) lie in the domain, \(T\) a multiple test blow-up of it, and \(\chi\colon T'\to T\) an isomorphism of towers over a
biholomorphism \(\Phi\) (Definition~\ref{hir:tower-isomorphism}). Then \(p\in W_o((\Phi^*\underline{\mathcal I})_{T'})\) iff \(\chi(p)\in W_o(\underline{\mathcal I}_T)\).
\item (Empty prefixes) For every tower \(P\), a point is in the regime \(o\) along \(PS\) iff it is in the regime \(o\) of \(\underline{\mathcal I}_P\) along \(S\). If
\(\underline{\mathcal I}\) lies in the domain and \(P\) has only empty centres, a point \(z\) of the last stage of \(P\) is in the regime \(o\) iff its image at the base is,
its run is \(\operatorname{length}P\) if so and \(0\) otherwise, and \(W_o(\underline{\mathcal I}_P)=p_P^{-1}W_o(\underline{\mathcal I})\).
\end{enumerate}
\end{lemma}
\begin{proof}
\leanok
(1) The nonmonomial part only depends on \(\operatorname{Res}_N\mathfrak a\) (Lemma~\ref{bcalc:monomial-part-etale}), and \(N\), \(E\) agree; after one blowup restricted-equivalent
marked ideals have equal transforms, and along every tower equal cosupports (Lemma~\ref{hir:restricted-equivalence-basic}). So the sets \(\{\nu\ge q\}\), which
determine the order, agree at every stage, and so do the cosupports; the regimes agree pointwise, the runs by induction along the tower, and the masks, which
are determined by \(\{\nu\ge o+1\}\).

(2) The transform of \(\underline{\mathcal I}|_U\) along \(T|_U\) is the restriction of \(\underline{\mathcal I}_T\) to the preimage of \(U\) (Lemma~\ref{hir:open-restriction-tower}),
whose set \(\{\nu\ge o+1\}\) is the trace of that of \(\underline{\mathcal I}_T\) (Lemma~\ref{bcalc:nonmonomial-order-local}(2), the stages lying in the domain,
Proposition~\ref{bcalc:inv-domain-stable}). On an open subset the closure of a trace is the trace of the closure, since the inclusion is continuous and open.

(3) By Lemma~\ref{bcalc:nonmonomial-order-local}(3) (with Lemma~\ref{mc:tower-isomorphism-transform}) the set \(\{\nu\ge o+1\}\) of \((\Phi^*\underline{\mathcal I})_{T'}\) is the
preimage under \(\chi\) of that of \(\underline{\mathcal I}_T\); \(\chi\) is a homeomorphism, so it maps closures to closures.

(4) \(\underline{\mathcal I}_{PS}=(\underline{\mathcal I}_P)_S\) (Lemma~\ref{hir:geo-transform-basic}). Along \(P\) with only empty centres, the cosupport at the last stage is the
preimage of the cosupport at the base (Lemma~\ref{mc:smooth-centre-invariant}, a blowup along \(\emptyset\) creating no cosupport) and \(\nu\) is that of the image
(Lemma~\ref{bcalc:nonmonomial-order-local}(2)); so the regime at every stage of \(P\) is that of the image at the base, and the run (Lemma~\ref{hir:regime-run}) is the
whole tower or \(0\). The stages lie in the domain, where \(W_o=\{\nu\le o\}\), and the last assertion follows.
\end{proof}

\paragraph{Admissible centres and finitely many orders in the companion phase.} Property (1) of the general step's value reads the level sets
of the maximal-order invariant of the companion ideal as centres of the marked ideal itself, and property (8) needs that only finitely many orders occur
over a compact set. The lemma below supplies both on the side of the order of the nonmonomial part.

\begin{lemma}[The companion phase: admissible centres and a bound on the order]
\label{bcalc:nonmonomial-order-phase}
\uses{bcalc:nonmonomial-order-tail,bcalc:nonmonomial-order-canonical,bcalc:nonmonomial-order-local,bcalc:companion-open,bcalc:companion-domain,bcalc:inv-domain-stable,hir:open-restriction-marked,hir:test-tower,hir:centre-function}
\source{wlodarczyk05-hironaka,bierstone-milman08-functoriality}
\upstreamlean{TranscendentalNakamaye.Pending.NonmonomialOrderPhase.isAdmissibleCentre_companion_iff}\lean{Hironaka.NonmonomialOrderPhase.isAdmissibleCentre_pullback_val_companion_iff}\upstreamlean{TranscendentalNakamaye.Pending.NonmonomialOrderPhase.exists_nmOrd_le_of_isCompact,TranscendentalNakamaye.Pending.NonmonomialOrderPhase.exists_nmOrd_le_of_isCompact_of_canonical}
\mixedleanok
Let \(n\ge1\) and \(\nu\) the order of the nonmonomial part.
\begin{enumerate}
\item Let \(\underline{\mathcal J}\) lie in the domain of Definition~\ref{hir:centre-function}, \(o\ge1\) with \(\nu\le o\) on \(\operatorname{cosupp}\underline{\mathcal J}\), and \(T\) a
multiple test blow-up of \(\mathcal O(\underline{\mathcal J},o)\). A set \(C\) is an admissible centre of \(\mathcal O(\underline{\mathcal J},o)_T\) iff it is an admissible centre of
\(\underline{\mathcal J}_T\) contained in \(\{\nu_{\underline{\mathcal J}_T}\ge o\}\); the same holds for the restrictions of both transforms to any open subset of the last stage.
\item Let \(\underline{\mathcal I}\) be well formed with \(\mathfrak a\) nonzero along \(N\). On every compact \(K\subset N\), \(\nu\) is bounded.
\item Let \(\underline{\mathcal I}\) lie in the domain, \(f\) a value function whose order dominates \(\nu\) (Lemma~\ref{bcalc:nonmonomial-order-canonical}), and \(K\) a compact
subset of the base. There is \(q\) such that along every multiple test blow-up \(T\) of \(\underline{\mathcal I}\) which is a canonical segment of \(f\), \(\nu\le q\) at every
point of \(\operatorname{cosupp}\underline{\mathcal I}_T\) over \(K\).
\end{enumerate}
\end{lemma}
\begin{proof}
\mixedleanok
(1) The companion ideal has the divisor of \(\underline{\mathcal J}\) (Proposition~\ref{bcalc:companion-domain}), so the transforms have the same divisor
(Lemma~\ref{bcalc:nonmonomial-order-tail}(1)), and \(\operatorname{cosupp}\mathcal O(\underline{\mathcal J},o)_T=\{\nu\ge o\}\cap\operatorname{cosupp}\underline{\mathcal J}_T\)
(Lemma~\ref{bcalc:nonmonomial-order-tail}(1)). An admissible centre is a closed submanifold in the cosupport with normal crossings with the divisor
(Definition~\ref{hir:test-tower}); both conditions on the divisor and on the cosupport therefore match. On an open set, restriction takes cosupports to their traces
(Lemma~\ref{hir:open-restriction-marked}; the transforms lie in the domain, Proposition~\ref{bcalc:inv-domain-stable}, and \(T\) is a multiple test blow-up of
\(\underline{\mathcal J}\), Lemma~\ref{bcalc:companion-open}(2)) and divisors to their traces.

(2) By Lemma~\ref{bcalc:nonmonomial-order-local}(1) every \(y\in K\) has a neighbourhood on which \(\nu\le\nu(y)\); finitely many of them cover \(K\), and the largest of the
finitely many values bounds \(\nu\) on \(K\).

(3) \(N\) is closed, so \(K\cap N\) is compact; let \(q\) bound \(\nu\) on it by (2). A point \(y\) of \(\operatorname{cosupp}\underline{\mathcal I}_T\) over \(K\) lies over
\(\operatorname{cosupp}\underline{\mathcal I}\cap K\subset N\cap K\), and \(\nu(y)\le\nu(p_T(y))\le q\) since the order never goes up along a canonical segment
(Lemma~\ref{bcalc:nonmonomial-order-canonical}(2), with \(T_0\) the empty tower).
\end{proof}

\paragraph{The monomial phase of the general step.} At a point where the nonmonomial part has order \(0\), the value of the general step is read off
the invariant \(\gamma_N\) of the current marked ideal (Lemma~\ref{hir:monomial-invariant-along}). The monomial case along \(N\) (Lemma~\ref{hir:monomial-step-along})
needs a marked ideal which is monomial along all of \(N\); the lemma below obtains its conclusions where only points of order \(0\) are involved.

\begin{lemma}[The monomial phase: admissible level sets, and the invariant off and over the centre]
\label{bcalc:nonmonomial-order-monomial}
\uses{bcalc:nonmonomial-order-zero,bcalc:nonmonomial-order-tail,bcalc:nonmonomial-order-local,hir:monomial-invariant-along,hir:monomial-step-along,hir:monomial-invariance-along,bcalc:centre-local,bcalc:restrict-open,hir:open-restriction-marked,bcalc:inv-domain-stable,bcalc:transform-wellformed,hir:giraud-along,hir:centre-function}
\source{wlodarczyk05-hironaka,bierstone-milman08-functoriality}
\lean{Hironaka.NonmonomialOrderMonomial.isMonomialAlong_pullback_val_ordMask,Hironaka.NonmonomialOrderMonomial.mem_ordMask_zero,Hironaka.NonmonomialOrderMonomial.isAdmissibleCentre_gammaAlong_level_of_ordZero}\upstreamlean{TranscendentalNakamaye.Pending.NonmonomialOrderMonomial.isAdmissibleCentre_maxLocus_gammaAlong_of_ordZero}\lean{Hironaka.NonmonomialOrderMonomial.gammaAlong_geoTransform_pullback_val,Hironaka.NonmonomialOrderMonomial.gammaAlong_geoTransform_eq_of_notMem_of_ordZero,Hironaka.NonmonomialOrderMonomial.preimage_maxLocus_gammaAlong_eq_maxLocusA,Hironaka.NonmonomialOrderMonomial.gammaAlong_geoTransform_lt_of_ordZero}
\mixedleanok
Let \(n\ge1\), \(\underline{\mathcal I}=(M,N,E,\mathfrak a,d)\) lie in the domain of Definition~\ref{hir:centre-function}, \(\nu\) be the order of the nonmonomial part
(Lemma~\ref{bcalc:nonmonomial-order-local}), \(\gamma=\gamma_N\) the invariant of Lemma~\ref{hir:monomial-invariant-along} (with the mark \(d\)), and \(W_0\) the open set of the
points off \(N\) or of order \(\nu=0\) (Lemma~\ref{bcalc:nonmonomial-order-zero}(3); the mask of order \(0\) of Lemma~\ref{bcalc:nonmonomial-order-tail}).
\begin{enumerate}
\item \(\underline{\mathcal I}|_{W_0}\) lies in the domain and is monomial along \(N\cap W_0\); a point of order \(0\) lies in \(W_0\).
\item (level sets) Let \(W\subset M\) be open and \(v\) a value of \(\gamma\) or \(\bot\), such that \(\nu=0\) and \(\gamma\le v\) at every point of
\(\operatorname{cosupp}\underline{\mathcal I}\cap W\). Then \(\{y\in\operatorname{cosupp}\underline{\mathcal I}\cap W:\gamma(y)=v\}\) is an admissible centre of \(\underline{\mathcal I}|_W\).
In particular, if \(\nu=0\) on \(\operatorname{cosupp}\underline{\mathcal I}\), the maximum locus of \(\gamma\) on \(\operatorname{cosupp}\underline{\mathcal I}\) is an admissible centre.
\item (off the centre) Let \(C\) be an admissible centre (possibly empty), \(\sigma\colon M'\to M\) the blowup along \(C\), \(\underline{\mathcal I}'\) the transform and
\(\gamma'\) its invariant along \(N'\). If \(z\in N'\), \(\sigma(z)\notin C\) and \(\nu(\sigma(z))=0\), then \(\gamma'(z)=\gamma(\sigma(z))\).
\item (over the centre) If \(\nu=0\) on \(\operatorname{cosupp}\underline{\mathcal I}\) and \(C\) is the maximum locus of \(\gamma\) on \(\operatorname{cosupp}\underline{\mathcal I}\), then
\(\gamma'(z)<\gamma(\sigma(z))\) at every \(z\in N'\) with \(\sigma(z)\in C\).
\end{enumerate}
\end{lemma}
\begin{proof}
\mixedleanok
(1) The restriction to an open set stays in the domain (Proposition~\ref{bcalc:inv-domain-stable}); it is monomial by Lemma~\ref{bcalc:nonmonomial-order-zero}(2), since
\(W_0\) misses \(\{\nu\ge1\}\). On the domain \(\{\nu\ge1\}\) is closed (Lemma~\ref{bcalc:nonmonomial-order-local}), so \(W_0\) is exactly the set of points off
\(\{\nu\ge1\}\) (Lemma~\ref{bcalc:nonmonomial-order-tail}).

Throughout, restriction to an open set \(U\) is read back as follows: \(\operatorname{cosupp}(\underline{\mathcal I}|_U)=\operatorname{cosupp}\underline{\mathcal I}\cap U\)
(Lemma~\ref{hir:open-restriction-marked}); \(\gamma\) of \(\underline{\mathcal I}|_U\) at \(p\) is \(\gamma(p)\) (Lemma~\ref{hir:monomial-invariance-along}(2)); the blowup restricts
to the blowup of \(U\) along \(C\cap U\) on \(\sigma^{-1}U\) (Lemma~\ref{bcalc:restrict-open}), and the transform of \(\underline{\mathcal I}|_U\) along it is
\(\underline{\mathcal I}'|_{\sigma^{-1}U}\) (Lemma~\ref{hir:open-restriction-marked}; the transform along an admissible centre is well formed,
Corollary~\ref{bcalc:transform-wellformed}). Hence \(\gamma'\) of the transform of \(\underline{\mathcal I}|_U\) at \(p\) is \(\gamma'(p)\).

(2) By Lemma~\ref{bcalc:centre-local}(2) it suffices that every \(x\in W\) has an open neighbourhood \(V\subset W\) on which the trace of the set is an admissible centre of
\(\underline{\mathcal I}|_V\). If \(\nu(x)\ge1\), \(x\) is not in the closed set \(\operatorname{cosupp}\underline{\mathcal I}\) (Lemma~\ref{hir:giraud-along}); on
\(V=W\setminus\operatorname{cosupp}\underline{\mathcal I}\) the trace is empty. If \(\nu(x)=0\), take \(V=W\cap W_0\): \(\underline{\mathcal I}|_V\) is monomial along \(N\cap V\)
by (1), so \(\gamma=\bot\) exactly at the points of \(N\cap V\) off the cosupport (Lemma~\ref{hir:monomial-invariant-along}(2)). For \(v=\bot\) the trace is empty. For
\(v\ne\bot\), \(\gamma\le v\) on all of \(N\cap V\), and the trace is \(\{y\in N\cap V:\gamma(y)=v\}\), an admissible centre by Lemma~\ref{hir:monomial-invariant-along}(3).
For the maximum locus: it is empty, or the level set of the value \(v\) of one of its points, on which \(\gamma\le v\) holds on the whole cosupport; take \(W=M\).

(3) Take \(U=W_0\ni\sigma(z)\). By (1) and the reading back above, Lemma~\ref{hir:monomial-step-along}(4) for \(\underline{\mathcal I}|_{W_0}\) and the restricted blowup gives
\(\gamma'(z)=\gamma(\sigma(z))\).

(4) Again \(U=W_0\), which now contains \(\operatorname{cosupp}\underline{\mathcal I}\). The trace of \(C\) on \(W_0\) is the maximum locus of \(\gamma\) of \(\underline{\mathcal I}|_{W_0}\) on
its cosupport: both are read off the same values at the same points, because every point of the cosupport lies in \(W_0\). By (2), \(C\) is admissible, and
Lemma~\ref{hir:monomial-step-along}(3) for \(\underline{\mathcal I}|_{W_0}\) gives \(\gamma'(z)<\gamma(\sigma(z))\).
\end{proof}

\paragraph{The monomial phase along canonical segments.} The next lemma applies Lemma~\ref{bcalc:nonmonomial-order-monomial} along canonical segments of
any value function whose first entry is the order of the nonmonomial part and which, at the points of order \(0\), is a strictly monotone function of \(\gamma\);
the value of the general step is one (it reads \((0,\gamma)\) at the points of order \(0\)).

\begin{lemma}[The monomial phase along canonical segments]
\label{bcalc:nonmonomial-order-monomial-canonical}
\uses{bcalc:nonmonomial-order-monomial,bcalc:nonmonomial-order-canonical,bcalc:nonmonomial-order-local,hir:monomial-invariant-along,bcalc:max-locus-local,hir:canonical-segments,bcalc:inv-domain-stable,hir:geo-transform-basic,hir:centre-function,hir:invariant-order}
\source{wlodarczyk05-hironaka,bierstone-milman08-functoriality}
\lean{Hironaka.NonmonomialOrderMonomialCanonical.OrdZeroValue,Hironaka.NonmonomialOrderMonomialCanonical.OrdZeroValue.apply,Hironaka.NonmonomialOrderMonomialCanonical.strictMono_toLex_zero_inl,Hironaka.NonmonomialOrderMonomialCanonical.ordZeroValue_toLex,Hironaka.NonmonomialOrderMonomialCanonical.nmOrd_eq_zero_of_mem_maxLocus,Hironaka.NonmonomialOrderMonomialCanonical.notMem_ordGe_one_of_nmOrd_eq_zero,Hironaka.NonmonomialOrderMonomialCanonical.maxLocus_eq_maxLocus_towerInvAlong,Hironaka.NonmonomialOrderMonomialCanonical.isAdmissibleCentre_level_of_nmOrd_eq_zero,Hironaka.NonmonomialOrderMonomialCanonical.eventually_le_of_nmOrd_eq_zero,Hironaka.NonmonomialOrderMonomialCanonical.towerInvAlong_step_of_nmOrd_eq_zero,Hironaka.NonmonomialOrderMonomialCanonical.value_step_of_nmOrd_eq_zero}\upstreamlean{TranscendentalNakamaye.Pending.NonmonomialOrderMonomialCanonical.towerInvAlong_le_of_canonical,TranscendentalNakamaye.Pending.NonmonomialOrderMonomialCanonical.value_le_of_canonical}
\mixedleanok
Let \(n\ge1\), \(\underline{\mathcal I}\) lie in the domain of Definition~\ref{hir:centre-function}, \(\nu\) be the order of the nonmonomial part and, for a tower \(T\),
\(\gamma_T\) the invariant of Lemma~\ref{hir:monomial-invariant-along} of \(\underline{\mathcal I}_T\) (with its mark). Let \(f\) be a value function whose order dominates
\(\nu\) (Lemma~\ref{bcalc:nonmonomial-order-canonical}), and \(\psi\) a strictly monotone map with \(f(T)(y)=\psi(\gamma_T(y))\) at every point \(y\) of
\(\operatorname{cosupp}\underline{\mathcal I}_T\) with \(\nu(y)=0\), for every multiple test blow-up \(T\) of \(\underline{\mathcal I}\). (For \(f=(\nu,g)\) ordered lexicographically
with \(g=\gamma\) in the first summand of a lexicographic sum at the points of order \(0\), \(\psi(w)=(0,w)\).) Let \(T\) be a multiple test blow-up of
\(\underline{\mathcal I}\).
\begin{enumerate}
\item If a point of order \(0\) lies on the maximum locus of \(f(T)\) on \(\operatorname{cosupp}\underline{\mathcal I}_T\), every point of the cosupport has order \(0\), and
that maximum locus is the maximum locus of \(\gamma_T\) on the cosupport.
\item (level sets) Let \(W\) be an open subset of the last stage and \(l\) a value with \(f(T)\le l\) on \(\operatorname{cosupp}\underline{\mathcal I}_T\cap W\), attained at a point
of the cosupport of order \(0\). Then \(\{y\in\operatorname{cosupp}\underline{\mathcal I}_T\cap W:f(T)(y)=l\}\) is an admissible centre of \(\underline{\mathcal I}_T|_W\).
\item (upper semicontinuity) Near a point \(y\) of \(\operatorname{cosupp}\underline{\mathcal I}_T\) with \(\nu(y)=0\), every point \(y'\) of the cosupport with \(\nu(y')=0\)
has \(f(T)(y')\le f(T)(y)\).
\item (one step) Let \(T\sigma\), \(\sigma\) the blowup along \(C\), be a multiple test blow-up of \(\underline{\mathcal I}\) and a canonical segment of \(f\), and \(z\) a point of
\(\operatorname{cosupp}\underline{\mathcal I}_{T\sigma}\) with \(\nu(\sigma(z))=0\). Then \(\sigma(z)\in\operatorname{cosupp}\underline{\mathcal I}_T\), \(\nu(z)=0\), and
\(\gamma_{T\sigma}(z)<\gamma_T(\sigma(z))\), \(f(T\sigma)(z)<f(T)(\sigma(z))\) if \(\sigma(z)\in C\), while \(\gamma_{T\sigma}(z)=\gamma_T(\sigma(z))\),
\(f(T\sigma)(z)=f(T)(\sigma(z))\) if \(\sigma(z)\notin C\).
\item (segments) Let \(T_0M\) be a multiple test blow-up of \(\underline{\mathcal I}\) and a canonical segment of \(f\), and \(z\in\operatorname{cosupp}\underline{\mathcal I}_{T_0M}\)
with \(\nu(p_M(z))=0\). Then \(p_M(z)\in\operatorname{cosupp}\underline{\mathcal I}_{T_0}\), \(\nu(z)=0\), \(\gamma_{T_0M}(z)\le\gamma_{T_0}(p_M(z))\) and
\(f(T_0M)(z)\le f(T_0)(p_M(z))\).
\end{enumerate}
\end{lemma}
\begin{proof}
\mixedleanok
The transforms along multiple test blow-ups stay in the domain (Proposition~\ref{bcalc:inv-domain-stable}), so Lemma~\ref{bcalc:nonmonomial-order-monomial} applies at every
stage; on the domain a point of \(N\) has \(\nu=0\) iff it is off \(\{\nu\ge1\}\) (Lemma~\ref{bcalc:nonmonomial-order-local}).

(1) A point \(y\) of the cosupport has \(f(T)(y)\le f(T)(c)\) for \(c\) on the maximum locus, so \(\nu(y)\le\nu(c)=0\). Hence \(f(T)=\psi\circ\gamma_T\) on the cosupport, and
a strictly monotone change of values does not change maximum loci (Lemma~\ref{bcalc:max-locus-local}; maximum loci only depend on the values on the set,
Lemma~\ref{hir:canonical-segments}).

(2) Let \(y_0\) be the point of order \(0\) with \(f(T)(y_0)=l\). Every \(y\in\operatorname{cosupp}\underline{\mathcal I}_T\cap W\) has \(f(T)(y)\le f(T)(y_0)\), hence \(\nu(y)=0\),
\(f(T)(y)=\psi(\gamma_T(y))\) and \(\gamma_T(y)\le\gamma_T(y_0)\); and \(f(T)(y)=l\) iff \(\gamma_T(y)=\gamma_T(y_0)\), \(\psi\) being injective. The set is therefore the level set
of \(\gamma_T\) at \(\gamma_T(y_0)\) of Lemma~\ref{bcalc:nonmonomial-order-monomial}(2).

(3) \(\gamma_T\) is upper semicontinuous on \(N_T\) (Lemma~\ref{hir:monomial-invariant-along}(1), which only uses that the ideal is nonzero along \(N_T\)), and
\(f(T)=\psi\circ\gamma_T\) at the points of the cosupport of order \(0\), \(\psi\) being monotone.

(4) By Lemma~\ref{bcalc:nonmonomial-order-canonical}(2), \(\sigma(z)\) lies in the old cosupport and \(\nu(z)\le\nu(\sigma(z))=0\). The centre \(C\) is admissible for
\(\underline{\mathcal I}_T\). If \(\sigma(z)\in C\), then \(C\) is the maximum locus of \(f(T)\) (it is not empty), and it contains \(\sigma(z)\), of order \(0\); by (1) the cosupport
has order \(0\) and \(C\) is the maximum locus of \(\gamma_T\), so Lemma~\ref{bcalc:nonmonomial-order-monomial}(4) gives \(\gamma_{T\sigma}(z)<\gamma_T(\sigma(z))\). If
\(\sigma(z)\notin C\), Lemma~\ref{bcalc:nonmonomial-order-monomial}(3) gives equality. Both points have order \(0\), so the values are \(\psi\) of the invariants.

(5) Induction on \(M\). For \(M=M'\sigma\), the image \(\sigma(z)\) lies in the cosupport at the last stage of \(T_0M'\), with
\(\nu(\sigma(z))\le\nu(p_{M'}(\sigma(z)))=\nu(p_M(z))=0\) (Lemma~\ref{bcalc:nonmonomial-order-canonical}(2)); (4) gives
\(\gamma_{T_0M}(z)\le\gamma_{T_0M'}(\sigma(z))\) and the induction hypothesis \(\gamma_{T_0M'}(\sigma(z))\le\gamma_{T_0}(p_M(z))\). The values follow as in (4), \(T_0\) being a
multiple test blow-up (Lemma~\ref{hir:geo-transform-basic}).
\end{proof}

\paragraph{The invariant of the monomial phase under the operations of the general step.} Properties (4)--(7) of the general step's value need, at the
points of order \(0\), that the invariant \(\gamma\) of the current transform behaves under the operations of the general step as the order \(\nu\) does
(Lemma~\ref{bcalc:nonmonomial-order-transport}); restriction to open sets is Lemma~\ref{hir:monomial-invariance-along}.

\begin{lemma}[The invariant of the monomial phase under re-basing, isomorphisms of towers and representatives]
\label{bcalc:nonmonomial-order-monomial-transport}
\uses{bcalc:nonmonomial-order-monomial,bcalc:nonmonomial-order-monomial-canonical,hir:monomial-invariant-along,hir:monomial-invariance-along,hir:exponent-along,hir:restricted-calculus,hir:restricted-equivalence-basic,hir:marked-equivalence-basic,hir:tower-isomorphism,mc:tower-isomorphism-transform,hir:geo-transform-basic,hir:open-restriction-tower,bcalc:inv-domain-stable}
\source{wlodarczyk05-hironaka,bierstone-milman08-functoriality}
\lean{Hironaka.NonmonomialOrderMonomialTransport.gammaAlong_congr_of_cosupport,Hironaka.NonmonomialOrderMonomialTransport.gammaAlong_congr_of_restrictEquiv,Hironaka.NonmonomialOrderMonomialTransport.towerInvAlong_congr,Hironaka.NonmonomialOrderMonomialTransport.towerInvAlong_isTowerIso,Hironaka.NonmonomialOrderMonomialTransport.towerInvAlong_append}\upstreamlean{TranscendentalNakamaye.Pending.NonmonomialOrderMonomialTransport.towerInvAlong_step_empty_of_nmOrd_eq_zero}
\mixedleanok
For a marked ideal \(\underline{\mathcal I}\) and a tower \(T\) over its base let \(\gamma_T\) be the invariant of Lemma~\ref{hir:monomial-invariant-along} of
\(\underline{\mathcal I}_T\) (along \(N_T\), with the members of \(E_T\) and the mark).
\begin{enumerate}
\item \(\gamma_{N,\mathfrak a,E,d}\) only depends on \(N\), \(E\), \(d\) and the cosupports \(\operatorname{cosupp}(\mathfrak a|_N,e)\), \(e\in\mathbb N\). In particular restricted-equivalent
marked ideals (Lemma~\ref{hir:restricted-equivalence-basic}) have the same invariant.
\item (representatives) If \(\underline{\mathcal I}\) and \(\underline{\mathcal I}'\) are restricted-equivalent, with \(N\) closed and coherent ideals, then
\(\gamma_T(\underline{\mathcal I})=\gamma_T(\underline{\mathcal I}')\) along every tower \(T\).
\item (isomorphisms of towers) Let \(\chi\colon T'\to T\) be an isomorphism of towers over a biholomorphism \(\Phi\) (Definition~\ref{hir:tower-isomorphism}),
\(T\) a multiple test blow-up of \(\underline{\mathcal I}\), \(N\) a closed submanifold and the ideal coherent. Then \(\gamma_{T'}(\Phi^*\underline{\mathcal I})(p)=\gamma_T(\underline{\mathcal I})(\chi(p))\).
\item (re-basing) \(\gamma_{PT_1}(\underline{\mathcal I})=\gamma_{T_1}(\underline{\mathcal I}_P)\) for all towers \(P\), \(T_1\).
\item (a trailing empty step at order \(0\)) Let \(n\ge1\), \(\underline{\mathcal I}\) lie in the domain of Definition~\ref{hir:centre-function}, \(T\) be a multiple test blow-up and
\(\pi\) the blowup along \(\emptyset\). At every point \(z\) of the strict transform \(N_{T\pi}\) with \(\nu(\pi(z))=0\), \(\gamma_{T\pi}(z)=\gamma_T(\pi(z))\).
\end{enumerate}
\end{lemma}
\begin{proof}
\mixedleanok
(1) The exponent \(e^N_H(\mathfrak a)(x)\) is defined by the condition that \(N\cap H\) lies in \(\operatorname{cosupp}(\mathfrak a|_N,e)\) near \(x\)
(Lemma~\ref{hir:exponent-along}), and \(\gamma\) is read off the members through the point and these exponents. A restricted inclusion
\(\operatorname{Res}_N\mathfrak a\subset\operatorname{Res}_N\mathfrak b\) gives \(\operatorname{cosupp}(\mathfrak b|_N,e)\subset\operatorname{cosupp}(\mathfrak a|_N,e)\) for every \(e\)
(Lemma~\ref{hir:restricted-calculus}); restricted-equivalent marked ideals have both inclusions and the same \(N\), \(E\), \(d\).

(2) For the empty tower this is (1). After one blowup the transforms of restricted-equivalent marked ideals are equal (Lemma~\ref{hir:marked-equivalence-basic}),
hence so are all later transforms.

(3) The transform of \(\Phi^*\underline{\mathcal I}\) along \(T'\) is \(\chi^*\underline{\mathcal I}_T\) (Lemma~\ref{mc:tower-isomorphism-transform}; the transforms keep \(N\) a closed submanifold
and the ideal coherent, Lemma~\ref{hir:open-restriction-tower}), and \(\gamma\) is invariant under biholomorphisms (Lemma~\ref{hir:monomial-invariance-along}(2)).

(4) \(\underline{\mathcal I}_{PT_1}=(\underline{\mathcal I}_P)_{T_1}\) (Lemma~\ref{hir:geo-transform-basic}).

(5) \(\underline{\mathcal I}_T\) lies in the domain (Proposition~\ref{bcalc:inv-domain-stable}), \(\emptyset\) is an admissible centre, and a point of order \(0\) is off \(\{\nu\ge1\}\)
(Lemma~\ref{bcalc:nonmonomial-order-monomial-canonical}); Lemma~\ref{bcalc:nonmonomial-order-monomial}(3) with \(C=\emptyset\) applies.
\end{proof}

\paragraph{The general value at the points of order zero.} At a point where the nonmonomial part has order \(0\) the general value
(Definition~\ref{hir:general-value}) is \((0,\gamma)\), with \(\gamma\) the invariant of the current transform. Its properties (4'), (5), (6) and (7) of
Definition~\ref{hir:centre-function} at such points therefore follow from those of the order \(\nu\) (Lemmas~\ref{bcalc:nonmonomial-order-local}
and~\ref{bcalc:nonmonomial-order-transport}) and of \(\gamma\) (Lemmas~\ref{bcalc:nonmonomial-order-monomial-transport} and~\ref{hir:monomial-invariance-along}); the
value function \(M\) of the companion phase plays no role there.

\begin{lemma}[The general value at the points of order zero]
\label{bcalc:general-value-order-zero}
\uses{hir:general-value,hir:general-value-basic,hir:general-admissible,hir:max-order-value,hir:centre-function,bcalc:nonmonomial-order-local,bcalc:nonmonomial-order-transport,bcalc:nonmonomial-order-monomial-transport,hir:monomial-invariance-along,hir:open-restriction-tower,hir:tower-isomorphism,mc:tower-isomorphism-transform,hir:restricted-equivalence-basic,hir:geo-transform-basic,bcalc:inv-domain-stable}
\source{wlodarczyk05-hironaka,bierstone-milman08-functoriality}
\lean{Hironaka.GeneralValueOrderZero.genInv_eq_towerInvAlong,Hironaka.GeneralValueOrderZero.genInv_eq_of_nmOrd_eq_zero,Hironaka.GeneralValueOrderZero.genInv_append_of_nmOrd_eq_zero,Hironaka.GeneralValueOrderZero.emptyPrefix_field_of_nmOrd_eq_zero,Hironaka.GeneralValueOrderZero.genInv_restrict_of_nmOrd_eq_zero,Hironaka.GeneralValueOrderZero.genInv_isTowerIso_of_nmOrd_eq_zero,Hironaka.GeneralValueOrderZero.genInv_congr_of_nmOrd_eq_zero,Hironaka.GeneralValueOrderZero.restrictOpen_field_of_nmOrd_eq_zero,Hironaka.GeneralValueOrderZero.towerIso_field_of_nmOrd_eq_zero}\upstreamlean{TranscendentalNakamaye.Pending.GeneralValueOrderZero.representative_field_of_nmOrd_eq_zero}
\mixedleanok
Let \(M\) be a value function (Definition~\ref{hir:max-order-value}) and \(\operatorname{inv}\) the general value of Definition~\ref{hir:general-value}. For a tower \(T\) let
\(\nu_T\) be the order of the nonmonomial part of \(\underline{\mathcal I}_T\) and \(\gamma_T\) the invariant of \(\underline{\mathcal I}_T\) of
Lemma~\ref{hir:monomial-invariant-along}.
\begin{enumerate}
\item If \(\nu_T(y)=0\), then \(\operatorname{inv}_{\underline{\mathcal I}}(T)(y)=(0,\gamma_T(y))\); two points of order \(0\), for any two marked ideals and towers, with the same
\(\gamma\) have the same value.
\item (Re-basing, every prefix) For all towers \(P\), \(T_1\) and every point \(z\) of the last stage of \(PT_1\) with \(\nu_{PT_1}(z)=0\),
\(\operatorname{inv}_{\underline{\mathcal I}}(PT_1)(z)=\operatorname{inv}_{\underline{\mathcal I}_P}(T_1)(z)\). No emptiness of \(P\) is needed.
\item (Open sets) Let \(n\ge1\), \(\underline{\mathcal I}\) lie in the domain of Definition~\ref{hir:centre-function}, \(U\subset X\) be open and \(T\) a multiple test blow-up. At
every point \(p\) over \(U\) with \(\nu_T(p)=0\), \(\operatorname{inv}_{\underline{\mathcal I}|_U}(T|_U)(p)=\operatorname{inv}_{\underline{\mathcal I}}(T)(p)\).
\item (Isomorphisms of towers) Under the same hypotheses, let \(\chi\colon T'\to T\) be an isomorphism of towers over a biholomorphism \(\Phi\)
(Definition~\ref{hir:tower-isomorphism}). At every point \(p\) with \(\nu_T(\chi(p))=0\),
\(\operatorname{inv}_{\Phi^*\underline{\mathcal I}}(T')(p)=\operatorname{inv}_{\underline{\mathcal I}}(T)(\chi(p))\).
\item (Representatives) If \(\underline{\mathcal I}\simeq\underline{\mathcal J}\) with \(N\) closed and coherent ideals, then along every tower \(T\), at every point \(y\) with
\(\nu_T(y)=0\) for \(\underline{\mathcal I}\), \(\operatorname{inv}_{\underline{\mathcal I}}(T)(y)=\operatorname{inv}_{\underline{\mathcal J}}(T)(y)\).
\item In particular, if \(T\) is a canonical segment of \(\operatorname{inv}_{\underline{\mathcal I}}\) (\(M\) a value function for the marked ideals of maximal order, \(n\ge1\),
\(\underline{\mathcal I}\) in the domain), properties (4'), (5), (6) and (7) of Definition~\ref{hir:centre-function} hold for \(\operatorname{inv}\) at the points of the cosupport of
order \(0\).
\end{enumerate}
\end{lemma}
\begin{proof}
\mixedleanok
(1) is the definition of the general value at order \(0\). In (2)--(5) the orders on both sides agree: in (2) because
\(\underline{\mathcal I}_{PT_1}=(\underline{\mathcal I}_P)_{T_1}\) (Lemma~\ref{hir:geo-transform-basic}); in (3) because the transform of \(\underline{\mathcal I}|_U\) along \(T|_U\) is the
restriction of \(\underline{\mathcal I}_T\) (Lemma~\ref{hir:open-restriction-tower}) and \(\nu\) is local (Lemma~\ref{bcalc:nonmonomial-order-local}(2)), the transforms staying
in the domain (Proposition~\ref{bcalc:inv-domain-stable}); in (4) by Lemma~\ref{bcalc:nonmonomial-order-local}(3); in (5) by Lemma~\ref{bcalc:nonmonomial-order-transport}(1).
So both points have order \(0\) and by (1) it suffices to compare \(\gamma\): in (2) by Lemma~\ref{bcalc:nonmonomial-order-monomial-transport}(4), in (3) by
Lemma~\ref{hir:monomial-invariance-along} applied to the restriction of \(\underline{\mathcal I}_T\), in (4) by Lemma~\ref{bcalc:nonmonomial-order-monomial-transport}(3) and
in (5) by Lemma~\ref{bcalc:nonmonomial-order-monomial-transport}(2). For (6), canonical segments of the general value are multiple test blow-ups
(Theorem~\ref{hir:general-admissible}), so (3) and (4) apply; (2) and (5) need no hypothesis on \(T\).
\end{proof}

\paragraph{An empty step in the middle of a sequence.} Włodarczyk's extensions of a resolution insert isomorphisms anywhere in a
sequence of blowups. Here a blowup along \(\emptyset\) appends the empty set to the divisor \(E\) (Lemma~\ref{hir:geo-transform-basic}), so an empty step in the middle of a
sequence moves every later member one position further, while the invariant \(\gamma\) of the monomial case records the members through a point by their
positions in \(E\) (Lemma~\ref{hir:monomial-invariant-geometry}). The lemma below computes the effect.

\begin{lemma}[The monomial invariant under an empty member inserted in the divisor]
\label{bcalc:monomial-invariant-empty-member}
\uses{hir:monomial-invariant,hir:monomial-invariant-geometry,hir:monomial-invariant-along,hir:monomial-resolution-along,hir:monomial-part-along,bcalc:nonmonomial-order-local,hir:geo-transform-basic,hir:restricted-equivalence-basic,mc:empty-blowup-inverse,bcalc:nonmonomial-order-monomial-transport,bcalc:max-locus-local,hir:general-value,bcalc:general-value-order-zero,hir:invariant-order}
\source{wlodarczyk05-hironaka,bierstone-milman08-functoriality}
\lean{Hironaka.MonomialInvariantEmptyMember.shiftIdx,Hironaka.MonomialInvariantEmptyMember.strictMono_shiftIdx,Hironaka.MonomialInvariantEmptyMember.shiftIdx_ne,Hironaka.MonomialInvariantEmptyMember.shiftIdx_of_lt,Hironaka.MonomialInvariantEmptyMember.shiftIdx_of_le,Hironaka.MonomialInvariantEmptyMember.shiftSet,Hironaka.MonomialInvariantEmptyMember.card_shiftSet,Hironaka.MonomialInvariantEmptyMember.sum_shiftSet,Hironaka.MonomialInvariantEmptyMember.shiftKey,Hironaka.MonomialInvariantEmptyMember.shiftKey_lt_shiftKey,Hironaka.MonomialInvariantEmptyMember.strictMono_shiftKey,Hironaka.MonomialInvariantEmptyMember.shiftVal,Hironaka.MonomialInvariantEmptyMember.strictMono_shiftVal,Hironaka.MonomialInvariantEmptyMember.shiftVal_bot,Hironaka.MonomialInvariantEmptyMember.shiftVal_coe,Hironaka.MonomialInvariantEmptyMember.shiftSet_ne_self,Hironaka.MonomialInvariantEmptyMember.shiftVal_ne_self,Hironaka.MonomialInvariantEmptyMember.maxLocus_shiftVal_comp,Hironaka.MonomialInvariantEmptyMember.key_shiftSet,Hironaka.MonomialInvariantEmptyMember.feasible_shiftSet,Hironaka.MonomialInvariantEmptyMember.max_image_of_monotone,Hironaka.MonomialInvariantEmptyMember.inv_shiftSet,Hironaka.MonomialInvariantEmptyMember.memberIdx_insertIdx_empty,Hironaka.MonomialInvariantEmptyMember.expDataAlong_insertIdx_empty,Hironaka.MonomialInvariantEmptyMember.gammaAlong_insertIdx_empty,Hironaka.MonomialInvariantEmptyMember.forall_get_insertIdx_empty_iff,Hironaka.MonomialInvariantEmptyMember.monomialPartAlong_insertIdx_empty,Hironaka.MonomialInvariantEmptyMember.insertEmpty,Hironaka.MonomialInvariantEmptyMember.insertIdx_append_of_le_length,Hironaka.MonomialInvariantEmptyMember.cosupp_insertEmpty,Hironaka.MonomialInvariantEmptyMember.nmPart_insertEmpty,Hironaka.MonomialInvariantEmptyMember.ordGe_insertEmpty,Hironaka.MonomialInvariantEmptyMember.nmOrd_insertEmpty,Hironaka.MonomialInvariantEmptyMember.gammaAlong_insertEmpty,Hironaka.MonomialInvariantEmptyMember.geoTransform_insertEmpty,Hironaka.MonomialInvariantEmptyMember.geoTransformAlong_insertEmpty,Hironaka.MonomialInvariantEmptyMember.le_length_E_geoTransformAlong,Hironaka.MonomialInvariantEmptyMember.towerInvAlong_insertEmpty,Hironaka.MonomialInvariantEmptyMember.nmOrd_geoTransformAlong_insertEmpty,Hironaka.MonomialInvariantEmptyMember.cosupp_geoTransformAlong_insertEmpty,Hironaka.MonomialInvariantEmptyMember.maxLocus_towerInvAlong_insertEmpty,Hironaka.MonomialInvariantEmptyMember.emptyStepBase,Hironaka.MonomialInvariantEmptyMember.length_E_emptyStepBase,Hironaka.MonomialInvariantEmptyMember.geoTransform_empty_eq_insertEmpty,Hironaka.MonomialInvariantEmptyMember.restrictEquiv_emptyStepBase_pullback,Hironaka.MonomialInvariantEmptyMember.geoTransformAlong_step_empty_append,Hironaka.MonomialInvariantEmptyMember.le_length_E_step_empty_append,Hironaka.MonomialInvariantEmptyMember.towerInvAlong_step_empty_append,Hironaka.MonomialInvariantEmptyMember.nmOrd_step_empty_append,Hironaka.MonomialInvariantEmptyMember.cosupp_step_empty_append,Hironaka.MonomialInvariantEmptyMember.maxLocus_towerInvAlong_step_empty_append,Hironaka.MonomialInvariantEmptyMember.towerInvAlong_step_empty_append_pullback,Hironaka.MonomialInvariantEmptyMember.genInv_insertEmpty_of_nmOrd_eq_zero,Hironaka.MonomialInvariantEmptyMember.genInv_step_empty_append_of_nmOrd_eq_zero,Hironaka.MonomialInvariantEmptyMember.genInv_step_empty_append_ne}
\leanok
For \(j\in\mathbf N\) let \(s_j(i)=i\) if \(i<j\) and \(s_j(i)=i+1\) if \(i\ge j\), and let \(\sigma_j\) act on the values \((n-|K|,\sum_Ka,K)\) of \(\gamma\) by
\(K\mapsto s_j(K)\) on the third entry (and fix \(\bot\)). For a marked ideal \(\underline{\mathcal I}=(M,N,E,\mathfrak a,d)\) and \(j\le\#E\) let
\(\underline{\mathcal I}^{(j)}\) be \(\underline{\mathcal I}\) with \(\emptyset\) inserted into \(E\) at position \(j\).
\begin{enumerate}
\item \(\sigma_j\) is strictly increasing; it moves every value whose set \(K\) has an element \(\ge j\). In particular maximum loci of \(\sigma_j\circ f\) and of \(f\) agree.
\item \(\gamma(\underline{\mathcal I}^{(j)})=\sigma_j\circ\gamma(\underline{\mathcal I})\); the monomial and nonmonomial parts, the order \(\nu\) of the nonmonomial part and the cosupport of
\(\underline{\mathcal I}^{(j)}\) are those of \(\underline{\mathcal I}\).
\item Inserting \(\emptyset\) commutes with the transforms: \((\underline{\mathcal I}^{(j)})_T=(\underline{\mathcal I}_T)^{(j)}\) along every tower \(T\). Hence
\(\gamma_T(\underline{\mathcal I}^{(j)})=\sigma_j\circ\gamma_T(\underline{\mathcal I})\), and \(\nu\), the cosupports and the maximum loci of \(\gamma\) along \(T\) are unchanged.
\item (An empty step in the middle) Let \(T_0\) be a tower over the base of \(\underline{\mathcal I}\), \(\sigma\) a blowup along \(\emptyset\) of its last stage, \(S\) a tower
over the source of \(\sigma\), and \(\underline{\mathcal I}'\) the transform of \(\underline{\mathcal I}_{T_0}\) by \(\sigma\) with its last (empty) member removed. Then
\(\underline{\mathcal I}_{T_0\sigma}=(\underline{\mathcal I}')^{(m)}\) with \(m=\#E+\operatorname{length}T_0\), and along \(T_0\sigma S\):
\(\gamma=\sigma_m\circ\gamma_S(\underline{\mathcal I}')\), \(\nu=\nu_S(\underline{\mathcal I}')\), the cosupport is that of \(\underline{\mathcal I}'_S\) and the maximum loci of \(\gamma\) agree. If \(N\)
and the members of \(\underline{\mathcal I}_{T_0}\) are closed, \(\underline{\mathcal I}'\) is restricted-equivalent to \(\sigma^*\underline{\mathcal I}_{T_0}\), and (for coherent ideals)
\(\gamma_S(\underline{\mathcal I}')=\gamma_S(\sigma^*\underline{\mathcal I}_{T_0})\).
\item (The general value) At a point of order \(0\), the general value of Definition~\ref{hir:general-value} of \(\underline{\mathcal I}^{(j)}\) along \(T\) is
\((0,\sigma_j(\gamma_T))\), and after an empty step in the middle as in (4) it is \((0,\sigma_m(\gamma_S(\underline{\mathcal I}')))\) while that of \(\underline{\mathcal I}'\) along \(S\) is
\((0,\gamma_S(\underline{\mathcal I}'))\). These two values differ whenever the set of \(\gamma_S(\underline{\mathcal I}')\) at the point contains a position \(\ge m\).
\end{enumerate}
So the values of the general step are not unchanged when an empty step is inserted in the middle of a sequence; they change by the strictly increasing
relabelling \(\sigma_m\), independent of the point, which keeps the maximum loci.
\end{lemma}
\begin{proof}
\leanok
(1) The first two entries are kept and \(K\mapsto s_j(K)\) preserves the colexicographic order since \(s_j\) is strictly increasing; if \(K\) has an element \(\ge j\), its
largest element \(k\) is sent to \(k+1\notin K\). Maximum loci of \(g\circ f\) for strictly increasing \(g\) are those of \(f\) (Lemma~\ref{bcalc:max-locus-local}).
(2) The members through a point of \(E^{(j)}\) are the images under \(s_j\) of those of \(E\) (the inserted member contains no point), with the same exponents
(Lemma~\ref{hir:monomial-invariant-geometry}). The feasible sets for \(s_j(S)\) are the images of the feasible sets for \(S\), with the same cardinalities and sums, so
the largest key is the image of the largest key (Lemma~\ref{hir:monomial-invariant}). The monomial part only involves the members through the point
(Proposition~\ref{hir:monomial-part-along}), so it does not change, nor do the nonmonomial part, \(\nu\) (Lemma~\ref{bcalc:nonmonomial-order-local}) and the cosupport, which
does not involve \(E\).
(3) The transform replaces each member by its strict transform and appends \(\pi^{-1}C\) at the end (Lemma~\ref{hir:geo-transform-basic}); the strict transform of
\(\emptyset\) is \(\emptyset\), so inserting \(\emptyset\) at position \(j\le\#E\) commutes with one transform, and by induction with the transforms along every tower (the
length of \(E\) only grows). The rest is (2) and (1).
(4) The transform by \(\sigma\) is \(\underline{\mathcal I}'\) with \(\sigma^{-1}\emptyset=\emptyset\) appended at position \(\#E(\underline{\mathcal I}_{T_0})=m\); apply (3) along \(S\). For closed
\(N\) and members the strict transforms along \(\emptyset\) are the preimages, and the ideals of \(\underline{\mathcal I}'\) and of \(\sigma^*\underline{\mathcal I}_{T_0}\) agree on
\(N'\) (Lemma~\ref{mc:empty-blowup-inverse}); restricted-equivalent marked ideals have the same \(\gamma\) along every tower
(Lemma~\ref{bcalc:nonmonomial-order-monomial-transport}(2)).
(5) By Lemma~\ref{bcalc:general-value-order-zero}(1) and (3), (4) (the orders agree, so both points have order \(0\)); the values differ by (1).
\end{proof}

\paragraph{Compact sets along sequences of blowups.} Local finiteness of the resolution invariant (property (8) of
Definition~\ref{hir:centre-function}, and its form along growing sequences of canonical segments) is proved phase by phase: at a stage \(T\) one counts the later nonempty centres meeting the
preimage of a compact set \(K\) of the base against a compact subset of the stage \(T\), inside an open set on which a lower-dimensional invariant lives.
The lemma below gives the topological part.

\begin{lemma}[Sequences of blowups are proper; the compact part of the cosupport]
\label{bcalc:tower-preimage-compact}
\uses{ana:blowup-compact,hir:tower-sequence,hir:test-tower,hir:geo-transform-basic,bcalc:inv-domain-stable,hir:giraud-along,hir:step-one-well-founded,hir:centre-function,hir:canonical-segments,bcalc:tower-restrict}
\source{wlodarczyk05-hironaka,bierstone-milman08-functoriality}
\upstreamlean{TranscendentalNakamaye.Pending.TowerPreimageCompact.exists_mem_nhds_isCompact_preimage,TranscendentalNakamaye.Pending.TowerPreimageCompact.isCompact_preimage_of_isBlowupAlong,TranscendentalNakamaye.Pending.TowerPreimageCompact.isCompact_preimage_proj,TranscendentalNakamaye.Pending.TowerPreimageCompact.countCentresMeeting_append}\lean{Hironaka.TowerPreimageCompact.exists_restrict_step}\upstreamlean{TranscendentalNakamaye.Pending.TowerPreimageCompact.isClosed_cosupp_geoTransformAlong,TranscendentalNakamaye.Pending.TowerPreimageCompact.isCompact_preimage_inter_cosupp}\lean{Hironaka.TowerPreimageCompact.proj_mem_cosupp_append,Hironaka.TowerPreimageCompact.countCentresMeeting_eq_inter_cosupp}\upstreamlean{TranscendentalNakamaye.Pending.TowerPreimageCompact.countCentresMeeting_append_eq_inter_cosupp,TranscendentalNakamaye.Pending.TowerPreimageCompact.countCentresMeeting_append_eq_restrict}
\mixedleanok
Let \(K\) be a compact subset of the base \(X\).
\begin{enumerate}
\item A blowup \(\pi\colon Z\to M\) along a smooth centre is proper: \(\pi^{-1}(K)\) is compact for every compact \(K\subset M\). Hence \(p_T^{-1}(K)\) is compact for the
projection \(p_T\) of every tower \(T\).
\item Let \(\underline{\mathcal I}\) lie in the domain of Definition~\ref{hir:centre-function} and \(T\) be a multiple test blow-up of it. Then
\(K_T=p_T^{-1}(K)\cap\operatorname{cosupp}\underline{\mathcal I}_T\) is compact, and along every multiple test blow-up \(TS\) the points of
\(\operatorname{cosupp}\underline{\mathcal I}_{TS}\) lie over \(\operatorname{cosupp}\underline{\mathcal I}_T\).
\item For towers \(T\), \(S\), the number of nonempty centres of \(TS\) meeting the preimage of \(K\) is that of \(T\) plus the number of nonempty centres of \(S\)
meeting the preimage of \(p_T^{-1}(K)\); if \(TS\) is a multiple test blow-up of \(\underline{\mathcal I}\) (closed submanifold \(N\), coherent ideal), the latter are the
nonempty centres of \(S\) meeting the preimage of \(K_T\).
\item If moreover \(W\) is an open set of the last stage of \(T\) containing \(K_T\), then \(K_T\) is a compact subset of \(W\) and the nonempty centres of \(TS\) meeting the
preimage of \(K\) are those of \(T\) plus those of the restricted tail \(S|_W\) meeting the preimage of \(K_T\). The restriction to \(W\) of \(S\) followed by a blowup is
\(S|_W\) followed by the restricted blowup, so restrictions of growing sequences are growing sequences.
\end{enumerate}
\end{lemma}
\begin{proof}
\mixedleanok
(1) Every point of \(M\) has a neighbourhood with compact preimage (the local part of Lemma~\ref{ana:blowup-compact}: on the centre a closed ball of a flattening
chart, off it a compact neighbourhood missing the centre, over which \(\pi\) has a holomorphic section). Finitely many of them cover \(K\), and \(\pi^{-1}(K)\) is a
closed subset of the union of their preimages (\(M\) is Hausdorff). For towers, induction on the length.
(2) \(\underline{\mathcal I}_T\) lies in the domain (Proposition~\ref{bcalc:inv-domain-stable}), so its cosupport is closed (Lemma~\ref{hir:giraud-along}) and
\(K_T\) is a closed subset of the compact \(p_T^{-1}(K)\). For one test blowup a point of the new cosupport lies over the old cosupport: over the centre, which is
admissible, and off it, where the blowup is an isomorphism (Theorem~\ref{hir:step-one-well-founded}); induction along \(S\).
(3) The projection of \(TS\) is \(p_T\circ p_S\) (Lemma~\ref{hir:tower-sequence}); induction along \(S\). A centre of \(S\) is admissible, so it lies in the cosupport of
its stage, which lies over \(\operatorname{cosupp}\underline{\mathcal I}_T\) by (2); so it meets the preimage of \(p_T^{-1}(K)\) iff it meets that of \(K_T\).
(4) \(K_T\) is compact in \(W\) since \(W\) carries the subspace topology. The count along \(S|_W\) of a subset of \(W\) is the count along \(S\)
(Lemma~\ref{hir:canonical-segments}), and the restriction of a tower is defined step by step (Lemma~\ref{bcalc:tower-restrict}).
\end{proof}

\paragraph{The phases of a canonical segment.} Along a canonical segment of the general value (Definition~\ref{hir:general-value}) the nonempty
centres are maximum loci, and their orders \(\nu\) of the nonmonomial part form the sequence of phases of the segment. Local finiteness is proved phase by
phase: the lemma below collects what the order contributes, for every value function whose order dominates \(\nu\)
(Lemma~\ref{bcalc:nonmonomial-order-canonical}).

\begin{lemma}[The phases of a canonical segment: the order of the nonmonomial part]
\label{bcalc:nonmonomial-order-global-phase}
\uses{bcalc:nonmonomial-order-canonical,bcalc:nonmonomial-order-local,bcalc:nonmonomial-order-regime,bcalc:nonmonomial-order-tail,bcalc:nonmonomial-order-phase,hir:regime-run,hir:step-one-phase,hir:monomial-invariant,hir:general-value,bcalc:general-value-order-zero,bcalc:inv-domain-stable,hir:test-tower,hir:centre-function,hir:invariant-order}
\source{wlodarczyk05-hironaka,bierstone-milman08-functoriality}
\upstreamlean{TranscendentalNakamaye.Pending.NonmonomialOrderGlobalPhase.nmOrd_le_append_of_le,TranscendentalNakamaye.Pending.NonmonomialOrderGlobalPhase.nmOrd_proj_eq_of_le,TranscendentalNakamaye.Pending.NonmonomialOrderGlobalPhase.ordRegime_and_length_le_regimeRun_of_le,TranscendentalNakamaye.Pending.NonmonomialOrderGlobalPhase.cosupp_subset_ordMask_of_le,TranscendentalNakamaye.Pending.NonmonomialOrderGlobalPhase.nmOrd_le_of_mem_centre,TranscendentalNakamaye.Pending.NonmonomialOrderGlobalPhase.pairwise_centreValues_nmOrd,TranscendentalNakamaye.Pending.NonmonomialOrderGlobalPhase.exists_centre_nmOrd_le_of_isCompact,TranscendentalNakamaye.Pending.NonmonomialOrderGlobalPhase.genInv_snd_of_fst_eq_zero,TranscendentalNakamaye.Pending.NonmonomialOrderGlobalPhase.not_strictAnti_zero_inl}\lean{Hironaka.NonmonomialOrderGlobalPhase.wellFoundedLT_genValue}
\mixedleanok
Let \(n\ge1\), \(\underline{\mathcal I}\) lie in the domain of Definition~\ref{hir:centre-function}, and \(f\) be a value function whose order dominates \(\nu\)
(Lemma~\ref{bcalc:nonmonomial-order-canonical}). Let \(T_0\) be a multiple test blow-up with \(\nu\le o\) on \(\operatorname{cosupp}\underline{\mathcal I}_{T_0}\).
\begin{enumerate}
\item (The phase is absorbing) Along every canonical segment \(T_0S\) which is a multiple test blow-up, every point \(z\) of the cosupport lies over the cosupport at
\(T_0\), and \(\nu(z)\le\nu(p_S(z))\le o\); if \(\nu(z)=o\), then \(\nu(p_S(z))=o\).
\item (Start of the run) Such a point \(z\) of order \(o\) is in the regime \(o\) (Lemma~\ref{bcalc:nonmonomial-order-regime}) and its run is at least the length of
\(S\): it reached the order \(o\) at the last stage of \(T_0\) or before. Moreover \(\operatorname{cosupp}\underline{\mathcal I}_{T_0}\subset W_o(\underline{\mathcal I}_{T_0})\),
the mask of Lemma~\ref{bcalc:nonmonomial-order-tail}.
\item (Phase entry; the phases never go up) If \(T\) followed by a blowup along \(C\) is a canonical segment, every point of \(C\) bounds the order on
\(\operatorname{cosupp}\underline{\mathcal I}_T\). Along a canonical segment which is a multiple test blow-up the orders of the points of the nonempty centres, in the order
of the steps, never increase, and every point of the last cosupport has order at most each of them.
\item (Over a compact set) For every compact \(K\subset X\) there is \(q\) such that along every canonical segment which is a multiple test blow-up, every point of a
nonempty centre meeting the preimage of \(K\) has order \(\le q\).
\item (The monomial phase) A value of the general step with first entry \(0\) is \((0,\gamma)\); no sequence of such values decreases strictly forever, since the values
of \(\gamma\) are well-founded (Lemma~\ref{hir:monomial-invariant}). If the values of the maximal-order invariant are well-founded, so are the values of the general
step.
\end{enumerate}
\end{lemma}
\begin{proof}
\mixedleanok
(1) The order never goes up along a canonical segment (Lemma~\ref{bcalc:nonmonomial-order-canonical}), and the image of a point of the cosupport lies in the cosupport.
(2) Induction along \(S\): after one more step a point \(z\) of order \(o\) lies over a point of the cosupport of order exactly \(o\) by (1), which is in the
regime with run at least the length so far (Lemma~\ref{hir:regime-run}). A point of the cosupport at \(T_0\) has order \(\le o\), so it is off the closed set
\(\{\nu\ge o+1\}\), i.e.\ in \(W_o\) (Lemma~\ref{bcalc:nonmonomial-order-tail}).
(3) A nonempty centre of a canonical step is the maximum locus, and a point of the maximum locus has the largest order on the cosupport
(Lemma~\ref{bcalc:nonmonomial-order-canonical}). Induction along the segment: the new entry is the order at a point of the new centre, which lies in the cosupport and so
has order at most every earlier entry; after the step the orders are at most those of the images, which are at most every entry, including the new one by phase entry.
(4) Lemma~\ref{bcalc:nonmonomial-order-phase} bounds \(\nu\) at the points of the cosupport over \(K\) along all such segments; a nonempty centre meeting the preimage of
\(K\) is the maximum locus, which lies in one level set of \(\nu\) and contains a point of the cosupport over \(K\).
(5) Definition~\ref{hir:general-value} and Lemma~\ref{bcalc:general-value-order-zero}(1). The pairs \((0,\gamma)\) are ordered as the \(\gamma\); lexicographic products and
ordered sums of well-founded orders are well-founded (as for the tuples of Lemma~\ref{hir:invariant-order}).
\end{proof}

\paragraph{Growing sequences: restriction and tails.} Property (8) in its form along growing sequences
(Definition~\ref{hir:centre-function-seq}) bounds the numbers of nonempty centres meeting a compact set along a growing sequence of towers (each extending the
previous one by one blowup). The phase argument passes, at a stage, to the rest of the sequence restricted to open sets of that stage.

\begin{lemma}[Restriction and tails of growing sequences]
\label{bcalc:growing-seq-tail}
\uses{hir:centre-function-seq,bcalc:tower-preimage-compact,bcalc:tower-restrict,hir:canonical-segments,hir:tower-sequence,hir:test-tower,hir:geo-transform-basic}
\source{wlodarczyk05-hironaka,bierstone-milman08-functoriality}
\lean{Hironaka.ResolutionInvariantSeq.GrowingSeq.restrict,Hironaka.ResolutionInvariantSeq.GrowingSeq.restrict_stage,Hironaka.ResolutionInvariantSeq.GrowingSeq.count_restrict}\upstreamlean{TranscendentalNakamaye.Pending.ResolutionInvariantSeq.GrowingSeq.tailTower,TranscendentalNakamaye.Pending.ResolutionInvariantSeq.GrowingSeq.stage_add_eq_append,TranscendentalNakamaye.Pending.ResolutionInvariantSeq.GrowingSeq.tail,TranscendentalNakamaye.Pending.ResolutionInvariantSeq.GrowingSeq.tail_stage,TranscendentalNakamaye.Pending.ResolutionInvariantSeq.GrowingSeq.count_add}\lean{Hironaka.ResolutionInvariantSeq.GrowingSeq.count_mono}\upstreamlean{TranscendentalNakamaye.Pending.ResolutionInvariantSeq.GrowingSeq.exists_bound_of_tail_restrict}
\mixedleanok
Let \(s=(T_i)_{i\ge0}\) be a growing sequence of towers over \(X\) (\(T_0\) empty, \(T_{i+1}\) is \(T_i\) followed by one blowup).
\begin{enumerate}
\item For an open \(U\subset X\), the restrictions \(T_i|_U\) form a growing sequence over \(U\); for \(K\subset U\) the numbers of nonempty centres meeting the
preimage of \(K\) along it and along \(s\) agree.
\item For a stage \(a\), write \(T_{a+i}=T_aS_i\); the towers \(S_i\) over the last stage of \(T_a\) form a growing sequence (the tail of \(s\) from \(a\)), and the number of
nonempty centres of \(T_{a+i}\) meeting the preimage of \(K\) is that of \(T_a\) plus the number of nonempty centres of \(S_i\) meeting the preimage of \(p_{T_a}^{-1}(K)\).
These numbers never decrease along \(s\).
\item Let the \(T_i\) be multiple test blow-ups of a marked ideal \(\underline{\mathcal I}\) of the domain of Definition~\ref{hir:centre-function}, \(a\) a stage, and
\(L_j\subset W_j\) (\(j\) in a finite set, \(W_j\) open in the last stage of \(T_a\)) with \(p_{T_a}^{-1}(K)\cap\operatorname{cosupp}\underline{\mathcal I}_{T_a}\subset\bigcup_jL_j\). If for
each \(j\) the numbers of nonempty centres meeting the preimage of \(L_j\) along the tail from \(a\) restricted to \(W_j\) are bounded, then the numbers of nonempty
centres meeting the preimage of \(K\) along \(s\) are bounded.
\end{enumerate}
\end{lemma}
\begin{proof}
\mixedleanok
(1) The restriction of a tower followed by a blowup is the restriction followed by the restricted blowup (Lemma~\ref{bcalc:tower-preimage-compact}(4),
Lemma~\ref{bcalc:tower-restrict}), and the restriction of the empty tower is empty; the counts agree by Lemma~\ref{hir:canonical-segments}.
(2) Choose the blowups of \(s\) after the stage \(a\); by induction \(T_{a+i}=T_aS_i\) (Lemma~\ref{hir:tower-sequence}), and the count splits by
Lemma~\ref{bcalc:tower-preimage-compact}(3). A blowup adds one or no centre.
(3) For \(i\le a\) the count is at most that of \(T_a\) by (2). For \(i=a+d\), the count of \(T_a\) plus the count of \(S_d\) for \(p_{T_a}^{-1}(K)\), which is the count for
\(p_{T_a}^{-1}(K)\cap\operatorname{cosupp}\underline{\mathcal I}_{T_a}\) (Lemma~\ref{bcalc:tower-preimage-compact}(3)), at most the sum over \(j\) of the counts for \(L_j\)
(counts are monotone and subadditive in the set, Lemma~\ref{hir:canonical-segments}), which are those along the restricted tails by (1).
\end{proof}

\paragraph{The count of centres by phases.} Along a canonical segment of the general value the values of the nonempty centres, in the order of
the steps, strictly decrease (Theorem~\ref{hir:step-one-phase}, Theorem~\ref{hir:general-admissible}); their first entries are the phases. The lemma below splits the count
of property (8) along growing sequences (Definition~\ref{hir:centre-function-seq}) by phases and settles the phase of order \(0\).

\begin{lemma}[The count of centres by phases; the monomial phase is finite]
\label{bcalc:general-phase-count}
\uses{hir:general-value,hir:general-admissible,hir:max-order-value,hir:step-one-phase,hir:centre-function-seq,bcalc:nonmonomial-order-global-phase,hir:monomial-invariant,hir:centre-function}
\source{wlodarczyk05-hironaka,bierstone-milman08-functoriality}
\lean{Hironaka.GeneralPhaseCount.exists_bound_of_isPrefix}\upstreamlean{TranscendentalNakamaye.Pending.GeneralPhaseCount.countCentresMeeting_le_length_filter,TranscendentalNakamaye.Pending.GeneralPhaseCount.length_filter_le_eq_sum,TranscendentalNakamaye.Pending.GeneralPhaseCount.countCentresMeeting_le_sum_phases}\lean{Hironaka.GeneralPhaseCount.centreValues_isPrefix_succ}\upstreamlean{TranscendentalNakamaye.Pending.GeneralPhaseCount.exists_bound_phase_zero,TranscendentalNakamaye.Pending.GeneralPhaseCount.exists_bound_of_phases}
\mixedleanok
Let \(n\ge1\), \(\underline{\mathcal I}\) lie in the domain of Definition~\ref{hir:centre-function}, \(M\) be a value function for the marked ideals of maximal order
(Definition~\ref{hir:max-order-value}) and \(\operatorname{inv}\) the general value. For a tower \(T\) let \(v(T)\) be the list of the values of the nonempty centres
(Theorem~\ref{hir:step-one-phase}).
\begin{enumerate}
\item Prefix-growing lists, each strictly decreasing, with entries in a set on which no sequence strictly decreases forever, have bounded lengths.
\item Let \(K\subset X\) and \(q\) be such that along every canonical segment every point of a nonempty centre meeting the preimage of \(K\) has order \(\le q\)
(for compact \(K\), Lemma~\ref{bcalc:nonmonomial-order-global-phase}(4)). Along a canonical segment \(T\), the number of nonempty centres meeting the preimage of \(K\)
is at most the number of entries of \(v(T)\) with first entry \(\le q\), which is the sum over \(o=0,\dots,q\) of the numbers \(\ell_o(T)\) of entries with first entry \(o\).
\item Along a growing sequence \((T_i)\) of canonical segments, \(\ell_0(T_i)\) is bounded: the monomial phase is finite.
\item Hence, if \(\ell_o(T_i)\) is bounded for each \(o=1,\dots,q\), the numbers of nonempty centres of \(T_i\) meeting the preimage of \(K\) are bounded.
\end{enumerate}
\end{lemma}
\begin{proof}
\mixedleanok
(1) Otherwise, for each \(k\) choose a list of length \(>k\) and take its \(k\)-th entry; by the prefix property all these entries are entries of one later list at the
positions \(0,1,2,\dots\), so they strictly decrease.
(2) Induction along \(T\): a step along a nonempty centre adds the value at a point of the centre, whose first entry is the order of that point
(Definition~\ref{hir:general-value}); if the centre meets the preimage of \(K\), it is \(\le q\). Partition the entries \(\le q\) by their first entry.
(3) The lists \(v(T_i)\) grow by prefixes and are strictly decreasing (Theorem~\ref{hir:step-one-phase} with the drop property of Theorem~\ref{hir:general-admissible});
so are their entries of first entry \(0\), which are \((0,\gamma)\) (Lemma~\ref{bcalc:nonmonomial-order-global-phase}(5)); \(\gamma\) has well-founded values
(Lemma~\ref{hir:monomial-invariant}), so (1) applies.
(4) By (2), the count is at most \(\sum_{o\le q}\ell_o(T_i)\), bounded by (3) and the hypothesis.
\end{proof}

\paragraph{The companion phase.} While the order of the nonmonomial part stays \(o\ge1\), Włodarczyk's Step 2a ``repeats the procedure'' for one
companion ideal. The lemma below makes this precise along canonical segments of the general value and bounds the phase of order \(o\) along growing
sequences (Definition~\ref{hir:centre-function-seq}) from property (8) of the value function for the marked ideals of maximal order.

\begin{lemma}[One companion ideal for a phase; the phases of positive order are finite]
\label{bcalc:general-companion-phase}
\uses{hir:general-value,hir:general-value-basic,hir:general-canonical-tail,hir:general-compare,hir:general-admissible,hir:max-order-value,hir:centre-function-seq,hir:step-one-phase,hir:tower-split,bcalc:nonmonomial-order-global-phase,bcalc:nonmonomial-order-regime,bcalc:nonmonomial-order-canonical,bcalc:nonmonomial-order-local,bcalc:nonmonomial-order-tail,bcalc:max-locus-local,bcalc:inv-domain-stable,bcalc:growing-seq-tail,bcalc:tower-preimage-compact,bcalc:general-phase-count,hir:canonical-segments,hir:centre-function}
\source{wlodarczyk05-hironaka,bierstone-milman08-functoriality}
\upstreamlean{TranscendentalNakamaye.Pending.GeneralCompanionPhase.exists_split_append,TranscendentalNakamaye.Pending.GeneralCompanionPhase.centresOrdGt_append,TranscendentalNakamaye.Pending.GeneralCompanionPhase.CentresOrdGe,TranscendentalNakamaye.Pending.GeneralCompanionPhase.genInv_eq_of_mem_cosupp_phase,TranscendentalNakamaye.Pending.GeneralCompanionPhase.isCanonical_restrict_phase,TranscendentalNakamaye.Pending.GeneralCompanionPhase.countPhase,TranscendentalNakamaye.Pending.GeneralCompanionPhase.countPhase_step,TranscendentalNakamaye.Pending.GeneralCompanionPhase.countPhase_step_of,TranscendentalNakamaye.Pending.GeneralCompanionPhase.countPhase_step_of_not,TranscendentalNakamaye.Pending.GeneralCompanionPhase.countPhase_le_countCentresMeeting,TranscendentalNakamaye.Pending.GeneralCompanionPhase.countPhase_append,TranscendentalNakamaye.Pending.GeneralCompanionPhase.countPhase_eq_zero_of_centresOrdGt,TranscendentalNakamaye.Pending.GeneralCompanionPhase.countCentresMeeting_le_sum_countPhase,TranscendentalNakamaye.Pending.GeneralCompanionPhase.countPhase_le_filter,TranscendentalNakamaye.Pending.GeneralCompanionPhase.countPhase_le_of_le,TranscendentalNakamaye.Pending.GeneralCompanionPhase.exists_bound_countPhase}
\upstreamleanok
Let \(n\ge1\), \(\underline{\mathcal I}\) lie in the domain of Definition~\ref{hir:centre-function}, \(M\) be a value function for the marked ideals of maximal order
(Definition~\ref{hir:max-order-value}), \(\operatorname{inv}\) the general value and \(\nu\) the order of the nonmonomial part. Let \(o\ge1\), and let \(T_0\) be a multiple
test blow-up with \(\nu\le o\) on \(\operatorname{cosupp}\underline{\mathcal I}_{T_0}\) and every centre in \(\{\nu\ge o+1\}\); let \(W\) be the open set \(\{\nu\le o\}\) of the last
stage of \(T_0\) and \(\underline{\mathcal O}\) the companion ideal of order \(o\) formed there on \(W\) (Definition~\ref{hir:general-value}).
\begin{enumerate}
\item If \(T_0S\) is a canonical segment of \(\operatorname{inv}\) and every point of every centre of \(S\) has order \(\ge o\), then \(S|_W\) is a canonical segment of
\(\underline{\mathcal O}\) for \(M\), and at every point \(q\) of the cosupport of the transform of \(\underline{\mathcal O}\) along \(S|_W\) the general value is
\((o,M_{\underline{\mathcal O}}(S|_W)(q))\).
\item For \(K\subset X\) let \(c_o(K,T)\) be the number of steps of \(T\) whose centre has a point over \(K\) of order \(o\). Then \(c_o(K,T)\) is at most the number of
centres meeting the preimage of \(K\), it is additive along concatenations, and it vanishes when every centre lies in \(\{\nu\ge o+1\}\). Along a canonical segment, with \(q\) as in
Lemma~\ref{bcalc:general-phase-count}(2), the number of centres meeting the preimage of \(K\) is at most \(\sum_{o\le q}c_o(K,T)\), and \(c_o(K,T)\le\ell_o(T)\).
\item Suppose that for every marked ideal of the domain of maximal order along its submanifold, every compact set and every growing sequence of canonical segments for
\(M\), the numbers of centres meeting the compact set are bounded (property (8) of \(M\), per sequence). Then for \(K\) compact and \(o\ge1\), along every growing
sequence \((T_i)\) of canonical segments of \(\operatorname{inv}\), \(c_o(K,T_i)\) is bounded.
\end{enumerate}
\end{lemma}
\begin{proof}
\upstreamleanok
(1) By induction along \(S\). A nonempty centre is the maximum locus of \(\operatorname{inv}\); its points have order \(\le o\) (the phase is absorbing,
Lemma~\ref{bcalc:nonmonomial-order-global-phase}(1)) and \(\ge o\), hence \(o\), and lie over \(W\). Over \(W\) the cosupport of the companion's transform is the set of points of
order \(o\) of the cosupport (Theorem~\ref{hir:general-canonical-tail}). At such a point the run of order \(o\) started at or before \(T_0\)
(Lemma~\ref{bcalc:nonmonomial-order-global-phase}(1)); splitting \(T_0S\) at the start of the run (Lemma~\ref{hir:tower-split}), the steps from there to \(T_0\) have centres in
\(\{\nu\ge o+1\}\), which miss the mask of the point (Lemma~\ref{bcalc:nonmonomial-order-tail}), so the companion value re-bases to \(T_0\)
(Theorem~\ref{hir:general-canonical-tail}): it is \((o,M_{\underline{\mathcal O}}(S|_W))\). Every other point of the cosupport has smaller order, hence smaller value, and the centre
has a point of order \(o\) over \(W\); so the trace of the maximum locus of \(\operatorname{inv}\) on the preimage of \(W\) is the maximum locus of \(M\)
(Lemma~\ref{bcalc:max-locus-local}).
(2) Induction along \(T\). For the sum: a centre meeting the preimage of \(K\) at a point \(c\) has \(\nu(c)\le q\) and is counted in \(c_{\nu(c)}\). For \(\ell_o\): the points of a
canonical centre have one value, whose first entry is their order (Lemma~\ref{hir:general-value-basic}, Lemma~\ref{bcalc:nonmonomial-order-canonical}).
(3) If at some stage \(\nu<o\) on the whole cosupport, this persists (Lemma~\ref{bcalc:nonmonomial-order-canonical}) and \(c_o\) is constant from there on. Otherwise every
canonical centre has order \(\ge o\). If at every stage some point of the cosupport has order \(>o\), no centre has order \(o\) and \(c_o=0\). Otherwise let \(a\) be the first stage
with \(\nu\le o\) on the cosupport: every centre before \(a\) lies in \(\{\nu\ge o+1\}\) (it is a maximum locus and some point has order \(>o\);
Lemma~\ref{bcalc:nonmonomial-order-local}), and every centre after \(a\) has order \(o\). By (1) the restricted tails of the sequence after \(a\)
(Lemma~\ref{bcalc:growing-seq-tail}) are canonical segments of the one companion ideal \(\underline{\mathcal O}\) at \(T_a\), which lies in the domain and has maximal order
(Theorem~\ref{hir:general-compare}). The part of the cosupport at \(T_a\) over \(K\) is compact and lies in \(W\), and the centres after \(a\) meeting the preimage of \(K\) are
those of the restricted tails meeting it (Lemma~\ref{bcalc:tower-preimage-compact}); the hypothesis bounds them, and \(c_o\) is at most the number of centres meeting the
preimage of \(K\).
\end{proof}

\paragraph{Admissible centres of the restricted marked ideal of Step 1ab.} In Step 1ab of Proposition~\ref{hir:maximal-order-step} the value at a
point of level \(s\ge1\) is that of the lower-dimensional invariant of the restricted marked ideal of Lemma~\ref{bcalc:member-restriction-domain} on a mask.
Admissibility of its maximum locus (Definition~\ref{hir:centre-function}(1)) is inherited from the inductive invariant through the lemma below, Włodarczyk's
Lemma ``le:S'' for a single centre.

\begin{lemma}[Admissible centres of the restricted marked ideal]
\label{bcalc:member-restrict-admissible}
\uses{bcalc:member-restriction-domain,mc:coefficient-restriction-drop-along,mc:coefficient-restriction-full-along,bcalc:nc-comparison,hir:drop-containing-members,hir:drop-members,hir:marked-equivalence-basic,hir:wellformed-tower,hir:step-one-compare,bcalc:member-restriction-rebase,hir:step-one-value,hir:step-one-lower-domain,hir:open-restriction-tower,hir:open-restriction-marked,hir:geo-transform-basic,mc:contact-independence-along}
\source{wlodarczyk05-hironaka}
\upstreamlean{TranscendentalNakamaye.Pending.MemberRestrictAdmissible.isAdmissibleCentre_memberRestrict_iff,TranscendentalNakamaye.Pending.MemberRestrictAdmissible.isAdmissibleCentre_memberRestrict_iff',TranscendentalNakamaye.Pending.MemberRestrictAdmissible.isAdmissibleCentre_lowerMarked_succ_iff,TranscendentalNakamaye.Pending.MemberRestrictAdmissible.isAdmissibleCentre_of_lowerMarked_succ}
\upstreamleanok
Let \(n\ge1\), \(\underline{\mathcal I}=(M,N,E,\mathfrak a,d)\) be well formed with \(d\ge1\), \(S\) a set of members of \(E\), and \(\underline{\mathcal S}\) the restricted marked ideal
\((M,N\cap H_S,E_S,\mathcal C_N(\mathfrak a),\mu_{\mathcal C})\) of Lemma~\ref{bcalc:member-restriction-domain}.
\begin{enumerate}
\item Along a multiple test blow-up \(T\) of \(\underline{\mathcal S}\), a set \(C\) of the last stage is an admissible centre of \(\underline{\mathcal S}_T\) if and only if it is an
admissible centre of \(\underline{\mathcal I}_T\) contained in the strict transform \((N\cap H_S)_T\), equivalently in \(\operatorname{cosupp}\underline{\mathcal S}_T\).
\item In Step 1ab, let \(T_0T_1\) be a multiple test blow-up of \(\overline{\underline{\mathcal J}}\), \(y_0\) a point of the last stage of \(T_0\) lying in its own mask \(W\), and
\(T_1|_W\) a multiple test blow-up of the lower marked ideal \(\underline{\mathcal S}(y_0)|_W\) (Definition~\ref{hir:step-one-value}). Then a set \(C\) of the preimage of \(W\) is an
admissible centre of the transform of \(\underline{\mathcal S}(y_0)|_W\) along \(T_1|_W\) if and only if it is an admissible centre of \(\overline{\underline{\mathcal J}}_{T_0T_1}\)
restricted to the preimage of \(W\) and lies on the strict transforms along \(T_1\) of the original members through \(y_0\).
\end{enumerate}
\end{lemma}
\begin{proof}
\upstreamleanok
(1) By Lemma~\ref{mc:coefficient-restriction-drop-along}, \(T\) is a multiple test blow-up of \(\underline{\mathcal I}\) with centres in the strict transforms of \(N\cap H_S\), and
\(\operatorname{cosupp}\underline{\mathcal S}_T=\operatorname{cosupp}\underline{\mathcal I}_T\cap(N\cap H_S)_T\); the same holds for the restricted coefficient ideal with the full divisor
(Lemma~\ref{mc:coefficient-restriction-full-along}), whose transform has the divisor of \(\underline{\mathcal I}_T\) (Lemma~\ref{hir:marked-equivalence-basic}), simple normal crossings
(Lemma~\ref{hir:wellformed-tower}). The members of the transformed divisor that are not members of the transform of \(E_S\) contain \((N\cap H_S)_T\)
(Lemma~\ref{hir:drop-containing-members}), so for \(C\subset(N\cap H_S)_T\) normal crossings with the two divisors are equivalent (Lemma~\ref{bcalc:nc-comparison}); admissibility
for \(\underline{\mathcal S}_T\) is admissibility for the transform with the divisor replaced (Lemma~\ref{hir:drop-members}). Both conditions put \(C\) in
\(\operatorname{cosupp}\underline{\mathcal I}_T\cap(N\cap H_S)_T\).
(2) On the mask the lower marked ideal is the restricted marked ideal of \(\overline{\underline{\mathcal J}}_{T_0}|_W\) at the members through \(y_0\)
(Lemma~\ref{bcalc:member-restriction-rebase}), and \(\overline{\underline{\mathcal J}}_{T_0}\) is well formed with mark \(\ge1\) (Lemma~\ref{hir:step-one-lower-domain},
Lemma~\ref{mc:contact-independence-along}); apply (1) on \(W\). The transform of \(\overline{\underline{\mathcal J}}_{T_0}|_W\) along \(T_1|_W\) is the restriction of
\(\overline{\underline{\mathcal J}}_{T_0T_1}\) (Lemmas~\ref{hir:open-restriction-tower} and~\ref{hir:geo-transform-basic}), and the cosupport of the lower transform consists of the
cosupport points of \(\overline{\underline{\mathcal J}}_{T_0T_1}\) on the strict transforms of the original members through \(y_0\) (Lemma~\ref{hir:step-one-compare};
Lemma~\ref{hir:open-restriction-marked} for the restricted cosupport).
\end{proof}

\begin{corollary}[Admissible centres of the restricted marked ideal on open sets]
\label{bcalc:member-restrict-admissible-local}
\uses{bcalc:member-restrict-admissible,hir:open-restriction-marked,bcalc:nc-comparison,mc:coefficient-restriction-drop-along,mc:coefficient-restriction-full-along,hir:drop-containing-members,hir:marked-equivalence-basic,hir:wellformed-tower,bcalc:member-restriction-rebase,hir:step-one-value,hir:step-one-lower-domain,hir:open-restriction-tower,hir:geo-transform-basic,mc:contact-independence-along}
\source{wlodarczyk05-hironaka}
\lean{Hironaka.MemberRestrictAdmissible.isAdmissibleCentre_pullback_val_of_relations,Hironaka.MemberRestrictAdmissible.memberRestrict_tower_relations,Hironaka.MemberRestrictAdmissible.isAdmissibleCentre_pullback_val_of_memberRestrict,Hironaka.MemberRestrictAdmissible.isAdmissibleCentre_pullback_val_of_lowerMarked_succ}
\leanok
Let \(\underline{\mathcal A}\), \(\underline{\mathcal B}\) be marked ideals on one manifold, \(\underline{\mathcal A}\) with coherent ideal and \(\underline{\mathcal B}\) well formed, such that
\(\operatorname{cosupp}\underline{\mathcal A}\subset\operatorname{cosupp}\underline{\mathcal B}\), every member of the divisor of \(\underline{\mathcal A}\) is a member of that of
\(\underline{\mathcal B}\), and every member of \(\underline{\mathcal B}\) missing from \(\underline{\mathcal A}\) contains \(\operatorname{cosupp}\underline{\mathcal A}\). Then for every open
set \(W''\), every admissible centre of \(\underline{\mathcal A}|_{W''}\) is an admissible centre of \(\underline{\mathcal B}|_{W''}\). The three relations hold for
\(\underline{\mathcal A}=\underline{\mathcal S}_T\), \(\underline{\mathcal B}=\underline{\mathcal I}_T\) in the situation of Lemma~\ref{bcalc:member-restrict-admissible}(1), and for the lower
marked ideal of Step 1ab and the restriction of \(\overline{\underline{\mathcal J}}_{T_0T_1}\) to the preimage of its mask in the situation of
Lemma~\ref{bcalc:member-restrict-admissible}(2); so the transfers of that lemma hold on every open set of the last stage.
\end{corollary}
\begin{proof}
\leanok
The restricted centre is a closed submanifold of \(W''\) for both; restriction to an open set takes cosupports to their traces
(Lemma~\ref{hir:open-restriction-marked}), so the trace of \(\operatorname{cosupp}\underline{\mathcal A}\subset\operatorname{cosupp}\underline{\mathcal B}\) gives the cosupport condition; the
restricted divisors are the traces of the divisors, that of \(\underline{\mathcal B}\) simple normal crossings, and the traces of the members of \(\underline{\mathcal B}\)
missing from \(\underline{\mathcal A}\) contain the centre, so normal crossings pass (Lemma~\ref{bcalc:nc-comparison}). The relations along \(T\): the cosupport of
\(\underline{\mathcal S}_T\) is \(\operatorname{cosupp}\underline{\mathcal I}_T\cap(N\cap H_S)_T\) (Lemma~\ref{mc:coefficient-restriction-drop-along}), the members of the transform of
\(E_S\) are members of \(E_T\), and the others contain \((N\cap H_S)_T\) (Lemma~\ref{hir:drop-containing-members}, with Lemmas~\ref{mc:coefficient-restriction-full-along},
\ref{hir:marked-equivalence-basic} and~\ref{hir:wellformed-tower}). For Step 1ab, the lower marked ideal on its mask is the restricted marked ideal of
\(\overline{\underline{\mathcal J}}_{T_0}|_W\) (Lemma~\ref{bcalc:member-restriction-rebase}), and the transform of \(\overline{\underline{\mathcal J}}_{T_0}|_W\) along \(T_1|_W\) is
the restriction of \(\overline{\underline{\mathcal J}}_{T_0T_1}\) (Lemmas~\ref{hir:open-restriction-tower} and~\ref{hir:geo-transform-basic}).
\end{proof}

\paragraph{The companion ideal along isomorphisms of towers.} Property (6) of the general step (Proposition~\ref{hir:general-step}) compares the
companion ideal formed at the start stage of a point, on its mask, with the same construction for \(\Phi^*\underline{\mathcal I}\) along an isomorphic
tower. An isomorphism \(T_0'T_1'\to T_0T_1\) over \(\Phi\) splits into \(\chi_0\colon T_0'\to T_0\) over \(\Phi\) and \(T_1'\to T_1\) over \(\chi_0\)
(Lemma~\ref{hir:tower-split-iso}(2)).

\begin{lemma}[The companion ideal along isomorphisms of towers]
\label{bcalc:companion-tower-iso}
\uses{bcalc:companion-rebase,bcalc:companion-domain,bcalc:nonmonomial-order-local,mc:tower-isomorphism-transform,mc:tower-iso-opens,mc:contact-marked-canonical,bcalc:inv-domain-stable,hir:canonical-segments,hir:tower-isomorphism,hir:tower-split-iso,hir:centre-function}
\source{wlodarczyk05-hironaka,bierstone-milman08-functoriality}
\lean{Hironaka.CompanionTowerIso.companionOf_isTowerIso}\upstreamlean{TranscendentalNakamaye.Pending.CompanionTowerIso.bound_isTowerIso,TranscendentalNakamaye.Pending.CompanionTowerIso.isGeoTestTower_companionOf_isTowerIso,TranscendentalNakamaye.Pending.CompanionTowerIso.inv_companionOf_isTowerIso}
\mixedleanok
Let \(n\ge1\), \(\underline{\mathcal I}\) in the domain of Definition~\ref{hir:centre-function}, \(o\ge1\), \(\chi_0\colon T_0'\to T_0\) an isomorphism of towers over a
biholomorphism \(\Phi\colon X'\to X\) with \(T_0\) a multiple test blow-up of \(\underline{\mathcal I}\), and \(W'\), \(W\) open sets of the last stages with
\(q\in W'\iff\chi_0(q)\in W\). Write \(\mathcal O_0=\mathcal O(\underline{\mathcal I}_{T_0},o)|_W\) and \(\mathcal O_0'=\mathcal O((\Phi^*\underline{\mathcal I})_{T_0'},o)|_{W'}\).
\begin{enumerate}
\item \(\mathcal O_0'=\psi^*\mathcal O_0\) for every biholomorphism \(\psi\colon W'\to W\) agreeing with \(\chi_0\).
\item If \(\nu\le o\) on \(\operatorname{cosupp}\underline{\mathcal I}_{T_0}\cap W\) (Lemma~\ref{bcalc:nonmonomial-order-local}), then \(\nu\le o\) on
\(\operatorname{cosupp}(\Phi^*\underline{\mathcal I})_{T_0'}\cap W'\).
\item Let \(T_1'\to T_1\) be an isomorphism of towers over \(\chi_0\). If \(T_1|_W\) is a multiple test blow-up of \(\mathcal O_0\), then \(T_1'|_{W'}\) is one of \(\mathcal O_0'\); if
moreover \(\operatorname{inv}\) is a resolution invariant and \(T_1|_W\) a canonical segment of \(\mathcal O_0\), then \(\mathcal O_0'\) lies in the domain, \(T_1'|_{W'}\) is a
canonical segment of it, and \(\operatorname{inv}_{\mathcal O_0'}(T_1'|_{W'})(p)=\operatorname{inv}_{\mathcal O_0}(T_1|_W)(q)\) at corresponding points of the cosupports.
\end{enumerate}
\end{lemma}
\begin{proof}
\mixedleanok
(1) \((\Phi^*\underline{\mathcal I})_{T_0'}=\chi_0^*\underline{\mathcal I}_{T_0}\) (Lemma~\ref{mc:tower-isomorphism-transform}), the companion ideal commutes with biholomorphisms
(Lemma~\ref{bcalc:companion-rebase}(2)), and restricting \(\chi_0^*\mathcal O\) to \(W'\) is pulling back \(\mathcal O|_W\) along \(\psi\), as \(\chi_0\) and \(\psi\) agree on \(W'\)
(Theorem~\ref{mc:contact-marked-canonical}); \(\mathcal O(\underline{\mathcal I}_{T_0},o)\) has a coherent ideal (Proposition~\ref{bcalc:companion-domain}, the stage in the domain by
Lemma~\ref{bcalc:inv-domain-stable}). (2) Cosupports correspond along \(\chi_0\) (Lemma~\ref{mc:tower-isomorphism-transform}) and so does \(\nu\)
(Lemma~\ref{bcalc:nonmonomial-order-local}(3)). (3) The isomorphism restricts to \(T_1'|_{W'}\to T_1|_W\) over the restriction \(\psi\) of \(\chi_0\) (Lemma~\ref{mc:tower-iso-opens});
multiple test blow-ups correspond along isomorphisms of towers (Lemma~\ref{mc:tower-isomorphism-transform}) and so do canonical segments
(Lemma~\ref{hir:canonical-segments}(2)); \(\psi^*\mathcal O_0\) lies in the domain (Lemma~\ref{bcalc:inv-domain-stable}); the values are those of
Definition~\ref{hir:centre-function}(6).
\end{proof}

\paragraph{Step 1ab's lower values along isomorphisms of towers.} Properties (5)--(7) of the value function of Step 1
(Proposition~\ref{hir:maximal-order-step}, Definition~\ref{hir:step-one-value-masked}) compare lower values computed from different but corresponding data.
At the levels \(s\ge1\) the lower marked ideal is the restricted marked ideal of Lemma~\ref{bcalc:member-restriction-domain}; for level \(0\) see
Theorem~\ref{mc:contact-lower-value}.

\begin{lemma}[Step 1ab's lower values along isomorphisms of towers]
\label{bcalc:member-restrict-tower-iso}
\uses{bcalc:member-restriction-domain,bcalc:member-restriction-rebase,mc:homogenized-along-etale,hir:step-one-value-masked,hir:step-one-nearby,mc:contact-lower-value,mc:contact-stage-kept,mc:tower-isomorphism-append,mc:tower-isomorphism-transform,mc:tower-iso-opens,mc:contact-marked-canonical,hir:biholomorphism-invariance,bcalc:inv-domain-stable,hir:canonical-segments,hir:centre-function,hir:step-one-lower-domain,hir:geo-transform-basic,mc:contact-independence-along}
\source{wlodarczyk05-hironaka}
\lean{Hironaka.MemberRestrictTowerIso.mem_of_preimage_mem_map_diffeo,Hironaka.MemberRestrictTowerIso.pullback_diffeo_memberRestrict_membersThrough,Hironaka.MemberRestrictTowerIso.lowerValueL_succ_eq,Hironaka.MemberRestrictTowerIso.lowerValueL_succ_isTowerIso}
\leanok
Let \(n\ge1\) and \(\underline{\mathcal J}\) be well formed with \(d\ge1\).
\begin{enumerate}
\item For a biholomorphism \(\chi\colon M'\to M\), a well-formed marked ideal \(\underline{\mathcal I}\) on \(M\) with \(d\ge1\), a list \(E_0\) of sets and \(p\in M'\), the
pullback along \(\chi\) of the restricted marked ideal of \(\underline{\mathcal I}\) at the members of \(E_0\) through \(\chi(p)\) is the restricted marked ideal of
\(\chi^*\underline{\mathcal I}\) at the members of \(\chi^{-1}E_0\) through \(p\).
\item (Anchor) Let \(T_0T_1\) be a multiple test blow-up of \(\overline{\underline{\mathcal J}}\) and \(p\) a point with lower data at level \(s+1\) at the split \(T_0T_1\)
(Definition~\ref{hir:step-one-value-masked}), \(\mathcal M\) the restricted marked ideal of \(\overline{\underline{\mathcal J}}_{T_0}\) at the original members through the image
\(y_0\) of \(p\). For every open \(W\ni y_0\) on which \(\mathcal M|_W\) lies in the domain of the inductive invariant \(R\), with \(T_1|_W\) one of its canonical segments, the
lower value of \(p\) is \(R_{\mathcal M|_W}(T_1|_W)(p)\).
\item (Property (6) at the levels \(s+1\)) Let \(\chi_0\colon T_0'\to T_0\) be an isomorphism of towers over a biholomorphism \(\Phi\), \(\chi\colon T_1'\to T_1\) one over
\(\chi_0\), and \(T_0T_1\) a multiple test blow-up of \(\overline{\underline{\mathcal J}}\). If \(p'\) has lower data at level \(s+1\) for \(\Phi^*\underline{\mathcal J}\) at the split
\(T_0'T_1'\) and \(\chi(p')\) has lower data at level \(s+1\) for \(\underline{\mathcal J}\) at the split \(T_0T_1\), the two lower values agree.
\end{enumerate}
\end{lemma}
\begin{proof}
\leanok
(1) The coefficient ideal commutes with biholomorphisms (Lemma~\ref{mc:homogenized-along-etale}); the submanifold and the divisor of the restricted marked ideal are
described by the members through the point (Lemma~\ref{bcalc:member-restriction-rebase}), and preimages under a bijection distinguish sets, so the members of
\(\chi^{-1}E_0\) through \(p\) are the preimages of those of \(E_0\) through \(\chi(p)\).
(2) The lower value is \(R\) of \(\mathcal M\) restricted to the mask along \(T_1\) restricted to it; on the two open sets the values agree at the points of the cosupport over
both (Lemma~\ref{hir:step-one-nearby}, on the restricted marked ideal, well formed by Lemma~\ref{bcalc:member-restriction-domain}, the stage
\(\overline{\underline{\mathcal J}}_{T_0}\) being well formed with mark \(\ge1\), Lemmas~\ref{hir:step-one-lower-domain} and~\ref{mc:contact-independence-along}).
(3) \(\overline{(\Phi^*\underline{\mathcal J})}=\Phi^*\overline{\underline{\mathcal J}}\) (Lemma~\ref{mc:contact-stage-kept}), so its transform along \(T_0'\) is \(\chi_0^*\) of that along \(T_0\)
(Lemma~\ref{mc:tower-isomorphism-transform}) and its original members are the preimages (Theorem~\ref{mc:contact-lower-value}); by (1) its restricted marked ideal at the
image \(y_0'\) of \(p'\) is \(\chi_0^*\mathcal M\), with \(\chi_0(y_0')=y_0\). On \(W'=\chi_0^{-1}(\text{mask of }y_0)\), \(\chi_0^*\mathcal M|_{W'}\) is the pullback of \(\mathcal M\) on
the mask along the restriction \(\psi\) of \(\chi_0\) (Theorem~\ref{mc:contact-marked-canonical}), it lies in the domain (Lemma~\ref{bcalc:inv-domain-stable}), and \(T_1'|_{W'}\) is
isomorphic over \(\psi\) to the restriction of \(T_1\) to the mask (Lemma~\ref{mc:tower-iso-opens}), hence a canonical segment with corresponding values
(Lemma~\ref{hir:canonical-segments}(2), Definition~\ref{hir:centre-function}(6)). The concatenations are isomorphic (Lemma~\ref{mc:tower-isomorphism-append}), so
\(T_0'T_1'\) is a multiple test blow-up of \(\overline{(\Phi^*\underline{\mathcal J})}\); conclude with (2) for \(\Phi^*\underline{\mathcal J}\) on \(W'\).
\end{proof}

\begin{lemma}[Step 1ab's lower values under restriction and change of representative]
\label{bcalc:member-restrict-restrict-repr}
\uses{bcalc:member-restrict-tower-iso,bcalc:member-restriction-domain,bcalc:member-restriction-rebase,hir:step-one-value-masked,hir:step-one-value-masked-basic,hir:step-one-nearby,hir:step-one-lower-domain,mc:contact-lower-value,mc:contact-stage-kept,mc:contact-representative,mc:coefficient-representatives-tower,mc:homogenized-along-basic,hir:marked-equivalence-basic,hir:restricted-calculus,mc:restriction-operations-along,mc:tower-isomorphism-append,mc:tower-isomorphism-transform,mc:tower-iso-opens,mc:contact-marked-canonical,bcalc:inv-domain-stable,bcalc:restrict-nested,bcalc:centre-local,bcalc:tower-restrict,hir:open-restriction-tower,hir:open-restriction-marked,hir:tower-split-iso,hir:wellformed-tower,hir:geo-transform-basic,hir:canonical-segments,hir:invariant-driver,hir:centre-function,mc:contact-independence-along}
\source{wlodarczyk05-hironaka}
\lean{Hironaka.MemberRestrictTowerIso.inv_pullback_isTowerIso,Hironaka.MemberRestrictTowerIso.lowerValueL_succ_restrict,Hironaka.MemberRestrictTowerIso.restrictEquiv_memberRestrict_membersThrough,Hironaka.MemberRestrictTowerIso.lowerValueL_succ_representative}
\leanok
Let \(n\ge1\) and \(\underline{\mathcal J}\) be well formed with \(d\ge1\), and \(T_0T_1\) a multiple test blow-up of \(\overline{\underline{\mathcal J}}\).
\begin{enumerate}
\item (Transport) If \(\mathcal A\) lies in the domain of the inductive invariant \(R\) with \(S\) one of its canonical segments, and \(\xi\colon S'\to S\) is an isomorphism of
towers over a biholomorphism \(\psi\), then \(\psi^*\mathcal A\) lies in the domain, \(S'\) is one of its canonical segments, and at every \(p\) with \(\xi(p)\) in the
cosupport along \(S\), \(p\) lies in the cosupport along \(S'\) and \(R_{\psi^*\mathcal A}(S')(p)=R_{\mathcal A}(S)(\xi(p))\).
\item (Property (5) at the levels \(s+1\)) Let \(U\) be open in the base, \(\chi_0\colon T_0'\to T_0|_U\) an isomorphism of towers over the identity of \(U\) and
\(\chi\colon T_1'\to T_1|_{T_0^{-1}U}\) one over \(\chi_0\). If \(p'\) has lower data at level \(s+1\) for \(\underline{\mathcal J}|_U\) at the split \(T_0'T_1'\) and \(\chi(p')\)
has lower data at level \(s+1\) for \(\underline{\mathcal J}\) at the split \(T_0T_1\), the two lower values agree.
\item For \(\underline{\mathcal J}\simeq\underline{\mathcal J}'\) restricted-equivalent with \(\underline{\mathcal J}'\) coherent, every tower \(T_0\) and every point \(y_0\) of
its last stage, the restricted marked ideals of \(\overline{\underline{\mathcal J}}_{T_0}\) and \(\overline{\underline{\mathcal J}'}_{T_0}\) at the original members through
\(y_0\) are restricted-equivalent.
\item (Property (7) at the levels \(s+1\)) For \(\underline{\mathcal J}\simeq\underline{\mathcal J}'\) as in (3), if \(p\) has lower data at level \(s+1\) at the split
\(T_0T_1\) both for \(\underline{\mathcal J}\) and for \(\underline{\mathcal J}'\), the two lower values agree.
\end{enumerate}
\end{lemma}
\begin{proof}
\leanok
(1) Domain (Proposition~\ref{bcalc:inv-domain-stable}), canonical segments (Lemma~\ref{hir:canonical-segments}(2)) and values (Definition~\ref{hir:centre-function}(6)) pass
along isomorphisms of towers; cosupports correspond (Lemma~\ref{mc:tower-isomorphism-transform}), \(S\) being a multiple test blow-up (Theorem~\ref{hir:invariant-driver}).
(2) \(\overline{(\underline{\mathcal J}|_U)}=\overline{\underline{\mathcal J}}|_U\) (Theorem~\ref{mc:contact-stage-kept}); multiple test blow-ups and transforms restrict to
open sets (Lemma~\ref{hir:open-restriction-tower}) and pass along the isomorphisms (Lemmas~\ref{mc:tower-isomorphism-append} and~\ref{mc:tower-isomorphism-transform}), and
the original members are the traces (Theorem~\ref{mc:contact-lower-value}). By Lemma~\ref{bcalc:member-restrict-tower-iso}(1) and
Lemma~\ref{bcalc:member-restriction-rebase}, the restricted marked ideal of \(\underline{\mathcal J}|_U\) at the image of \(p'\) is \(\chi_0^*(\mathcal M|_{T_0^{-1}U})\), with
\(\mathcal M\) that of \(\underline{\mathcal J}\) at the image \(y_0\) of \(\chi(p')\). Shrink the mask \(A\) of \(y_0\) to \(W_1=A\cap T_0^{-1}U\)
(Lemma~\ref{bcalc:restrict-nested}, values by Lemma~\ref{hir:step-one-nearby}), view \(W_1\) inside \(T_0^{-1}U\) (Lemmas~\ref{bcalc:centre-local}
and~\ref{bcalc:restrict-nested}), transport by (1) along the restriction of \(\chi_0\) to \(W'=\chi_0^{-1}W_1\) (Lemma~\ref{mc:tower-iso-opens},
Theorem~\ref{mc:contact-marked-canonical}), and conclude with the anchor (Lemma~\ref{bcalc:member-restrict-tower-iso}(2)) for \(\underline{\mathcal J}|_U\) on \(W'\).
(3) \(\overline{\underline{\mathcal J}}\simeq\overline{\underline{\mathcal J}'}\) (Theorem~\ref{mc:contact-representative}), so the coefficient ideals of their transforms along
\(T_0\) are restricted-equivalent (Lemma~\ref{mc:coefficient-representatives-tower}), the divisors of the transforms agree (Lemma~\ref{hir:marked-equivalence-basic}) and
so do the original members (Theorem~\ref{mc:contact-lower-value}); the submanifolds and divisors of the two restricted marked ideals are described by the members through
\(y_0\) (Lemma~\ref{bcalc:member-restriction-rebase}), and the restricted inclusions pass to the smaller submanifold (Lemma~\ref{hir:restricted-calculus}).
(4) \(T_0T_1\) is a multiple test blow-up of \(\overline{\underline{\mathcal J}'}\) as well (Lemma~\ref{hir:marked-equivalence-basic}). On the intersection \(W\) of the two
masks both restricted marked ideals lie in the domain with \(T_1|_W\) a canonical segment (Lemma~\ref{bcalc:restrict-nested}) and \(p\) in the cosupport
(Lemma~\ref{hir:step-one-nearby}); their restrictions to \(W\) are restricted-equivalent by (3) (Theorem~\ref{mc:contact-representative}), so the values agree
(Definition~\ref{hir:centre-function}(7)); conclude with the anchor on both sides.
\end{proof}

\paragraph{The pushforward of \(\mathcal O(-KD)\).} For the blowup \(\pi\) along a smooth centre \(C\) with exceptional divisor \(D\), a function whose
pullback vanishes to order \(K\) along \(D\) vanishes to order \(K\) along \(C\): \(\pi_*\mathcal O(-KD)=\mathcal I_C^K\). The statement below is local on
the blowup: the hypothesis is at a single point of \(D\). It is the blowup half of the step ``the nonmonomial part stays nonmonomial'' in the
proof of Lemma~\ref{hir:nonmonomial-transform}.

\begin{theorem}[The pushforward of \(\mathcal O(-KD)\)]
\label{bcalc:exceptional-pushforward}
\uses{bcalc:monomial-lift,bcalc:exceptional-divisor,bcalc:strict-ideal,hir:polarization,hir:order-power,hir:point-power-pullback,mc:normal-vanishing,hir:normal-power-pullback,hir:hadamard-bridge,mc:monomial-local-equation}
\source{wlodarczyk05-hironaka,bierstone-milman08-functoriality}
\lean{Hironaka.ExceptionalIdealPushforward.eq_zero_of_symmetric_of_apply_diag_on_open,Hironaka.ExceptionalIdealPushforward.eventually_mem_idealPow_pointIdeal,Hironaka.ExceptionalIdealPushforward.mem_idealPow_of_mem_at,Hironaka.ExceptionalIdealPushforward.liftPoint,Hironaka.ExceptionalIdealPushforward.liftDir,Hironaka.ExceptionalIdealPushforward.liftPoint_apply,Hironaka.ExceptionalIdealPushforward.liftDir_apply,Hironaka.ExceptionalIdealPushforward.monomial_liftPoint_add,Hironaka.ExceptionalIdealPushforward.mem_idealPow_vanishingIdeal_of_pullback_mem}
\leanok
Let \(\pi\colon Y\to X\) be the blowup along \(C\), \(D=\pi^{-1}(C)\), \(z\in D\) and \(f\) holomorphic at \(\pi(z)\). If \(f\circ\pi\in\mathcal O(-D)^K\) at \(z\),
then \(f\in\mathcal I_C^K\) at \(\pi(z)\).
\end{theorem}
\begin{proof}
\leanok
Let \(x=\pi(z)\). Take flattening coordinates \(\Phi\) at \(x\), centred at \(0\), with \(C=\mathbf C^k\times0\), put \(F=f\circ\Phi^{-1}\), and take a monomial
chart \(j\ge k\), \(W\) around \(z\) with its lift \(\ell\), \(\pi\circ\ell=\Phi^{-1}\circ\sigma_j\) (Lemma~\ref{bcalc:monomial-lift}); \(\ell\) is read on
\(\mathbf C^n\) as in Lemma~\ref{mc:monomial-local-equation}. Since \(\Phi(x)=0\), the point \(W(z)\) has \(W_j(z)=0\) and vanishing tangential coordinates.

(1) The germs of \(\mathcal O(-D)=\pi^*\mathcal I_C\) at \(z\) vanish on \(D\) near \(z\) (Theorem~\ref{bcalc:strict-ideal}), so a germ of \(\mathcal O(-D)^K\) at \(z\)
lies in \(\mathfrak m_u^K\) at every \(u\in D\) near \(z\): induction on \(K\), a germ of \(\mathfrak a^{K+1}\) being a combination of products of a germ of \(\mathfrak a^K\)
and a germ of \(\mathfrak a\). Pulling back along \(\ell\) (Lemma~\ref{hir:point-power-pullback}), \(G=F\circ\sigma_j\in\mathfrak m_\lambda^K\) at every \(\lambda\) near
\(W(z)\) with \(\lambda_j=0\), so \(D^iG(\lambda)=0\) for \(i<K\) (Lemma~\ref{hir:order-power}).

(2) For \(w\in\mathbf C^k\times0\) and a direction \(h\) with \(h_j\neq0\), let \(\lambda(w,h)\) have the tangential coordinates of \(w\), the coordinates
\(h_i/h_j\) for the normal \(i\neq j\), and \(\lambda_j=0\), and let \(v(h)\) have the tangential coordinates of \(h\), \(v_j=h_j\), and \(0\) otherwise. Then
\(\sigma_j(\lambda+tv)=w+th\) for every \(t\). With \(h_0=W(z)+e_j\) one has \(\lambda(0,h_0)=W(z)\), so for \(w\) near \(0\) and \(h\) in an open neighbourhood
\(V\) of \(h_0\), \(\lambda(w,h)\) lies where (1) applies and
\[
 D^aF(w)(h,\dots,h)=\tfrac{d^a}{dt^a}\Big|_{t=0}F(w+th)=\tfrac{d^a}{dt^a}\Big|_{t=0}G(\lambda+tv)=D^aG(\lambda)(v,\dots,v)=0\qquad(a<K).
\]

(3) The diagonal \(h\mapsto D^aF(w)(h,\dots,h)\) is a polynomial vanishing on the open set \(V\), hence everywhere (identity principle), and the
symmetric form \(D^aF(w)\) is zero by polarization (Lemma~\ref{hir:polarization}). So every derivative of order \(<K\) of \(F\) vanishes on
\(\mathbf C^k\times0\) near \(0\); \(F\) lies in the \(K\)-th power of the ideal of the normal coordinates near \(0\) (Lemma~\ref{mc:normal-vanishing}), and
pulling back along \(\Phi\) gives \(f\in\mathcal I_C^K\) at \(x\) (Lemma~\ref{hir:normal-power-pullback}), as in the Hadamard bridge
(Lemma~\ref{hir:hadamard-bridge}).
\end{proof}

\begin{corollary}[A controlled transform inside \(\mathcal O(-D)\) raises the order]
\label{bcalc:controlled-not-exceptional}
\uses{bcalc:exceptional-pushforward,bcalc:controlled-transform-def,bcalc:controlled-transform,hir:bridge-consumers,hir:cosupport,bcalc:isomorphism-off-centre,bcalc:strict-submanifold}
\source{wlodarczyk05-hironaka,bierstone-milman08-functoriality}
\upstreamlean{TranscendentalNakamaye.Pending.ExceptionalIdealPushforward.mem_cosupport_succ_of_controlledTransform_le,TranscendentalNakamaye.Pending.ExceptionalIdealPushforward.exists_not_mem_exceptionalDivisor}\lean{Hironaka.ExceptionalIdealPushforward.eventually_le_vanishingIdeal_of_isCoherent}\upstreamlean{TranscendentalNakamaye.Pending.ExceptionalIdealPushforward.frequently_le_vanishingIdeal_of_controlledTransform_le,TranscendentalNakamaye.Pending.ExceptionalIdealPushforward.exists_not_mem_vanishingIdeal_geoStrictTransform}
\mixedleanok
Let \(\mathfrak a\) be coherent, \(C\subset\operatorname{cosupp}(\mathfrak a,\mu)\), \(\pi\) the blowup along \(C\) and \(z\in D\). If \(\sigma^c(\mathfrak a,\mu)\subset\mathcal O(-D)\)
at \(z\), then \(\operatorname{ord}_{\pi(z)}\mathfrak a\ge\mu+1\). Hence, at every point of \(D\) over a point of \(C\) where \(\operatorname{ord}\mathfrak a=\mu\), some germ of
\(\sigma^c(\mathfrak a,\mu)\) is not in \(\mathcal O(-D)\). Likewise for a closed set \(H\) with strict transform \(H'=\overline{\pi^{-1}(H\setminus C)}\): if
\(\sigma^c(\mathfrak a,\mu)\subset\mathcal I_{H'}\) at a point \(z\in H'\), then \(\mathfrak a\subset\mathcal I_H\) at points of \(H\setminus C\) arbitrarily close to
\(\pi(z)\); so if \(\mathfrak a\not\subset\mathcal I_H\) at the points of \(H\setminus C\) near \(\pi(z)\), some germ of \(\sigma^c(\mathfrak a,\mu)\) at \(z\) is not in
\(\mathcal I_{H'}\). These are the two clauses used in the proof of Lemma~\ref{hir:nonmonomial-transform}.
\end{corollary}
\begin{proof}
\mixedleanok
By Corollary~\ref{hir:bridge-consumers}, \(\pi^*\mathfrak a=\mathcal O(-\mu D)\,\sigma^c(\mathfrak a,\mu)\). For \(a\in\mathfrak a_{\pi(z)}\), \(a\circ\pi\in\pi^*\mathfrak a\) at \(z\) is a
combination of products \(pq\) with \(p\in\mathcal O(-\mu D)\) and \(q\in\sigma^c(\mathfrak a,\mu)\subset\mathcal O(-D)\) at \(z\), so \(a\circ\pi\in\mathcal O(-D)^{\mu+1}\) at
\(z\). By Theorem~\ref{bcalc:exceptional-pushforward}, \(a\in\mathcal I_C^{\mu+1}\subset\mathfrak m_{\pi(z)}^{\mu+1}\) at \(\pi(z)\) (as \(\pi(z)\in C\)), so
\(\operatorname{ord}_{\pi(z)}\mathfrak a\ge\mu+1\) (Definition~\ref{hir:cosupport}). The second statement is the contrapositive.

For \(H'\): \(\sigma^c(\mathfrak a,\mu)\) is coherent (Corollary~\ref{hir:bridge-consumers}), so the containment in \(\mathcal I_{H'}\) at \(z\) holds at every
point near \(z\) (each of finitely many generators vanishes on \(H'\) near \(z\)). Points \(z'\in\pi^{-1}(H\setminus C)\) accumulate at \(z\in H'\). At such a
\(z'\), for \(a\in\mathfrak a_{\pi(z')}\), \(a\circ\pi\) is a combination of products \(pq\) with \(q\in\sigma^c(\mathfrak a,\mu)\subset\mathcal I_{H'}\), so it vanishes on
\(H'\) near \(z'\). The section \(s\) of \(\pi\) over \(X\setminus C\) with \(s(\pi(z'))=z'\) (Lemma~\ref{bcalc:isomorphism-off-centre}) maps the points
\(y\in H\setminus C\) near \(\pi(z')\) into \(\pi^{-1}(H\setminus C)\subset H'\) near \(z'\), so \(a(y)=(a\circ\pi)(s(y))=0\): \(a\in\mathcal I_H\) at \(\pi(z')\), and
\(\pi(z')\to\pi(z)\).
\end{proof}

\paragraph{Along a submanifold.} For a marked ideal \((M,N,E,\mathfrak a,d)\) with \(N\ne M\) only the restrictions of germs to \(N\), and of the
transformed germs to the strict transform \(N'=\overline{\pi^{-1}(N\setminus C)}\), matter (Definition~\ref{hir:geo-transform}). A germ \(g\) at \(z\)
\emph{lies in \(\mathfrak b\) on \(N'\)} if near \(z\) it agrees on \(N'\) with a holomorphic combination of germs of \(\mathfrak b\) at \(z\). The
retraction of Lemma~\ref{bcalc:chart-retraction} reduces the statements of this subsection along \(N'\) to their versions on the whole space; it is the
blowup half of the normal form \(\mathfrak a^\flat=r^*(\mathfrak a|_N)\) of Lemma~\ref{hir:along-calculus}.

\begin{lemma}[The chart retraction]
\label{bcalc:chart-retraction}
\uses{bcalc:nested-flattening,bcalc:monomial-lift,bcalc:snc-chart,bcalc:strict-submanifold,hir:order-along-flat,hir:hadamard-bridge-along,mc:monomial-local-equation}
\source{wlodarczyk05-hironaka,bierstone-milman08-functoriality}
\lean{Hironaka.ExceptionalPushforwardAlong.monomial_planeRetract,Hironaka.ExceptionalPushforwardAlong.mdifferentiableAt_flatMap,Hironaka.ExceptionalPushforwardAlong.mdifferentiableAt_flatInv,Hironaka.ExceptionalPushforwardAlong.flatMap_monomialChart,Hironaka.ExceptionalPushforwardAlong.lift_mem_geoStrictTransform,Hironaka.ExceptionalPushforwardAlong.exists_chartRetraction}
\leanok
Let \(\pi\colon Y\to X\) be the blowup along \(C\), \(N\supset C\) a closed submanifold and \(z\in N'\) with \(\pi(z)\in C\). There are a holomorphic map
\(r\) from a neighbourhood of \(\pi(z)\) to \(N\) with \(r=\mathrm{id}\) on \(N\), and a holomorphic map \(r'\) from a neighbourhood of \(z\) to \(N'\) with
\(r'(z)=z\) and \(\pi\circ r'=r\circ\pi\).
\end{lemma}
\begin{proof}
\leanok
Take nested flattening coordinates \(\Phi\) at \(\pi(z)\) (Lemma~\ref{bcalc:nested-flattening}): with \(\psi=\Phi\circ e\) and \(\varphi=\psi^{-1}\), near
\(\pi(z)\) we have \(\psi(\pi(z))=0\), \(C=\{\psi\in\mathbf C^{k_C}\times0\}\) and \(N=\{\psi\in\mathbf C^{k_N}\times0\}\). Let \(P\) be the linear retraction
of \(\mathbf C^n\) onto \(\mathbf C^{k_N}\times0\) and \(r=\varphi\circ P\circ\psi\) (Lemma~\ref{hir:order-along-flat}); it is holomorphic near \(\pi(z)\), maps into
\(N\), and fixes the points of \(N\). Let \(j\ge k_C\), \(O\), \(W\) be a monomial chart through \(z\) with lift \(\ell_j\) (Lemma~\ref{bcalc:monomial-lift}):
\(\sigma_j(W)=\psi\circ\pi\) on \(O\) and \(\ell_j(W(u))=u\). By Lemma~\ref{bcalc:snc-chart} applied to \(N\), \(z\in N'\) gives \(j<k_N\) and
\(W(z)\in\mathbf C^{k_N}\times0\). For \(j<k_N\), \(\sigma_j\circ P=P\circ\sigma_j\): \(\sigma_j\) multiplies the coordinates \(w_t\), \(t\ge k_C\), \(t\ne j\), by
\(w_j\), which \(P\) does not change. Put \(r'=\ell_j\circ P\circ W\) near \(z\) (Lemma~\ref{mc:monomial-local-equation} reads \(\ell_j\) on \(\mathbf C^n\)). Then
\(r'(z)=\ell_j(W(z))=z\), and \(\pi\circ r'=\varphi\circ\sigma_j\circ P\circ W=\varphi\circ P\circ\sigma_j\circ W=\varphi\circ P\circ\psi\circ\pi=r\circ\pi\).
Finally \(\ell_j\) maps every point \(q\in\mathbf C^{k_N}\times0\) of its domain into \(N'\): if \(q_j\ne0\), then \(\sigma_j(q)\in\mathbf C^{k_N}\times0\) and
\(\sigma_j(q)\notin\mathbf C^{k_C}\times0\), so \(\pi(\ell_j(q))\in N\setminus C\); if \(q_j=0\), \(\ell_j(q)\) is the limit of \(\ell_j(q+te_j)\), \(t\to0\),
\(t\ne0\), and \(q+te_j\in\mathbf C^{k_N}\times0\) as \(j<k_N\). As \(P(W(u))\in\mathbf C^{k_N}\times0\), \(r'\) takes values in \(N'\).
\end{proof}

\begin{theorem}[The pushforward of \(\mathcal O(-KD)\) along the strict transform of a submanifold]
\label{bcalc:exceptional-pushforward-along}
\uses{bcalc:chart-retraction,bcalc:exceptional-pushforward,bcalc:exceptional-divisor}
\source{wlodarczyk05-hironaka,bierstone-milman08-functoriality}
\lean{Hironaka.ExceptionalPushforwardAlong.comp_mem_idealPow,Hironaka.ExceptionalPushforwardAlong.comp_mem_exceptionalDivisor,Hironaka.ExceptionalPushforwardAlong.mem_idealPow_vanishingIdeal_along_of_pullback}
\leanok
Let \(\pi\) be the blowup along \(C\), \(N\supset C\) a closed submanifold, \(z\in N'\) with \(\pi(z)\in C\), and \(f\) holomorphic at \(\pi(z)\). If \(f\circ\pi\)
lies in \(\mathcal O(-D)^K\) on \(N'\) at \(z\), then \(f\) lies in \(\mathcal I_C^K\) on \(N\) at \(\pi(z)\).
\end{theorem}
\begin{proof}
\leanok
Take \(r,r'\) of Lemma~\ref{bcalc:chart-retraction} and put \(f^\flat=f\circ r\), holomorphic at \(\pi(z)=r(\pi(z))\) and equal to \(f\) on \(N\). A germ \(q\) of
\(\mathcal O(-D)=\pi^*\mathcal I_C\) at \(z\) is a combination \(\sum b_i\,(c_i\circ\pi)\) with \(c_i\in\mathcal I_C\) at \(\pi(z)\); so
\(q\circ r'=\sum(b_i\circ r')(c_i\circ r\circ\pi)\), and \(c_i\circ r\in\mathcal I_C\) because \(r\) fixes the points of \(C\subset N\). Hence \(q\circ r'\in\mathcal O(-D)\) at
\(z\), and by induction on \(K\) (using \(r'(z)=z\)) \(g\circ r'\in\mathcal O(-D)^K\) at \(z\) for \(g\in\mathcal O(-D)^K\) at \(z\). If \(f\circ\pi=\sum a_ig_i\) on
\(N'\) near \(z\) with \(g_i\in\mathcal O(-D)^K\), then near \(z\), since \(r'\) takes values in \(N'\),
\(f^\flat\circ\pi=f\circ\pi\circ r'=\sum(a_i\circ r')(g_i\circ r')\in\mathcal O(-D)^K\) at \(z\). By Theorem~\ref{bcalc:exceptional-pushforward},
\(f^\flat\in\mathcal I_C^K\) at \(\pi(z)\), and \(f=f^\flat\) on \(N\) near \(\pi(z)\).
\end{proof}

\begin{corollary}[The transform along \(N'\) at the points of \(D\) and of strict transforms]
\label{bcalc:controlled-not-exceptional-along}
\uses{bcalc:exceptional-pushforward-along,bcalc:controlled-not-exceptional,hir:geo-transform,hir:geo-transform-wellformed,hir:hadamard-bridge-along,ana:restricted-colon,bcalc:isomorphism-off-centre,bcalc:exceptional-divisor,hir:cosupport}
\source{wlodarczyk05-hironaka,bierstone-milman08-functoriality}
\lean{Hironaka.ExceptionalPushforwardAlong.isInSpanOnAt_mul_idealPow_succ,Hironaka.ExceptionalPushforwardAlong.restrictLe_pullback_geoTransform,Hironaka.ExceptionalPushforwardAlong.mem_cosupport_succ_of_geoTransform_le,Hironaka.ExceptionalPushforwardAlong.exists_not_isInSpanOnAt_exceptionalDivisor,Hironaka.ExceptionalPushforwardAlong.frequently_le_vanishingIdeal_of_geoTransform_le,Hironaka.ExceptionalPushforwardAlong.exists_not_mem_vanishingIdeal_of_geoTransform}
\leanok
Let \(\underline{\mathcal I}=(M,N,E,\mathfrak a,d)\) with \(N\) a closed submanifold and \(\mathfrak a\) coherent, \(C\subset\operatorname{cosupp}\underline{\mathcal I}\),
\(\pi\) the blowup along \(C\), and \(\mathfrak a'\) the transformed ideal (Definition~\ref{hir:geo-transform}). Then:
(0) every germ of \(\pi^*\mathfrak a\) at a point of \(N'\) lies in \(\mathcal O(-dD)\,\mathfrak a'\) on \(N'\);
(1) if at \(z\in N'\) over \(C\) every germ of \(\mathfrak a'\) lies in \(\mathcal O(-D)\) on \(N'\), then \(\operatorname{ord}_{\pi(z)}(\mathfrak a|_N)\ge d+1\); hence where
\(\operatorname{ord}(\mathfrak a|_N)=d\), some germ of \(\mathfrak a'\) at every point of \(D\cap N'\) over it is not in \(\mathcal O(-D)\) on \(N'\);
(2) for a closed set \(A\subset N\) (for instance \(A=H\cap N\) for a member \(H\) of \(E\)) and \(z\in A'=\overline{\pi^{-1}(A\setminus C)}\): if
\(\mathfrak a'\subset\mathcal I_{A'}\) at \(z\), then \(\mathfrak a\subset\mathcal I_A\) at points of \(A\setminus C\) arbitrarily close to \(\pi(z)\); so if \(\mathfrak a\not\subset\mathcal I_A\) at
the points of \(A\setminus C\) near \(\pi(z)\), some germ of \(\mathfrak a'\) at \(z\) does not vanish on \(A'\) near \(z\).
These are the two clauses of Corollary~\ref{bcalc:controlled-not-exceptional} along \(N'\), as used in Lemma~\ref{hir:nonmonomial-transform} for \(N\ne M\).
\end{corollary}
\begin{proof}
\leanok
(0) By the Hadamard bridge along \(N\) (Lemma~\ref{hir:hadamard-bridge-along}), \(\mathfrak a\subset\mathcal I_C^d\) on \(N\); pulling back
(Lemma~\ref{hir:geo-transform-wellformed}) and restricting to \(N'\subset\pi^{-1}(N)\), \(\pi^*\mathfrak a\subset\pi^*\mathcal I_C^d=\mathcal O(-dD)\) on \(N'\)
(Lemma~\ref{bcalc:exceptional-divisor}). As \(\mathcal O(-dD)\) is invertible along \(N'\) (Lemma~\ref{hir:geo-transform-wellformed}) and \(\mathfrak a'\) is the colon of
\(\pi^*\mathfrak a\) by it along \(N'\), Lemma~\ref{ana:restricted-colon} gives (0).
(1) Let \(a\in\mathfrak a\) at \(\pi(z)\). By (0), \(a\circ\pi\) is, on \(N'\), a combination of products \(pq\) with \(p\in\mathcal O(-dD)\) and \(q\in\mathfrak a'\) at \(z\), and
each \(q\) is on \(N'\) a combination of germs of \(\mathcal O(-D)\); so \(a\circ\pi\) lies in \(\mathcal O(-D)^{d+1}\) on \(N'\). Theorem~\ref{bcalc:exceptional-pushforward-along}
gives \(a\in\mathcal I_C^{d+1}\subset\mathfrak m_{\pi(z)}^{d+1}\) on \(N\) (as \(\pi(z)\in C\)), that is \(\operatorname{ord}_{\pi(z)}(a|_N)\ge d+1\) (Definition~\ref{hir:cosupport}).
(2) The ideal \(\mathfrak a'\) is coherent (Lemma~\ref{hir:geo-transform-wellformed}), so the containment in \(\mathcal I_{A'}\) at \(z\) holds at every point near \(z\);
points \(z'\in\pi^{-1}(A\setminus C)\subset N'\) accumulate at \(z\). At such \(z'\), by (0), \(a\circ\pi\) agrees on \(N'\) with a combination of germs of
\(\mathcal O(-dD)\,\mathfrak a'\), which vanish on \(A'\subset N'\); so \(a\circ\pi\) vanishes on \(A'\) near \(z'\). The section of \(\pi\) over \(X\setminus C\)
(Lemma~\ref{bcalc:isomorphism-off-centre}) maps the points of \(A\setminus C\) near \(\pi(z')\) into \(\pi^{-1}(A\setminus C)\subset A'\), so \(a\) vanishes on \(A\) near
\(\pi(z')\), and \(\pi(z')\to\pi(z)\).
\end{proof}

\begin{lemma}[The transform of an ideal with an exceptional factor, along \(N'\)]
\label{bcalc:transform-exceptional-factor-along}
\uses{hir:marked-operations-along-step,ana:restricted-calculus,hir:geo-transform,hir:geo-transform-wellformed,bcalc:exceptional-divisor}
\source{wlodarczyk05-hironaka,bierstone-milman08-functoriality}
\upstreamlean{TranscendentalNakamaye.Pending.ExceptionalPushforwardAlong.restrictLe_geoTransform_of_pullback}
\upstreamleanok
Let \(\underline{\mathcal I}=(M,N,E,\mathfrak a,d)\) with \(N\) a closed submanifold, \(\mathfrak a\) coherent, \(C\) an admissible centre and \(\mathfrak a'\) the
transformed ideal. If \(\pi^*\mathfrak a\) and \(\mathcal O(-eD)\,\mathfrak b\) have the same restriction to \(N'\), for a coherent \(\mathfrak b\) and \(e\ge d\),
then \(\mathfrak a'\) and \(\mathcal O(-(e-d)D)\,\mathfrak b\) have the same restriction to \(N'\). With Lemma~\ref{bcalc:member-pullback}
(\(\pi^*\mathcal I_H=\mathcal O(-D)\,\mathcal I_{H'}\) where \(C\subset H\), \(=\mathcal I_{H'}\) elsewhere) this computes the transform of a monomial marked ideal along
\(N'\) (Lemma~\ref{hir:monomial-blowup}).
\end{lemma}
\begin{proof}
\upstreamleanok
By Lemma~\ref{hir:marked-operations-along-step}, \(\pi^*\mathfrak a\) and \(\mathcal O(-dD)\,\mathfrak a'\) have the same restriction to \(N'\). As
\(\mathcal O(-eD)\,\mathfrak b=(\mathcal O(-(e-d)D)\,\mathfrak b)\,\mathcal O(-dD)\) (powers add, products are associative and commutative for coherent ideals),
\(\mathfrak a'\,\mathcal O(-dD)\) and \((\mathcal O(-(e-d)D)\,\mathfrak b)\,\mathcal O(-dD)\) have the same restriction to \(N'\). The ideal \(\mathcal O(-dD)\) is invertible
along \(N'\) (Lemma~\ref{hir:geo-transform-wellformed}), so it cancels (Lemma~\ref{ana:restricted-calculus}); both ideals are coherent
(Lemma~\ref{hir:geo-transform-wellformed} for \(\mathfrak a'\)).
\end{proof}

\begin{lemma}[Strict transforms of the traces of the members]
\label{bcalc:trace-strict-transform}
\uses{bcalc:snc-simultaneous,bcalc:snc-chart,bcalc:monomial-lift,bcalc:strict-submanifold,hir:snc}
\source{bierstone-milman08-functoriality}
\lean{Hironaka.ExceptionalPushforwardAlong.coordSubspace_union,Hironaka.ExceptionalPushforwardAlong.geoStrictTransform_inter_eq}
\leanok
Let \(\pi\) be the blowup along \(C\), \(N\supset C\) closed, \(E\) a simple normal crossings divisor transverse to \(N\) with which \(C\) has only normal
crossings, and \(H\) a member of \(E\). Then \(\overline{\pi^{-1}((H\cap N)\setminus C)}=H'\cap N'\). In particular, in Corollary~\ref{bcalc:controlled-not-exceptional-along}~(2)
with \(A=H\cap N\), \(A'\) is the trace on \(N'\) of the member \(H'\) of the transformed divisor.
\end{lemma}
\begin{proof}
\leanok
\(\subset\): strict transforms are monotone. \(\supset\): let \(z\in H'\cap N'\); then \(\pi(z)\in H\cap N\), as \(H'\subset\pi^{-1}(H)\) and \(N'\subset\pi^{-1}(N)\). If
\(\pi(z)\notin C\), \(z\in\pi^{-1}((H\cap N)\setminus C)\). If \(\pi(z)\in C\), take coordinates at \(\pi(z)\) adapted to \(C\subset N\) and \(E\)
(Lemma~\ref{bcalc:snc-simultaneous}): \(C=\mathbf C^{k_C}\times0\), \(N=L_{T_N}\), \(T_N=\{t\ge k_N\}\), \(H=L_{\{s\}}\), so \(H\cap N=L_{\{s\}\cup T_N}\), and a
monomial chart \(j\) through \(z\) (Lemma~\ref{bcalc:monomial-lift}). By Lemma~\ref{bcalc:snc-chart}, \(z\in H'\) gives \(j\ne s\) and \(W_s(z)=0\), \(z\in N'\)
gives \(j\notin T_N\) and \(W(z)\in L_{T_N}\); so \(j\notin\{s\}\cup T_N\) and \(W(z)\in L_{\{s\}\cup T_N}\), that is \(z\in(H\cap N)'\) by the same lemma.
\end{proof}

\paragraph{The trace of a member on a submanifold.} In the cases of the algorithm along \(N\ne M\) a member \(H\) of \(E\) enters through its trace
\(H\cap N\): the exponents along \(N\) are defined with \(\mathcal I_{H\cap N}\) (Lemma~\ref{hir:divisor-exponent}), while the transform and the monomial part are
written with the ideals \(\mathcal I_H\) restricted to \(N\). For \(E\) transverse to \(N\) the two agree on \(N\), also for the exceptional divisor on the strict
transform; a retraction onto \(N\) preserving the members compares them.

\begin{lemma}[A retraction onto \(N\) preserving the members]
\label{bcalc:member-retraction}
\uses{bcalc:snc-simultaneous,hir:order-along-flat,hir:hadamard-bridge-along,bcalc:chart-retraction,hir:snc}
\source{wlodarczyk05-hironaka,bierstone-milman08-functoriality}
\lean{Hironaka.MemberTraceIdealAlong.exists_memberRetraction}
\leanok
Let \(E\) be a simple normal crossings divisor transverse to \(N\) (Definition~\ref{hir:snc}) and \(x\in N\). There is a map \(\rho\), holomorphic on a
neighbourhood of \(x\), with \(\rho(x)=x\), \(\rho(y)\in N\) and \(\rho(y)=y\) for \(y\in N\) near \(x\), and \(\rho(y)\in H\) for every member \(H\) of \(E\) and
every \(y\in H\) near \(x\).
\end{lemma}
\begin{proof}
\leanok
If \(n=0\), the coordinates of the transversality at \(x\) make \(N\) a coordinate subspace of \(\mathbf C^0\), so \(N\) is a neighbourhood of \(x\), and
\(\rho=\mathrm{id}\). If \(n\ge1\), take coordinates \(\psi=\Phi\circ e\) at \(x\) adapted to \(N\subset N\) and \(E\) (Lemma~\ref{bcalc:snc-simultaneous}, with \(C=N\)):
near \(x\), \(N=\{\psi\in\mathbf C^{k}\times0\}\) and each member \(H_i\) through \(x\) is \(\{\psi_{s_i}=0\}\) with \(s_i<k\). Let \(P\) be the linear retraction
of \(\mathbf C^n\) onto \(\mathbf C^k\times0\), which does not change the coordinates \(w_t\), \(t<k\) (Lemma~\ref{hir:hadamard-bridge-along}), and
\(\rho=\psi^{-1}\circ P\circ\psi\), defined and holomorphic near \(x\) since \(P(\psi(x))=\psi(x)\) (Lemma~\ref{hir:order-along-flat}; the differentiability
of \(\psi\) and \(\psi^{-1}\) as in Lemma~\ref{bcalc:chart-retraction}). It maps into \(N\), fixes \(N\), and for \(y\in H_i\) near \(x\),
\(\psi_{s_i}(\rho(y))=(P\psi(y))_{s_i}=\psi_{s_i}(y)=0\), so \(\rho(y)\in H_i\). A member not through \(x\) is closed, so no point near \(x\) lies on it.
\end{proof}

\begin{lemma}[The ideal of the trace of a member]
\label{bcalc:member-trace-ideal}
\uses{bcalc:member-retraction,ana:ideal-sheaves,hir:snc,hir:wellformed-tower}
\source{wlodarczyk05-hironaka,bierstone-milman08-functoriality}
\upstreamlean{TranscendentalNakamaye.Pending.MemberTraceIdealAlong.comp_mem_idealPow_of_comp_mem,TranscendentalNakamaye.Pending.MemberTraceIdealAlong.mem_idealPow_of_forall_mem}\lean{Hironaka.MemberTraceIdealAlong.comp_mem_vanishingIdeal,Hironaka.MemberTraceIdealAlong.isInSpanOnAt_of_retraction,Hironaka.MemberTraceIdealAlong.mem_vanishingIdeal_of_isInSpanOnAt,Hironaka.MemberTraceIdealAlong.vanishingIdeal_mem_inter}\upstreamlean{TranscendentalNakamaye.Pending.MemberTraceIdealAlong.isInSpanOnAt_idealPow_vanishingIdeal_iff}\lean{Hironaka.MemberTraceIdealAlong.isInSpanOnAt_vanishingIdeal_iff,Hironaka.MemberTraceIdealAlong.isInSpanOnAt_vanishingIdeal_iff_mem}\upstreamlean{TranscendentalNakamaye.Pending.MemberTraceIdealAlong.isInSpanOnAt_idealPow_vanishingIdeal_iff_of_isWellFormed}\lean{Hironaka.MemberTraceIdealAlong.isInSpanOnAt_vanishingIdeal_iff_mem_of_isWellFormed}
\mixedleanok
Let \(E\) be a simple normal crossings divisor transverse to \(N\), \(H\) a member of \(E\), \(x\in N\) and \(K\ge0\). Then \(\mathcal I_H^K\) and
\(\mathcal I_{H\cap N}^K\) have the same restriction to \(N\) at \(x\): a function agrees on \(N\) near \(x\) with a holomorphic combination of germs of the one
iff it does so with germs of the other. In particular a germ \(g\) holomorphic at \(x\) lies in \(\mathcal I_H\) on \(N\) iff it vanishes on \(H\cap N\) near \(x\):
\(\mathcal I_{H\cap N}=\mathcal I_H+\mathcal I_N\) at the points of \(N\). This holds at every stage of a multiple test blow-up of a well-formed marked ideal, whose
transforms are well formed (Lemma~\ref{hir:wellformed-tower}).
\end{lemma}
\begin{proof}
\mixedleanok
\(\mathcal I_H\subset\mathcal I_{H\cap N}\), and powers are monotone at \(x\) (by induction on \(K\), as membership in \(\mathfrak l^{K+1}\) at \(x\) only involves
germs at \(x\) of \(\mathfrak l^K\) and \(\mathfrak l\)); this gives one direction. Conversely take \(\rho\) of Lemma~\ref{bcalc:member-retraction}. For
\(q\in\mathcal I_{H\cap N}\) at \(x\), \(q\circ\rho\) is holomorphic at \(x=\rho(x)\) and vanishes on \(H\) near \(x\), since \(\rho\) maps the points of \(H\) near \(x\)
into \(H\cap N\) and tends to \(x\); so \(q\circ\rho\in\mathcal I_H\) at \(x\). By induction on \(K\) (writing a germ of \(\mathfrak l^{K+1}\) as \(\sum a_i\,p_iq_i\) and
composing each factor with \(\rho\)), \(s\circ\rho\in\mathcal I_H^K\) at \(x\) for \(s\in\mathcal I_{H\cap N}^K\) at \(x\). If \(g=\sum a_is_i\) on \(N\) near \(x\) with
\(s_i\in\mathcal I_{H\cap N}^K\), then, \(\rho\) being the identity on \(N\), \(g=\sum a_i\,(s_i\circ\rho)\) on \(N\) near \(x\). For the last sentence: if \(g\) is
holomorphic and \(g=\sum a_is_i\) on \(N\) with \(s_i\) vanishing on \(H\cap N\), then \(g\) vanishes on \(H\cap N\); conversely \(g\) itself is a generator of
\(\mathcal I_{H\cap N}\) at \(x\).
\end{proof}

\begin{corollary}[The exceptional divisor on the strict transform]
\label{bcalc:exceptional-trace-ideal}
\uses{bcalc:member-trace-ideal,bcalc:snc-transform,bcalc:strict-ideal,bcalc:exceptional-divisor,hir:nonmonomial-step,hir:snc-basic}
\source{wlodarczyk05-hironaka,bierstone-milman08-functoriality}
\lean{Hironaka.MemberTraceIdealAlong.exceptionalDivisor_mem_iff_vanishingIdeal}\upstreamlean{TranscendentalNakamaye.Pending.MemberTraceIdealAlong.isInSpanOnAt_idealPow_exceptionalDivisor_iff}\lean{Hironaka.MemberTraceIdealAlong.isInSpanOnAt_exceptionalDivisor_iff_mem}\upstreamlean{TranscendentalNakamaye.Pending.MemberTraceIdealAlong.isInSpanOnAt_exceptionalDivisor_of_mem_vanishingIdeal}
\mixedleanok
Let \(\pi\) be the blowup along \(C\), \(N\supset C\) a closed submanifold, \(D=\pi^{-1}(C)\) and \(z\in N'\) with \(\pi(z)\in C\). Then \(\mathcal O(-D)=\mathcal I_D\) at
\(z\), and \(\mathcal O(-KD)\) and \(\mathcal I_{N'\cap D}^K\) have the same restriction to \(N'\) at \(z\). In particular a germ vanishing on \(N'\cap D\) near \(z\)
lies in \(\mathcal O(-D)\) on \(N'\): \(\mathcal I_{D\cap N'}\subset\mathcal O(-D)+\mathcal I_{N'}\) at \(z\), and a holomorphic germ lies in \(\mathcal O(-D)\) on \(N'\) iff
it vanishes on \(N'\cap D\) near \(z\).
\end{corollary}
\begin{proof}
\mixedleanok
\(\mathcal O(-D)=\pi^*\mathcal I_C\subset\mathcal I_D\) (Lemma~\ref{bcalc:strict-ideal}); conversely \(\mathcal I_D\subset\mathcal O(-D)\) at the points of \(D\)
(Lemma~\ref{hir:nonmonomial-step}, its first step, with \(E=\emptyset\), with which every closed submanifold has normal crossings, Lemma~\ref{hir:snc-basic}). So
\(\mathcal O(-D)^K=\mathcal I_D^K\) at \(z\) by monotonicity of powers at a point. The one-member divisor \((D)\) is simple normal crossings and transverse to \(N'\)
(Theorem~\ref{bcalc:snc-transform} with \(E=\emptyset\)), and Lemma~\ref{bcalc:member-trace-ideal} applies to it at \(z\in N'\).
\end{proof}

\paragraph{Normal crossings divisors.} Let \(E=(H_1,\dots,H_s)\) be an ordered simple normal crossings divisor on \(X\)
(Definition~\ref{hir:snc}) and \(N\supset C\) a closed submanifold. The \emph{transformed divisor} is
\(E'=(H_1',\dots,H_s',\pi^{-1}(C))\) with \(H'=\overline{\pi^{-1}(H\setminus C)}\), as in the transform of a marked ideal
(Definition~\ref{hir:geo-transform}). Coordinates \(\Phi\) at \(x\in C\) are \emph{adapted} to \(C\subset N\) and \(E\) if
\(C=\mathbf C^{k_C}\times0\), \(N=\mathbf C^{k_N}\times0\) and each member of \(E\) through \(x\) is a coordinate hyperplane \(\{w_{s_i}=0\}\) with
\(s_i<k_N\), the indices \(s_i\) being distinct. This is the simultaneous normal crossings of \(N\), \(C\) and \(E\) at \(x\) (for
\(E=\emptyset\) it is Lemma~\ref{bcalc:nested-flattening}; in general it follows from the normal crossings of \(C\) with \(E\) and the
transversality of \(E\) to \(N\), Lemma~\ref{bcalc:snc-simultaneous}).

\begin{lemma}[Strict transforms of coordinate subspaces]
\label{bcalc:snc-chart}
\uses{bcalc:monomial-lift,bcalc:strict-submanifold,ana:local-holomorphic-inverse,hir:snc}
\source{bierstone-milman08-functoriality,wlodarczyk05-hironaka}
\lean{Hironaka.BlowupSNCTransform.mem_geoStrictTransform_iff_of_chart,Hironaka.BlowupSNCTransform.monomial_mem_coordSubspace,Hironaka.BlowupSNCTransform.monomial_mem_coordSubspace_iff,Hironaka.BlowupSNCTransform.isClosed_coordSubspace,Hironaka.BlowupSNCTransform.mem_coordSubspace_singleton,Hironaka.BlowupSNCTransform.sub_mem_coordSubspace_iff,Hironaka.BlowupSNCTransform.injective_fderiv_of_monomialInverse,Hironaka.BlowupSNCTransform.exists_recentred_coords,Hironaka.BlowupSNCTransform.isCoordinateFamilyAt_of_preimage}
\leanok
Let \(\Phi\) flatten \(C\) onto \(\mathbf C^{k}\times0\) and a set \(A\) onto a coordinate subspace \(L_T=\{w:w_t=0,\ t\in T\}\), and let \(j\ge k\),
\(O\), \(W\) be a monomial chart of the blowup (Lemma~\ref{bcalc:monomial-lift}). For \(u\in O\),
\[
 u\in\overline{\pi^{-1}(A\setminus C)}\iff j\notin T\ \text{and}\ W(u)\in L_T.
\]
Moreover \(W-W(z)\), read in a chart, is a coordinate system centred at any \(z\in O\); and coordinates adapted to a family of sets
at \(g(z)\) pull back to coordinates adapted to the preimages at \(z\) along a holomorphic map \(g\) with invertible differential at \(z\).
\end{lemma}
\begin{proof}
\leanok
On \(O\), \(\pi(u)\in C\iff W_j(u)=0\) and \(\pi(u)\in A\iff\sigma_j(W(u))\in L_T\), where \(\sigma_j(w)_t=w_tw_j\) for \(t\ge k\), \(t\ne j\), and
\(\sigma_j(w)_t=w_t\) otherwise. Hence \(\sigma_j\) maps \(L_T\) into itself, and for \(w_j\ne0\): \(\sigma_j(w)\in L_T\iff j\notin T\) and \(w\in L_T\).
So \(\pi^{-1}(A\setminus C)\cap O=\{u:W(u)\in L_T,\ W_j(u)\ne0\}\) if \(j\notin T\), and \(=\emptyset\) if \(j\in T\). In the first case the closure in
\(O\) lies in the closed set \(\{W\in L_T\}\), and contains it: if \(W(u)\in L_T\) and \(W_j(u)=0\), the points \(\ell_j(W(u)+te_j)\), \(t\ne0\) small, lie in
\(\pi^{-1}(A\setminus C)\) (as \(j\notin T\)) and tend to \(u\). The differential of \(W\) is injective (Lemma~\ref{bcalc:monomial-lift}), so
\(W-W(z)\) is a local biholomorphism at \(z\) (Lemma~\ref{ana:local-holomorphic-inverse}). For the last claim compose the adapted coordinates
at \(g(z)\) with a local holomorphic inverse of \(g\) read in charts, as in Lemma~\ref{ana:regular-transport}.
\end{proof}

\begin{lemma}[Simultaneous normal crossings coordinates]
\label{bcalc:snc-simultaneous}
\uses{hir:snc,hir:snc-basic,ana:local-holomorphic-inverse,ana:regular-transport}
\source{bierstone-milman08-functoriality,wlodarczyk05-hironaka}
\lean{Hironaka.SimultaneousNormalCrossings.isAdaptedAt_of_hasNormalCrossings,Hironaka.SimultaneousNormalCrossings.exists_model_coords,Hironaka.SimultaneousNormalCrossings.exists_perm_nested,Hironaka.SimultaneousNormalCrossings.exists_perm_nested_aux,Hironaka.SimultaneousNormalCrossings.card_filter_le_lt,Hironaka.SimultaneousNormalCrossings.selMap,Hironaka.SimultaneousNormalCrossings.selMap_apply,Hironaka.SimultaneousNormalCrossings.apply_eq_zero_of_eventually_smul,Hironaka.SimultaneousNormalCrossings.norm_zeroCoords_le,Hironaka.SimultaneousNormalCrossings.eq_sum_single_of_support,Hironaka.SimultaneousNormalCrossings.eq_of_eventually_mem_iff,Hironaka.SimultaneousNormalCrossings.pos_of_isBlowupAlong}
\leanok
Let \(E\) be a simple normal crossings divisor transverse to \(N\), let \(C\subset N\), and suppose that \(C\) has only normal crossings
with \(E\). If \(n>0\), then \(C\subset N\) and \(E\) admit adapted coordinates at every point of \(C\). The hypothesis \(n>0\) holds whenever
\(C\) is the centre of a blowup and \(C\neq\emptyset\): in dimension \(0\) every point is open, while a centre has empty interior. (In
dimension \(0\) no member of \(E\) passes through a point, but the index family \((s_i)\) of adapted coordinates is required for every member.)
\end{lemma}
\begin{proof}
\leanok
Two members \(H_i,H_j\) of \(E\) through \(x\) that agree near \(x\) are equal: in the coordinates of \(E\) at \(x\) they are the hyperplanes
\(\{w_a=0\}\), \(\{w_b=0\}\), with \(a\ne b\) if \(i\ne j\), and the points with coordinates \(\tau e_b\), \(\tau\ne0\) small, lie in the first and not in
the second.

Let \(x\in C\). Take coordinates \(u\) at \(x\), centred at \(0\), in which \(C=L_Y=\{u_t=0,\ t\in Y\}\) and each member \(H_i\) through \(x\) is
\(\{u_{t'_i}=0\}\) (normal crossings of \(C\) and \(E\)), and coordinates \(v\) in which \(N=L_{S_0}\) and \(H_i=\{v_{t_i}=0\}\) with
\(t_i\notin S_0\) (transversality). By the first paragraph \(i\mapsto t'_i\) is injective on the members through \(x\). Let \(\xi=v\circ u^{-1}\)
near \(0\) and \(L=d\xi(0)\), an isomorphism. As \(\xi(L_Y)\subset L_{S_0}\) (since \(C\subset N\)), the rows \(L_s\), \(s\in S_0\), vanish on
\(L_Y\). For a member with \(t'_i\in Y\) (that is, \(H_i\supset C\) near \(x\)), \(\xi_{t_i}\) vanishes on \(\{u_{t'_i}=0\}\), so
\(L_{t_i}=c_i\,du_{t'_i}\), and \(c_i\ne0\) as \(L_{t_i}\ne0\). Let \(Y'=Y\setminus\{t'_i:t'_i\in Y\}\). Every \(b\in\mathbf C^{S_0}\) is
\((L_sw)_{s\in S_0}\) for some \(w\) supported on \(Y'\): choose \(w\) with \((Lw)_s=b_s\) for \(s\in S_0\) and \((Lw)_r=0\) otherwise; then
\(c_iw_{t'_i}=(Lw)_{t_i}=0\) for \(t'_i\in Y\) (as \(t_i\notin S_0\)), and replacing \(w\) by its projection to the coordinates in \(Y\)
changes \(w\) by a vector of \(L_Y\), hence does not change the \(L_sw\). So the vectors \((L_se_y)_{s\in S_0}\), \(y\in Y'\), span
\(\mathbf C^{S_0}\); choose \(X\subset Y'\) such that those with \(y\in X\) form a basis, so \(|X|=|S_0|\), and a bijection
\(\beta:X\to S_0\). Put \(F(u)_t=\xi_{\beta(t)}(u)\) for \(t\in X\) and \(F(u)_t=u_t\) otherwise. If \(dF(0)w=0\), then \(w\) is supported on
\(X\) and \(L_sw=0\) for all \(s\in S_0\), so \(w=0\) by the choice of \(X\). By the inverse function theorem
(Lemma~\ref{ana:local-holomorphic-inverse}) \(F\) is a coordinate system on a ball \(B\) around \(0\), and on \(B\):
\begin{itemize}
\item \(\xi(u)\in L_{S_0}\) if and only if \(F_t(u)=0\) for all \(t\in X\), as \(\beta\) is onto \(S_0\);
\item \(u\in L_Y\) if and only if \(F_t(u)=0\) for all \(t\in Y\): if \(u\in L_Y\) then \(\xi(u)\in L_{S_0}\); conversely, let \(u'\) be \(u\) with
the coordinates in \(Y\) put to \(0\); then \(u'\in B\cap L_Y\), so \(\xi(u')\in L_{S_0}\), and \(F(u')=F(u)\), whence \(u=u'\in L_Y\) by injectivity
of \(F\) on \(B\);
\item each member through \(x\) is \(\{F_{t'_i}=0\}\), as \(t'_i\notin X\) (either \(t'_i\notin Y\supset X\) or \(t'_i\in Y\setminus Y'\)).
\end{itemize}
Finally a permutation \(\sigma\) of the coordinates with \(\sigma(t)\ge n-|X|\iff t\in X\) and \(\sigma(t)\ge n-|Y|\iff t\in Y\) (it exists as
\(X\subset Y\)) gives coordinates, composed with \(u\) (Lemma~\ref{ana:regular-transport}), in which \(C=\mathbf C^{n-|Y|}\times0\),
\(N=\mathbf C^{n-|S_0|}\times0\), and each member through \(x\) is the hyperplane with index \(s_i=\sigma(t'_i)<n-|S_0|\), these indices being distinct.
\end{proof}

\begin{theorem}[The transformed divisor]
\label{bcalc:snc-transform}
\uses{bcalc:snc-chart,bcalc:nested-flattening,bcalc:exceptional-divisor,bcalc:strict-submanifold,hir:snc,hir:snc-basic}
\source{bierstone-milman08-functoriality}
\lean{Hironaka.BlowupSNCTransform.IsAdaptedAt,Hironaka.BlowupSNCTransform.exists_coords_over_centre,Hironaka.BlowupSNCTransform.isSNCDivisor_geoTransformE,Hironaka.BlowupSNCTransform.isTransverse_geoTransformE,Hironaka.BlowupSNCTransform.isAdaptedAt_nil,Hironaka.BlowupSNCTransform.isSNCDivisor_preimage_centre,Hironaka.BlowupSNCTransform.isTransverse_preimage_centre,Hironaka.BlowupSNCTransform.length_map_append_singleton,Hironaka.BlowupSNCTransform.get_map_append_singleton,Hironaka.BlowupSNCTransform.eq_of_not_lt}
\leanok
Suppose that \(E\) is a simple normal crossings divisor and that \(C\subset N\) and \(E\) admit adapted coordinates at every point of \(C\). Then
\(E'\) is a simple normal crossings divisor; if moreover \(E\) is transverse to \(N\), then \(E'\) is transverse to
\(N'=\overline{\pi^{-1}(N\setminus C)}\). In particular \((\pi^{-1}(C))\) is a simple normal crossings divisor transverse to \(N'\).
\end{theorem}
\begin{proof}
\leanok
Let \(y\in Y\). If \(\pi(y)\notin C\), then \(\pi\) is a local biholomorphism at \(y\), near which \(H'=\pi^{-1}(H)\), \(N'=\pi^{-1}(N)\) and
\(\pi^{-1}(C)=\emptyset\); the coordinates of \(E\) (and \(N\)) at \(\pi(y)\) pull back (Lemma~\ref{bcalc:snc-chart}). If \(\pi(y)\in C\), take adapted
coordinates \(\Phi\) at \(\pi(y)\) and the monomial chart \(j\ge k_C\), \(O\), \(W\) at \(y\). By Lemma~\ref{bcalc:snc-chart}, near \(y\):
\(\pi^{-1}(C)=\{W_j=0\}\); for a member \(H\) through \(\pi(y)\), \(H'=\{W_{s_i}=0\}\) if \(s_i\ne j\) and \(H'\) misses \(O\) if \(s_i=j\) (then \(H\supset C\)
has the centre as its trace in codimension one); \(N'=\{W_t=0,\ t\ge k_N\}\) if \(j<k_N\), and \(N'\) misses \(O\) otherwise. A member not through
\(\pi(y)\) has a strict transform not through \(y\) (\(H'\subset\pi^{-1}(H)\), \(H\) closed). In the coordinates \(W-W(y)\) the members of \(E'\) through \(y\)
are therefore the hyperplanes with the indices \(s_i\) (distinct, and \(\ne j\)) and \(j\): \(E'\) is simple normal crossings at \(y\). If \(y\in N'\), then
\(j<k_N\), all these indices are \(<k_N\), and \(N'\) is cut out by the coordinates \(t\ge k_N\): \(E'\) is transverse to \(N'\) at \(y\). The last
statement is the case \(E=\emptyset\), where adapted coordinates are nested flattening coordinates (Lemma~\ref{bcalc:nested-flattening}).
\end{proof}

\begin{corollary}[The transformed divisor of an admissible blowup]
\label{bcalc:snc-transform-nc}
\uses{bcalc:snc-simultaneous,bcalc:snc-transform,hir:snc}
\source{bierstone-milman08-functoriality}
\lean{Hironaka.SimultaneousNormalCrossings.isSNCDivisor_geoTransformE_of_hasNormalCrossings,Hironaka.SimultaneousNormalCrossings.isTransverse_geoTransformE_of_hasNormalCrossings,Hironaka.SimultaneousNormalCrossings.isTransverse_univ_of_isSNCDivisor}
\leanok
Let \(E\) be a simple normal crossings divisor and let the centre \(C\) of the blowup \(\pi\) have only normal crossings with \(E\). Then
\(E'=(H_1',\dots,H_s',\pi^{-1}(C))\) is a simple normal crossings divisor. If moreover \(E\) is transverse to a closed set \(N\supset C\), then
\(E'\) is transverse to \(N'=\overline{\pi^{-1}(N\setminus C)}\).
\end{corollary}
\begin{proof}
\leanok
A simple normal crossings divisor is transverse to the whole manifold (cut out by no coordinate). Lemma~\ref{bcalc:snc-simultaneous}, applied
to \(C\subset X\), respectively to \(C\subset N\), gives adapted coordinates at every point of \(C\) (\(n>0\) as \(C\) is a centre, and there is
nothing to prove if \(C=\emptyset\)); conclude by Theorem~\ref{bcalc:snc-transform}.
\end{proof}

\begin{corollary}[The transform of a marked ideal is well formed]
\label{bcalc:transform-wellformed}
\uses{bcalc:snc-transform,bcalc:snc-simultaneous,hir:geo-transform,hir:geo-transform-wellformed,hir:marked-ideal,hir:hadamard-bridge-along}
\source{bierstone-milman08-functoriality}
\lean{Hironaka.GeoTransformWellFormed.isWellFormed_geoTransform_of_isAdmissibleCentre,Hironaka.GeoTransformWellFormed.isWellFormed_geoTransform}\upstreamlean{TranscendentalNakamaye.Pending.GeoTransformWellFormed.isWellFormed_geoTransform_of_E_eq_nil}\lean{Hironaka.GeoTransformWellFormed.isWellFormed_geoTransform_of_adapted}\upstreamlean{TranscendentalNakamaye.Pending.GeoTransformWellFormed.isWellFormed_geoTransform_of_E_eq_nil'}
\mixedleanok
Let \(\underline{\mathcal I}=(M,N,E,\mathfrak a,d)\) be a well-formed marked ideal and \(C\) an admissible centre. Then the transform of
\(\underline{\mathcal I}\) by the blowup along \(C\) is well formed.
\end{corollary}
\begin{proof}
\mixedleanok
Since \(C\subset\operatorname{cosupp}\underline{\mathcal I}\), \(\mathfrak a|_N\subset\mathcal I_C^d|_N\) (Lemma~\ref{hir:hadamard-bridge-along}). Then \(N'\) is a closed
submanifold and the transformed ideal is coherent (Lemma~\ref{hir:geo-transform-wellformed}). As \(\underline{\mathcal I}\) is well formed and \(C\)
is admissible, \(E\) is a simple normal crossings divisor transverse to \(N\supset C\) with which \(C\) has only normal crossings, so \(C\subset N\)
and \(E\) admit adapted coordinates at every point of \(C\) (Lemma~\ref{bcalc:snc-simultaneous}; \(n>0\) as \(C\) is a nonempty centre, and there
is nothing to prove if \(C=\emptyset\)), and \(E'\) is a simple normal crossings divisor transverse to \(N'\) by Theorem~\ref{bcalc:snc-transform}.
\end{proof}

\paragraph{Dropping members that contain the centre.} Step~1a of the algorithm replaces \(E\) by \(E_S\), the divisor without the members
that contain the submanifold \(S\) (Włodarczyk 2005, Step~1a; Lemma~\ref{hir:drop-containing-members}). The submanifold, the ideal and the
cosupport of a marked ideal and of its transforms do not depend on \(E\); admissibility depends on \(E\) only through the normal crossings of the
centre with \(E\), and the members containing the centre do not change it.

\begin{lemma}[Sub-divisors]
\label{bcalc:sub-divisor}
\uses{hir:snc,hir:marked-ideal}
\source{bierstone-milman08-functoriality}
\lean{Hironaka.NormalCrossingsComparison.hasNormalCrossings_of_forall_mem,Hironaka.NormalCrossingsComparison.isTransverse_of_forall_mem,Hironaka.NormalCrossingsComparison.isSNCDivisor_of_sublist}\upstreamlean{TranscendentalNakamaye.Pending.NormalCrossingsComparison.isWellFormed_withE}
\mixedleanok
Let \(E\) be a finite sequence of closed subsets of \(X\) and \(E'\) a finite sequence each of whose members is a member of \(E\).
(1) If \(C\) has only normal crossings with \(E\), it has only normal crossings with \(E'\); if \(E\) is transverse to \(N\), so is \(E'\).
(2) If \(E\) is a simple normal crossings divisor and \(E'\) is a subsequence of \(E\), then \(E'\) is a simple normal crossings divisor.
(3) If \(\underline{\mathcal I}=(M,N,E,\mathfrak a,d)\) is well formed and \(E'\) is a subsequence of \(E\), then \((M,N,E',\mathfrak a,d)\) is well formed.
\end{lemma}
\begin{proof}
\mixedleanok
(1) The coordinates at a point \(x\) of \(C\) (of \(N\)) given by the hypothesis serve for \(E'\): every member of \(E'\) through \(x\) is a member of
\(E\) through \(x\), hence a coordinate hyperplane (whose coordinate is not among those cutting out \(N\)). (2) A subsequence is indexed by a strictly
increasing, hence injective, map of the indices; composing the hyperplane indices of \(E\) at \(x\) with it keeps them injective on the members through
\(x\). For an arbitrary sequence of members this fails: a member through \(x\) repeated twice would need two distinct hyperplanes. (3) \(N\) and
\(\mathfrak a\) are unchanged; apply (1) and (2) to \(E'\subset E\).
\end{proof}

\begin{theorem}[Normal crossings and the members containing the centre]
\label{bcalc:nc-comparison}
\uses{bcalc:sub-divisor,bcalc:snc-simultaneous,hir:snc,hir:transform,ana:local-holomorphic-inverse}
\source{bierstone-milman08-functoriality,wlodarczyk05-hironaka}
\lean{Hironaka.NormalCrossingsComparison.hasNormalCrossings_iff_of_forall_mem,Hironaka.NormalCrossingsComparison.hasNormalCrossings_of_forall_notMem,Hironaka.NormalCrossingsComparison.IsPartialNCAt,Hironaka.NormalCrossingsComparison.isPartialNCAt_of_hasNormalCrossings,Hironaka.NormalCrossingsComparison.IsPartialNCAt.insert,Hironaka.NormalCrossingsComparison.mem_complexCoordinatePlane_pred_iff}\upstreamlean{TranscendentalNakamaye.Pending.NormalCrossingsComparison.hasNormalCrossings_filter_iff,TranscendentalNakamaye.Pending.NormalCrossingsComparison.hasNormalCrossings_drop_iff}\lean{Hironaka.NormalCrossingsComparison.isAdmissibleCentre_withE_iff}
\mixedleanok
Let \(E\) be a simple normal crossings divisor on \(X\), \(C\subset X\), and \(E'\) a finite sequence of members of \(E\) such that every member of
\(E\) that is not a member of \(E'\) contains \(C\). Then \(C\) has only normal crossings with \(E'\) if and only if it has only normal crossings with
\(E\). In particular, if \(A\) is a set of indices with \(C\subset H_i\) for \(i\in A\), then \(C\) has only normal crossings with \(E\) if and only if
it has only normal crossings with the subsequence \(E\setminus A\) of the members with index not in \(A\); and if \(E\) is the divisor of a marked
ideal \(\underline{\mathcal I}=(M,N,E,\mathfrak a,d)\), then \(C\) is an admissible centre for \(\underline{\mathcal I}\) if and only if it is one for
\((M,N,E',\mathfrak a,d)\).
\end{theorem}
\begin{proof}
\mixedleanok
``If'' is Lemma~\ref{bcalc:sub-divisor}~(1). Conversely, let \(x\in C\). Call holomorphic coordinates \(u\) at \(x\), centred at \(0\), \emph{adapted to
a set \(I\) of members} if \(C\) is a coordinate subspace \(L_Y=\{u_t=0,\ t\in Y\}\) near \(x\) and every member of \(I\) through \(x\) is a coordinate
hyperplane \(\{u_{p_i}=0\}\). Coordinates of the normal crossings of \(C\) with \(E'\) are adapted to the set of members of \(E\) that are members of
\(E'\). Let \(u\) be adapted to \(I\) and let \(H_0\notin I\) be a member containing \(C\), so \(x\in H_0\). In simple normal crossings coordinates \(v\) of
\(E\) at \(x\) the members through \(x\) are \(\{v_{t_i}=0\}\) with distinct indices, so \(t_i\ne t_0\) for the members of \(I\) through \(x\).
Read \(\xi=v\circ u^{-1}\) near \(0\); \(\xi(0)=0\). Since \(C\subset H_0\), \(\xi\) maps \(L_Y\) into \(\{v_{t_0}=0\}\), and \(\xi_{t_i}\) vanishes on
\(\{u_{p_i}=0\}\) for the members of \(I\) through \(x\). This is the situation of the proof of Lemma~\ref{bcalc:snc-simultaneous}, with
\(S_0=\{t_0\}\) (the hypersurface \(H_0\) in the role of \(N\)) and the members of \(I\) in the role of the members through \(x\): the construction there
gives \(X\subset Y\) with one element \(y_0\), different from every \(p_i\), and the coordinates \(F\) obtained from \(u\) by replacing \(u_{y_0}\) by
\(\xi_{t_0}\) (a local biholomorphism by Lemma~\ref{ana:local-holomorphic-inverse}, followed by a permutation of the coordinates), in which near \(0\)
\(L_Y\) is a coordinate subspace, \(\xi^{-1}(\{v_{t_0}=0\})\), that is \(H_0\), is a coordinate hyperplane, and the coordinates \(u_{p_i}\) are kept.
As \(\xi(0)=0\), \(F(0)=0\). So \(F\circ u\) is adapted to \(I\cup\{H_0\}\). After finitely many steps, one for each member of \(E\) that is not a member
of \(E'\), the coordinates are adapted to all members of \(E\): \(C\) has only normal crossings with \(E\).

For the index form take \(E'=E\setminus A\): a member \(H_i\) of \(E\) that is not a member of \(E\setminus A\) has \(i\in A\), so it contains \(C\). The
cosupport of a marked ideal does not depend on \(E\), so admissibility (Definition~\ref{hir:transform}) is compared through the normal crossings alone.
\end{proof}

\begin{lemma}[Dropping members along one blowup]
\label{bcalc:nc-comparison-transform}
\uses{hir:geo-transform,hir:snc}
\source{wlodarczyk05-hironaka}
\lean{Hironaka.NormalCrossingsComparison.geoStrictTransform_mono,Hironaka.NormalCrossingsComparison.forall_mem_transformE}\upstreamlean{TranscendentalNakamaye.Pending.NormalCrossingsComparison.sublist_transformE}\lean{Hironaka.NormalCrossingsComparison.forall_notMem_transformE,Hironaka.NormalCrossingsComparison.geoTransform_withE}
\mixedleanok
Let \(E\) and \(E'\) be finite sequences of subsets of \(X\) such that every member of \(E'\) is a member of \(E\) and every member of \(E\) that is not a
member of \(E'\) contains \(S\). Let \(\pi\colon Y\to X\) be a map, \(C\subset X\), and let \(E^\pi=(H_1',\dots,H_s',\pi^{-1}(C))\), with
\(H'=\overline{\pi^{-1}(H\setminus C)}\), be the transformed sequence, \(E'^\pi\) that of \(E'\). Then every member of \(E'^\pi\) is a member of \(E^\pi\)
(and \(E'^\pi\) is a subsequence of \(E^\pi\) if \(E'\) is one of \(E\)), and every member of \(E^\pi\) that is not a member of \(E'^\pi\) contains
\(S'=\overline{\pi^{-1}(S\setminus C)}\). The transform of \((M,N,E',\mathfrak a,d)\) (Definition~\ref{hir:geo-transform}) is the transform of
\((M,N,E,\mathfrak a,d)\) with \(E^\pi\) replaced by \(E'^\pi\).
\end{lemma}
\begin{proof}
\mixedleanok
\(A\mapsto\overline{\pi^{-1}(A\setminus C)}\) is monotone. A member of \(E^\pi\) that is not a member of \(E'^\pi\) is not \(\pi^{-1}(C)\), which belongs to
both, so it is \(H'\) for a member \(H\) of \(E\); \(H\) is not a member of \(E'\) (otherwise \(H'\) would be a member of \(E'^\pi\)), so \(S\subset H\) and
\(S'\subset H'\). The submanifold, the ideal and the mark of the transform do not involve \(E\).
\end{proof}

\begin{lemma}[Generators of the ideal of a submanifold have multiplicity one]
\label{bcalc:multiplicity-one}
\uses{bcalc:giraud-strict,bcalc:strict-submanifold,hir:derivative-ideal,ana:local-holomorphic-inverse}
\source{wlodarczyk05-hironaka}
\lean{Hironaka.GeneratorMultiplicityOne.surjective_fderiv_toFun}\upstreamlean{TranscendentalNakamaye.Pending.GeneratorMultiplicityOne.exists_chartDeriv_ne_zero_of_vanishingIdeal_eq_span,TranscendentalNakamaye.Pending.GeneratorMultiplicityOne.exists_chartDeriv_ne_zero_controlledTransform}
\mixedleanok
Let \(N\) be regular of dimension \(k<n\) at \(y\), and suppose \(\mathcal I_N=(g)\) at \(y\) with \(g\) holomorphic. Then \(dg(y)\ne0\): some partial
derivative of \(g\) in a chart does not vanish at \(y\). Consequently, in the situation of Corollary~\ref{bcalc:giraud-strict} with \(N\) of
dimension \(k<n\) at each of its points, the controlled transform \(u'\) satisfies \(du'(z)\ne0\) at every \(z\in N'\) (Włodarczyk's Lemma on
tangent directions, part (2)).
\end{lemma}
\begin{proof}
\mixedleanok
Let \(\Phi\) flatten \(N\) onto \(\mathbf C^{k}\times0\) at \(y\). The coordinate \(\Phi_k\) vanishes on \(N\), so \(\Phi_k=a\,g\) near \(y\) with \(a\) holomorphic;
and \(g(y)=0\) since \(g\) itself lies in \(\mathcal I_N\) and \(y\in N\). Hence \(d\Phi_k(y)=a(y)\,dg(y)\). The differential of \(\Phi\) is invertible at
\(y\) (it has a holomorphic inverse, Lemma~\ref{ana:local-holomorphic-inverse}), so \(d\Phi_k(y)\ne0\) and therefore \(dg(y)\ne0\). For the
consequence, \(N'\) is regular of dimension \(k\) at \(z\) (Theorem~\ref{bcalc:strict-submanifold}) and \(\mathcal I_{N'}=(u')\) at \(z\)
(Corollary~\ref{bcalc:giraud-strict}).
\end{proof}
